% source: https://learning.edx.org/course/course-v1:HarvardX+TinyML2+1T2022/home
% Applications of TinyML
% google log in

% Bilder TinyML/Applications

% Reihenfolge
%1. TinyML/Fundamentals
%2. TinyML/Applications
%3. TinyML/deployment
%4. TinyML/MLOps for Scaling TinyML


\chapter{TinyML Applications}


\section{Welcome to TinyML applications}

VIJAY JANAPA REDDI: Welcome to course 2 where
we're going to focus on Tiny Machine Learning applications.
Now one of the exciting things about this particular course
is that we're going to learn the nuts and bolts behind TinyML applications.
We're going to take three different application scenarios
and try and understand how we build an end-to-end system for TinyML.
Now some of you might be joining us directly in course 2.
So one of the key things that I'm going to do right
at the beginning of the course is to kind of do a quick recap, a fast recap,
just to make sure that you're comfortable with the material
that we will be covering in course 2.
For those of you who are coming from course 1, it might be a quick recap.
Bear with me.
Let's start off with the fundamentals.
Again, what is tiny machine learning?
Tiny machine learning is a fast-growing subfield within machine learning.
One of the key and most interesting things about tiny machine learning
is not just that it focuses on machine learning algorithms,
but it, specifically, brings about the interesting aspects
of embedded hardware and software that make it quite unique,
as compared to any other machine learning class that's out there.
One of the key things about this particular domain
is the fact that much of the machine learning that we're going to be doing
is going to be processing sensor data.
And the key thing is that we're going to be processing all that sensor data
on device, not offloaded to the cloud.
We're going to do it on the physical metal that's right in front of us.
One of the key requirements in TinyML is that we've
got to do all of this with ultra-low power consumption.
Now, if you showed up to me, and you said you know TinyML,
one of my questions to you is going to be what's the rough power
budget for a TinyML device.
The answer has got to be that it's, typically, in the order of milliwatts.
One of the reasons we want to do this in such low power consumption
is because we want to run machine learning all
the time, always-on machine learning.
And there's several interesting application use cases.
We're going to touch on a few of them and go pretty deep into them
in this particular course.
The key thing, at the end of the day, is that you've
got to package up all of this stuff, and you
need to run these machine learning models
on these tiny little microcontroller devices
and do so on battery-operated power, which
means you've got to be able to take a small, little coin cell battery
and then be able to power that device off of that battery, which presents us
a remarkable challenge.
And there are lots of fun things in trying to tackle those challenges.
You've heard about a couple of these in the previous course,
and I'm going to take you much deeper by focusing on the application
aspects in this particular course.
In course 1, we learned the language of TinyML , which
was the fundamentals of TinyML, which is great if you don't know anything
about machine learning.
That's why we intended it to be the basics.
Now, if you knew about machine learning, hopefully, it was a good reminder.
And, if you knew about machine learning, hopefully,
you picked up a couple of new, interesting things
about what's tiny in the context of machine learning.
The whole resource requirements of embedded devices
plays extreme challenges for us.
So, hopefully, you picked up a little bit of that.
Again, we're going to go deeper into that in course 2.
Course 2 is where we try and understand how
to engineer some TinyML applications.
Now, you might be wondering, what's an example of a TinyML application?
Remember, I showed you this little snippet in course 1.
When you say "OK, Google," the machine wakes up.
How does it know how to do that?
How does it take an analog signal, my voice or speech,
and then know how to convert that into some sort of digital signal
that actually triggers the machine to wake up?
Alternatively, you might have said "Hey, Siri" on your watch, which
is running nearly all day.
Or what about saying "Alexa," which makes up and then
responds back to you when you say "Alexa."
Pretty remarkable things. These things are already being deployed here.
Obviously, when we're talking about TinyML,
we're talking about untethering these devices from the wall and, actually,
being able to power them on coin cell kind of batteries
and allow them to operate for a very, very long time.
I'm talking about weeks, months, sometimes,
even years if you remember what I said in course 1.
All right, that's about applications that we already
have here and right now.
But I also mentioned, in course 1, that we're
going to be looking at more forward-looking applications, things
that have to do with robotics, things that have to do with smart glasses
where you're embedding small, little sensors that are constantly
perceiving the environment and doing a lot of machine learning processing.
There are also many different kinds of medical applications
that we touched upon.
We're going to try and understand what the fundamental building
blocks behind all of these kinds of different use cases are.
Specifically, we're going to do that by looking deeper in this course
into the sensory aspects.
Remember, I said that the exciting things about TinyML
is the fact that the input data is coming from sensors that
are external to the computing system.
Now they could be variety of different kinds of sensors.
We're talking about motion sensors, acoustic sensors,
environmental sensors, touch screen sensors, image sensors,
biometric sensors.
Oh, there's a world of sensors out there.
I cannot even count them all.
In this course, we're going to focus on acoustic sensors.
We're going to focus on image sensors.
We're going to focus on motion sensors.
There's a very good reason that I've picked
these three different kinds of sensors in order to teach you all about TinyML.
I'll get into that later.
But the key thing is that these three sensors
are going to be our input signals.
We're not yet going to deal with the physical signal in this course.
We're going to deal with the data that's coming in from the signals.
In course 3, we'll get more into the physical aspects of the sensors.
So the key nugget is we have data that's coming in from the sensors.
We need to understand a little bit about that.
I'll go into that a little bit more deeper soon.
The next step is to understand what it means
to take the input data that's coming in from these sensors
and be able to build a machine learning model.
One of the interesting properties about each of those sensors that I touched on
is that the data that they present the neural network is different.
So it's not like you're going to be rinsing and repeating
the exact same thing for all three sensors.
You are actually going to learn something interesting
that's unique about each of those sensors
and understand how the neural network actually
needs to adapt in order to be able to cope with the input signal.
And, finally, we're going to learn how to assess
the performance of the neural network.
You might be thinking, hey, I just did some classification.
So what I care about is accuracy.
Well, that's a little naive because accuracy is a rather delicate thing.
It's one thing for your model to be able to predict something accurately
in your Google Colab.
It's a whole different thing for your model
to be able to perform well in the real world.
And one of the things that I really want to do in this course
is try and get across the importance of looking at the end-to-end TinyML
application design.
Most of the courses that are out there, most of the books that are out there,
often, focus just on the middle portion.
They tend to focus on the neural network design only.
I'm not saying that that's not important.
It is very important to focus on that, but, look, at the end of the day,
if you want to build an application, you
got to understand the whole pipe.
And that's what's going to make you different
because one of the key nuggets that I'm going to focus on
in this particular class is taking you through the machine learning lifecycle
for building and deploying a TinyML application.
With that, I'll see you in the next video
where we'll go a little bit deeper into some of these aspects.
First and foremost, I'm going to start off
with a quick recap of what we covered in course 1
so that we can level the playing field for everybody.
I'll see you in the next video.

\section{Building Blocks (from Course 1)}

VIJAY JANAPA REDDI: Hi.
In this video, I'm going to focus on the building blocks
that we need from Course 1 to be successful in Course 2.
Specifically, I'm going to touch on the language that we
use in Course 1 and the terminologies that we introduce
and the concepts that you practiced, so that it
can be a refresher for the things that we're going to be doing in Course 2.
I'm going to try and make a couple of connections
about how that wealth of material that you picked up in Course 1
is actually going to be applicable here in Course 2.
Now remember, in the last video I said, that we're
going to be looking at three different sensors,
and then we're going to be looking at the unique characteristics
about these sensors, and then looking at the networks
that we're going to be training for them.
And then understand how we assess those networks,
in the context of a TinyML application use
case, which means looking beyond just the typical model accuracy,
but trying to understand the nuances that are there about how
our application might really perform.
And is looking at just a neural network model performance enough,
or the accuracy enough, or are the other metrics
that we really have to consider?
All that said, let's first start with the basics
of doing a total recap on the language that we introduced in Course 1.
Course 1, almost exclusively focused on neural networks.
We wanted to make sure that you understood
all you needed to know to be successful in Course 2
from the machine learning perspective.
What did we touch upon?
I'm going to throw out a couple of terms here.
This is not the full gamut of terms that we touched upon.
So, I want to make sure that some of these key terms that we
touch on actually resonate and remind you about the knowledge you picked up.
Or my hope is that if you're coming straight into Course 2,
that you have familiarity with these terms,
because it's a fundamental thing at this point.
Because I'm going to be moving up a little faster, in terms of how
we cover material in this course.
First and foremost, I hope you know what a neural network is
and where it stands in the big context of AI.
Remember, AI goes into machine learning, further narrows into deep learning.
And that's where your neural networks come.
So it's a very specific type of artificial intelligence
that we're talking about.
It is not the only solution.
Hopefully, you remember a little bit of that.
Then we talked about training these networks using gradient descent.
What exactly is gradient descent?
How do you actually optimize a network to learn by using loss functions?
Specifically, what are you training in the network?
Kernels and filters.
What's the difference between them?
Do you know the difference between a CNN or convolutional neural network
versus a more traditional deep neural network?
I'm kind of referring to MLPs or multilayer perceptrons.
We have touched on both of these kinds of networks in Course 1.
And in this course, we are going to be using
both CNNs and traditional MLP/DNNs.
So you need to be familiar with both.
Furthermore, we did a lot of training in Course 1,
and then we did a couple of inferences on the MNIST data set and the MNIST
fashion data set.
We're also going to be doing some more training.
We're going to be doing training on data that corresponds to sensors.
And then we're going to learn how to do inferences on them,
and then learn how to assess the inference capabilities of the TinyML
application.
Well, before you can train a network, you have to extract key features.
We're going to be doing a lot of feature engineering in this particular course.
Feature engineering, that specifically looks at the traits of the sensors,
so each of the sensor requires different kinds of feature engineering.
We're going to talk about that.
At this point, do you recognize the word features.
What do they really mean?
Think about the images that you looked at in Course 1.
What exactly are we trying to extract?
What is it that the convolutional kernels
are doing as the image is moving through the different stages of the network?
Why do you have multiple convolutional layers?
Why not just have one big convolutional layer?
They're extracting different features.
I hope that resonates with you.
Then we talked about overfitting.
When you're training a network over the data that you're given,
the network can become overfitted to that particular training
data, which can be a real problem when you deploy it.
Because in deployment, or during deployment,
you are not necessarily going to get that exact kind of data.
You might get data that is a little off.
But your network might not actually generalize well.
So you need to avoid overfitting.
And we learned how to do that.
We also talked about how to get data into the network.
How do you do data augmentation?
Why is data augmentation extremely critical?
In this course, we're going to go really deep into data engineering,
because we're going to look at things that are well beyond data augmentation.
And it's extremely crucial because it's garbage
in garbage out with neural networks.
And we need to know how to preprocess the data.
We touched a very little bit of preprocessing in the Course 1.
And in Course 2 we're going to do a lot of preprocessing.
We're going to understand how you preprocess data differently
for different types of sensors.
So hopefully, you're comfortable with the notion of preprocessing.
Finally, when we talk about the data set that we're using,
the same data can be used for multiple different things.
You usually start with the large data.
And you take that large data and you break it up into pieces.
You get your training data.
Then you get your validation data, and you get your test data.
Hopefully, you know what each of these means.
We will actually be recapping some of these things further on.
So hopefully that will help you connect the dots back to Course 1.
Then we built different kinds of machine learning models, not necessarily just
neural networks.
You learned about regression as one of the first things.
Regression is all about fitting a line and trying
to minimize the error between a scattered set of points.
Then we learned about classification.
How can we get convolutional neural nets to be
able to tell us if it's a cat or a dog, or what
the digit is in an MNIST data set.
All of this put together, one of the key things that I wanted to emphasize,
and hopefully you got across, was that it's extremely important for us
to be responsible in order to deploy TinyMLs safely into the real world.
Look, all too often many of us make this mistake
that we get very excited about the actual mechanics-- the nuts and bolts,
and how to build things.
And then we forget the consequences of putting this technology out
into the real world.
Organizations, big organizations, world leading machine learning organizations,
like Google, Amazon, Facebook, Microsoft and so forth, they're all increasingly
expecting their engineers to be responsible AI engineers.
This is the reason that in every single one of these courses,
we repeatedly emphasize to you that you need to be aware about responsible AI
design.
You need to be aware of responsible AI development.
You need to be aware about responsible AI deployment.
We talked about responsible AI design in course 1.
We're going to talk about responsible AI development in Course 2.
And Susan Kennedy is going to lead that effort in this course.
Now, one of the things that we talked about in Course 1 was look,
we wanted to make models really accurate.
And so the models ended up getting really big.
We got into VGG-16 and so forth, as models were evolving in order
to improve the accuracy.
This is from the ImageNet accuracy evaluation.
This plot that has kind of faded in the background,
is actually looking at the accuracy on the y-axis.
And the x-axis is showing you the computational requirements.
Now in Course 1, we talked about how we ended up with VGG and said,
that, wow, the models are getting way too big.
And that there was a real need to actually start
building smaller models that are also accurate and try and deploy them.
This transition happened as we started building machine
learning models that are getting deployed in the cloud,
to then trying to deploy them in smartphone kind of devices.
And now, as we're looking towards the next big wave,
in terms of deploying things onto smaller networks,
we need to think even more carefully.
So one of the key things that I introduced at a very high level--
I only introduced the language, I did not explain the mechanics,
was how can we take that network, a big network, and shrink it down?
One of the ways we can do that is by pruning the network.
We can prune the synapses.
That is, we can cut out some of the connections.
Or we can also prune the neurons all together,
which is eliminate some of the neurons, so that way
you don't have to do the weight calculations on them,
which reduces the amount of computational complexity
of the network.
And it also reduces the total number of neurons or weights
that you have in the network.
And in doing so, you shrink the model size,
and you also reduce the amount of computational power needed.
We only touched on the language of this.
We did not actually explain the mechanics.
In this course, I'm going to open up that box.
And I'm going to step us into that box, so we
understand what it means to actually be able to do these optimizations.
What are the optimizations?
So we'll spend a lot of time on this.
It's extrememly important as you try and make anything tiny, is quantization.
The idea that instead of using many bytes for representing values.
For example, if you use a float32, that's
going to require 4 bytes of memory storage.
Instead of the 4 bytes, if you can compress those values down
into an int8 byte, which is a single byte, which can represent values
from negative 128 to 127, then that gives you a 4x reduction
in the network straight out.
But of course, there's a trade off.
You make trade offs in accuracy.
We only introduce the term quantization.
In this course, we're actually going to learn
how to program for quantization, specifically, that's one.
Two, we're also going to open that up and understand what is actually
happening when you do quantization.
Quantization is not as simple as it seems.
There are a lot of nuances around it.
So I want to give you a firsthand experience about that.
And it's extremely important for TinyML engineers to know this,
because quantization is going to be used over and over
and over and over in something like TinyMLs.
So you need to master this.
We also touched on this very high level notion
that, look, when you're training models you tend to use TensorFlow.
But as you start moving towards smaller devices, like smartphones and embedded
microcontrollers, we have to use different frameworks.
We have to use something like TensorFlow Lite.
The question is why.
Because you're trying to deploy them on resource constrained things.
I gave you a very high level intuition about who
uses TensorFlow versus TensorFlow Lite.
Person who is doing researching into the machine learning space
might use TensorFlow.
Someone who is actually trying to deploy the model
is going to care about TensorFlow Lite.
We're going to open up that box, and we're
going to go pretty deep into understanding
what are the fundamental differences.
This is going to help you be a better engineer,
because when you're out there, trying to build TinyML applications
and then deploy them, you're going to need to know the differences.
Very often, most people don't know the difference
between TensorFlow and TensorFlow Lite.
In fact, you can build a model and you can run a model completely
in TensorFlow.
But in the real world, you would never do that.
Why?
Well, you'll find out in this course.
So the other thing that I touched a lot on when you're deploying things
on to small little embedded devices there
are these big trade offs between large computing
systems and the small embedded microcontrollers.
I really hope you took away some key principles down there,
and not brush it away as, oh, this is some sort of high-level concept
that I don't need to worry about.
Because if you did do that, it's going to come back and bite you
when you're trying to go out for an engineering interview.
The reason I say that is because there are
massive differences between a traditional large computing
system like a laptop versus what happens in an embedded desktop.
One of the things that I'm going to ask you right now
is, do you know why the yellow bar here says, that the smaller embedded
system has lower code portability?
Think about it for a second.
And if you've taken Course 1, I want to make sure you actually
know how to answer that.
Why is it that in a larger system you have universal code portability
versus in a smaller system you don't have universal code portability?
If you don't know this, go back and look at the material in Course 1.
It's important to understand this, especially as we go into Course 3,
it's going to become an even more important.
And then, looking at the green rows below.
Why are embedded system's cost low?
Why do they consume very little power?
Why is it that the engineering effort for them
is very low compared to the larger systems?
I'll give you a hint.
Cost. Think about the fact that a traditional embedded system might not
have the ability to do very simple floating point math.
For instance, when I said that pi can be approximated using 22 by 7,
where the first three digits are 3.14.
I don't necessarily need to have a floating point arithmetic to get that.
I could use fixed point arithmetic to do that.
Well, if I take that floating point arithmetic out,
it automatically drops the power consumption
inside my processor in the embedded system.
If I take that floating arithmetic out, it
lowers the cost, because I don't need to put much engineering effort into it.
So embedded systems, as much as they're pervasive and they're about 250 billion
of them all over the place, there are trade offs you make.
And those tight trade-offs, those tight corners we cut,
that's the reason that it makes deploying machine running on them hard.
But then again, they're everywhere.
And that's what's exciting, and that's what opens up the space.
So please take a moment to kind of reflect on some of these differences
that we have between embedded systems and non-embedded systems.
It's extremely crucial.
There's a lot of material out there as well.
I'll throw out a couple of pointers that are actually
helpful for a quick recap on what the differences are if you need it.
So as a quick recap in Course 1, we talked
about the fundamentals of TinyML, tried to understand the language.
That was the key thing we focused on-- understanding the language
and then doing the hands-on programming with using Google Colab.
In this course, we're going to focus on taking the language that we learn
and the hands on building blocks about the neural networks,
and we're going to learn how to put it all together in order
to solve applications, which I will talk about in greater
depth in the next video.
With that, I'll see you soon.

\section{What You'll Learn in This Course}

VIJAY JANAPA REDDI: Hi.
In this particular video, I'm going to focus
on what are the critical nuggets you're going to learn in course 2.
Recall that I keep showing you this diagram
that we are going to focus on sensors, neural network models,
and talk about different kinds of metrics.
What I want to do now is actually take each one of these
and kind of talk a little bit more in depth
about what you're going to pick up.
Once we get into the meat of the course, there's
going be a lot of information coming at you.
Trust me.
It's going to be coming at a good pace.
You're going to pick up a lot of good knowledge,
but what I want to try and make sure is you sort of understand
how things are laid out for you.
Why are we picking these three different types of sensors?
TinyML is all about sensors, and I showed you
earlier the many different types of sensors.
So why did I pick these three particular types of sensors
to teach you about TinyML?
Acoustic sensors, image sensors, and motion sensors,
let me ask you a question.
What is different about each of these sensors?
What does an acoustic sensor hear?
And what does it generate as its output?
An acoustic sensor is, basically, picking up sounds,
for instance, if we use microphones.
That's an audio signal.
Then how do you take an audio signal and then try and preprocess it
in a way that you can feed it into a machine learning model?
You'll be quite surprised at how we do this actually.
Then, an image sensor, an image sensor is a camera sensor.
We're going to be using a traditional image sensor in the sense
that it's actually going to generate pictures.
If it's a picture, you might be using a traditional convolutional-based neural
network.
Ah, but you might say, well, I already dealt with pictures in course 1
or have dealt with pictures in the past.
Ah, well, you're still going to learn something new.
The reason is because we're going to be looking at how to actually reduce
the amount of time that's actually needed
to process this input that's coming in.
The third one is accelerometer.
Where do you find them?
Ah, we find them in industrial applications,
and motion sensors would generate time series data.
We'll talk about this a bit more, but the key thing
is that each of these sensors are quite different.
And that's what's really exciting about them.
The industry, largely, bins these sensors
into three Vs. The first one being voice for microphones or acoustics,
second one image, dealing with vision aspects, the third one being motion,
dealing with vibration kind of data, using different kinds of sensors.
So you see these three axes down here?
This, effectively, allows us to teach you very many interesting things
because we're looking at very different dimensions of how
the sensors are actually used in the TinyML application space.
So we'll try and understand what the modalities of these different sensors
are and what kind of data that they generate.
Some of them generate more data.
Some of them generate less data.
Some of them generate time series data.
Some of them generate images.
So there are very different nuances about these things,
and we'll learn about this.
And that's what's going to make it very interesting to look
at the signal characteristics of these different sensors.
Then there's an aspect about the network itself.
What is it that our network has to learn from these different kinds of sensors?
Because the input is, actually, different,
you might have to do different things within the network.
And that's what's going to be the interesting thing.
For acoustic, you're going to do something different.
For image, you're going to do something different.
For the accelerometer kind of data, you're going to do something different.
In fact, in the last one, it's actually going to be quite interesting.
You'll be surprised about the kind of network we're going to look at.
So you should ask yourself, are all neural networks the same in order
to deal with sensors?
That's the key question that we're going to try and answer in this course.
Then, of course, what is the metric that we're going after?
Is there a single metric that matches all of those different sensor inputs
and how the network processes them and generates the output?
Or is there a unique sort of a metric for every different kind of sensor
input?
Should we consider just a single metric when we
we're evaluating a particular sensor?
Or should we actually consider multiple metrics?
So we're going to learn quite a bit about that.
And, hopefully, in doing that, you will pick up the nuances
that it's not always as simple as it seems on the surface.
And that's what's going to make you different when you say,
hey, I'm actually a TinyML engineer.
And, specifically, when we're talking about metric assessment,
we're going to talk about the metric assessment
in the context of the resource constraints
that we're facing with today, right?
That is in the context of the embedded device.
So we'll be training the models so that they're actually fitting
into the rough constraints that we're going to face.
Again, we're going to look at that end-to-end flow.
It's extremely important.
And I'm actually going to talk about that end-to-end flow much,
much more in the next coming couple of videos.
At the end of the day, my goal is really quite simple.
I want to help you understand how to apply TinyML.
What does that mean?
That means you have to understand the intersections between the application
space that you're focusing on, what the machine learning models
need to look like, and how do you build that machine learning model
in the context of the embedded system.
I touched on this in course 1.
And, a quick recap for the newbies, like one
of the key things when we're looking at applications
and then trying to use ML is we want to understand
what are some interesting applications out there
where machine learning is, actually, the right tool.
Please take away one thing from this course
is that machine learning is not like a hammer
that you can just knock everything out with.
It's good sometimes.
It's not so good sometimes.
I'm going to try and teach you that later
on when we actually deal with motion sensors,
as to when machine learning is good and what type of machine learning
is actually good for a particular type of application.
Then, of course, the next thing is now let's say
you do end up deciding that, OK, machine learning is the right problem--
or a right solution for this particular application problem
I'm trying to solve.
Then how do I build that machine learning model
with the unique constraints of the embedded hardware?
So we'll look at training-based methods.
We'll look at different types of networks in this course.
And then we'll also be looking at what are the trade-offs you're
going to make in order to try and squeeze your networks down
into the embedded hardware.
Now we're not going to focus just on the models.
One of the key things that I really want you to take away from this course
is I want to focus you on also the software components that
allow the model to run efficiently and then how that works in the hardware.
We'll give a light preview into that software aspect,
a bit deeper than we did in course 1, but, as we go into course 3,
we'll go much deeper into it.
So you will understand, again, the full end-to-end pipe.
Then, of course, there's the question about what
kind of interesting use cases are really where embedded systems and TinyML
applications flourish out of.
There are certain embedded hardwares where you simply don't know or don't
want to be deploying ML because it's not the right kind of substrate
in order to enable the application.
We'll be looking at three different applications in great depth,
and you will see what the fit is and what the lack of the fit is.
Now that's what course 2 is going to be all about.
It's all about understanding these applications
from the perspective of the sensors.
And one of the things I want you to understand
is that we're not actually going to physically deploy anything
yet onto a TinyML hardware.
The core nugget of this course is to take the machine learning fundamentals
that we all know and then apply them for the application space
and try and understand what are the trade-offs we make
and how do we optimize the models so that they fit into this.
This is one of the critical stages.
You don't always just jump directly to the hardware.
You iterate on this stage quite a bit before you actually
push things out onto the hardware, which will happen in course 3.
That's where you will actually be getting a kit.
Now there is a reading that I have put out
where we actually outline how you can go ahead
and pre-order the kit because this is something
you want to start doing now because it's going to take a little a while for you
to get your kit from the store.
With that said, you will soon be deploying your models
that you trained in this course onto the hardware.
So stay tuned for that.

\section{What Resources are Needed for this Course}

  \includegraphics[width=\textwidth]{TinyML/2 Applications/3-1-5-1.png}


In Course 2, you will only need a laptop with a web browser (we suggest Google Chrome for the best performance using Google Colab), a microphone, a camera, and an internet connection. All of the model training and testing will be done in Colab! We hope this makes the course accessible to everyone!

\section{Preview of TinyML Applications}

VIJAY JANAPA REDDI: Hi.
In this video, I want to give you a sneak preview
of the kind of TinyML applications that we will be dealing with.
I have been telling you so far that we are
going to be looking at acoustic sensors, image sensors, and motion sensors
and then train the networks, and then evaluate them.
What I want to do is help you understand why
we have chosen these different sensors and where
they're actually going to be deployed.
I want to help you understand the real world implications of how
to think about taking sensors, training models, and deploying them
in the field with a real world perspective, OK?
What do I mean by that?
The sensors that we're going to be talking about--
each one of them generally fits in one of these buckets
where there are vast applications of TinyML that people are envisioning.
Let's start with the home scenario.
And then we'll work into the office and industrial applications.
In the home, the most common one of these
is what we call a keyword spotting.
Sometimes it's also referred to as hot word detection.
And by this point of time, you've heard me use this example several times.
You know exactly what I'm talking about when you say, hey Siri or Alexa
or OK Google.
It's all about trying to weight that machine up
and trying to do that on that device.
I'm going to talk quite extensively about keyword spotting.
In this particular course, we're going to go pretty deep into it.
And then we will look at some of their applications
as well that involve sound or speech.
What happens when you do keyword spotting?
You got to understand that the input is coming in.
That's the first thing that happens.
You have to capture the input from the microphone,
which results in sort of a WAV file or an MP3 file or whatever that is.
Then you want to process that.
You want to be able to feed that into a neural network, which
will then be actually be able to make an intelligent decision.
Specifically in the context of keyword spotting,
it's about understanding if you actually said
a word that corresponds to the machine.
Did I say Alexa in order to wake up Amazon's Alexa?
Did I say, hey Siri?
Or did I say he's serious, right?
It's about detecting that keyword, and that
might seem like a very simple thing.
But there are a lot of nuances in it because, for instance, I
might, in passing, just be making a random comment that says, well,
sometimes when I'm saying Google or so forth and just move along.
And the machine has got to be able to pick the keyword up
in that transient streaming audio.
Or it's just sometimes I might just literally be looking in the machine
and say, OK, Google, right?
These are differences that have implications in terms of how
the network can actually handle it or cannot handle it.
We'll learn about this.
Finally, there's the aspect of just generating
the output, which in the case of the Alexa example
was simply dim the lights.
Now, this is the full end-to-end pipeline.
When you look at that end-to-end pipeline,
there's several fundamental technical questions.
It seems like a natural process.
But when you can break it down into stages,
that's when you really get to appreciate the nuances that are involved.
So let me step you through a few of these.
The first question is, how do you capture the data feed into the network?
Sure, you're using a microphone.
But remember, you have to train the network, right?
So you have to understand how you're actually
going to be able to get that data that's from the microphone,
convert it into a signal that the network would actually
be able to understand.
Then there's, of course, the question about how
do you design the neural network itself to take in the speed signal.
What are the input vectors looking like, right?
Then you've got to think about, well, that data
set that you're using to train the network--
what does that look like?
Is it just some arbitrary set of sounds?
Let's say, for instance, let's use the example, OK Google.
So when I'm doing that, think about what's involved.
OK, Google-- that's me saying it.
How do you see it?
How does your mother say it?
How does your sister say it?
My mom says it very differently.
Sometimes, in fact, the machine cannot pick up my mom's way of saying OK
Google because she has a strong Indian accent,
which the machine does not understand.
But why is it that the network does not understand that?
Well, it's because the data set that was used
in order to train that OK Google neural network model was not
fully broad and diverse.
And that's a big problem.
So we need to understand how that data set's actually constructed.
And we're actually going to be doing some of that in this particular course.
And then how do you pre-process the data that's going into that network?
So let's say you have this ideal data set now that you're using this.
And then you have to take this data set.
And then you have to somehow pre-process it
so that it extracts out all the key features.
And you have to do this feature engineering that's involved in there.
And then both, you do the feature engineering
in terms of the training data set, as well as the actual input
signal that's coming in.
So we'll look into what those pre-processing steps are.
And then, of course, you've got to take the output that's
coming out of the network.
It's not that the network is simply saying, oh, this is a cat or a dog.
You actually have to do some amount of intelligent post-processing
on that in order to be able to make that signal that's
coming out of the neural network be useful in the bigger
context of an application.
Now, one other thing that you have to worry about is that how can
you make sure when you have that trained network
that you're in then deploying out into the field.
The small little tiny model that you have is not biased--
in other words, that it would actually work for me
and that would actually work for my mother as well.
So you have to understand to make sure that they're
no biases in the data set before you ship it out into the real world.
And, of course, there's the ultimate question
of how do you deploy this onto the tiny little microcontroller.
Deploying a tiny little model onto a microcontroller
is not simply taking the model that you've trained in Google Colab
and then simply tossing it into the microcontroller.
There's actually several steps involved in terms
of converting the model from the format that it
is in whatever training environment due to deploying environment.
We'll talk about this much more in course three.
But in this particular course, I'm going to try and give you a couple of hints
about what is actually involved.
Similarly, let's look at a different sensor
that is in the context of the office environment.
In the office context, we're going to be using an image sensor.
So it's just a regular camera.
It doesn't have to be in the context of an office.
But I wanted to pick the office deployment because it is not directly
something that we're dealing with.
But think of it, for instance, as a sensor--
a camera sensor-- that is detecting occupancy in a given office building,
right?
If there is no one there in an office building room,
then it should automatically activate the system
in order to turn the lights off so that we can conserve energy.
That's a very popular use of a TinyML application for vision sensor.
There are many other ones of course, but that's
one that's actually really interesting.
So what's actually involved there?
Now, here, you have an image sensor.
It's a camera module, right?
It's going to feed in pixel data into the network
that you're going to evaluate.
Let's take the example that we'll be using.
We'll be using a small little or do cam sort
of a camera that's going to be feeding the network with images.
And what we want to be able to do is pick up some interesting thing.
For example, maybe we'll do person detection
where you're trying to detect the person in that particular image.
Of course, this goes back to something you might have already done
in course one or elsewhere.
You might have trained a convolution neural network in order
to be able to pick up a person and put a bounding box on top of them.
What we're going to specifically focus on with a twist in this class
is how do you accelerate that training time so you can try and accelerate
the time to deployment from collecting data to actually getting on a inference
small?
How do you accelerate that whole pipe, because it's extremely important
to be able to build TinyML applications with a very quick turnaround time.
So we're going to touch on things called transfer learning, for instance, where
we're actually going to learn how to train
the network differently and then deploy it much faster so you don't have to wait
for two days for the network to train.
You'll just wait a couple of minutes, and you are kind of done.
But, of course, the whole thing that comes down to
is even if you come up with ways to train the model faster,
there's the ultimate question about trying to shrink the model down,
which we talked about in course one where we said that, hey, look,
if you have a MobileNet model, which is small.
It's one of the smallest models that we have.
But it's still about 17 megabytes in size and recall
that our embedded microcontrollers tend to have
about 256 kilobytes of main memory, which is barely anything, right?
It's a significant order of magnitude different from what
the smallest models today look like.
So how do we shrink them down so we can actually
put them on a microcontrollers?
That's a key thing.
Some of the nuances that we're going to learn about here
are what are the computing operations that we are actually using?
And how can we, perhaps, change them out so we can actually
do it much more effectively?
We'll learn about some of the numerical representations-- and sort
of using the floating point, how can we use end date?
If we do that quantization process, then what
are the trade-offs that we have to make there?
Then how do we compress the networks down just
from the pure networking size?
Some concepts we'll touch on and go pretty deep.
Like, for example, quantization, we're going to learn quite a bit
in this course.
But in course three, we'll touch on some of the ones
which we'll touch on much more likely in here.
But over the course two and course three,
we'll cover the whole gamut for you.
Now, again, we want to focus on the end-to-end flow in terms
of going from, getting, collecting your data,
all the way to training your model and being able to do inferences with it.
Now, in this context, the pipeline of going
from collecting your data to actually ending up within a model
is largely the same for any kind of TinyML application.
What is interesting, though, is what is new here?
The same sort of questions that we already talked about previously
for the microphone sensor are exactly going to repeat here,
but there are two new ones that come up.
How do you ensure that the neural network that we have trained
is actually resilient for lighting conditions or subtle differences,
for instance, in the office environment?
How do you make it stronger and robust and generalizable?
We'll also talk about how to train that network really fast.
So in a way while we're introducing a new application,
we're also introducing new machine learning concepts in place.
So your learning is going to be much more than just building a TinyML okay
application.
The third application that we're going to be touching on
is something that you often find in industry applications.
It's extremely important to understand this use case
because people are saying it's going to be pervasive.
TinyML has huge amount of industry applications.
And we're going to focus specifically on gyroscope and external meter--
that sort of data.
Specifically, we're going to be looking at things that you would actually
find on machines.
So this little square that's up on the top left
is a very simple three-axis compass, but we'll actually
be using an accelerometer in our assignment.
But the key idea here is that we want to be looking or hearing or listening
to characteristics about how a machine is performing in the field
so that we can do predictive maintenance.
This is something I mentioned previously.
We want to be able to predict if machines
are going to fail in the future.
Why?
If we can predict, then the downtime goes down, right?
If their neural network is able to say, hey, I
think this machine is vibrating a little differently than it normally
would, well, that's important, because that
can raise a flag to the operator who can send someone out
to quickly see if the machine is going to break.
And if it looks like that something is actually wrong, they can replace it.
So without ever having the machine to actually stop working,
which would then shut down the industry plant, which would have consequences.
So predictive maintenance is seen as something that is extremely valuable.
And it really comes down to being able to do all this sensing in real time
because you have so many of these different engineering
parts that are out there.
So we'll look much deeper into this.
And what's specifically interesting about this particular kind
of applications, so we'll look at time series data.
Time series data is where you have a continuous signal that's
constantly coming through.
And when the machine is working normally,
like the blue or the green lines when you're looking at different times,
you'll see that the machine is humming.
Let's just assume that the amplitude there
is just how the machine is vibrating.
And it seems to be going along just fine.
And the neural network looks at the signal.
And then it'll say it'll give a thumbs up and say everything is OK.
That's what you want the neural network to say.
Is everything OK or not?
Now, if everything is OK, nothing needs to be done.
But occasionally, you might see that the machine will actually
have a blip of something-- something about the way it's vibrating,
for instance, is off.
That, somehow, the neural network will pick up and say, hey,
this is an abnormal behavior, right?
Now, how would the network do this?
This type of data is actually not supervised learning.
When you're dealing with image data, for instance, it's supervised learning,
because you're given a set of data.
And you're asked to find similar kinds of features in new data
that's given to you.
So if I give you a cat--
a whole bunch of cats-- and then train you on a whole bunch of cats
and I give you a picture that looks like a cat, then you know that, OK, well,
it's got all the features.
So this is probably a cat.
That's how you train a neural network in supervised learning.
But here, you don't know what you don't know.
You're looking for an abnormality which is not
going to be the common case when you're looking
to train the machines, abnormal or out-of-activity behavior
or rare scenarios.
So how do you get the neural network to actually learn this?
Well, that's the really fun thing about this particular scenario.
It's a fundamentally different concept.
We're going to introduce this notion of what's called an autoencoder.
I'm not going to open up the autoencoder box now.
When we get into the application, we'll dive real deep into it
that you will have a lot of fun, because each of these sensors
is actually touching on something different.
An autoencoder quite different in the way it behaves.
And we can learn about that.
And, again, this particular application goes back
to the same sort of fundamental questions
that we were talking about earlier in terms of how do
you capture the time series data now.
How do you design the autoencoder neural network?
How is this autoencoder different from the other kinds of networks?
So we're trying to cover not only different scenarios of sensors,
but different ways of how these things are actually
executed across the entire pipe.
And that's the really fun part about learning
TinyML in this particular course.
I really hope that you will have a really fun time in terms
of understanding the nuances between the different sensors, the way
you would train the network, the kind of data
that you would try and feed into the network,
and what happens if you have good data versus bad data--
how does the network actually behave?
So that hopefully will give you a very practical take on TinyML applications.
And that's just a preview of what's coming up, so stay tuned.

\section{The Role of Sensors in TinyML Applications}

\subsection{The Role of Sensors in TinyML Applications}

Sensors are important for machine learning applications as they act as a vehicle between algorithmic resources and characteristics of the physical environment. Today there is a continued push towards low-cost and high-performance sensors, especially for applications such as air quality monitoring and predictive maintenance. There are two main types of sensing: in situ and remote sensing. 

In situ sensors come in numerous forms and each have a particular sensing mechanism. For example, the sensing mechanism may involve chemical reactions, light, or measuring voltages from material changes caused by stress. In situ sensors are often characterized as physical sensors and are constrained to measure local attributes at the location of the sensing element. Remote sensing methods work at a range, and are almost always optical techniques. For example, temperature can be measured at a distance using an infrared thermometer, and gas concentrations in the atmosphere can be measured by columnar measurements made by satellites in low Earth orbit. In situ sensors can use optical mechanisms, but are often distinguished because they only measure local attributes.

In TinyML, we are typically interested in sensor data in the form of an image or time series, enabling us to perform spatial or time series analysis using machine learning models. The utility of sensors for tiny machine learning allows real-time feedback and monitoring of the local environment, either at the point of the sensing element or at a distance using remote sensing methods. For example, an array of in situ sensors each with gas sensing elements can be used to monitor pollution levels across a university campus, or surrounding an industrial facility. Similarly, an example of a remote sensing method might be a camera used to monitor local weather conditions based on the appearance of overlying clouds - perhaps useful in locations with limited internet access. 

For a sensor outputting a time series, the sensor measurement can be used as a single feature in a machine learning model. A segment of the time series might also be used for classification or anomaly detection purposes. For an image-based method, the bitmap can be flattened and each pixel passed as a feature in a neural network. 

Sensors can be used individually or can be combined to produce synergistic effects. This provides even more powerful capabilities, especially when combined with online machine learning. Perhaps the most utilized example of this for TinyML currently is the combination of an accelerometer and gyroscope. These devices are often present in smartphones and other embedded devices, and can be used together to differentiate between complex 6-axis movements. These allow your smartphone to know when it has been picked up, or to know when the holder has fallen over. They can also be used in more complex devices such as drones to estimate the local wind speed and direction via how it perturbs from normal operation. Other sensors apart from accelerometers and gyroscopes are also common. Vibration sensors are becoming increasingly common in industrial applications for predictive maintenance. These sensors are used with the aim to improve machinery uptime by detecting anomalous vibrations and repairing them before a catastrophic failure occurs.

Using sensors for embedded machine learning presents several additional challenges. Firstly, sensors have a finite working life, and the designer should bear this in mind during development. For example, we might not want a complex sensing system to become unusable if a single sensor becomes defective, and we might want it to default  to a specific action in the event of a failure. Sensors also tend to exhibit noise and drift, which is often infeasible to remove or model statistically. TinyML systems should be robust enough that they are not adversely impacted by this noise or drift during regular operations. In a similar vein, sensors that are not properly isolated electrically can have the potential to distort other measurements. Finally, the more sensors that are present in a given system, the more complex and power-hungry it becomes. To be an adept user of TinyML, it is not enough to know about machine learning and embedded systems, but also mindful of the constraints imposed by electrical and mechanical engineering challenges.

To overcome the above challenges, the following rules of thumb are recommended:

\begin{description}
  \item[1 - Collect rich data.] In the machine learning world, data is power. A sufficiently large amount of data should be collected to ensure that our TinyML models are able to effectively approximate the distribution of the data. More data never hurts!

  \item[2 - Don't over-engineer your system.] While it may seem attractive to add fifteen sensors to a system, we should follow Occam’s razor when engineering our system that “entities should not be multiplied without necessity”. Being intelligent and yet parsimonious with our sensor selection will minimize the possibility of running into hardware complications.

  \item[3 - Plan to cover all possible sources of variation.] The statistical power of our model will always be limited by the data at hand. Consequently, it is helpful to cover as much ground as possible without conflicting with \#2. This can involve collecting data that covers edge cases as well as using additional sensors that are expected to provide some additional explanatory power to the model.

  \item[4 - Use the highest workable sampling rate.] Although it is highly expensive to transmit and store high-frequency data, we limit these issues when we are using TinyML systems since inferences are made on-device and streamlined to minimize storage. It is much easier to downsample data when there is too much than to upsample it when there is not enough. Thus, when collecting data to use for model training, it is recommended to obtain the highest resolution data possible to provide flexibility in downstream data engineering processes.
\end{description}

Regardless of whether you are using sounds, images, vibrations, electrical signals, or other sensor data, these signals can be combined and used to train machine learning models to help model, classify, or predict events at the edge in real-time using inexpensive microcontrollers, which is a very powerful tool indeed.

\section{The Kit for Course 3}

\subsection{Looking Ahead to the Resources You Will Need for Course 3}

  \includegraphics[width=\textwidth]{TinyML/2 Applications/3-1-10.png}

In Course 3, you will need everything from Course 2, a laptop with a microphone and camera and an internet connection, as well as the course kit as we will be deploying our models onto microcontrollers. You can buy the course kit directly from Arduino at this link:

\url{https://store.arduino.cc/usa/tiny-machine-learning-kit}

The course kit includes a proprietary shield which will make the camera module much more reliable. However, if you would like to purchase individual components instead of the full kit you will need to purchase the following components:

\begin{itemize}
  \item \href{https://docs.arduino.cc/hardware/nano-33-ble-sense?queryID=undefined}{Arduino Nano BLE 33 Sense WITH HEADERS}
  \item \href{https://www.robotshop.com/en/arducam-640x480-03-mp-lens-ov7675-cmos-camera-module-w--adapter-board.html}{ArduCam OV7675}
  \item Jumper wires and optionally a breadboard for connecting the camera to the Arduino
\end{itemize}

Keep in mind that if you don't use the Arduino kit, we have found that the camera connection is often unreliable over jumper wires, and it will be hard for your fellow classmates and others to provide guidance if you experience any issues.

The tinyML course kit from Arudino showing the Arduino Nano 33 BLE sense, the OV7675 camera, the tinyML shield, the usb cable, and the box that contains all of it.

  \includegraphics[width=\textwidth]{TinyML/2 Applications/akx00028_00.default_1.jpg}


\section{ML Lifecycle}

VIJAY JANAPA REDDI: Hi.
In this video, I want to focus on what I call as the ML lifecycle.
What do I mean by the lifecycle?
Typically, when you take any application where
you're trying to embed machinery, there's
a certain process you go through from the beginning in terms
of collecting the data in order to be able to train the model before you
can actually deploy it.
And I want to make sure you understand this full pipe so you
have a big picture before we go into the details in one of these stages
later down in the course.
Having that big picture is extremely important
because it allows you to understand how much effort
needs to be put in various stages.
Because when we are talking about our TinyML applications,
for those three sensors that we touch upon
in terms of the speech, the image, and the vibration stuff,
we are going to need to have deep appreciation for this full pipeline.
And it's extremely important that we have that
in order to be able to build a successful TinyML application.
All right, let's start off by understanding what I mean by lifecycle.
Typically when you take a machine learning course,
you end up focusing very heavily on the neural network.
You try and understand, what's my input data?
What's my output from the neural network?
How do I do the inference?
And mostly, how do I end up training the network?
All of this is a very small portion of the big puzzle of deploying
machine learning into the field.
It's important that you keep that in mind,
because this is just the ML code box.
This code box is a small, little portion of what it takes
to build a full-on TinyML application.
The boxes that I'm showing you here are general.
They apply to any kind of machine learning pipeline.
And I'm showing you this because you need
to have an appreciation for all the components that are involved.
And if you do that, then you will know how much effort you're putting
where and why it's actually worth.
So this whole thing boils down to having an AI infrastructure.
An AI infrastructure includes the components
that support AI, as well as non-AI components.
And to build any application successfully,
you need to have a vertical introspection of the whole pipe.
So the first element of this AI infrastructure
is going to be data engineering.
We're going to be doing data engineering quite a bit, in fact, in this course,
across all three sensors.
And We will also try to understand what are the nuances in data engineering.
Then we'll talk about model engineering.
Model engineering is all about building a model
that is based on the data engineering you've done previously.
Then there's the model deployment, where you actually
take the model you are training, you try and evaluate it.
Is that a model that is a good fit for deploying?
And of course, finally after you've deployed it,
you end up doing product analytics, where
you tend to have dashboards where you are constantly
monitoring how the machine learning system is actually
performing in the field.
And then you circle back up on top.
That's typically how it works.
These are the four fundamental building blocks,
I would say, in any AI infrastructure system.
I want you to keep this in mind, because I'm going
to step through each one of these now.
What is data engineering?
Data engineering, as you try and build any Tiny or non-Tiny application,
is all about defining what are the data requirements.
If you want to be able to say, OK, Google,
what are the data requirements there?
Is it simply about, I need to be able to detect these two words and that's it?
Or are there many more nuances?
For instance, do you need to think about,
what does the data need to look like in a noisy environment versus a very
clean, quiet house when you're actually collecting the data to train the model?
You probably never thought of that, huh?
There are a lot more nuances like that in data engineering.
Then you got to go about the process of actually collecting the data.
Where are you going to get this data from?
Are you and I going to stand in front of the microphone
and keep saying, OK, Google, OK, Google, OK, Google, OK, Google, OK, Google.
Well, that's still a male voice.
So no matter how you do it, you've got to systematically think about where
your data is going to come from.
The scale, the quality, the features, all of that nuances
is actually going to matter quite a bit.
We'll be looking into this much deeper.
Then you got to go about the process of actually labeling that data
so your network knows how to actually learn from that data.
Then you've got to make sure you actually clean the data that you got,
because not all data is really good.
Because sometimes, a mic doesn't work, like [SILENCE]
I wasn't actually saying anything, but that
actually happens in the real world.
So you got to make sure you actually clean the data.
Then of course, you've got to train the data.
But you've got to prep the data and pre-process it in a way
that it's actually ready to be fed into the neural network models.
Then you might want to augment the data.
Remember, we learned about data augmentation
for images, where you can skew things.
What does that do?
It allows your network to be able to take the exact same image
but look at it in slightly twisted ways.
So effectively, rather than collecting those sorts
of twisted data, what you're able to do is actually
build on top of existing data.
So if I take an image and then I rotate it,
well, I didn't have to go and get a rotated image.
I could just reuse that image.
So that's augmentation.
So that's part data engineering, and we'll
be doing that for all three different application use cases.
And then of course, you tend to repeat the cycle,
because machine learning is all about data.
So you're going to go and get more data.
And then you repeat this process over and over.
Then let's look at model engineering.
You've got to train the machine learning model.
You've got to improve the training speed.
We'll look at transfer learning, for instance, where,
instead of training the network from scratch,
you take a trained, pre-trained model, and then
make subtle tweaks in order to accelerate its ability
to train much faster.
Then you've got to define the right set of metrics.
What is a good accuracy metric?
It depends on the application use case.
You need to understand this.
It's not about just building a model that can tell a cat from a dog
and being able to put it on a smaller microcontroller.
No one cares if you got the wrong metrics by which you train the model.
So you got to understand the application use case.
Then we've got to evaluate the model against those metrics
and see how well it does.
And then we've got to repeat the cycle in terms of optimizing the model
training in order to get that model to be
able to meet our metric requirements.
And if you're a model engineer, typically,
you've got to keep up with the state-of-the-art,
which means reading a lot of literature.
That's the job of a machine or an engineer.
It's not just about building stuff.
You actually got to know the field really well.
We'll try and give you a little bit of the sampling.
I want to make sure that you're a practical engineer when you get out
of this program.
Now the third step is model deployment.
Deployment is when you actually physically
put the model onto an embedded device, which we
will be doing physically in Course 3.
But the general idea is you need to understand, how do you
take a TensorFlow model and convert it onto a format that can actually
fit on a tiny device?
And what is actually learned in that conversion process?
We'll open that box later on, of course.
Then you've got to evaluate it for performance capabilities.
Does it meet the real-time requirement that you need?
Is it OK if I say, hey, Siri, and then it takes, like, three seconds.
But you got to understand what it means in terms of understanding
that performance optimization.
What's that latency that an end user is willing to tolerate, for instance?
Then you're putting things onto a small, little microcontroller.
So you've got to be very conscious about energy.
Remember, these things are going to be running on battery devices.
So what that means is that always optimizing for performance
might not necessarily give you the best energy efficiency.
So you need to think about that a little bit.
Then there's the aspect of security and privacy.
When you're deploying the models, what data is this thing collecting?
Is that data kept safely or not?
So we got to engineer the system accordingly.
Then of course, you're performing inference using some sort of APIs.
They got to be fluid.
And we'll look into this much more in Course 3,
but I'm trying to give you a little flavor of these things that
are coming up.
And then you finally might want to do some on-device fine-tuning
of the model, where you actually train the model
and you actually put it onto the device.
But then even after you put it onto the device,
you're trying to train it by fine-tuning it
I'll give you a little sampling of this when we actually
do keyword spotting later on.
And then of course, you tend to do product analytics,
where you build dashboards that allow you to look at data,
you visualize them, and you evaluate that data.
And then you try and figure out, what's the value
added that the ML application deployment is actually giving you?
And then finally, you try and repeat the cycle so you can actually
get advancements and improvements.
So we're going to focus on this full vertical stack in TinyML.
It's extremely important.
It might seem non-technical at this point in time.
I want you to keep the concepts in mind, because later on, as we go in,
we're going to focus on the exact mechanics.
How do you do data engineering for keyword spotting?
How do you do data engineering for anomaly detection?
How do you do data engineering for a vision-based pipeline?
And similarly for model engineering, model deployment, and so forth.
Now, by and large, because we're not building a product that we're actually
going to ship out into the market, the focus
is going to be on the top three bars, data engineering, model engineering,
model deployment.
With that said, let me walk you through the full cycle of the lifecycle of ML.
It starts with data collection being the critical first step.
Then we're going to go into doing some data pre-processing.
After that, we're going to have to see, is the data actually good enough,
or do we need to go back and collect more data?
If the data is good enough, then we're going
to start doing some labeling on it, saying it's a cat or a dog
or whatever it is.
Once we get that, then we're going to prep
the data to actually be ready to go and do the training process pipeline.
Once we do that, we start training the networks.
Then we start evaluating whether the models are actually performing
well given the metrics we have defined.
And then we're going to start thinking about how we actually
deploy this into the real world.
And once we deploy, we get analytics and we continue this sort of circle
over and over and over and over.
And that is the typical lifecycle of a machine learning system.
This applies not just for TinyML.
It applies in general for any kind of machine learning application space.
But that said, it's the exact same thing you're going to need for TinyML.
So I'm trying to make sure that you understand the big picture.
And we will go through most of these phases
for every one of those applications, because there's subtle differences,
so you can actually have a deep appreciation.
And we're going to do that both in Course 2 and Course 3.
You will get some critical nuggets in Course 2.
But that said, I hope you keep this in mind
as you go through into the deeper aspects of each of these boxes.
And whenever you feel lost, when you're deep inside one of the course modules,
I want you to be able to come back up and look at this lifecycle
and say, ah, that's right.
I'm spending time in this particular box.
And they're having me do all these things because this
is one of the critical steps.
And if any one of these steps is a weak link, the whole system will fail.
Because if your data is poor right at the beginning,
and you do all this hard work to deploy it,
when the machine is actually doing its inferences, it's garbage in,
garbage out.
Your model, if it's not trained well on the right kind of data,
it does no good.
For instance, if the keyword spotting is supposed to work for my mom,
but it doesn't, why?
Because, well, the model was not trained well for Indian accents specifically.
And then now, if I try and buy all of these devices
to put around the house for my mom, well, too bad.
None of them really work well.
So even if I have a fancy device there which works for me,
it doesn't work for her because the model didn't have her kind of data.
So it has serious implications.
So that's why I want you to keep this in mind.
And with that said, in the next video, we're
going to focus on one of the critical pipes
for every single application, which I will
call the machine learning workflow.
And I'll see you in that video.

\section{ML Lifecycle}

\subsection{Machine Learning Lifecycle}

  \includegraphics[width=\textwidth]{TinyML/2 Applications/MLLifeCycle}


The development of a machine learning model follows a multi-step design methodology, commonly referred to as the machine learning lifecycle. Diagrams typically illustrate this lifecycle using a varying set of discrete steps. This methodology is also applicable to tiny machine learning applications, albeit with different requirements to conventional machine learning models. Some of the most important steps in the machine learning lifecycle are as follows:

\textbf{Design Requirements.} The first stage of the machine learning lifecycle is  where the application requirements are established. For tiny machine learning applications, hardware constraints (e.g., memory and storage footprint, latency) are often the most pressing requirements. Other requirements can include temporal resolution of data, the number of inferences per second, or minimum model accuracy. This step, whilst being the least well-defined, is the most crucial step, as it sets the scene for the remainder of the lifecycle. The more clearly defined the design criteria are, the easier it is to evaluate the model in later stages. Data requirements, such as the variable of interest and potential feature variables, can also be specified during this stage.

\textbf{Data Collection/Curation.} In this stage, design requirements are known and data begin to be collected that are suspected to aid in predicting the variable of interest. This can often involve active data collection or curation of information from external sources, such as pre-existing datasets (e.g., CIFAR10, ImageNet). This is often one of the most time consuming stages, as it can be difficult and time-consuming to curate a sufficiently large dataset, simultaneously ensuring to minimize dataset bias. No data analysis is performed during this stage. In Course 2, we will learn the challenges associated with data collection for different types of sensors that serve as inputs for tiny machine learning applications.

\textbf{Exploratory Data Analysis/Data Preprocessing.} During this phase, the collected data is analyzed to determine which features are most informative at predicting the variable of interest. Feature engineering is common at this stage, extracting new information from available data. An example of this would be to extract the day of the week or whether it is a weekend from a time variable. Preprocessing involves ensuring the data follows standard modeling assumptions (e.g., are normally distributed) or manipulating data to meet these assumptions (e.g., log transformations on highly skewed variables). Data imputation is also typically done to appropriately fill in or remove missing data, and invalid or duplicate data is handled accordingly. When we develop our keyword spotting model, we will learn how to pre-process the audio signal.

\textbf{Model Development.} All of the previous stages focused on the design aspects, as well as the procurement and analysis of data. These stages often make up the dominant proportion of the time in the machine learning lifecycle. Once a viable dataset is available, the next stage is creating a suitable machine learning model. There are a plethora of machine learning techniques that have various pros and cons. One of the jobs of a machine learning practitioner is selecting an appropriate model for the task at hand. For example, basic linear regression may be suitable in environments where interpretability is key, but are constrained to the set of linear functions. This may present bias in situations where the distribution of data is highly non-linear in feature space, which can result in poor model accuracy. Conversely, neural networks may offer high performance as they are able to model the non-linear data distribution more effectively, but are less interpretable to the user. Keeping in mind the design requirements is important during this stage to know on what grounds to evaluate different models. One thing we will do in Course 2 is to understand the different model development approaches. Training a neural network from scratch is not always necessary.

\textbf{Model Validation.} It is commonplace to generate multiple models and to compare their performance on an unseen data set, known as the test set. The presence of a test set helps to ensure that the model has effectively modeled the distribution of data and has not overfit to points in the training set. The design criteria help to determine what metrics should be used in comparing the performance of various models. For example, if accuracy is the dominant metric, the model with the highest accuracy should be chosen. This stage can also be used to determine whether specific features improved or hindered model performance (e.g., by analyzing the relative importance of features or by successively removing variables and retraining the model, known as ablation study). But there is more to model validation than just looking at the model metrics. We will discover how to think about metrics in a new light, specifically, from an application deployment perspective. 

\textbf{Model Deployment.} The last stage of the lifecycle involves the deployment of the model in the production environment. For tiny machine learning applications, this is one of the most important stages. The model size must be reduced sufficiently to fit on the embedded device through various means such as model compression (i.e., model distillation), type conversion (e.g.., changing model weights from floats to integers), and lossless compression (e.g., Huffman encoding). The model must then be compiled into a format compatible with the end device. We will learn about “quantization” and use it to optimize the model size for performance and latency to make sure we are building sufficiently ‘tiny’ models. We will use TensorFlow’s quantization API, as well as understand what’s going on behind the scenes. 


\bigskip

After all of these stages, continuous monitoring of the model in the production environment is typically necessary to ensure that it is working effectively. If, after all of these steps, the performance of the model is unsatisfactory in the production environment, the model must be updated. This can be done using new design requirements, data, or modeling techniques. This is the essence of the machine learning lifecycle.

\section{ML Workflow}

VIJAY JANAPA REDDI: In this video, I'm going
to teach you about the Machine Learning Workflow, which
is different from the Machine Learning Lifecycle that I touched upon earlier.
This Machine Learning Workflow is going to become our bread and butter
pipeline that we will be focusing on throughout our course.
Earlier on I showed you that an AI infrastructure
supports Data Engineering, Model Engineering, Model Deployment,
and Project Analytics.
We are specifically going to focus on the first three
Data Engineering, Model Engineering, and Model Deployment.
We're going to look at this notion of Data Engineering, Modeling Engineering,
and Model Deployment from the perspective
of building a keyword spotting model to doing person detection using
image sensors or camera, as well as looking at times series data
based on an Accelerometer.
These three sensors are going to become our bedrock
for building the applications, but the way we build these applications
is going to be in a very principled way.
The reason I want to emphasize the principled way
is because I want you to really internalize
the steps of Data Engineering, Model Engineering, and Model Deployment
so you know how to systematically build applications, not just throw
a couple of things together and be able to say,
ah, I've got a TinyML application.
That's not what a good engineer does.
A good engineer is systematic in the way they do the things,
and that's what I'm going to try and teach you.
Now, what's important is the diagram that's at the bottom of the slide.
Over the entire course, we're going to be using that bottom row.
And that I want you to pay a lot of attention to.
And in this slide, I want you to pay attention
to how that bottom row evolves with respect to the layers on the top.
The layers on top-- the Data Engineering, Model Engineering,
Model Deployment, Product Analytics-- those are the general principles.
Out of those, we're going to pick a key set of things
in order to emphasize in this course.
First and foremost is data engineering.
We're going to learn how to collect data and how to pre-process the data.
For example, we're going to learn what it
means to collect keyword spotting data from you, from me, and everybody.
And what happens if we don't collect the data properly?
Will it still work?
Will it not work?
We try and understand some of those nuances.
Assuming that you have collected the data properly,
then we go into the second stage of Model Engineering.
We try and understand, how do you actually train a model so that you can
actually get it to detect the words?
Or how do you actually train the network?
Or how you design and train the network if you're trying to detect something
like a time series signal?
Where there is a bunch of signals that are constantly
coming, and trying to find that abnormality that happens rarely.
It's not only about just training a given network,
it's also engineering the network and designing it in such a way
that you can actually get the network to learn the relative characteristics that
are of importance for you.
Then the final stage is actually the Model Deployment.
How do you go from a trained model to preparing the model so that it's
actually ready to be optimized for deployment,
both in terms of performance, energy, latency, and various other metrics that
actually matter?
Finally, we'll get to the last stage of doing inferences, which
is about the Product Analytics part.
That said, we're going to focus on the top three ones
and really less on Product Analytics.
Ultimately though, one of the key things that I want to get across for you
is really understanding what it means to package a TinyML application together.
You see, most of the time when you're taking a Machine Learning
course, what you end up focusing on is really just the neural network part.
You don't tend to think about the glue that holds it together.
You don't really think too much about the data.
You don't really tend to think about the optimizations in terms
of deploying things.
You tend to just kind of focus on, did I get my network
to classify a dog versus a cat correctly?
Right?
That's not really what building a true real world application is about.
There is much more to it, and I want to train you to be a TinyML engineer who
understands that big picture.
Yet another point that is important is that it
is important to understand not just the neural network models,
you also want to understand the software stack that actually supports running
those neural networks efficiently.
If I give you a really small neural network model,
but then I give you a big, bloated system that
is actually going to take the network and then run it, what good does it do?
If you're talking about it-- especially in the context of a TinyML device--
maybe your model is extremely efficient.
But the code that's trying to run the network is really inefficient.
Well, what was the point of then building that really efficient model?
So you need to understand both TensorFlow as a training and inference
framework, and how do you make that TensorFlow ecosystem itself really thin
and mean so that you can deploy it efficiently
along with an efficient model?
This is something that is often overlooked by many engineers.
But it's extremely crucial that you understand this.
And that's one of the unique value adds that I'm
going to try and bring to you in this particular course.
We're going to learn about the nuance differences between TensorFlow
and TensorFlow Lite.
Especially as we go into Course 3, I'm going
to try and teach you just internals of TensorFlow Lite Micro itself.
And why and how it has been engineered.
That's a lot to cover.
So even though we're going to cover all of this stuff,
in this particular course, I want to emphasize that we're
covering the first three parts.
The blue parts, the yellow parts, and the green parts.
Stuff about data, stuff about training the model, and stuff about optimizing
the model for deployment.
In the next course, we'll focus on how do we take that small neural network
model that we've optimized and how do we deploy it?
Deploying is not simply taking some graph
in then being able to toss it onto an embedded device.
There is actually a lot of embedded systems aspects
that are involved in there.
And in Course 3, well go deeper into that aspect.
For now though, let's try and understand what
it means to build a TinyML application, starting with collecting data
all the way to actually optimizing the model.
Specifically in the context of those three application scenarios
that we're going to look at.

\section{ML Workflow Details}

\subsection{ML Workflow / Pipeline}

  \includegraphics[width=\textwidth]{TinyML/2 Applications/3-2-5-1.png}


\subsection{TinyML Application Pipeline}

The figure above shows the workflow methodology we will be following for building the tinyML applications in this course. These are broken down into three main stages:

\begin{itemize}
  \item Step 1: Collect \& Preprocess Data
  \item Step 2: Design and Train a Model
  \item Step 3: Evaluate, Optimize, Convert and Deploy Model
\end{itemize}

\subsection{What's Underneath the Hood?}
Often, when learning about machine learning, we forget that models are just a piece of the bigger picture. There is more to machine learning than just the models. These are the underlying ML runtimes and hardware that enable these models to run efficiently. The ML runtimes we will be dealing with are TensorFlow, TensorFlow Lite, and TensorFlow Lite Micro. 

\begin{figure}
  \includegraphics[width=\textwidth]{TinyML/2 Applications/3-2-5-2.png}

  \caption{Machine learning models, runtimes, and hardware all work together and are related to each other.}
\end{figure}  

One of the goals of our course is to dive into understanding the fundamental differences between machine learning runtimes so that we can deploy the models we prepare efficiently. To this end, it is important to understand the various frameworks we use across all three applications in this course: Keyword Spotting, Visual Wake Works, and Anomaly Detection. 

So coming up next is material that highlights the fundamental differences between TensorFlow and TensorFlow Lite (for microcontrollers). You will learn not only the high-order bits but also the mechanics that implement the API calls you learn to make with the code. 

% 1.3: Machine Learning on Mobile and Edge IoT 

\section{TensorFlow: Where We Left Off}

Before we dive into our TinyML applications' underlying mechanics, you must understand the software ecosystem that we will be using. The software that wraps around the models we train is just as important as the models themselves. Imagine working hard to train a tiny machine learning model that occupies only a few kilobytes of memory, only to realize that the TensorFlow framework used to run the model is itself several megabytes large. 

To avoid such oversight, we first introduce the fundamental concepts around TensorFlow's usage versus TensorFlow Lite (and soon TensorFlow Lite Micro). We will be using TensorFlow for training the models. We will be using TensorFlow Lite for evaluation and deployment because it supports the optimizations we need. It has a minimal memory footprint, which is especially essential for deployment in mobile and embedded systems. To this end, in the following two modules, Laurence and I are going to first introduce the concepts, programming interfaces, and the mechanics behind them before diving deep into the applications. Every application reuses these fundamental building blocks. 

\subsection{Recap of the Machine Learning Paradigm}

In the previous course you had an introduction to machine learning, exploring how it is a new programming paradigm that changes the programming paradigm. Instead of creating rules explicitly with a programming language, the rules are learned by a neural network.

\begin{figure}
    \includegraphics[width=\textwidth]{TinyML/1 Fundamentals/MachineLearning}
    \caption{Answers and data are fed into a machine learning algorithm and rules are produced.}
\end{figure}  


Using TensorFlow you could create a neural network, compile it and then train it, with the training process, at a high level looking like this:


\begin{figure}
    \includegraphics[width=\textwidth]{TinyML/1 Fundamentals/MachineLearningLoop}
    \caption{sequence of making a guess, measuring accuracy, and optimizing guess. Then repeat.}
\end{figure}  



The neural network would randomly initialize, effectively making a guess to the rules that match the data to the answers, and then over time it would loop through measuring the accuracy and continually optimizing.

An example of this type of code is seen here:

\begin{lstlisting}
tensorflow
data = tf.keras.datasets.mnist

(training_images, training_labels), (val_images, val_labels) = data.load_data()
training_images  = training_images/255.0
val_images = val_images /255.0

model = tf.keras.models.Sequential([tf.keras.layers.Flatten(input_shape=(2828)),
               tf.keras.layers.Dense(20, activation=tf.nn.relu),
                tf.keras.layers.Dense(10, activation=tf.nn.softmax)])

model.compile (optimizer='adam',
               loss='sparse_categorical_crossentropy',
               metrics=['accuracy'])
model.fit(training_images, training_labels, epochs=20 validation_data=(val_images, val_labels))
\end{lstlisting}


With the bottom line being that instead of creating a program or an app, you train a model, and you use this model in future to get an inference based on the rules that it learned.

\subsection{Model Training vs. Deployment}

With TensorFlow there are a number of different ways you can deploy these models. As you’ve been studying this specialization, you’ve been running inference within the colabs that you were training with. However, you can’t share your model with other users that way.

The overall TensorFlow ecosystem can be represented using a diagram like this -- with the left side of the diagram showing the architecture and APIs that can be used for training models. 

\begin{figure}
    \includegraphics[width=\textwidth]{TinyML/2 Applications/3-3-1-1.png}
    \caption{An architecture diagram of the TensorFlow ecosystem. Data (tf.data and other datasets) and model design (Keras and other estimators) lead into training which can be done on different types of hardware (CPUs, GPUs, TPUs, etc.). Analysis can be done on training in Tensorboard. Serialization of the training results creased saved model, which can be saved to the model repository (TesorFlow Hub). Various services like browsers, microcontrollers, and the cloud can serve up the experience using various TensorFlow technologies like TFLite Micro.}
\end{figure}  


In the center is the architecture for saving a model, called TensorFlow SavedModel. You can learn more about it at: 

\url{https://www.tensorflow.org/guide/saved_model}

On the right are the ways that models can be deployed.

\begin{itemize}
  \item The standard TensorFlow models that you’ve been creating this far, trained without any post-training modification can be deployed to Cloud or on-Premises infrastructures via a technology called TensorFLow Serving, which can give you a REST interface to the model so inference can be executed on data that’s passed to it, and it returns the results of the inference over HTTP.  You can learn more about it at \url{https://www.tensorflow.org/tfx/guide/serving}

  \item TensorFlow Lite is a runtime that is optimized for smaller systems such as Android, iOS and embedded systems that run a variant of Linux, such as a Raspberry Pi. You’ll be exploring that over the next few videos. TensorFlow Lite also includes  a suite of tools that help you convert and optimize your model for this runtime. \url{https://www.tensorflow.org/lite}
  
  \item TensorFlow Lite Micro, which you’ll explore later in this course, is built on top of TensorFlow Lite and can be used to shrink your model even further to work on microcontrollers and is a core enabling technology for TinyML. \url{https://www.tensorflow.org/lite/microcontrollers}
  
  \item TensorFlow.js provides a javascript-based library that can be used both for training models and running inference on them in JavaScript-based environments such as the Web Browsers or Node.js. \url{https://www.tensorflow.org/js.}
\end{itemize}

Let's now explore TensorFlow Lite, so you can see how TensorFlow models can be used on smaller devices on our way to eventually deploying TinyML to microcontrollers in Course 3. Over the rest of this lesson we’ll be exploring models that will generally run on devices such as Android and iOS phones or tablets. The concepts you learn here will be used as you dig deeper into creating even smaller models to run on tiny devices like microcontrollers.


\section{Introduction to TensorFlow Lite}


Now that you've seen an overview of what's
involved in the process of converting a model
to run at the edge using the TensorFlow Lite converter,
next up you'll explore the code to train a model in TensorFlow
and then use the conversion tools.
Colab supports a TensorFlow Lite interpreter, so using that,
you'll be able to test your converted model.
You'll see that later in a hands-on lab, but first,
let's explore how to convert a model using code.
If we draw a high-level architecture of the major components within TensorFlow,
we'd get something like this.
The components on the left are primarily for training models.
You've used the Carus APIs already in creating models,
and you've had a look at how data sets are
stored using Tf.data and TF data sets.
When it comes to sharing your models with the world
so that people can use them, there are a number of deployment options,
including TensorFlow Serving, which gives you a REST interface
to pass data to and get inferences back.
There's TensorFlow JS, which has a JavaScript library that
lets you train models or run inference on them in JavaScript environments.
And then there's TensorFlow Lite and TensorFlow LiteMicro,
which provide a runtime for inference on small systems like phones,
embedded systems, and microcontrollers.
And this is the focus of what we'll be looking at next,
starting with the understanding that for models to run on these devices,
they'll need to be modified somewhat due to the fact
that the devices running them are relatively constrained.
They may have lower powered processors, rely on battery life,
and stuff like that.
As a result, models will need to be shrunk to run on them.
Fortunately, you can do this using tools that
act on the trained model instead of you needing
to learn a whole new way of training models.
We'll explore that next.
Recall the flow from earlier in the course, where you start with a model
that you train in TensorFlow.
You then convert that model to TensorFlow Lite format
using tools that are part of the API.
Your converted model can then be optimized to run on smaller systems.
It's then deployed at the edge on systems such as iOS, Android,
Linux-based embedded systems like the Raspberry Pi,
and of course microcontrollers.
Once it's deployed to the edge, you can then make inferences with the model.
But the first step once you have your model
is, of course, to convert it, as we saw earlier.
So let's recap the conversion process, where
the primary task that the converter will do
is to compress the model into a format that's
suitable for running at the edge.
It will first do this by pruning the model.
And as the name suggests, this means simplifying the model
by removing excess content that may not have
a major impact on the final result. So for example, if removing these synapses
does not result in a drastic loss of quality,
the conversion tool will remove them.
And if removing these neurons also doesn't end up
with a dramatic loss in model accuracy, it will also remove them.
The results?
The tool removes many of the memory-consuming elements
of the network without a large loss in accuracy.
Similarly, quantization can take place in the conversion.
And this is the process of recasting the data from high-memory data
storage, such as a float32, to a lower memory consumer, such as int8,
without a significant loss in accuracy.
So for example, the distribution on the right
here can end up having a very similar distribution when it's recast
into int8, but only using a quarter of the memory.
OK, so now that we've seen an overview of what's
involved in the process of converting a model
to run at the edge using the TensorFlow Lite converter, next up,
you'll explore the code to train a model in TensorFlow
and use the conversion tools.
Colab supports a TensorFlow Lite interpreter
so that using that, you'll be able to test the converted model.
We'll see that later in the hands-on lab.
But first, let's explore how to convert a model using code.

\section{Using the TFLite Converter}

SPEAKER 1: OK so here's the Colab where we'll first
start exploring how to convert a model into TensorFlow Lite.
We'll start just by making sure that Colab is running TensorFlow version 2.
Here you'll see it's initializing the VM.
It's connected to the VM, and we can see the RAM and disk that we have.
And TensorFlow 2 is available.
So now I'm going to import what I need, and that's just TensorFlow kera
sequential, keras dense, and numpi.
And now I'm going to create the model that you're probably familiar with
by now where I have just a set of x's and y's, and it's going to fit
the x's to the y's.
We'll take a moment to do that.
It's 500 epochs, but they're pretty fast.
So here we can see 500 epochs, tried it for 10, we got 18.98.
And it learned-- for y equals wx plus b it learned--
1.998 and minus 0.99, which is very close to x minus 1, but remember we
only had six data points.
OK so now we're going to start the process of converting this
into a lite model.
So the first thing I'm going to do is I'm going to save the model.
I'm specifying an export directory, which is just \PYTHON{saved\_model/1} and then
I'm using \PYTHON{tf.saved\_model.save} and putting the model into that export
directory.
If I click the directory folder here, in Colab,
I can actually look at \PYTHON{saved\_model}.
We can see 1, which I had created.
And then in here are the details of the saved models.
There is the \FILE{.pb} file as well as any assets or variables that it might have
needed.
The variables will just have some data in them,
and that's actually the saved weights of the 1.99 and minus 0.99.
That kind of stuff.
So now that I have the saved model, I can then
use the \PYTHON{tf.lite.TFLiteConverter}, and say I'm doing it from a saved model.
And that saved model is coming from the export directory that I had created
earlier on, which is just \PYTHON{saved\_model/1}.
So I can invoke the converter.
And now I can take a look at the TF Lite model file.
I'm just saying \PYTHON{pathlib.Path(model.tflite)},
and I'm going to write it out to that.
And if we go back and we look here, now we can see model.tflite has been saved.
And this 764 is actually the size of the file.
It's 764 bytes.
Now I don't need to have an embedded system, or a mobile device,
or a tiny email system if I want to test it.
Because within Python we actually have the standalone interpreter for TFLite.
So I'm just going to say interpreter equals \PYTHON{tf.lite.Interpreter},
and I'm going to say use the tflite model that we just created.
And I actually double-clicked there.
Sorry about that, which is why we got that warning.
And so now it's going to allocate the tensors to that,
specifying the input and output details, and printing them.
So we can see the name of the input details as \PYTHON{dense\_input},
and it's going to take in an array of one by one, which is an int32,
and then the output is this.
So now if I say, OK, on my input, I'm going
to specify it's going to be an array of type flow at 32.
I'm going to pass in 10.0, print that out, set the tensor of the inputs
to this value to predicts, invoke the interpreter,
read the results by reading \PYTHON{get\_tensor} on the \PYTHON{output\_details},
and then printing those results.
Might seem a little bit convoluted, but remember
the TensorFlow Lite model tries to be as encapsulated as possible.
It has an input tensor.
It has an output tensor.
So you need to get data, and format it into what
it expects for you to pass it to it.
And then when it has an output tensor, you get that and you pass it back.
So here we can see we passed in the value 10,
and we got back the value 18.98676.
And the eagle eyes amongst you may have noticed that that's
exactly the same as we got here.
So we have the same model, but it's now operating in TensorFlow Lite.
So give it a try, and see how it works for you.

\section{Using the TFLite Converter in Colab}

Now that you've seen how we can use the converter to create a TFLite model, you'll get a chance to explore the code for yourself in this next Colab. Try to take your time and walk through the code to make sure you understand the Converter!

\url{https://colab.research.google.com/github/tinyMLx/colabs/blob/master/3-3-4-TFLiteConverter.ipynb}


\section{How to use TFLite Models}

Previously you saw how to train a model and how to use TensorFlow’s Saved Model APIs to save the model to a common format that can be used in a number of different places.

Recall this architecture diagram:

\begin{figure}
    \includegraphics[width=\textwidth]{TinyML/2 Applications/3-3-1-1.png}
    
    \caption{An architecture diagram of the TensorFlow ecosystem. Data (tf.data and other datasets) and model design (Keras and other estimators) lead into training which can be done on different types of hardware (CPUs, GPUs, TPUs, etc.). Analysis can be done on training in Tensorboard. Serialization of the training results creased saved model, which can be saved to the model repository (TesorFlow Hub). Various services like browsers, microcontrollers, and the cloud can serve up the experience using various TensorFlow technologies like TFLite Micro.}
\end{figure}  



So, after training your model, you could use code like this to save it out:

\PYTHON{export\_dir = 'saved\_model/1'}

\PYTHON{tf.saved\_model.save(model, export\_dir)}


This will create a directory with a number of files and metadata describing your model. To learn more about the SavedModel format, take a little time now to read \url{https://www.tensorflow.org/guide/saved_model}, and also check out the colab describing how SavedModel works at \url{https://colab.research.google.com/github/tensorflow/docs/blob/master/site/en/guide/saved_model.ipynb}, and in particular explore ‘The SavedModel format on disk' section in that colab.

Once you had your saved model, you could then use the TensorFlow Lite converter to convert it to TF Lite format:

\PYTHON{converter = tf.lite.TFLiteConverter.from\_saved\_model(export\_dir)}

\PYTHON{tflite\_model = converter.convert()}

This, in turn could be written to disk as a single file that fully encapsulates the model and its saved weights:


\PYTHON{import pathlib}

\PYTHON{tflite\_model\_file = pathlib.Path('model.tflite')}

\PYTHON{tflite\_model\_file.write\_bytes(tflite\_model)}

To use a pre-saved tflite file, you then instantiate a \PYTHON{tf.lite.Interpreter}, and use the \PYTHON{model\_content} property to specify an existing model:

\begin{lstlisting}
interpreter = tf.lite.Interpreter(model_content=tflite_model)
\end{lstlisting}

Or, if you don’t have the existing model already, and just have a file, you can use the \PYTHON{model\_path} property to have the interpreter load the file from disk:

\begin{lstlisting}
interpreter = tf.lite.Interpreter(model_path=tflite_model_file)
\end{lstlisting}

Once you’ve loaded the model you can then start performing inference with it. Do note that to run inference you need to get details of the input and output tensors to the model. You’ll then set the value of the input tensor, invoke the model, and then get the value of the output tensor. Your code will typically look like this:

\begin{lstlisting}
# Get input and output tensors.
input_details = interpreter.get_input_details()
output_details = interpreter.get_output_details()
to_predict = # Input data in the same shape as what the model expects
interpreter.set_tensor(input_details[0][], to_predict)

tflite_results = interpreter.get_tensor(output_details[0]['index'])
\end{lstlisting}

\ldots and a large part of the skills in running models on embedded systems is in being able to format your data to the needs of the model. For example, you might be grabbing frames from a camera that has a particular resolution and encoding, but you need to decode them to 224x224 3-channel images to use with a common model called mobilenet. A large part of any engineering for ML systems is performing this conversion. 

To learn more about running inference with models using TensorFLow Lite, check out the documentation at: 

\url{https://www.tensorflow.org/lite/guide/inference#load_and_run_a_model_in_pytho}

\section{Running Models with TFLite}

SPEAKER: Earlier, we saw an example of a Colab that was using TenserFlow Lite,
and it was a very simple example, where we
were doing a linear regression, where y equals 2x minus 1.
Now let's take a look at a little bit more sophisticated of an example.
This one is going to be image classification using
the famous cats versus dogs model.
I'll have the same code here just to ensure that TenserFlow 2 is running.
And now, I'm going to run this piece of code.
And what this piece of code is just doing,
is it's downloading the cats versus dogs data.
It's setting up my training batches and my validation batches
and all that kind of stuff.
So it's preparing the data for us.
So that we can train the model.
In order to keep things fast, I'm going to do transfer learning instead
of training a model from scratch.
And we'll see that in just a moment, where I download the model
from something called TensorFlow Hub.
And the TensorFlow Hub model will then be
used with transfer learning to create a new model off
of the cats versus dogs data.
So the data is now ready, and we can see we have like 23,000 samples.
And this means they're in two classes.
So now this piece of code is where I'm going to create the model.
And I'm going to create a very simple model, because I'm
doing transfer learning, where I'm downloading,
it's something called MobileNetV2.
And MobileNetV2, I'll pull the feature extractor from this.
And then that feature extractor will be used as the main layers in my model.
And then I'm just going to add my own classifier
at the bottom of that with two classes, because I have cats versus dogs,
as opposed to the number of classes that MobileNet supports.
So I'll start training this model.
And instead of it taking a few hours over many epochs,
I'm just going to train over five epochs,
because of the transfer learning.
So this'll just take a couple of minutes.
OK, now that it's finished training, we can
see that we have accuracy of 99.36% on our training sets,
but 98.67% on the validation setts.
So we look like we're in pretty good shape.
So what am I going to do?
Well, exactly the same as the process that we did for the y
equals 2x minus 1.
I'm going to specify a save model.
And in this case, I'm just calling it exp saved model.
And then, I'm calling tf.saved model.save
and saving my model to that cats versus dogs save model.
And if I look at my folder, we can see exp saved model is there.
And the same thing, I've a saved model.pb,
and I have assets and variables.
There's going to be a lot more in this, because it's a much bigger model.
But we can see the same kind of structure.
And this is how saved model actually works.
So I close this.
And now go back.
Now, I can use the TF Lite converter to convert this model
to one that will work with TF Lite.
So I'm going to do the same thing.
I'm using TF Lite convertor from saved model,
and giving it the cats versus dogs save model.
And then going to do the conversion.
And we can see here, now that my file is much bigger.
It's 8.9 megabytes.
It's eight million bytes, as we can see here.
But still, that's pretty good for a mobile model which is complete.
And it's like self complete for doing computer vision.
So now here, if I want to explore the model and take a look at it,
I'm going to run it.
And I can look at the validation sets.
And using the validation set and like I grabbed like 100 sample images
from that, and it tested them.
And we could see that out of 100 predictions, they got 99 correct.
Now remember, this is the mobile model.
So this is us after we've converted the model into TensorFlow Lite.
It's still getting that level of accuracy.
And here in the Colab, if you want to give it a shot,
there's some things that you can just plot out what it looks like,
and all of those kind of things.
And if you want to visualize some of the outputs, I've put some code in here.
So you can visualize them along with the predictions.
So let's take a look at some of these, where it got them right or wrong.
That one is correct.
That one's correct.
And all of these are correct, because it's a pretty good model.
It's getting it right 99% of the time.
This one is kind of interesting, because it's two cats, but it's still detected.
I don't know if there are any images with both a cat and a dog.
I'd be curious to see what it would do there.
But this one, for example, is quite interesting
because there's a human arm and a human leg, from the look of it in the image,
but it still wasn't fooled, and it detected it as a dog.
So give this model a try for yourself.
Remember, this is the one that has been converted into mobile,
so quantization, optimization, all those kind of things, have been done to it.
But it's been done without any kind of loss of accuracy.
And we're seeing here, I'm not seeing any of them being wrong.
I always find it interesting to find one that's wrong,
so we can see why it might be wrong.
And again, this one is really cool because it's a human carrying the dog,
but it's still detected the dog.
This one was only 91% certain it was a cat, but it still was a cat.
And this one's kind of funny, but it's still a dog, as we can see.
So that's it.
Try out the Colab.
Give it a shot for yourself, and see what you can come up with.

\section{Running Models with TFLite in Colab}

Now that you’ve seen how to test TensorFlow Lite models in Colab you’ll get a chance to explore this yourself! Make sure to start to pay attention to the file structure of Colab as you’ll need to understand how to upload and download data and models from Colab later in the course.

\url{https://colab.research.google.com/github/tinyMLx/colabs/blob/master/3-3-7-RunningTFLiteModels.ipynb}

\section{TFLite Optimizations and Quantization}

LAURENCE MORONEY: Recall the flow from earlier
in the course where you start with a model that you train in TensorFlow.
You then convert that model to TensorFlow Lite format
using tools that are part of the API.
Your converted model can then be optimized to run on smaller systems,
and it's then deployed to the Edge.
Once it's deployed to the Edge, you can make inferences with that model.
For you to get the most out of your model,
it will need to be optimized for the Edge.
And optimization can take place in a number of ways.
You can set up optimizations on the converter
when you're converting the model, or you can do post-conversion optimization.
And in this video, we'll look at a couple of these.
First of all, let's consider optimization
at the point of model conversion.
There are a couple of optimizations that you can explore at the converter level,
but let's start with the easiest.
And that's using it's built-in, automatic optimization
through setting a property.
So for example, consider this code where we have a saved model for cats and dogs
at the directory indicated.
And we'll convert it to a TensorFlow Lite model by invoking the converter,
getting a tflite file, and then saving that out.
So we could introduce optimizations like this,
by setting the converter optimizations property here.
And in this case, I'll set it to tf.lite.Optimize.DEFAULT.
There are a few options for what you could do here, so let's explore them.
Optimize for size, as its name suggests, performs optimizations
to make the model as small as possible.
Optimize for latency will perform optimizations
that reduce inference time as much as possible.
They're both worth experimenting with.
And because Optimize.DEFAULT will try to optimize for both of these--
and it usually works quite well at doing so--
you'll usually get by just by keeping it at a default like this.
But it is good to know about the others.
As an example, when I was using the cats versus dogs model
we're discussing here, my model size was about 9 megabytes.
After using the default optimization, it ended up a 2.3 megabytes.
Experiments have shown that models can be made up to about four times smaller
with a two to three time speed-up.
Depending on your model type, there can be a loss in accuracy,
so be sure to test thoroughly to see if the losses are acceptable.
For example, my cats versus dogs model dropped from about 99%
accurate to about 96\% accurate when using this optimization.
The second method we can explore is quantization
through representative data.
This changes the weights in the model from 32-bit floating
points to 8-bit integer.
And this can have a huge impact on model size and latency
with a relatively small impact on accuracy.
And it's particularly powerful for hardware
that's optimized for 8-bit integers.
So starting with the same code that we just saw,
we can set our converter optimizations.
But we can then add code to create a representative data set of the values
that we want to optimize around.
Let's look at this piece by piece.
We'll start by taking a sample of our test data
from the data set that we originally used to train and test the model.
The goal is to have a representative sample of the type of data
that the model would expect to see, so that the model can be calibrated
for the ranges of data expected.
This is then wrapped in a function that returns its input values.
This type of function is typically called a Data Generator,
so here, the function name is Representative Data Generator.
Now that you have a Data Generator that can
provide a good example of what type of data
the model is expected to run inference on,
we can tell the converter to use that as the representative
data set for the model.
And then we can tell the model that the \PYTHON{supported\_ops}
should be \PYTHON{INT8} by setting the \PYTHON{supported\_ops} property like this.
If you want to compress the \PYTHON{INT8} --
which is recommended for hardware that supports it --
you can use this property value, \PYTHON{tf.lite.OpsSet.TFLITE\_BUILTINS\_INT8}.
Note that this interface is constantly evolving
and the API surface might change, but the underlying
concept should remain the same.
Now that we've explored coding optimization and quantization
at the converter level, it's your turn to try them out in the Colab.
Play with the different options, and explore
the impact of the optimizations.

\section{TFLite Optimizations and Quantization in Colab}

In this next Colab you will get to explore quantization in TFLite through a hands-on example of recognizing pictures of cats vs. dogs. We will explore both how the model size changes and how the accuracy changes. You’ll also get to look at individual predictions of the images to see which pictures it gets correct and incorrect!

Note: Changes have been made to the TFLite Interpreter since Laurence filmed the previous video that further optimize it for mobile use at the expense of speed in Colab. As such, you'll find that while Laurence was able to achieve speeds of 16 iterations per second for Model 2 in the previous video, you may only see speeds of 1-2 iterations per second.

\url{https://colab.research.google.com/github/tinyMLx/colabs/blob/master/3-3-10-TFLiteOptimizations.ipynb}


\section{Quantization Aware Training}

LAURENCE MORONEY: We've been looking at the flow for training, converting,
optimizing, deploying, and then using a model
at the Edge in the last few videos.
The one thing we haven't really explored is
what we can do to optimize the model for Edge inference right at the beginning.
And that's when you're training the model.
So we'll do that here.
Looking into the model optimization tool kit, and how it can help
you quantize models at training time before you convert them.
There are a couple of optimizations you can explore at the converter level,
and let's start with the easiest.
And that's using it's built-in automatic optimization
through setting a property.
We'll start with this code, and it should be quite familiar by now.
And it's a convolutional neural network that
classifies the famous mnist 28 by 28 handwriting digits
data set into 10 categories.
And those are the digits 0 through 9.
As an experiment, I trained this for just one epoch,
and I got these results.
Accuracy of 92.44% on the training set and 96.95% on the validation set.
Let's use this as our baseline.
Now, if I want to quantize the model, I can use the model optimization tool
kit.
And let's look at this step by step.
We'll start by importing the tool kit as tfmot.
If you don't have it already, you'll need to pip install it.
Model training quantization can be done using the quantize model
API, which is available in tfmot.quantization.keras namespace.
You can use this by simply passing at your model,
and it will return a new model to you, and we'll call this q aware model,
as it's a quantization aware model.
And you treat this model like any other.
For example, here we compile it with an atom optimizer,
sparse categorical cross entropy as its loss function,
and we'll track accuracy metrics.
If we return to our baseline-- the non-optimized model showing the 92.44%
accuracy and 96.95\% validation accuracy--
and then we'll do the same with our optimized model, we get the following.
And these figures are very similar with equivalent accuracy.
And this is just after one epoch.
Later you can try this for yourself, training for more epochs
to see the impact.
Know at this point your model is already quantized with INT8 weights and UINT8
activations.
So you don't need to do things like having a representative data set
to convert your model.
Here you had an introduction to how the model optimization tool kit works
for quantization while you're training, so before you convert to .tflite.
Next up, you'll explore the code, and you can try it out for yourself,
and you can examine the impact on the model size and model accuracy.

\section{Quantization Aware Training Colab}

There are two primary forms of quantization -- post training quantization that you had seen previously -- where, as part of the conversion process, your model’s internal weights and ops get converted to int8 and uint8. It’s much easier to use this, and it’s a good way to get started.

As you want to further optimize your model, you may want to explore quantization aware training. You got a taste of that in the previous video seeing how you could pass your simple model to an api and get a quantized model back. It turns out that there is a lot more fine-grained control available to you. To get some exposure to some of those controls, please work on the below colab provided by the TensorFlow team on Quantization-aware Training (QAT). After you are done, if you would like to go deeper, including a comprehensive guide on all the APIs available in the toolkit, check out and read through the \href{https://www.tensorflow.org/model_optimization/guide/quantization/training}{model optimization site on TensorFlow.org}.

In particular, note the results posted by the Google teams when comparing accuracy of models before and after quantizing like this -- the effects on accuracy are negligible!

\begin{tabular}{lll}
Model & Non-quantized Top-1 Accurac &>8-bit Quantized Accuracy\\
MobilenetV1 224 & 71.03\% & 71.06\% \\
Resnet v1 50  & 76.3\% & 76.1\% \\
MobilenetV2 224 & 70.77\% & 70.01\% \\
\end{tabular}

Colab link for QAT for you to work through:

\url{https://colab.research.google.com/github/tinyMLx/colabs/blob/master/3-3-12-QAT.ipynb}

\section{Assignment: Quantization in TFLite}

Now that we've explored how to do quantization in TFLite, it's your turn. In this assignment we are going to quantize a ``rock, paper, scissors'' hand gesture recognition model. We'll set up the problem for you and you’ll be tasked to perform conversion and set up the interpreter for inference.

\url{https://colab.research.google.com/github/tinyMLx/colabs/blob/master/3-3-14-Assignment.ipynb}

\section{Assignment Solution}

For this assignment you needed to define the converter, run it and then define the TFLite interpreter.

You can instantiate the converter, leveraging the saved model file as noted in the hint, using the handy saved model function shown below:

\begin{lstlisting}
converter = tf.lite.TFLiteConverter.from_saved_model(

ROCK_PAPER_SCISSORS_SAVED_MODEL)
\end{lstlisting}


You then can run the converter as follows:

\begin{lstlisting}
tflite_model = converter.convert()
\end{lstlisting}

It's that simple! TFLite does all of the heavy lifting under the hood for you.

Defining the interpreter (again leveraging a saved model -- this time our TFLite model) is also a simple one line command:

\begin{lstlisting}
interpreter = tf.lite.Interpreter(model_path=TFLITE_MODEL_FILE)
\end{lstlisting}

And that's it! From there you should have been able to test out the model's accuracy and visualize which images it got correct and incorrect. We hope you also took us up on the challenge of improving the model. Let us know in the comments to the Colab what worked best. We’re curious to see how accurate you all can make this model! As always you’ll find the link to the assignment again below in case you want to explore it a bit more!

\url{https://colab.research.google.com/github/tinyMLx/colabs/blob/master/3-3-14-Assignment.ipynb}

% 1.4: Machine Learning on Mobile and Edge I

\section{Why are 8-Bits Enough for ML?}

When neural networks were first being developed, the biggest challenge was getting them to work at all! That meant that accuracy and speed during training were the top priorities. Using floating-point arithmetic was the easiest way to preserve accuracy, and GPUs were well-equipped to accelerate those calculations, so it’s natural that not much attention was paid to other numerical formats.

These days, we actually have a lot of models being deployed in commercial applications. The computation demands of training grow with the number of researchers, but the cycles needed for inference expands in proportion to the number of users. That means that inference efficiency has become a burning issue for a lot of teams in organizations that are deploying ML solutions, TinyML included. 

That is where quantization comes in. It’s an umbrella term that covers a lot of different techniques to store numbers and perform calculations on them in more compact formats than 32-bit floating-point. We are going to focus on an eight-bit fixed point in this course, for reasons we will go into more detail on later.

\subsection{Why does Quantization Work?}

Neural networks are trained through stochastic gradient descent; applying many tiny nudges to the weights. These small increments typically need floating-point precision to work (though there are research efforts to use quantized representations here too), otherwise, you can get yourself into a pickle with such things as “vanishing gradients.” Recall that activation functions (Course 1, ding ding ding!) restrict the large range of values into a rather confined numerical representation so that any large changes in the input values do not cause the network to have catastrophically different behavior. 

Taking a pre-trained model and running inference is very different. One of the magical qualities of deep networks is that they tend to cope very well with high levels of noise in their inputs. If you think about recognizing an object in a photo you’ve just taken, the network has to ignore all the CCD noise, lighting changes, and other non-essential differences between it and the training examples it’s seen before, and focus on the important similarities instead. This ability means that they seem to treat low-precision calculations as just another source of noise, and still produce accurate results even with numerical formats that hold less information. For instance, does your brain recognize the below picture? The answer has got to be a yes, unless, of course, you weren’t watching the videos!

\begin{figure}
    \includegraphics[width=\textwidth]{TinyML/2 Applications/BlurryPicture}
    
    \caption{blurry picture of mobile phone}
\end{figure}  


You can run many neural networks with eight-bit parameters and intermediate buffers (instead of full precision 32-bit floating-point values), and suffer no noticeable loss in the final accuracy. OK, so fine, sometimes you might suffer a little bit of loss in accuracy but often the gains you get in terms of performance latency and memory bandwidth are justifiable.  

\subsection{Why Quantize?}

Neural network models can take up a lot of space on disk, with the original AlexNet being over 200 MB in float format, for example. Almost all of that size is taken up with the weights since there are often many millions of these in a single model. Because they’re all slightly different floating-point numbers, simple compression formats like “zip” don’t compress them well. They are arranged in large layers though, and within each layer, the weights tend to be normally distributed within a certain range, for example, -3.0 to 6.0.

The simplest motivation for quantization is to shrink file sizes by storing the min and max for each layer and then compressing each float value to an eight-bit integer representing the closest real number in a linear set of 256 within the range. For example with the -3.0 to 6.0 range, a 0 byte would represent -3.0, a 255 would stand for 6.0, and 128 would represent about 1.5. This means you can get the benefit of a file on disk that’s shrunk by 75\%, and then convert back to float after loading so that your existing floating-point code can work without any changes.

Another reason to quantize is to reduce the computational resources you need to do the inference calculations, by running them entirely with eight-bit inputs and outputs. This is a lot more difficult since it requires changes everywhere you do calculations, but offers a lot of potential rewards. Fetching eight-bit values only require 25\% of the memory bandwidth of floats, so you’ll make much better use of caches and avoid bottlenecking on RAM access. You can also typically use hardware-accelerated Single-Instruction Multiple Data (SIMD) operations that do many more operations per clock cycle. In some cases, you’ll have a digital signal processor (DSP) chip available that can accelerate eight-bit calculations too, which can offer a lot of advantages. (Don’t worry if you don’t fully understand these hardware capabilities, we will discuss them more in Course 3.)

Moving calculations over to eight-bit will help you run your models faster, and use less power, which is especially important on mobile devices. It also opens the door to a lot of embedded systems that can’t run floating-point code efficiently, so it can enable a lot of applications in the TinyML world.

\section{Post Training Quantization (PTQ)}

VIJAY JANAPA REDDI: Hi.
In this video, I'm going to take you through the realm of quantization.
Lawrence introduced us to a variety of different APIs calls,
like the code that you're looking here, where in TensorFlow Lite,
you can ask the system to optimize the network for different characteristics.
Specifically, you can ask TensorFlow Lite
to optimize the network in terms of size or latency,
or there's some sort of default setting.
What I want to do is try and explain to you what
is really going on behind the scenes.
Lawrence gave you a little bit of intuition
about the fundamental concept.
We're going to take that concept much further so you have a deep appreciation
for why we want to do quantization.
What exactly is quantization?
Quantization is an optimization that works
by reducing the precision of the numbers that
are used to represent a model's parameters, which
by default, in TensorFlow, happen to be in 32-bit floating-point numbers.
Now, if you can reduce the precision from 32-bit floating-point numbers
into something smaller, then you will be able to reduce the model size.
And you will be able to increase portability of that model.
Why?
Because the model is smaller so it can fit onto smaller constrained devices.
Not all embedded systems have a large amount of memory.
And if you also make the model smaller, there's
a chance you might be able to increase the speed of the model,
the inference performance of the model.
We'll talk about how that actually happens in this video.
When we're talking about reducing the precision, what we mean
is that we're going to take a value that's
expressed in float 32-bits, which takes 4 bytes,
and compress that value down into an int8.
An int8 is only 1 byte.
So immediately, if you do that, you get a 4x reduction.
You go from 4 bytes down 1 byte.
That sounds wonderful.
Now nothing's a free lunch.
There are critical trade-offs that we have to make when we do this.
We'll learn about them.
But first and foremost, let's do a quick recap
on why we actually want to quantize.
Given that quantization is so important, it
is extremely crucial we appreciate it.
The three fundamental reasons-- size, latency, and portability--
let's step through each one of these.
The first thing is about size.
What about size is actually critical?
The first thing that comes to mind is when you have a neural network model,
you've got to store that model someplace.
That storage size becomes critical.
Now if you remember, the storage size on our microcontroller
is a 1 megabyte size.
Now if you bounce back a couple of videos earlier,
I had mentioned that a MobileNet model for instance, which
is something that is often used in vision tasks
on small little mobile devices or even embedded devices,
tends to be in the order of several megabytes.
But we only have 1 megabyte of storage.
You can't even store it on the device.
So we've got to think about how we can compress the models down.
If you can take a 32-bit model and reduce it down into an 8-bit model
where each weight, for instance, is expressed in 8 bits instead of 32 bits,
then there's an automatic 4x reduction in memory.
So you might be able to fit that onto your device.
That's the first reason.
There's a second reason.
The second reason is primarily from the RAM size.
Now you've got a model that-- let's say you're able to store on the device
right now.
Now you compress that model down and you want to be able to run it.
Well, even during execution time, there are implications of that.
For example, the graph on the right here is showing you the RAM size
collected from our junior device.
And what you will see is that the TensorFlow stack, which
is the orange color, and the red, the model size,
effectively are the neural network parts.
And the rest of it, there's a whole bunch of stuff
that's actually just dynamic memory and the normal memory usage of the program.
So what I'm trying to get across is that even though we are talking about having
a RAM of 256 kilobytes, you cannot assume that all of that is purely
dedicated for your machine learning model or even the machine-learning
inference engine, the TF Micro.
You have to make assumptions for the fact
that an application is comprised of machine-learning tasks
and non machine-learning tasks.
So there's a bunch of goo that actually holds the application together.
And that also requires memory usage.
So the lesser amount of RAM you can actually take,
the better it is because it frees up the memory
for other aspects of the application, which is crucial in the embedded system
world.
One of the reasons we can actually do this massive reduction in terms
of going from floating-point values down to 8-bit values
is because the weights that are in the neural network
tend to be very concentrated around certain critical ranges.
For example, here is the network distribution weights for AlexNet.
What you are seeing is that it's a histogram, which
shows that most of the weight values, which is really all these bars,
are all kind of clustered together into two ginormous lumps.
It's not a very broad distribution.
It's a rather narrower distribution.
So therefore, it's possible that you can actually take these values
and then be able to express them in a limited range, which
an int8, for instance, might be able to offer.
That's the critical observation.
We'll get into this more in depth in this video later.
That's for the size.
Now let's talk about latency.
Where does latency improvement come from if you do quantization?
Think about this.
If you're operating on a floating-point value versus if you're
operating on an integer value, a big difference
is the fact that floating-point arithmetic
as implemented in any silicon is often slower, much slower than an integer
arithmetic.
Integer arithmetic might only take a cycle or two cycles,
as we would see, as compared to a floating-point arithmetic, which
might actually take over 10 cycles to 15 cycles or so.
So it's almost an order of magnitude different in terms
of the computational performance requirement.
So reducing a model from using floating-point values
to int8 values that are easier to physically compute in silicon actually
gives you a massive improvement in performance.
Another benefit that you get by doing this translation
from floating-point integer values is that you can cut down the power.
Now think about where that reason might come from.
Why would the power go down if I'm changing
from floating point to integer?
Well, it's the same reason as the silicon aspect.
Floating point arithmetic, which takes longer to compute,
is also more power hungry.
And so by switching to a cheaper way of doing arithmetic inside the processor,
we're able to reduce the power consumption down.
This is quite substantial.
So here, I'm showing an example for three different models that
are operating in either floating point, which is the top gray bar,
or integer, which is the green bar right next to it.
And what you're seeing is in all three models, the green bar is much shorter.
The reason is because the arithmetic is actually running faster.
The latency for the model is much, much better.
What this ends up translating to is you get better experience if you are,
for instance, deploying a Tiny ML application
onto a resource-constrained device, it's actually
able to perform many more inferences per second, which is awesome.
So what a quantized model is giving us is
that it's making the models run faster, and it's also shrinking them down
so they can actually fit on the limited memory that we have.
So overall, this is awesome.
Now another key reason is portability.
Often when people talk about quantization,
they only tend to focus on the size and latency and tend to ignore portability.
But portability in embedded systems is extremely crucial.
Why?
This goes back to the fundamental message
that I tried to convey to you in course one that I'm repeating here again.
Embedded systems are not easily portable.
The systems often make very aggressive trade-offs.
For instance, one piece of an embedded silicon
might support floating point arithmetic.
Another might not support it.
So then imagine, you are training the model here and preparing it
for inference.
You want to deploy it on some embedded microcontrollers.
So you do your optimizations and then you're ready to deploy it.
And it happens to be in floating point values.
Well, that's only going to work on a subset of the embedded systems.
It's not going to work on the other subset.
So now suddenly, you have a model that's not generally
applicable across the ecosystem.
And recall why embedded engineers make this trade-off.
They make this trade-off because they want
to keep the cost of the embedded system low.
And also they want to keep the power consumption on the embedded system
low, which means by removing functionality like the floating point
unit, they can effectively improve the overall system power
efficiency, which is number one and it's king in the embedded system world.
So you can effectively quantize the network
into integer, which is a common baseline numeric that
is actually available across all of the different systems.
Then you get better portability, and you get better efficiency.
Great.
So now let's go deeper and ask a question
of how exactly do we do the quantization.
I don't mean what function calls you call.
I mean what happens behind the function calls.
OK.
So when you take a neural network, calculations are all over the place.
Hopefully, you will recognize some of this.
One of the first things you will notice is
that we have neurons that have weights.
Then, these neurons also have biases.
For every single computation you're doing,
there's a bias that's associated with that network.
There are many different bias vectors all over the network.
Then you have activations that are actually going out the neuron.
So you can perform quantizations at various different levels.
One of the things that we're going to focus on
is on doing quantization with weights because weights
play a major role in terms of storage, for instance,
because those are the things that you actually physically
store inside your flash memory.
Now earlier, I showed you this graph on the left, which
is showing you the weight distribution, the values for all the weights
across all those neurons across the entire network.
And what you're seeing is that there's limited range.
That's a first point.
Second is that you can actually quantixze these values
without losing the general representation of that space.
So the values in the right graph are quantized values.
Quantization effectively means, again, binning the floating point ranges
into little buckets.
We'll talk about that more.
But effectively, if you look at the graph on the left,
and you look at the graph on the right, you discretizing the data.
But can you still help the same trend?
Yes.
You still have those two big lumps.
And you have that tailing effect in the second lump.
So how can we go further with this then?
Effectively, what's happening is that when
you are taking those original weight values, which are in floating point
representation, you are then binning them into an 8-bit encoding scheme.
Remember, because we are using 8-bit integer values,
we have 2 to the 8, or 256 potential values,
including 0 all the way up to 255.
That's 256 values.
So we discretize that long floating point range.
And we took those values and we bucket them.
We hash them into this 8-bit range that we have.
Then when we are ready to do the computation again--
now that the network is actually 8 bits in the middle stage, you save it.
Then when you start running it, what you want to do is you take the 8-bit values
and you reconstruct them into a 32-bit floating point value.
And you can run it.
This is the first step of the quantization scheme,
because now imagine that you've taken a floating point values
and you've compressed the weights into 8-bit encoding,
so it's shrunk the model down by 4x.
And so you're able to fit it.
At this point, let's assume that you have floating point arithmetic
capability.
And so therefore, you would effectively decompress it and then
be able to run it.
Now, one of the things that's going to happen
when you do this is that certain values that are inside the floating point
original range will get bucketed into some bit encoding.
For example, negative 5.4 gets bucketed into 0, which later you
decompress into the negative 5.4 value, because 0 happens to map for that.
However, sometimes you can have a situation where two floating point
ranges are so close together, the two floating point values
are so close to each other, they fall into the exact same integer bucket.
The unfortunate side effect of this is that when you decompress the value,
it's going to come out into the exact same floating point value, which
means now you've lost a little bit of that resolution
that you previously had in the original.
This is a trade-off you will be making.
There's also the aspect that, again, you have the issue that one of the values
may not precisely map into an exact value.
So this generally gives you a flavor for the issue that you might be seeing.
The first red box being the most important one.
Now let's go a bit further in this quantization.
When we're talking about post-training quantization, which
is on the left-hand side, you can go further into this
and talk about weight compression and calculating the inference
much more effectively.
Now remember, when we're talking about weight compression
on the left-hand side, what we're doing is we're taking the integer value,
converting that into a floating point value when we're about to run it.
And that's exactly what the pseudocode here is showing.
When we have compressed volume, we are storing it on the disk.
And then we're about to run the model.
For every time we get the quantized value, which is an integer value,
we convert that or decompress that value as you
can see in the below box into a floating-point value.
The arithmetic map here is very simple.
Take a moment to look at it and you will understand.
It's effectively looking at what the range is and scaling the range back
into its original value.
That's it.
I can give you some pointers that you can
look at later in the reading section.
In the end, once the decompress function finishes,
you now have a floating-point value back from the stored integer value.
And then you perform the multiplication.
One great benefit is that the size is now smaller.
But recall here, as you're doing the input n times the decompress,
that multiplication is still a floating-point multiplication.
So we have addressed the first part, which is we have compress network.
But we're decompressing it into floating points, so we can actually run it.
Now there's the other alternative where you actually
want to speed up the inference.
If you want to speed up the inference, you've got to get rid of that multiply
from being a floating-point value.
And you've got to convert that multiply in the output box
into an integer value.
That's where the second leg of post-training quantization comes in.
You take the input and you quantize the input itself into a lower value.
And you take the weight and you also do the quantize on the weight value.
And effectively, this way, what you have done
is you've converted the floating-point value into an integer value.
That's ideally what you want to get to.
In doing so, now you can actually use integer multiplication,
instead of using floating-point multiplication,
which again is much slower than the integer multiplications.
Now at a very high level, this should cover the core nugget
that post-training quantization has two legs to it.
You have one way of just compressing the network
and this is storing it as integers.
The second is rather than kind of decompressing it back
into floating point, you can actually use integer arithmetic on it.
But then you would also have to quantize inputs.
Now hopefully, that's the key takeaway that you're
going to remember from this part.
Now what are the trade-offs?
I said earlier, there's no free lunch.
When you quantize, yes, the network is smaller.
Yes, you can run the network faster.
But then you end up losing accuracy.
That's the critical trade-off you're going to make.
What this graph is showing you is that you're
able to take a point on accuracy versus latency
and move that point to the left, which means you have
decreased the latency of the network.
But then you see that it's also a little lower, which
means you have dropped the accuracy.
And the same thing happens for a variety of different models.
This is actually coming from the original MobileNet's model.
And they are looking at what happens when
you do quantization across a whole range of different kinds of MobileNet models.
Now what else happens?
We have quantized the weights.
Now weights are just one portion of the network as I was saying earlier.
There are many other portions that you can potentially
quantize within the neural network.
You can quantize the activations.
You can quantize the channels.
There are many, many different places you can actually do quantization.
And they offer different trade-offs.
There is no one silver bullet answer.
But that said, weights are often the place
you tend to quantize because they are all over the network.
So in summary quantization is great because if you use 8-bit integer
values, you get faster arithmetic.
At the end of the day, ideally, you not only
want to store the model as integer 8-bit.
You also want to compute in 8-bit.
And that gives you faster arithmetic which improves both latency.
And then by doing the quantization, you can also
cut down the memory because your models are now a quarter of the size.
Finally, overall, for microcontrollers, it's
extremely important that we do integer arithmetic because they don't always
have floating-point hardware.
So you almost always are going to be quantizing.
In the context of a tiny amount, it is extremely important
that you really take quantization to heart.
And with that said, I'm going to leave you recalling
the critical aspects of quantization, specifically
from post-training quantization.
And in the next video, I'm going to pick up on quantization aware training
in order to address the accuracy loss that we saw as a result of making
the accuracy latency trade-off.

\section{PTQ Weight Distribution}

In this Colab we are going to explore Post Training Quantization (PTQ) in more detail. In particular we will use Python to get a sense of what is going on during quantization (effectively peeking under the hood of TensorFlow). We will also visualize the weight distributions to gain intuition for why quantization is often so successful (hint: the weights are often closely clustered around 0).

\url{https://colab.research.google.com/github/tinyMLx/colabs/blob/master/3-4-3-PTQ.ipynb}

\section{Quantization Aware Training (QAT)}

VIJAY JANAPA REDDI: More on quantization, specifically
a quantization-aware training.
Now, previously, when we talked about post-training quantization,
one of the key things that I was trying to emphasize
are the positive aspects of post-training quantization.
Take a minute and think about what are the reasons
that we want to do quantization?
It is about size.
It's about latency.
And it's about portability--
all wonderful things to achieve, which we get with post-training quantization.
However, there is a trade-off, if you remember.
The trade-off that we make is that when we do quantization,
our accuracy, which here is being shown for floating-point values,
is going to drop when we do post-training quantization.
As you can see, across all three random sample networks,
there is an accuracy drop.
In MobileNets, you see an accuracy drop that is as much as 1.46%.
Now, the question you've got to ask yourself at this point,
is this accuracy drop OK?
Is it worth it in order to reduce the size of the model?
Or is it not something that is actually acceptable?
This question is really a function of the application use case.
You've got to ask yourself, for your particular application,
is this accuracy drop actually tolerable?
For example, if I'm looking at a cat or a dog in a picture on my camera,
maybe it's not the end of the world.
But if I'm a doctor looking at an image that's
using machine learning to try and detect cancer cells,
now, is an accuracy drop there the same as being
able to have an accuracy drop on that cat versus dog?
I think you know where I'm going with this.
Accuracy drops in medical applications are extremely sensitive.
And that's going to be the case in whatever deployment you have.
So is there anything we can do to improve
the accuracy while keeping the benefits that we
get from post-training quantization?
And the answer is yes.
But before I give it to you, pause and think for a moment.
Where do you think the error is coming from when
we do post-training quantization?
We're taking a floating-point value.
We're converting it into integer 8-bits.
That's where the conversion problem is.
The conversion problem effectively is chopping off
some bit of information that the floating-point values have,
which is a much wider range.
Now, is there anything we can do?
Yes, the answer is.
Can we try and recoup that error?
What if you could tell the network, hey, you
are going to have this error if I do quantization
on you while the network is training.
Can it automatically learn that error to become resilient, effectively,
against that error?
That's the critical observation here.
By making the observation that when you do post-training quantization
you're getting an error, and then saying,
oh, I know I'm going to do post-training quantization later
and there's going to be some amount of error, which is not so great.
Then how about I can train the network to automatically learn
that I'm going to be doing this sort of a quantization later
and then inject that error into the training, into the training pipeline,
so that the network naturally learns to become resilient against that error?
And so when it's doing its forward pass and back propagation to fix the error,
it'll naturally learn over that time to become resilient against that error.
And the accuracy will go back up.
Now, that is a critical observation about quantization-aware training.
So quantization-aware training effectively
emulates the inference-time quantization behavior
by creating a model the downstream tools will use
to produce the actual quantized models.
The quantized models, naturally, will be using lower precision.
And so knowing that we'll be using lower precision
and introducing that into the training pipeline
while its training will effectively allow the network to become resilient.
Let's take an observation about where that accuracy drop is coming from.
Remember, we took an original continuous signal.
That's in the floating-point values up on top-- nice curvy line.
Then we discretize that signal because we're now
representing it using integer values, which
means now you can see that the lines are more choppy,
which is effectively what we're doing when we go into integer 8.
And that's because we're going from floating-point integer 8 values, which
doesn't have the same sort of resolution.
Now, in addition to that, there are many other conversions
that happen behind the scenes.
It's not just the weights.
You might be about how you're actually doing your multiply accumulates
for the weight computations, where you might
have to accumulate things into an integer 32 value
after you do the additions.
And then you bend them back into int8 and so forth.
There's a lot more nuance, which I'll cover in the reading materials.
But for now, N I'm trying to get across is the critical point that there are
many sources of accuracy drops, the biggest one being the one that we
talked about, floating-point to 32-bit integer--
floating-point 32-bit to int8.
So in order to introduce the error into the training graph, what we want to do
is we want to modify the original graph that we typically
have in floating-point values so that when we're doing the computations,
we are implicitly quantizing the values.
And when we do the calculations, the natural output that goes out
when we're doing the weights times the inputs plus the biases,
that computation effectively should have the error
that we are going to have in the post-training quantization.
And by injecting that error directly into the training pipeline, what
happens is a network is only seeing that error be propagated through.
And so when it does its back propagation and it does its gradient descent
and the loss functions are working on it,
it's naturally going to try and become resilient to that issue.
And so effectively, by exposing the training pipeline to the error,
we're going to be able to naturally get the neural network
to learn to improve its accuracy by itself.
And that will allow us to then later on be able to use the network effectively
by using the int8 multiply accumulates.
So as an interesting result here, let's take a look
at what happens when we are comparing post-training quantization
versus quantization-aware training.
The first thing you will notice is that in MobileNets, we won.
Our accuracy for post-training quantization, which was 69.5%,
got better with quantization-aware training because it is now 71.06%.
However, in the second one, if you see, our post-training quantization
had 70.2\%, whereas our quantization-aware training is only
at 70.01\%, which means quantization-aware training is not
guaranteed to always be able to do better.
It'll mostly do better.
That's the general wisdom out there.
This is part of the black box magic with machine learning.
You have to play around with it a little bit.
Some of you probably already witnessed this
when you're trying to choose what the loss function should be
or how you should optimize your learning rate.
Well, this is yet another one of those hyperparameters
that you will have to think about.
Another interesting thing I want to point out in this particular table
is that if you look at MobileNet--
we want the first line--
an interesting observation is that the original floating-point value actually
had an accuracy of 71.03%.
But if you look at the quantization-aware training cell,
it is 71.06\%, which means that the quantization-aware model is actually
outperforming the floating-point baseline, which
is a really interesting observation.
And this, again, goes back to the point that, hey,
you cannot control exactly what's happening in quantization because when
you're doing quantization-aware training,
the network is learning a little differently from the way it originally
learned.
So key takeaway, though, is that in general, quantization-aware training
has a lot more benefits.
And it almost always tries to do better than the post-training quantization
because you're learning to be resilient against the error.
And this is why Lawrence was earlier showing you,
hey, you can take certain tools that TensorFlow offers
and you can extract the operators that allow you to do model quantization
original graph.
So when you go to quantize model, you take the original model
and you give it this graph.
And what that API is doing is that it's extending that original network
graph with the ability to actually mimic the quantized behavior that
would be happening during the inference time, during the training time.
And so when you train the network, the network
is implicitly learning to be resilient.
And therefore, what you find is that when
you're running a quantization-aware trained INT8 model,
its accuracy is very close to the FP32 model.
So in summary, what I hope you take away is
that you have learned what is quantization.
Then I hope you learned that there's post-training quantization.
And then it's quantization-aware training.
And that why quantization is especially important for us
in the TinyML context given size, portability, and latency.
And that what happens in post-training quantization
is that you are introducing a little bit of error
as a result of mapping everything from floating-point into 32,
but you get the benefits of the quantized smaller model.
And you can also get the benefits of improved latency
if you use the numerical methods involved for doing int8 calculations,
where your multiply accumulates are actually
happening in int8 instead of FP32.
And hopefully you learned what is quantization-aware training
and that quantization-aware training is something
that happens during the training phase, where you're actually
injecting the error in order to reduce the quantization loss.
And with that said, that effectively wraps up
the whole thing about quantization.
Hopefully you understand when Lawrence was talking about with respect
to quantization APIs.
I've now kind of opened up for you.
And so you're able to peek a little bit into the box.
There's a lot more nuance in quantization.
You can go pretty deep into post-training quantization
and quantization-aware training.
But the nuggets that we touched on here, these
are the fundamental building blocks.
And this is the high-level intuition that you need to be successful.
Anything beyond this is going to be really just kind of specializing
in the quantization screen.
Now, that's great for researchers.
But if you want to be a TinyML engineer, you
don't necessarily have to push too much harder beyond what we have covered.
And with that, I'll see you in the next video, my friends.

\section{Inference Engine: TF vs. TFLite}

VIJAY JANAPA REDDI: Hi.
In this video, we're going to learn about the inference engine,
or the machine learning framework that is actually
going to support our models that we're working very hard on making small
and so forth.
It's extremely important that you understand these differences.
Because as a TinyML engineer, you're going to need it.
Because every little bit counts in terms of extracting efficiency
from the system.
In this course, as we go through this machine learning workflow,
we're going to be using both TensorFlow and TensorFlow Lite.
In course 3, we'll actually be using TensorFlow Lite from microcontrollers
specifically.
In this video, I'm going to try and focus
on the big differences between TensorFlow and TensorFlow Lite.
I want to make sure you understand what are the technical differences
between these software stacks.
Because this is going to be fundamental to actually
how you go about building your TinyML applications.
Let's start off with TensorFlow.
TensorFlow is a machine learning framework that runs in the cloud,
as you are doing with your Google Colabs.
It relies on big, powerful processors.
That might be CPUs, GPUs, and these custom-built Google
TPUs, which are used for efficiently running neural network processing
graphs.
However, when you're doing the models, you're training them,
you need big frameworks.
But if you're deploying them onto small devices, for example, like your phone,
you don't necessarily have all of that massive computational horsepower
and capabilities.
Why?
Because these devices have less memory.
They have less compute power.
And also, they're only focused on inference.
So then why do you need a big training framework,
a system that's capable of both training and inference, when you only
do one half of that world on something like a mobile device
or an embedded device?
That's the fundamental rationale behind TensorFlow Lite.
Let's make something nice, and clean, and mean.
Now that boils down to some critical differences
in terms of how these two systems are effectively implemented.
Let's talk about them.
From a network topology, think about what's
happening when you are training the network versus when
you're deploying the network.
When you're training the network, are you typically
playing around with the network, like what operations to use,
what max pulling you use, maybe doing drop-out as an optimization.
These kinds of things require you to change the network structure.
Therefore, TensorFlow has a lot of capability for that.
But when you're only doing inference, are you
ever changing the network structure?
The answer is no.
Because you've already designed the model, and you're deploying it.
Similarly for weights, when you're training the network,
are your weights constantly changing?
The answer is yes, right?
Because we're doing a forward pass.
And then we're coming back, and we're doing backpropagation,
and updating the weights using gradient descent on our loss functions.
So our weights are constantly getting updated.
Versus in TensorFlow Lite, are we doing any weight updates?
The answer is no, because we're deploying it.
Think about the memory.
I've emphasized this many times, that when you're deploying things
onto mobile devices and embedded systems,
memory is always extremely critical.
So we tend to focus on that.
So we want to strip away anything that is in the framework that
actually ends up eating the memory.
Because not only do we want our models to be small,
we actually need the frameworks that is supporting the model also to be small.
Do we need distributed compute?
Do we need multiple computers to work together?
Well, when you're training a complicated model, the answer is yes.
You need multiple computers to actually do the computation.
But when you're doing inference calculations,
no, you're typically running the model on a single device.
Therefore you don't need all of this extra capability that's
in the software framework to actually be able to do your inference efficiently.
Bottom line-- TensorFlow is really designed for machine
learning engineers and researchers who are trying to architect the model.
TensorFlow Lite is designed for application developers.
And by the way, you're in a TinyML applications course right now.
So let's go a little bit further to try and understand
the critical differences between TensorFlow and TensorFlow Lite
in these three buckets.
Because this has implications.
In addition to just talking about TensorFlow Lite,
I'll also talk about TensorFlow Lite Micro, which
is an extension of the TensorFlow Lite ecosystem, which
Lawrence hinted at earlier.
But let's start off with the model size.
From a modeling perspective, as I've said earlier, TensorFlow,
you're going to be doing training.
But in TensorFlow Lite and TensorFlow Lite for microcontrollers,
you don't need to be doing any training.
So you can strip away those parts.
Also, can TensorFlow do inference?
The answer is, absolutely, yes.
But the question is whether you can do it efficiently or not.
It's not necessarily optimized for that.
It's really optimized for the training side.
However, TensorFlow Lite and TensorFlow Lite Micro
are extremely optimized for inference.
In course 3, we'll go much deeper into this.
Also, how many ops a particular framework supports depends on the task
it's designed to do.
As a machine learning researcher or engineer,
you're often exploring many different kinds of operators when
you're trying to train the network.
Therefore TensorFlow has over 1,400 ops right now.
That's a lot of ops.
But many of those ops are actually not used by the widely popular models.
Therefore TensorFlow Lite and TensorFlow Lite for microcontrollers
only have a handful of ops.
There is almost 130 ops in TensorFlow Lite
as compared to TensorFlow Lite for microcontrollers,
which has only 50 ops.
This is an order of magnitude different.
Just from an efficiency perspective, you want
to focus on where things really matter, not just
provide general capability, which is absolutely
needed in TensorFlow for being able to train the models effectively.
And the hope is that many of those critical ops that are being trained
are actually the ones actually are useful.
And it turns out only about 130, 150 ops are typically useful.
And also there's the whole notion of native support for quantization.
At this point, you know why quantization is so crucial to us.
TensorFlow doesn't really support any specific quantization,
because it's not really needed.
Versus TensorFlow Lite and TensorFlow Lite for microcontrollers, quantization
is sort of natively supported.
The optimization APIs are all naturally there.
Let's look at the software side.
First thing is, do you need an operating system in order to run the frameworks?
On TensorFlow, the answer is yes.
You need something like Linux, or Windows, or so forth,
because it relies on things like dynamic memory allocation.
Same thing happens for TensorFlow Lite, which
typically runs on smartphones, where you tend
to have something like Android or iOS, which can actually do dynamic memory
allocation.
But in TensorFlow Lite for microcontrollers, the answer is no.
Because if you had taken course 1, you would
know that we don't have nice operating system support for microcontrollers,
because we want the microcontroller to be lean and efficient.
So you don't necessarily need to have a mid-level operating system support.
Also, you need the ability to map models efficiently
in embedded devices and mobile devices.
Memory mapping is a very quick and effective way
of bringing the model into memory, which is something that TensorFlow
and TensorFlow Lite differ upon.
Finally, the ability to delegate to hardware
accelerators-- we'll talk about this in much more detail down
the road in course 3.
But for now, the ability exists in both TensorFlow, and TensorFlow Lite.
And TensorFlow Lite Micro is soon going to have the ability
to deploy onto custom micro neural processing
units on embedded microcontrollers.
Final element is the hardware differences.
Think about the binary size.
I've always said that it's extremely important to keep the binary size small
because of flash, for instance, on microcontrollers, is extremely small.
Similarly, on mobile devices the memory is also very limited.
But if you look at the binary size of TensorFlow, it's 3 megabytes,
versus on microcontrollers and mobile devices,
it's only in the order of, like, a few kilobytes.
That's because TensorFlow Lite and TensorFlow Lite for microcontrollers
has been extremely optimized to keep the whole software
infrastructure itself very small.
So the base footprint, once the model is actually up and running,
as compared to TensorFlow, is an order of magnitude
smaller for TensorFlow Lite for microcontrollers and TensorFlow Lite.
And finally, these systems are optimized for different hardware execution units.
TensorFlow is optimized for x86, the GPUs, and the Google TPUs.
TensorFlow Lite, which is optimized for mobile devices,
is optimized for Cortex A, or application class processors,
and x86 processors.
Whereas TensorFlow Lite for microcontrollers
is really custom optimized for embedded systems,
which are like the Cortex m class microcontrollers
from Arm or the digital signal processors, all of which--
this particular thing will be something that we
talk about in great detail and course 3 so
that we can actually close the loop in terms
of understanding how we go from building models, to building applications,
to actually deploying them very efficiently.
Now with all that said, the key message that I'm trying to get across here
is that it's important not only to optimize the models,
but it's also important to optimize the infrastructure, or the runtime that
actually runs that model during inference time.
I hope you took away that there are big differences between TensorFlow
versus TensorFlow Lite ecosystem.
Even within the TensorFlow Lite ecosystem,
which is custom designed for the inference world,
is also split up further based on mobile devices
versus the embedded ecosystem.
The reason is because even when you do inference,
it depends on the kind of system you are running inference on.
For example, embedded systems do not have an operating system.
They also have extremely limited amount of memory.
So there are some big differences.
So hopefully this video gives you a little bit of insight
into why we actually use the converter, and what happens behind the converter
calls in terms of optimization, and why the converter
itself is actually designed to be lean and mean.
With all that said, I'm very much looking forward to the next video,
where we're finally ready to get started with our TinyML applications.
Because everything we have covered so far
lays the foundations that we will be using for building our applications.
I'll see you in the next video.


\section{Conversion and Deployment}

Now that we have talked about the difference between TensorFlow and TensorFlow Lite, let's dive a little under the hood to understand how models are represented and how conversion is processed.

\subsection{Tensorflow's Computational Graph}

Tensorflow represents models as a computational graph. To better understand what that means, let's use a simple example of the program shown below:

\begin{lstlisting}
var1 = 6
var2 = 2
temp1 = var1 + var2
temp2 = var1 - var2
var3 = 4
temp3 = temp2 + var3
res = temp1/temp3
\end{lstlisting}

This program can be represented by the following computational graph which specifies the relationships between constants variables (often represented as tensors), and operations (ops):

\begin{figure}
    \includegraphics[width=\textwidth]{TinyML/2 Applications/3-4-7.png}
    
    \caption{Constants like the numbers 6, 2, and 4 are connected to variables (like the temp variables mentioned above) and can then have operations like addition, subtraction, and division apply to the variables to combine all of the values to produce the final mathematical output.}
\end{figure}  



The benefit of using such a computational graph is twofold:

\begin{enumerate}
  \item Optimizations such as parallelism can be applied to the graph after it is created and managed by a backend compiled (e.g., Tensorflow) instead of needing to be specified by the user. For example, in the case of distributed training, the graph structure can be used for efficient partitions and learnable weights can be averaged easily after each distributed run. As another example, it enables efficient auto-differentiation for faster training by keeping all of the variables in a structured format.

  \item Portability is increased as the graph can be designed in a high level language (like we do using Python in this course) and then converted into an efficient low level representation using C/C++ or even assembly language by the backend compiler.
\end{enumerate}

Both of these reasons become increasingly important as we increase the size of our models --- even for TinyML many models have thousands of weights and biases!

If you'd like to learn more about this we suggest checkout this great article which we used for inspiration for this section of the reading.

\subsection{Tensorflow Checkpoints}

As you train your model you will notice that Tensorflow will save a series of .ckpt files. These files save a snapshot of the trained model at different points during the training process. While it may seem like a lot of files at first (e.g., you will notice later that that /train directory will get quite full when you train your custom keyword spotting model in the next section), you do not need to be able to understand their contents. The most important thing to keep in mind about the checkpoint files is that you can rebuild a snapshot of your model. However, at a high level:

\begin{itemize}
  \item .ckpt-meta files contain the metagraph, i.e. the structure of your computation graph, without the values of the variables
  \item .ckpt-data contains the values for all the weights and biases, without the structure. 
  \item .ckpt-index contains the mappings between the -data and -meta files to enable the model to be restored. As such, to restore a model in python, you'll usually need all three files.
\end{itemize}

If you'd like to take Google’s hands on crash course on all things checkpoint files \href{https://colab.research.google.com/github/tensorflow/docs/blob/master/site/en/guide/checkpoint.ipynb}{check out this Colab!}

\subsection{Freezing a Model}

You may also notice some .pb files appear as well. These are in the Google ``\href{https://developers.google.com/protocol-buffers}{Proto Buffers}'' format. A Proto Buffer represents a complete model (both the metadata and the weights/biases). A \FILE{.pb} file is the output from ``freezing'' a model --- a complete representation of the model in a ready to run format!

Google differentiates between checkpoint files and frozen model files (which it refers to as “saved models”) as follows: 

\begin{itemize}
  \item Frozen/saved models are designed to be deployed into production environments from that format
  \item Checkpoint files are designed to be used to jumpstart future training
\end{itemize}

If you'd like to learn more about freezing/saving models you can check out this Colab made by Google.

\subsubsection{Optimizing a Model (Converting to TensorFlow Lite)}

As we described earlier in the course, the Converter can be used to turn Tensorflow models into quantized TensorFlow Lite models. We won’t go into detail about that process here as we covered it in detail earlier but suffice it to say that that process often reduces the size of the model by a factor of 4 by quantizing from Float32 to Int8.

TensorFlow Lite models are stored in the \href{https://www.tensorflow.org/lite/api_docs/cc/class/tflite/flat-buffer-model}{FlatBuffer} file format \FILE{as.tflite} files.

The primary difference between FlatBuffers and ProtoBuffers is that ProtoBuffers are designed to be converted into a secondary representation before use (requiring memory to be allocated and copied at runtime), while FlatBuffers are designed to be per-object allocated into a compressed usable format. As such FlatBuffers offer less flexibility in what they can represent but are an order of magnitude more compact in terms of the amount of code required to create them, load them (i.e., do not require any dynamic memory allocation) and runt them, all of which are vitally important for TinyML!


% 1.5: Keyword spotting

\section{Keyword Spotting - Introduction}

VIJAY JANAPA REDDI: Here we are.
Finally, we are ready to get started on the TinyML applications,
kicking things off with keyword spotting.
We have got a lot of foundations that we have laid out so far.
And we're going to start building applications
on those strong foundations.
One of the key things that I want to get across to you as we kick off
keyword spotting as an application is understanding
that keyword spotting falls in the general realm of speech recognition.
It is undoubtedly one of the most successful examples of TinyML.
It has to operate with ultra low power efficiency.
It has to run continuously and it has to be performed on-device, all of which
we'll be looking at in much greater detail.
We'll learn about keyword spotting data sets.
We'll learn about how to train the models.
We'll learn about how we deploy them and how to systematically evaluate
and think about them.
But before we go there, I want to help place
keyword spotting in the general realm of that speech recognition ecosystem.
Typically, when you're dealing with machines
that have a full-on understanding or have a full-on conversation with you,
there are large ASR models, automatic speech recognition models.
They're large.
They're power hungry.
And keyword spotting is not one of those.
It falls in the realm of overall speech recognition.
But it's a critical front end in order to wake a more complex beast that
is sitting behind the scenes.
You want things that are at the front of the pipeline
to be lean and mean and efficient, like keyword spotting.
What do I specifically mean?
For instance, when you say, Okay, Google, that's
a set of keywords that wakes the machine up,
where you can then perform full-on automatic speech recognition.
In this particular scenario, the aspect of "show me the calendar"
is the full-on speech recognition that's going on there, where there is speech
through text converting the words that I'm saying into text, which are then
used by the machines in order to figure something out and then
be able to come back to you and then do text to speech,
where the machine speaks back to you.
Now, this is in the context of OK, Google.
But the same fundamental thing effectively happens when you say,
``Alexa--''.
The machine wakes up.
I give it a command in general English.
And it processes it and does something.
Keyword spotting here is really, OK, Google or Alexa.
We're going to go pretty deep into that.
These two examples that I've just shown you
are examples that we have used repeatedly in the course so far.
And we have also talked about them in the course one
as saying that they're everywhere.
However, keyword spotting has a lot of applications.
Just think about your household.
Think about the numerous different devices you have in your house.
Just look around your room right now and take, for a moment,
and think, which of these devices could actually
use keyword spotting where you give a command and the machine
actually wakes up to be able to do something?
For example, if you have an oven, you might be able to say, bake at 350.
All you need the machine to do is listen to the word "bake"
and when it hears the word "bake", you know it's the oven.
And then you can give it a simple follow-on command.
Now, keyword spotting does not necessarily
mean that it's just one word.
It could be a few words.
It's a subset of the general English vocabulary, for instance.
Now, this applies, obviously, to any language of your choosing.
But the key thing is that you can take any device that is right around you
and you can think of ways to give it the capability
to be able to respond to keywords that will get it to do something.
Now, in the house, there are a variety of different places.
I'll give it to you as a homework in order
to just go out and have fun and try and think
of how would a particular device around your house wake up.
Think about it in the context of cars.
Most automotive companies today are looking at keyword spotting as a way
to make the user experience much more smoother in the car.
For instance, you might be able to get your car to say, ``Play music.''
Or you might get it to say, ``Roll the window down.''
As you can see, those are very specific and small,
short sentences that I'm using.
And the key thing that you want the machine to be able to
is spot those keywords in the general context of things that
are going on in the speech paradigm.
Another example could be when you are wearing smart glasses in the future--
augmented reality, for instance.
You don't necessarily have to say, OK Google.
Maybe you're having a conversation and you say something in passing.
But the machine is listening for those critical words.
And then as they happen, perhaps it projects an interesting piece
of information on your lens.
Those are the kind of applications that people are really excited about when
it comes to keyword spotting.
That said, think about situations where you're not explicitly
giving commands to the machine.
In the example, it's like, when we say, OK, Google
or when we say Alexa, the key thing that we're doing
is that we're looking at the device and giving it a very specific command.
But there are cases where the machine is constantly listening
and it's able to pick up these keywords on the fly in streaming audio.
A very good example here is Now Playing by Google,
where your phone will be able to tell you
what song is playing in the background if it happens to detect it.
Now, how does that work?
You want your phone to use as little energy as possible.
So you want it to be off most of the time.
Otherwise, you're going to be charging it multiple times throughout the day.
So how is it that we're able to build these phenomenal things
and still have the battery last all day or sometimes even two,
three days, while the machine is still able to sense
and listen to things around us?
You're going to find out in this particular series of videos
that have to do with keyword spotting.
In this series of videos that we have around keyword spotting,
we'll start off by understanding what are the fundamental challenges around
keyword spotting?
It's extremely important to have a deep appreciation
for how to look at the problem space.
Remember, one of the things that I want to make
sure I teach you is that we look at things from an end-to-end pipeline
perspective, not just let me just train my model.
That's not useful.
Trust me.
I'm telling you this from experience.
Once we understand the broader challenges
of how to deploy keyword spotting effectively,
then we will move on to looking at the keyword spotting pipeline.
Remember, I said that we will walk all the way
from data sets all the way to how we deploy the model.
We will start by looking at what are the characteristics of the data
that we need to have?
We will be actually using Pete's own speech data
set, which has become the standard, effectively,
for keyword spotting today.
Then we'll learn how to train a small little model that
can actually fit in a microcontroller.
We'll do the training in Google Colab, of course.
And then in course 3, we'll learn how to deploy it.
Then, of course, we've got to make this model small, right?
So we have to understand how to optimize them.
I'm glad you learned about TensorFlow Lite and the whole thing
about quantization and why it's so critical
because it's going to be used repeatedly in order to shrink our models down.
So hopefully you remember some of those fundamentals that I taught you.
And then finally, we'll actually get into the process
of trying to deploy it.
But that's, of course, in course 3.
The last part that's critical here is that we
will continue to use Google Colab in order to train the model.
We'll look at various metrics and trade-offs that we have,
mostly from a model standpoint.
But I also want to help you understand how
to evaluate this particular application from an end-to-end scenario.
With that, stay tuned, get excited because we're
going to have a lot of fun with keyword spotting.

\section{How (else) would you use KWS?}

What are some creative ways  you could think of using keyword spotting as a use case? Sure, we can wake up the machine to respond to us. But what are some things you think we can do if we are able to detect words in just an ongoing conversation? Get creative!

\section{Keyword Spotting Challenges/Constraints}

VIJAY JANAPA REDDI: All right.
Back at Keyword Spotting.
Specifically in this video, we're going to try and understand
some of the challenges and constraints that we have around how
to deploy keyword spotting effectively.
We're going to focus on bandwidth, latency, privacy, security,
and various kinds of challenges that we have.
I'm not going to go specifically into every single one of those things that
have extensive depth.
Instead, what I'm going to try and do, is try and get you
to think about how we should deploy notion of keyword
spotting in a responsible manner.
That's what the focus of this particular video is going to be.
So let's start.
What are the challenges and constraints?
Broadly, I would bend them into performance characteristics
like latency and bandwidth from a system standpoint, from a model standpoint,
in terms of the accuracy and the personalization that we would expect,
from the safety, security, and privacy standpoint
of preserving our data carefully, and then looking at the resource
constraints that we have on a microcontroller,
and how we deploy things effectively.
Let's go through each and every single one of these,
systematically, starting off with latency and bandwidth.
Latency, the notion here, is that we will
want to provide the results quickly.
And we want to respond when a user says some sort of keyword,
we want to respond very actively.
Imagine you say, "OK Google," and then you have to wait about three seconds.
Do you have any idea how many contact switches
your brain is doing if you have to wait three seconds?
The brain can only focus on things for about 100 milliseconds.
So if there's a gap of about 100 to 200 milliseconds,
your brain automatically switches off and goes on to doing other things.
So think about that aspect of responsiveness.
It's about understanding how we, as humans, operate in general,
and then trying to provide a quality of experience
from a system standpoint that meets that expectation.
And therefore, latency is crucial.
When we use a keyword, we want to see some sort of responsiveness.
Think, for example, when you have an OK Google or an Alexa device sitting around.
Why do you think they tend to have lights or some way
to communicate that the machine has actually heard you?
It is a way for us to be able to take a device
and be able to show the human, who's interacting with this device,
that the machine has actually heard them, that it's actually awake
and that it's doing something.
It gives us a sense of satisfaction in our brain.
That's a means to mask latency, effectively.
Fundamental point I'm getting across, is that latency is extremely crucial
because it matters tremendously in deploying
keyword spotting applications.
The next one is about bandwidth.
One of the things that you want to do, is
that you want to minimize the amount of data movement.
In any computing system, you've got data and you've got compute.
Guess which one is the most difficult thing to actually deal with?
It's actually data movement that's really costly.
Computing is actually very cheap.
The notion of, for instance, when you do an addition doing a 1 plus 1,
from a pure logic standpoint, it's extremely cheap in terms of energy.
But when you have to move data inside the computing processor, from point A
to point B, oh, that is significantly much more expensive.
So you take that and you want to minimize that data movement, not just
within a chip, but across the network.
Because you don't want to send every single word that you say all the way up
into the cloud where you have the big computers process that data,
and then bring it back.
Because that's going to be significantly slow.
So what you want to be able to do is minimize that data movement.
Because it's not only about the slowness,
it's also about the energy that's going to cost in order to move that data.
The more local you keep the data, the better
it is from a battery and lifetime sort of perspective.
The next aspect is accuracy.
Remember, when you want to enable keyword spotting,
you want to wake the machine up.
Well the machine's got to be continuously listening to us.
And it's only got to show that light or some sort of cue
that it's actually heard you when it's the right time.
Because you don't want false positives.
And neither do you want false negatives.
So accuracy is extremely critical.
We're going to try and understand what are
the fundamental metrics that are used in order to assess keyword spotting.
So when you train a model, what are the right metrics to use,
and how do you deploy them?
We'll go pretty deep into that over a series of courses
corresponding to keyword spotting.
Then there's the aspect of personalization.
You got to remember that the machine is simply listening to sounds.
And it's trying to find interesting patterns
in the sounds that happen to resemble a word that might match the keyword.
When we do spectrograms, we will understand this much more deeply
from a visual standpoint.
But in doing that, what you will you realize
is that our sound can sometimes be completely the same
as some sort of background sound.
That sounds crazy, but it's true.
So somehow we've got to be able to teach our neural network model to only
be able to respond to the user's voice and not necessarily any background
noise that might seemingly look like that.
So it's got to be very personalized.
You can push this personalization to quite an extreme.
Specifically, you might want the machine, for instance your phone,
might want to unlock only when it hears you
say, "unlock", not when my wife, or my friends, or my daughter,
comes along and says "unlock".
Well in that particular case, it's got to learn not only the keyword,
it's got to learn to recognize that it's you in addition to actually
being able to recognize the keyword.
You can certainly complement that and do that perhaps through a camera.
So detecting that it's me, has a facial recognition,
and then be able to couple that with the audio signal.
Sure, you can do that.
But I'm talking about getting a device to respond to you just purely based
on audio signals.
We'll talk about that more.
Now, the other thing that you've got to really worry about
is, when you're dealing with data because your audio signal is data.
And if you're going to do this processing up in the cloud,
for instance, if you record the sound and you send it up into the cloud
in order to do the keyword spotting in the cloud.
Because for instance maybe we're not able to it on a tiny device.
Well then you've got to realize that all this data might actually
get compromised when it's being sent to the cloud.
So you want to keep that data secure.
The other aspect is about privacy.
You will be surprised at how little we think about privacy in the real world.
But our data, when something bad happens, we get very upset about it.
And so when data is going up to the cloud,
there's a remote chance that that data privacy could be violated.
And people can use it in ways that you were not expecting.
And so therefore security and privacy is something that we got to protect now.
Susan Kennedy, who you met in course one,
and if you haven't you will meet her again in this particular course,
is going to help us understand what it means
to be secure and private about our data and how
to do things in a responsible way for keyword spotting.
Then, of course, there's the aspect of just overall efficiency.
Remember that tiny ML devices are operating with battery so we got
to do things with extreme efficiency.
So they have to consume very little energy.
And the second thing is that memory on a tiny ML device
is very constrained as well, which means we
have to be extremely good about training our neural network model for keyword
spotting to be resource-constrained or resource-aware.
So the network has to be really small.
So we will understand how to do this soon.
There are many other constraints that come up, countless of them,
but these are the big ones.
So hopefully you can kind of recognize why they are actually important.
So how do companies deal with keyword spotting today?
Because we're going to try and mimic a little bit of that.
And at the same time I'm going to try and teach you
something that's a little different so that you're prepared for what's
coming down the road in the future.
Today the fundamental idea when you open up
the box in say your Google device, or your Google Home device,
or Alexa device, you've got a microphone that's
there that is continuously listening.
I'm sorry.
Yes, the device is continuously listening.
But of course it's only listening for the keyword.
It's not constantly recording the data.
It's listening.
That's it.
So at some point, what happens is that you might say, "Alexa".
And at that critical junction, what ends up
happening is that the neural network will pick up those sounds
and it'll process them.
We'll figure out how to do that.
You will actually go through that process.
It actually processes that data.
And if the keyword is detected, and if the network is confident enough
that it is really a keyword and it gives a strong thumbs up,
then the remaining part of the data.
For instance, when I say, "Alexa, dim the lights",
``Alexa'' is the keyword part.
``Dim the lights'' is a command that usually is offloaded to the cloud.
Because that requires you to do natural language processing because you now
understand what exactly you are trying to do there.
And that's a very simple example.
You can actually have more complex things.
For instance, when you say ``Play my favorite songs'',
that's a rather complicated process.
Because I need to maintain history about what kind of songs you like,
and then go pick some songs and then be able to start playing them.
So that's a very complex process.
But you could think of the fact that you still have to start with "Alexa".
And all of that back-end processing for more complicated tasks
is usually done in the cloud and now, that's how things typically work today.
Now you might be thinking for instance, hey, I've just got a single device.
And sure, I've got this cloud where I've got this big neural network that
can do full on speech recognition.
Well you're only thinking about one microphone
and using the database as a single user.
But in reality, you have to remember that there are thousands of users
out there.
Billions of Alexa words go out to Amazon on a regular day.
That's a lot of people trying to invoke the thing.
Imagine if every single one of these keywords themselves
were actually being processed up in the cloud.
That would pretty much, I think meltdown the internet.
Not quite, but you get the concept.
That's going to be a huge volume of data going up
into the cloud for very small little bits of information.
Now you might be thinking about the end point, which
is the microphone on your device, and the network, that's
the neural network that's got to run and all the computing horsepower that
is up there in the cloud.
Yes.
But there's also a lot of computing infrastructure
in the middle that is going to be boiling up a lot of energy.
And you have to keep that in mind because applications are using up
a significant amount of energy.
And that results in a lot of carbon dioxide production that's really
bad for the environment as a whole.
And this is becoming an important problem to solve.
Green AI is actually emerging as a key topic.
So you don't necessarily want all these devices and their keywords
also going up into the cloud.
Instead what you really want to be able to do,
is to come up with an architecture that allows
us to provide that sort of keyword spotting capability
such a way that it scales billions of users.
Now of course, as I have said this in Course I,
we can go out and build massive data centers,
put them right next to resources, like water,
that can be used to cool and power these data centers,
and build custom processors like the Cloud
TPU that consume 300 watts of power.
But that's not what we want.
Because if we have to scale to billions of users
and really bring the whole planet up with these kinds of AI capabilities,
it's not going to work.
Because we're going to have a serious energy problem.
Why do you think in fact, Google went ahead
to build this custom machine learning processor that consumes so much power.
It's specifically because you need to be able to run this extremely efficiently.
So yes, they ended up building a custom processor.
But if they did not do this, they would be
using so much more power per processor in order
to be able to handle OK Google, for instance.
And what are we trying to get to?
We're trying to get to something that would run in the order of milliwatts
or microwatts with extreme efficiency, which
is mind blowing really complicated.
So we have to somehow figure out how to be
able to run that keyword spotting sort of a task on a tiny ML device.
And that's what we're going to figure out.
Now, the key thing that keyword spotting really enables,
is that by running the keyword spotting application directly at the end point,
we're effectively acting as a net or a filter,
in curbing the amount of data that is actually going up to the cloud, thereby
effectively improving the overall system efficiency of how we provide speech
recognition capability to the planet.
So keyword spotting allows us to wake the machine up and then
provide some complicated services of back end
rather than having even the keywords actually going all the way up
to the network.
So it's a very powerful system level optimization.
And this is one of the reasons why keyword spotting has
become ubiquitous on so many devices.
And of course, as I said earlier in the previous video,
you can do more than keyword spotting.
You can actually expand and do keyword spotting
with the limited vocabulary set.
You can potentially do, "Dim the lights" on the microcontroller.
That's an extension of the very specific,
more than just one word to wake the machine up.
And that's doable.
And we'll get there.
So what I'm trying to get across to you in this constraints and challenges
video, is that you really need to understand the full end-to-end picture
and why keyword spotting has become so important.
There are of course, latency, bandwidth, security, privacy, personalization,
accuracy, memory constraints, battery constraints.
There are all kinds of things that you need
to think about when you're building the model at the end of the day
and you're trying to deploy it.
But you also need to be aware of the full end-to-end picture of how
this particular tiny ML application is deployed in the bigger context.
And hopefully you got a little sampling of that.
We will soon start going into the machine learning pipeline that's
involved to build the application.
And we will be building with some of these constraints
and challenges in mind.
And I'll see you in that video.

\section{Keyword Spotting Application Architecture Overview}

Keyword Spotting Application Architecture
In this reading, we will go over the high-level architecture of the application before we go deep into building the keyword spotting application (KWS).

\begin{figure}
    \includegraphics[width=\textwidth]{TinyML/2 Applications/3-5-5.png}
    
    \caption{Machine learning workflow and pipeline. Begins with collection and preprocessing of data then the design and training of a model. These steps are done in TensorFlow. Then a model is optimized, converted and deployed using TensorFlow Lite. Finally inferences are made using TensorFlow Lite Micro.}
\end{figure}  


There are several steps in the ML workflow and this is what each of the stages entails for KWS.

\subsection{Step 1 Data collection}

For Keyword Spotting we need a dataset aligned to individual words that includes thousands of examples which are representative of real world audio (e.g., including background noise).

\subsection{Step 2: Data Preprocessing}

For efficient inference we need to extract features from the audio signal and classify them using a NN. To do this we convert analog audio signals collected from microphones into digital signals that we then convert into spectrograms which you can think of as images of sounds.

\subsection{Step 3: Model Design}

In order to deploy a model onto our microcontroller we need it to be very small. We explore the tradeoffs of such models and just how small they need to be (hint: it’s tiny)!

\subsection{Step 4: Training}

We will train our model using standard training techniques explored in Course 1 and will add new power ways of analyzing your training results, confusion matrices. You will get to train your own keyword spotting model to recognize your choice of words from our dataset. You will get to explore just how accurate (or not accurate) your final model can be!

\subsection{Step 5: Evaluation}

We will then explore what it means to have an accurate model and why your training/validation/test error may be different from the accuracy experienced by users.

\subsection{Additional Topics:}

We will also consider some additional topics that are unique to keyword spotting: post processing and cascade architectures.

\section{Keyword Spotting Datasets}

VIJAY JANAPA REDDI: All right, here we are
getting started with building our keyword application.
One of the first things we have to do when we want to build any application
is we've got to start with data sets.
So, in this particular video, I'm going to go deep into understanding data
sets.
Specifically, I'm going to talk about how
data sets, which are right at the beginning of the keyword
spotting pipeline, influence how the model will do.
So we need to understand what are the kind of things that our data set really
needs to have in order to be effective for building a TinyML model that
enables keyword spotting or that does keyword spotting.
What does it mean to build a good data set?
Take a look at this particular GIF here.
We've got Alexa, and then the machine comes back.
And I say, ``dim the lights.''
Let's keep it really simple.
Keyword spotting allows you to have Alexa,
and you might also be able to do dim the lights, potentially,
because it's a limited set of words on device,
but let's keep it really simple, the vanilla version of keyword spotting
where it wakes up and listens to a command.
So it's Alexa, just that.
Now first thing you have to think about is who
are the users for this Alexa device.
Is it you and only you?
Is it you and me?
Is it you, me, and my parents who are currently
in India who have a very different way of speaking English
because they tend to have a bit of a stronger accent?
Everybody is capable, perfectly capable, of saying Alexa,
but the crucial difference is how the machine interprets that sound
or hears that audio signal.
We'll look at that in the form of spectrograms
later and see there are subtle differences for the exact same word.
Because the machine is just looking at audio signals,
it needs to be able to generalize, which means that the data set that we're
looking at, the data that we're going to train a neural network model on,
needs to be really diverse.
So we need to understand who the users are in order
to build the proper data set.
So ask yourself, when you're looking at all the different devices
that I ask you to look around and see how
we can empower them with keyword spotting, who
are all the people who will interact?
Maybe it's in your bedroom, and you're the only person that interacts.
Or maybe it's a very general device that's in your living room,
and it's got to be able to understand your parents.
It's got to be able to understand your grandparents.
It's got to be able to understand you.
So you've got to think about who that user set is, and it's
extremely crucial because you've got to have a rich enough data in there to be
able to train that model well.
Now the second question is, what do they need?
What is the particular model that you're trying to build?
What is it trying to process?
What kind of a task are you trying to do?
Because you've got to make sure you capture
the right aspects into that data set.
So, for example, when I earlier said bake at 350,
I used that as an example to say that I want my oven in the kitchen to be
able to listen for the keyword bake.
And then, when it hears the temperature, set the temperature at the right level.
Well, that's a very particular task.
So there I know I need to know that the keyboard is going to be bake.
You've got to think about these keywords.
That might seem really simple, but, later on,
when we get into the nuances of actually doing
the preprocessing on the audio signal, you'll
realize that the keyword you choose is, actually, quite critical
because some words are not easy to pick up by using machine learning models
because it's very hard to tell the difference.
So you might actually have a lot of false positives.
So you need to pick the right keyword.
Now that's crucial because, again, you're building the data set.
So, if you chose a wrong keyword, well, that's
going to get you into serious trouble later down the road.
Then what particular task are you trying to solve?
Of course, in the example of bake, well, I'm
trying to get my oven to actually bake my bread or bake my cake.
Well, in that particular case, you know the particular task.
So you need to clearly define that particular task
because that outlines the requirements.
And I'll show you an example soon.
And how do you interact with the system?
In this particular example, for instance, I'm looking straight at you.
And I'm talking directly at you or, well, the camera and the microphone.
I'm talking directly to the microphone.
What happens if, for instance, the device is actually
located in some corner?
And it's supposed to be listening to me.
So, in that particular case, I better think
about how to actually create the data because the data might, actually,
be a very weak audio signal that it picks up
something that is reflecting off the walls
and then bounces it into that corner there.
Well, that data set needs to look quite different from when I'm actually
directly standing in front of the machine
and giving it a command, which might be the case,
for instance, if you're building a keyword spotting model for your phone.
Maybe you take your phone out, and you say take a picture, right?
Or you say camera, and then it opens up the camera app.
Well, that proximity has implications because that audio signal
sound is actually going to be pretty large there,
as compared to some sort of device that's
sitting in some corner of the room, and it's got to listen,
which means you've got to think about all these things
because all of the stuff has to go into the data set so that your model can
be trained properly.
So I'm hinting at a lot of things here.
And there are a couple of readings that are coming up
that will take you through more systematically
and how you need to ask these questions yourself in order
to be able to get the right data set.
Then, of course, there are aspects of the real world
that make it really hard.
Specifically, you're not in some sort of a quiet room that's
pristine at all times.
A device might be-- in my office, right now, it's nice and quiet,
but I have kids.
They can get really loud.
Well, when that happens, for instance, all hell breaks loose.
Point that I'm trying to make is that there can be background noises,
for instance, that will happen.
So your data set needs to be having all of these kinds of background noise
and like real-world characteristics so that, when
you deploy the product, the sounds that you're listening for,
the model has learned to be a little resilient against all
of these kinds of issues, such as background noise.
Now, all of this said, let's take one specific example, speech commands.
And let's go deep into it to try and understand
how those characteristics that I mentioned earlier
have been addressed through one example by Pete
who you will see in the next course.
And, in fact, he's one of the great guys, and I love him.
And he's actually inspirational.
And it's one of the reasons that I'm actually here because both of us
tagged up to say let's teach TinyML to the world.
So, getting back on point, Pete actually put together
this data set a little while back.
As you can see the date, it's in 2018.
Wow, he's one of the guys who was thinking
about this kind of a thing for a while and realized that, back then, there was
no limited vocabulary speech data set.
There was no notion of keyword spotting data set.
There was nothing there.
So he thought hard, talked to many people,
and tried to compose a data set that tries to address some of the issues
that I mentioned to you earlier.
Let's go through and see what's inside this data
set because you will be using this data set in doing your Colabs.
The speech commands data set, basically, has a whole bunch of individual words.
It's not sentences.
It's not dim the lights.
It's something just like a single word.
So it's composed of a whole collection of words,
which I'll show you in a second.
And it's got many, many recordings for each word.
Think about why do you need that.
Well, it's for that exact reason I was saying.
Like there are many people who might be interacting with the device.
So you've got to capture the fact that people are diverse.
They speak differently.
They see things differently.
They emphasize things differently.
So, by having a broad distribution of examples in your data set,
you are, effectively, going to try and get the model
to be able to recognize all of them.
And how did he build this?
Well, you need a lot of data.
If you need a lot of data, well, you've got to enlist people.
Typically, companies will pay labellers, as they say.
They'll hire an organization that actually does the labelling.
Later, when we do data set engineering as another module
in this particular course, we'll talk about some
of the issues that will come about this, but,
for now, what Pete did was that he actually got the community
to contribute those example words.
Those 1,000 and 4,000 examples of each word,
well, the community contributed it, went to a website,
recorded the audio in for the given word.
And one of the key things is that the data set actually
includes background noise in order to kind of be
reflective of real-world situations.
Again, I said, for instance, this video, my family is really quiet because they
know I'm actually doing a recording.
But, in reality, it's like a little zoo down here
with my kids screaming in the background at all times.
Well, point is that you need to have some
of that sort of background, realistic audio in there
so that we can train the network.
And the speech commands data set actually
has that sort of background noise included in there.
And so the data set is actually limited.
Remember, I said keyword spotting is all about a limited set of keywords.
Now, in the particular case of speech commands,
it's got 25 different IoT kind of keywords.
I'll show you what those are.
And it's got 10 different unknown words.
There's a little bit more description about this
in the readings that will follow.
Well, what are the data sets actually looking like?
Well, here's a distribution of words that are there in the data set.
On the left side, you're seeing all the different words that are there.
And the bars are, basically, showing the number of examples
that are there in the data set.
It's not a perfect data set, but it's a good enough data set.
It's not perfect in the sense that some of the words
don't have that many samples.
For example, the words at the top, like for bird and so forth,
there aren't that many samples, as compared to yes and no.
Now this is not by mistake.
So, often, when you think about the particular use case,
you will tend to emphasize words that are more relevant to you.
And that's what's actually going on here for instance.
Here you can see that some of these words
are actually the more often used words.
So think of it like, for instance, if I have a robot where
I want to say yes, no, left, right, or I have some sort of machine where
I want to say yes or no to.
Well, these are quite pronounced in the data set.
And, in fact, these are the words that we'll actually
be using because there's a large sample of them.
And, by having a large sample, that data is actually quite diverse.
Now I'm just emphasizing this as one example data sets,
typically, have these characteristics, but you need
to understand what that data set is.
And, hopefully, one of the things that you're getting out
of this particular video right now is you
understand that, as you look in TinyML, you need to have specific data sets.
Now, typically, for instance, if you are familiar with machine learning
or very familiar with machine learning, you
have probably used ImageNet or MS COCO or KITTI or any of those big data sets.
Well, you're talking about TinyML here.
So we're not talking about having data that's, broadly, able
to like look at a very large range.
So like ImageNet, which is a humongous, beautiful database of images
for image classification, has over 14 million different images
that are, broadly, classified over a thousand different classes.
Well, in TinyML, you are much more focused.
You know what your application area is, and then
you try and construct the data set accordingly.
That is a crucial difference in data set engineering for TinyML.
So keep that in mind.
And, therefore, when you look at speech commands, that's what motivated.
There are other speech data sets that are out there,
but they are all around sentences.
That's not what you need for TinyML and keyword spotting.
You need a data set that, specifically, focuses on commands or keywords.
Therefore, that inspired Pete to actually put the speech commands
data set together, trying to address some of the questions
that we raised earlier.
Now some of the things that you've got to think about, generally, in terms
of data sets is about quality control.
You've got to keep in mind that sounds are imperfect
just because of the way sounds kind of propagate through.
For instance, right now, when I'm talking to you,
this sound is actually propagating all the way
to the microphone, which is getting encoded
into some sort of compressed format.
Then that's getting pushed over to the HDMI cable
onto my computer, which is then picking it up and then doing a decode on that,
decompressing the signal.
Then it's passing through filters in order to clean the signal a little bit.
Then it's recording it, and, voila, you hear me.
So what I'm trying to say is that there are lots of complicated processes
that are involved.
So sound can be quite delicate.
So, when you're looking at recording data,
you've got to think about is that sound clip that you're
recording quite representative of the real audio signal or not.
And you've also got to capture the noisy backgrounds
aspect that I mentioned to you earlier.
It's very, very important.
We'll actually do this in the Colab later on.
And, with that said, I'm going to stop here and let
you think for a moment about why we focused on data sets.
Ask yourself, what did you learn in this video?
Data sets are extremely crucial and extremely fundamental.
You get them wrong, everything you do down the application pipeline
is a waste of effort.
So think carefully.
And, specifically, think about what is unique about building TinyML data sets,
as I hinted earlier in this particular video.
With that said, I'll see you in the next video.
And, hopefully, we can push forward in terms of building the models.


\section{Wake Words Dataset Creation}

As we just learned, there is a lot of work that goes into designing a good dataset for keyword spotting. And this is true for all applications of TinyML. This is because the requirements for different applications (and thus the data that needs to be collected) can vary drastically even for the same kinds of inputs. In fact, even for the same keyword deployed into different environments!

This makes it impossible to come up with hard and fast rules for how to design a good dataset. But/and here are a couple of things you might want to consider when designing an audio data collection scheme.

\subsection{Who are the anticipated end users?}

It is important to understand who the end users might be as they may engage with the application differently and may present different kinds of challenges for how you need to design the application and what data you will need to collect. For example:

\begin{itemize}
  \item What languages will they speak? 
  \item What accents will they have? 
  \item Will they use slang or formal diction?
  \item Will the users speak clearly?
\end{itemize}

\subsection{In what environments will the users employ the application?}

Imagine a Keyword Spotting device. Is it placed in a quiet room where it can hear you clearly? Or is it placed right next to your TV in the living room? The environment can greatly affect the quality and type of audio that can be collected, as well as the amount of additional sounds that must be accounted for. For example:

\begin{itemize}
  \item How much background noise do we expect?
  \item Who/how many people will be talking in the environment?
  \item How far will the users be from the device and sources of noise?
  \item Will the user be stressed vs. calm vs. panicked?
  \item Will the user use different volumes of voice (whispered/normal/loud/shouted)?
  \item How likely is it that these keywords may be triggered unintentionally during conversation?
\end{itemize}

What are the goals for using the application? How will this impact the requirements of the ML model's performance?

\bigskip

Depending on the goals of the application, model errors can either have very little consequences, or they can be catastrophic. This can directly impact how robust of a data collection scheme you need to implement. For example consider the difference between these three scenarios:

\begin{itemize}
  \item The user is trying to turn the thermostat up/down
  \item The user is trying to arm/disarm a security system
  \item The user is trying to play their favorite playlist from a smart speaker
\end{itemize}

Given all of these previously mentioned factors here a short checklist to consider when implementing your final data collection scheme:

\begin{itemize}
    \item What custom keywords will you select?
    \item What background noise samples do you need to collect to augment your dataset so that the trained model is robust?
    \item What other words do you need to collect to ensure the model learns the difference between them and your keywords? It is important to collect other words so that the model can learn your particular word and not simply the general sounds of humans talking!
    \item How much data will you need to collect? We all know more is better (usually), but you also live in the real world with time constraints so how much do you think will be enough?
    \item In what environments will you collect these samples?
\end{itemize}


\section{Keyword Spotting Data Collection / Pre-Processing}

VIJAY REDDI: In this video, we're going to focus on keywood spotting data
pre-processing, which is part of the data collection kind of pipeline,
but it's also useful when we're actually training models.
So think about what happens, for instance, when we say, Okay Google.
The first thing that happens is that you have your microphone, which
is a sensor, which is picking up that audio signal, which is actually
an analog signal and converting that into a digital signal that can then
be processed effectively for your neural network.
So we will be understanding the sequence that we're going through
and what are all the nuances you have to watch out for.
So this is going to be very much at the beginning of this pipeline
that we will be walking through.
We've already talked about collecting the data
and how to build a good rich data set.
Well, that's the sounds.
But now, before you can actually get around to training the model,
you actually have to do preprocessing on that data, which
we learn in this particular video, and then use
that to actually train the network on.
Let's step through and understand what exactly
is going on when we are collecting an audio signal
from a sensor like the microphone.
All we hear, really, are just sound waves that are vibrating.
And your microphone is effectively picking up those vibrations
and converting them into an electrical signal, which
is in the form of voltage distortions you're going to have.
And that is going to be sampled at a certain rate,
because you're effectively picking up signals periodically.
With the microphone for instance, you will
see that you will be picking up an audio signal at around 16 kilohertz, which
means 16,000 samples are getting picked up continuously.
And that data of 16,000 samples per second and somehow
has got to be fed into the network.
So this leads into an interesting set of questions.
Here, while I am focusing very much on the audio signal,
I want you to keep in mind that what I'm trying to really teach you
is that when we have a sensor, we have to think about what it's data
input looks like and what its data feeding into the network
needs to look like.
So for the microphone, we have 16 kilohertz signal that's coming.
That means we have 16,000 data points.
First question-- do we need all of that data?
Something to think about.
What level of sampling is actually good something that we should keep in mind?
Do we need to feed all of that data into the network?
Or can we try and extract some critical features?
That's where creative thinking comes in.
We need to understand what aspects of the signal
that we really want to pick up.
Let's take this example here of no said loudly.
And what we're looking at is a time domain signal.
x-axis is effectively just time and then y-axis is just
the amplitude of that signal.
What's interesting here is that you have a wide signal,
and the important part the note is somewhere in the middle of the signal.
Everything is a signal.
So the question is, how do we first and foremost figure out where that no is?
So somehow, we actually have to identify the no.
Doesn't matter whether it's said loudly or quietly--
you have to identify that crucial region.
The other thing here is that it's the same word, no,
said loudly versus said quietly.
And yet what you were seeing here is that by just looking at the raw
waveforms that are coming in, we cannot tell the difference.
The left hand and the right hand look completely different,
even though it's the exact same word being said.
If I show you this figure, do you know what these words are?
Of course, no, you can't.
You can't possibly know.
Well, but can you tell us something interesting here?
Well, the bottom two are effectively the no from the previous slide.
The top ones are yes.
And again, I'm trying to make the same critical point here
that the same word can look totally different in its raw format.
This is why preprocessing the signal is actually extremely critical.
Because we want to extract the salient features.
Because the two no's and the two yeses--
there is some hidden information in there that clearly corresponds to yes
and no, but looking at these raw waveforms,
we cannot tell that the left hand no and the right hand no and the left hand yes
and the right hand yes are actually exactly the same.
So how do we go about distilling this?
Some of the interesting challenges that we're going to get into
are really, that when you have this basic waveform that's coming out,
you got to first and foremost figure out when does the word really start.
That's called alignment.
Keep that in mind.
So you've got to align or extract the critical part of the signal that
actually corresponds to us in order to figure out whether it's
a yes or no or any other word.
That's the first step you're going to do.
We need to figure out how to actually extract the aligned part of the signal,
extract the critical features that are in there.
So let's step through those things next.
Specifically, how can we differentiate the uniqueness in the signals?
If we go back and we look at this figure,
that 16 kilohertz analog signal is getting a digital representation,
which we're looking at in the kind of waveforms that we're seeing earlier.
Now what we want to be able to do is to pull out the interesting tidbits
between these two variants.
So if you look at any signal, it's fundamentally a composition
of other primitive signals.
So the signal on the right, for instance,
that we're looking at that has an interesting shape,
is really a composition for instance of a signal that's on top and the bottom.
Now we totally just drew this, but do you get the idea.
Is that a complex signal can be deconstructed
into a set of different signals that each have a critical frequency.
So the frequency on the top and the frequency on the bottom one
are completely different.
And the resulting combination of those signals
also looks different, which is why you have the right hand waveform.
So what I'm hinting at is if we are looking at a keyword like this,
what does the signal decomposition look like?
What are the frequencies, the fundamental frequencies,
that are all sitting inside the signal?
If we can extract that, then perhaps regardless of how loud or how quiet
it is we can actually get to the sense of the word.
That's where Fourier transforms come in.
Fourier transform is a very simple process
that tries to extract a time domain signal into a frequency domains.
So it does this conversion so you're trying
to look at what are the frequencies that are present in the signal.
That's the first step.
Why do we do this?
Because by extracting out the critical frequencies,
it's like trying to find the fingerprint, effectively.
It's one of the steps, but it's a crucial step.
If you're not familiar with Fourier transforms--
and FFT stands for how you do a fast Fourier transform--
don't worry about it, I have a reading that will actually give you pointers
and explains the exact process that's involved.
Because if I have to teach you FFT, we'll
have to start a whole new course and just signal processing.
But let's stay focused on the critical building
blocks that correspond to building a keyword spotting ML application.
So the first step is really extracting the frequencies,
which FFT allows us to do.
Now one of the key things you'll realize is when you say,
no, it's not a snapshot.
It is actually a set of frequencies that are emitting over time.
No.
That takes maybe half a second--
I don't know, something like that.
When I say, yes, that's a set of frequencies
that are changing over time.
And that is the pattern that we want to try and extract or essentially capture,
because that pattern is a unique fingerprint.
And how do we do that?
Well, this is where we have to do the FFT, which we're extracting
those critical frequencies out.
Well, that, we're going to have to do that over the entire sound.
So once we align or chop out the critical bit of no in here,
we start at its beginning, and we have what's
called a sliding window, where we're doing that FFT to extract out
the frequencies along that time domain.
And in doing that, we look something like this,
where we're overlapping a window, as it's kind of sliding every time.
It's about 25 to 30 milliseconds.
That's typically the window size.
The reading explains why it's about that size.
It's effectively a very good measure for capturing the critical features,
and to be able to see the transience that are happening over time
and the frequency response to it.
So if we kind of step through this, what we can do
is we can generate interesting plot that actually shows us something about what
frequencies are pronounced and how they're
changing over time in their strength on the right hand
side, which is called a spectrogram.
A spectrogram is effectively a visual representation of the signal,
and all of the different frequencies that are involved,
and how it changed over time.
That's all there is.
It's a beautiful picture.
Now if I go back and look at these original figures that we're looking at,
these audio signals, and I convert them into spectrograms so beautiful,
what you will see is that the top and the bottom ones look quite distinct.
Specifically, if we look at the top yes ones,
then you see that they have a distinctive signature that stands out
from the ones that are below for the no's.
The no's don't have any strong amplitude of visibility there.
The strength of that color or the vibrancy of that color
is telling you how strongly those frequencies are actually
pronounced there.
And you can clearly see a distinction here,
which is then allowing us to sort of get a clue that these are not
the same words.
By looking at the signal, we're able to say that,
well, yes and no look quite different.
Moreover, if you look at the horizontal side
where you have no loud and no that is said quietly, look at them.
They look similar.
Even though they're actually in the raw form, they were different.
That's the beauty about converting an audio signal into a frequency domain,
and then looking at it visually through a spectrogram.
Can we do even better?
This is often typical in machine learning.
You're always iterating on the data to try and extract more salient features.
You're doing feature engineering.
The answer is, yes, of course.
Well, a little bit of a machine running.
Well, it's simply about trying to mimic how our ear is actually hearing sound.
So our ear does not hear all the sounds equally
with the same frequencies-- don't respond well in the ears.
The lower band frequencies are actually much more crisper to us.
We can tell apart the low frequencies far more easily
than the high frequency ranges.
And so we want to effectively use that thing,
and we call that process a Mel filterbanks.
Mel filterbanks are basically allowing us to understand
how the human perception is of sound.
And this is something that was developed a very long time back.
We will be using these Mel filterbanks to effectively bend
the frequencies to extract out a much stronger signal in the spectrogram.
So again, in the readings, I explain what are these MFCCs
that we'll be talking about.
We also explain what a spectrogram is.
But as you can see on the left hand side is the original spectrogram,
the right hand side is what's called a MFCC, which
has been processed effectively.
And what you can clearly see are that for the yes and the no,
they look different.
And by extracting these different frequencies
that we're seeing over time, what the model gets a chance to look at
is how the things are changing over time on the x-axis.
And it's that sort of a pattern that the network
is going to try and look for or learn to look for in identifying
the keywords, yes and no.
We can certainly do a bit more of engineering.
For instance, we can clean the signal.
Sometimes, you want to get rid of that quiet yes and the loud yes,
those amplitude differences.
Sometimes, you can normalize the signal in order
to make sure that the loudness and the quietness
is actually kind of ripped apart, so that way, you can really
focus on a normalized clean signal, and the amplitudes are all
the same for all of the sentences or the words that are inside your data set.
By doing that, it actually helps your neural network learn.
Now this is a very typical part of any machine learning pipeline.
It's not specific to keyword spotting, so that's
why I'm not really poking at it too much.
But this is something that we actually do in our particular pipeline
as well as we prep the audio signal in order to feed it into the network.
So in this video, what I want you to take away
is that we preprocess signals that are coming out of the sensors.
We do not pick the raw data and feed it straight into the model--
no, no.
Instead, what we do is we preprocess that data.
And by doing the preprocessing, we are looking for critical features
that we want to extract out so that it's easier for our models
to be able to pick those features up or learn to pick those features up
as we train the network.
That's the critical message that I'm trying to convey to you here.
Moreover, what's interesting is that when you do the preprocessing,
we are not looking at the signal in the original format.
We're actually completely converting it on a different representation.
We went from an audio time squeeze signal into a frequency signal.
And we're looking at the frequency as a picture.
Think about that.
That's important.
So when we're dealing with sensors, think of useful conversions
that we would go through.
We will give you a couple of examples to bootstrap you
on the kinds of different preprocessing steps that are done.
We will actually be using them in order to train our models, as well.
So you will get quite familiar with this sort of a process.
Generally, though, whatever we are touching in the pipeline
is actually going to be useful for any audio sensor that you take.
So I've tried to structure it so that you get the fundamentals that can
be repurposed for this type of sensor.
And with this, we'll continue in the next video.

\section{Spectrograms and MFCCs}

In the last video you explored how audio is preprocessed for Keyword Spotting.

You saw that the raw analog audio signal is first transformed into a digital signal.

\begin{figure}
    \includegraphics[width=\textwidth]{TinyML/2 Applications/3-5-9-1.png}
    
    \caption{A taw smooth analog signal is converted to a chunked up digital figure. The orginal (16kHz) signal is shown to be wave-like in visual appearance and the analog to digital conversion produces bars that align to the contours of the original wave.}
\end{figure}  


You then saw how the digital signal in the time domain (time on the x-axis and amplitude on the y-axis) is transformed into the frequency domain (frequency on the x-axis and amplitude on the y-axis) through the use of a fourier transform. If you’d like to dive deeper into understanding how the fourier transform works we suggest you watch this fantastic video on the 3Brown1Blue youtube channel.

Of course for the purposes of TinyML we need to worry about how we can compute this transformation efficiently -- both in the memory needed to store the program, and in the memory and latency needed for the computation. The solution is the use of an approximate fast fourier transform (FFT) of which there are multiple possible implementations. TensorflowLite Micro uses the \PYTHON{KISS\_FFT} library, which includes more documentation about the design choices made by the library to be compact in memory on their GitHub page.

By applying the FFT through a sliding window procedure we can analyze the intensity (in color) of each of frequencies over time. This is a spectrogram.

\begin{figure}
    \includegraphics[width=\textwidth]{TinyML/2 Applications/3-5-9-2.png}
    
    \caption{A visual depiction of the sliding window of an FFT and the resulting spectogram. On the left there is a time-domain digital signal (time on the x-axis vs. amplitude on the y-axis) overlayed with many red boxes showing the slices of sliding windows that are used to construct the spectogram. On the right is the spectogram which is shown to again have time along the x-axis but now frequency on the y-axis and then the color is the amplitude. It now looks much more like a paint splatter as compared to custered bars.}
\end{figure}  



A spectrogram is a better input for a deep learning system than the standard input audio signal as we have already extracted the frequencies for the neural network and so it doesn’t need to learn how to do that on its own. Even visually, it is easier to see the difference between the ``yes'' and ``no'' spectrograms.

\begin{figure}
    \includegraphics[width=\textwidth]{TinyML/2 Applications/3-5-9-3.png}
    
    \caption{4 spectrograms are shown for ``yes'' and ``no'' both for a loud and quiet recording. All four spectograms have the same area highlighted visually. In both yes spectrograms data exists in the highlighted area while both no spectograms do not have data. Showing how the spectrogram can islate the words well.}
\end{figure}  



We then discussed how the human ear can tell low frequency signals apart easier than high frequency signals according to the Mel Scale. We can therefore create a Mel Filter Bank to take advantage of this phenomenon and group together larger clusters of higher frequency signals to provide more useful data to our machine learning algorithm. The output of this process is a Mel Frequency Cepstral Coefficient (MFCC). For more information on how this works here are three great posts: \href{http://practicalcryptography.com/miscellaneous/machine-learning/guide-mel-frequency-cepstral-coefficients-mfccs/}{post1}, \href{https://jonathan-hui.medium.com/speech-recognition-feature-extraction-mfcc-plp-5455f5a69dd9}{post2}, \href{https://jonathan-hui.medium.com/speech-recognition-phonetics-d761ea1710c0}{post3}.

\begin{figure}
    \includegraphics[width=\textwidth]{TinyML/2 Applications/3-5-9-4.png}
    
    \caption{The ``yes'' and ``no'' loud spectograms are shown again as well as the corresponding yes and no loud MFCCs. Visually one can see how the MFCCs are much sharper in their color gradients and contrasts. This is because the MFCC clustering makes the differences more stark.}
\end{figure}  



Of course there are many different ways you could solve this learning problem and it is likely that very large neural networks would perform quite well on this task without this preprocessing. However, for the case of TinyML this preprocessing is vital and drastically improves the final model performance without taking up a lot of memory or latency! 

In Course 2, we will use these methods via Google Colab but in Course 3 you will learn how to deploy these pre-processing methods directly onto the embedded hardware platform.

\section{Spectrograms and MFCCs Colab}

Now that you understand how the audio preprocessing pipeline works, in this Colab you’ll get to explore it for yourself. Record your own audio files and explore how each step in the processing pipeline changes your audio signal and reveals features!

This Colab, along with the others in this section of the course, will have a lot of code in it. We have tried to provide comments throughout to explain the code. At the same time, since machine learning is such a growing and changing field, a core skill for ML practitioners is to learn how to read documentation and leverage new libraries. Therefore we strongly suggest you google the various libraries used if you’d like to learn more about some of the functions we call! For example, at the end of this Colab we use the Librosa library to compute the MFCCs which has some great documentation which you can find here.

\url{https://colab.research.google.com/github/tinyMLx/colabs/blob/master/3-5-10-SpectrogramsMFCCs.ipynb}

\section{A Keyword Spotting Model}

VIJAY JANAPA REDDI: All right.
Now we're getting to the part of actually training the model itself.
Remember that we first came from collecting the data set
and learning about speech commands, then we
learned about how to pre-process the input data that we have,
which is in the form of raw audio files, which are then getting converted
into a representation that is actually useful for our neural network
to be able to ingest.
Now we're going to focus on designing the model,
and we're going to think about how we should design the model for TinyML
specifically.
Now, because we're dealing with an image it should not come to you as a surprise
that we might just want to use a convolutional neural network
at this point.
You've used convolutional neural networks previously
in order to classify things such as the MNIST digits, or the fashion data set,
or in terms of telling the horses versus humans.
You've used it a few different ways.
Well, can we just reuse it?
Well, what do you think?
We've got an image, which we create into a big tensor
that we then feed into the network.
And the network is effectively using a series of filters via convolutions
to extract out interesting features, and then
classify those features to get some interesting outputs.
That's all we're doing.
If that's the case, if we're just dealing with an image,
why don't we just replace that image with our spectrograms?
Voila, spectrogram has interesting features.
Perhaps our network can actually piggyback on the exact same setup
that we have, and then be able to extract interesting features out of it
in order to do the classification.
Then we might be able to classify as whether the word, which came in
from the microphone, is an analog signal that
got converted into a digital signal, which then got
processed in through FFTs and then came through here
as a spectrogram is a yes or no, or pick a favorite word.
Well, can we reuse that old model that we touched on?
Yes, of course.
Yes, we can.
Well, what has this network got?
The first part and the second part are what I'm going to focus on.
The first part, the five convolutional layers plus pooling layers,
are all designed to extract features out of the image, right?
And then we're using a dense layer in order
to help us do classification at the bottom.
So the two major portions.
Can I use this on the spectrogram?
What do you think?
Of course, the answer is yes.
Well, all right.
Before we use it, hmm, what are some of the constraints for a TinyML
application?
Remember, I've been beaten into you that it's extremely crucial for us
to be cognizant about the end-to-end system when you're doing TinyML.
TinyML needs to work on a real device.
So, if we look at a real device, we only got 256 kilobytes of RAM.
And we only have one megabyte of flash memory or storage on the device.
Well, that might prove to be a little bit of a problem for our network.
Why?
Because if we look at the model and we're
converting a TensorFlow model that we already saved into a TFLite, model--
ding, ding, ding.
I hope you remember why we learned about TF versus TFLite.
We keep reusing it now.
Any time you want to try and do something for deployment,
we're going to have to deal with really simple, clean inference engines that
are efficient.
So therefore, we're going to be using it.
This is why we kept talking to you about TensorFlow versus TensorFlow Lite.
OK?
So we do the conversion, which simplifies the whole model down
to its leanest and meanest form.
And then we print the size.
Voila.
It's about 6,800 kilobytes.
You think we have a problem?
The answer is yes.
We only have about 256 kilobytes of RAM, and we only
have about one megabyte, or 1,024 kilobytes, of flash for storage.
This model won't even physically sit where you can't even
loaded onto the device for saving it.
Ah, don't you worry, I know how to quantize the network,
is probably what you're thinking.
Very well, sure.
Pick your favorite quantization approach.
Post-training quantization of our training.
At this point, you're a ninja at these things.
Well, let's say we do that quantization process,
and then we end up with a model that's 1,700 kilobytes.
Hmm.
Does it fit?
Well, no, because it's 1,700 kilobytes.
We only have one megabyte of flash storage,
and at best we have 256 kilobytes of working RAM during execution.
So we've got a bit of a problem here because our network
is too big even though it's really simple, as we
kept saying in the beginning.
So can we do something different here?
We need something smaller.
The key thing that I want you to take away in this particular video
is realize that when you're designing tiny ML applications,
they're very application specific.
They are really domain-specific machine learning tasks.
They just need to perform one task, and they just need to do it well.
That's typically the critical key, or the critical secret sauce
for actually enabling models that are really small
that can run efficiently on embedded microcontrollers.
So there comes the TinyConv model.
It's still a convolutional model because we're still
going to use the spectrograms of the original it got.
It just happens to be a really tiny model.
Look, you can come up with many different kinds of models.
I'm just going to toss one out here for you.
It's a very simple network in that it's got
a one conv networks, one conv layer, followed by a dense layer
for doing the classification.
It gives you a probability of what the word is, a yes or a no,
and is it some sort of unknown word.
That's all it does.
Really simple.
You take the input data, which, again, is the pre-processed input data that we
looked at in the previous videos.
And then it goes through the convolutional operations, and then
the fully connected layers, and out comes your output classification.
It's simple.
It's simple by design.
This is not an unfortunate thing.
Or it's by design simple because all you're doing
is trying to detect very simple words.
You not doing a full-on speech recognition then
these very complicated, deep networks.
It's about simplicity, and that's the beauty about this thing.
Now, certainly you can improve this network quite a bit.
We'll get there at some point.
But for now, I just want to give you a baseline to start off with.
If we look at this network, it's only got about 16,652 parameters,
which is very, very small.
One of the things, by the way, I want you to take away
on this slide is learn how to calculate the number of parameters that
are involved because those parameters or weights are effectively
the key sense of the network, which lead to the size of the network.
In the future, you're going to need to be able to calculate these things.
So here's a very simple calculation, so that's
what I'm showing you the numbers.
So if you look at the amount of parameter model size that it leads to,
well, it's about 70 kilobytes.
That's pretty small.
That's pretty good.
If I quantize it down, well, that's going to be only about,
like, 17 kilobytes.
So that's very exciting.
So, that said, can I fit it on the network, onto my device?
The answer is absolutely yes, and that's the exciting part about it.
What's nice about this TinyConv model is not only that it's actually small,
but it also leads to really fast inference.
It's only going to be fast because we have few number of parameters
in the network, and we only have like very simple convolutional layers
followed by a dense layer.
So there are not too many multiply accumulates that you have to do.
So in addition to thinking about the memory size being small,
you also got to keep in mind that your microcontroller is
a very simple processor because it's going
to be low power, low cost, and extremely efficient.
It has to be simple, which means your network calculations also
need to be simple, because you want to keep that inference fast.
So the key takeaway here is that I've given you
a baseline network that will allow you to do some sort of classifications
on keywords.
The crucial reason that it can be so small
is because it's highly specialized.
So that's the thing I really want you to take away.
I don't care if you don't remember the network specifics.
What I really care about is if you remember
that the network can be small because you're performing highly specialized
tasks.
So again, this is where you really want to understand
what is the specific task you're trying to perform,
and try and specialize in networks for just that task,
because that's what allows them to be small and effectively deployable
on microcontroller kind of devices.
And that's the essence of designing a model for keyword spotting on TinyML.


\section{Keyword Spotting in Colab}

Now that you have seen how the Keyword Spotting application works we thought it might be fun to explore how well our default pre-trained model works! This will also start to get you familiar with the (relatively large amount of) code we are going to use to train a custom model later as that Colab will build directly on top of this one. And for fun we’ve added in a section where you can upload your own audio samples to see if the model can figure out if you are saying “yes” or “no” or something else “unknown.”

As mentioned in the last Colab, this Colab, along with the others in this section of the course, will have a lot of code in it. We have tried to provide comments throughout to explain the code. At the same time, since machine learning is such a growing and changing field, a core skill for ML practitioners is to learn how to read documentation and leverage new libraries. Therefore we strongly suggest you google the various libraries used if you’d like to learn more about some of the functions we call! For example in this Colab we use the pickle library to load in a python object from a web file that we download through the use of wget. We also make extensive use of the TensorFlow Micro example code which can be found here.

\url{https://colab.research.google.com/github/tinyMLx/colabs/blob/master/3-5-13-PretrainedModel.ipynb}

\section{Forum}

What was your experience like with the pre-trained model? Did its performance surprise you? Did it work well for you when you provided your own audio? Why do you think that was the case?

\section{Intro to Training in Colab}

VIJAY JANAPA REDDI: All right.
We're one step closer now to training our model.
So in this particular video, I'm going to touch
on some of the critical aspects that we just
need to keep in mind when we get around to doing hands-on training
on the Google Colab environment.
Specifically, I want to emphasize that we will be going through all the bits
and pieces that we talked about so far.
We talked about pre-processing.
We talked about the model design.
Then we talked about the data set itself.
And then we're going to learn a little bit
about how to evaluate the model that we're
getting, all in the Colab perspective but with the resource constraints
in mind.
We don't necessarily need the microcontroller right now.
We know what it's going to do.
We know what the model is going to be like in order
to be able to deploy it on the microcontroller
because we know the static requirements of the model, for instance.
So we can still do everything we need to do before we physically deploy
the model, which we will do in Course 3.
But deployment has got its own set of challenges and interesting things
you can learn about.
For now, I want to focus you very much on the model aspects.
So we're going to get around to actually training our model.
At this point, we're very much still in the training pipe
so we're going to be using TensorFlow.
We're going to be leveraging or reusing a lot of the knowledge
that you have already picked up in Course 1,
or hopefully you're already familiar with the notions of gradient descent,
classification, and all these kinds of things
if you're coming straight into Course 2.
Specifically, remember in the previous course
we were computing a classification exercise.
That's exactly what we're doing.
I already mentioned this in the previous video
that, instead of using a particular image,
we're just going to swap it with a spectrogram image.
We talked about cats versus dogs and horses versus humans.
Now it's going to be spectrograms.
That's all.
We're going to try and classifier keywords versus other words
that are there that we're hearing.
And we can continue to use all the methods that we touched on previously.
Specifically, we're going to use gradient descent over our training
data.
Hopefully, you remember what gradient descent is
and what the loss function does.
We're going to reuse all of that knowledge that you've built up.
Now, one thing I want you to keep in mind
is that we need big data for training our model.
So for example, on this side what you're looking at
is the training run that Brian did.
Brian is our head TA for this course.
And what he's plotted down here on the x-axis--
in the y-axis is effectively showing you what
the accuracy is over the number of epochs that are going by.
And so what you're seeing is that the network learns really fast right
at the beginning, but then getting that last bit of improvement
gets really, really hard.
And you can see that pretty much the network gets to around 80%
accuracy by around maybe 3,000 steps or so.
And after that, it's taking forever to actually slowly improve the accuracy.
But every bit of accuracy really matters.
I'll give you a hint.
Typically, state-of-the-art accuracies are in the order of around 93% to 94%
right now in terms of precisely being able to say whether it was a keyword
that matches what the machine is supposed to respond to or not.
This model is by no means perfect.
That said, everything we're doing here is still useful for building some
of those state-of-the-art models.
Now, we're going to be training over a couple of thousand samples.
Now, that data set might seem big, but it's actually really small.
Just keep that in mind.
One of the reasons I'm harping about this data
set and the length of training time is because,
while it's going to take a little while to actually train that network even
though it is quite small, so you might want to take a little break.
Go get yourself a little cup of coffee is what I'm trying to say here.
On a serious note, it will take a little while.
And especially as we do some of the other colabs that
are involving vision processing, it's going to take much, much longer.
We've got some techniques to cut it down.
But the fundamental point is this.
Don't sit around and wait in front of the computer,
because it's not going to be as simple as the previous networks
that we've looked at.
But when you do get around to actually doing the inference,
it's going to be really fast, so you'll be able to enjoy
and it's going to be worth the effort of training the model.
Personally, I think it's very exciting the first time you actually
experience dealing with something that is sound
and where you've done all this complex transformations to try and get
the machine to learn an interesting signal that
is encoded in a visual representation.
So I hope you have a little bit of fun, and feel
free to share your thoughts in the forum or post some discussion topics.

\section{Training in Colab}

\subsection{Training Keyword Spotting in Colab}

The code for training your model in Colab is very similar to the code you used to test out the pre-trained model. However, we’d like to point out a couple of things that will hopefully make training in Colab easier (and more intuitive).

\subsection{Timeouts}

After 12 hours your Colab instance will automatically recycle so make sure to download anything you want to keep before then. Also after about 90 minutes of no activity Colab may also recycle the instance as it may deem it “stale.” Therefore, during training it is best practice to check in every hour and click around on the Colab to make sure it knows you are still there! But/and beyond that you can absolutely minimize the window and continue with other work.

\subsection{File Structure}

The training process will create three directories of note.

\begin{itemize}
  \item /data: will hold all of the data from the Speech Commands dataset. If you want to take a look at particular audio files used for training this is where you can inspect and download them. We will revisit this in Course 3 when you design your own dataset as you’ll need to ensure that it matches this format / file structure.

  \item /train: will hold all of the model checkpoint files. These files describe the state of the model after X steps of training (usually denoted in the file name and the training script auto-saves them every 100 or so steps). If you e.g., trained for 100 steps and wanted to train for another 100 more, you can use the latest checkpoint file at step 100 to start from there instead of re-training from scratch for 200 steps.

  \item /model: will hold all of the final processed model files. There you can find both the frozen Tensorflow model, the TensorflowLite model, and in course 3 you will also find the TensorFlowLite Micro model there. If you’d like to download and save your trained model to compare against another trained model at another time you should make sure to download it from there!
\end{itemize}


\subsection{The Training Script}

The training script is designed to automate a bunch of the training process. It first sets up the optimizer using a \PYTHON{sparse\_softmax\_cross\_entropy} loss function. It then runs the optimizer (by default stochastic gradient descent) on the training data and automatically preprocesses the audio files as we've discussed in previous sections, such as computing the spectrogram and then the MFCC of the audio stream, before passing them to the model/optimizer.

If you want to take a look at the source code it can be found here: 

\url{https://github.com/tensorflow/tensorflow/blob/master/tensorflow/examples/speech_commands/train.py}. You'll notice that there are even more flags than we set in the colab. If you’d like to customize the training process even more you can!


\section{Monitoring Training in Colab}

\subsection{Monitoring Training in Colab}

While the Keyword Spotting training script is running (for ~2 hours with the default settings and a GPU runtime), you may want to check-in and analyze how the training is performing. We’ve included two ways that you can do this into the Colab.

\subsection{TensorBoard}

TensorBoard is an application that is designed to help you visualize the training process. It will output graphs of the accuracy over time that look something like the below (this is a screenshot of the staff training a Keyword Spotting model):

\begin{figure}
    \includegraphics[width=\textwidth]{TinyML/2 Applications/3-5-17-1.png}
    
    \caption{A screenshort of tensorboard showing the accuracy of training over time. The model starts with around 30\% accuracy initially and rapidly increases its accuracy to almost 80\% by the 2,000 epoch and then shows drastically decreasing marginal returns on more training ending up only around 8\% by the 14,000 epoch.}
\end{figure}  



If you'd like to learn more about TensorBoard, Google has put together a nice intro Colab that explore some of its features:

\url{https://colab.research.google.com/github/tensorflow/tensorboard/blob/master/docs/get_started.ipynb}

\subsection{DEBUG Mode}

Unfortunately, the staff has found that it sometimes doesn't start showing data for a while ($\approx$15 minutes) and sometimes doesn't show data until training completes (and instead shows No dashboards are active for the current data set.). Therefore, we have also set the training script to run in DEBUG mode. By doing so it will print out information about the training process as it progresses. In general you will see printouts at each training step showing the current accuracy that look something like this:

\begin{figure}
    \includegraphics[width=\textwidth]{TinyML/2 Applications/3-5-17-2.png}
    
    \caption{A screenshot of the training results for steps 25-32 of training a keyword spotting model. Where the accuracy bounces around between 31\% and 42\%.}
\end{figure}  



Remember because we are using a variant of stochastic gradient descent, it is natural for the error to go up and down a bit during training. The key thing to look for is that the larger trend is for the accuracy to be going up on average! For example, if I scroll down a bit on this training run we can see that ~500 steps later we now have much higher accuracy on average (so learning seems to be going well):

\begin{figure}
    \includegraphics[width=\textwidth]{TinyML/2 Applications/3-5-17-3.png}
    
    \caption{A screenshot of the training results for steps 552-560 of training a keyword spotting model. Where the accuracy bounces around between 70\% and 83\%.}
\end{figure}  



Every 1,000 training steps you will also see a Confusion Matrix printed out that looks something like this:


\begin{figure}
    \includegraphics[width=\textwidth]{TinyML/2 Applications/ConfusionMatrix}
    
    \caption{Confusion Matrix.}
\end{figure}  



What this is describing is the number of samples that were correctly and incorrectly classified by class for the validation set at this point in the training process. Labeling the axis we can see how this works a little better:

\begin{figure}
    \includegraphics[width=\textwidth]{TinyML/2 Applications/3-5-17-4.png}
    
    \caption{A visual depction of a 4x4 confusion matrix showing that the actual labels are on the left and the predicted labels are on the right. Along the diagonal (highlighted in green) are the number of times the actual label matched the predicted label. On the off diagonals we see all of the mistakes. Since this example is taken from a fully trained model we find that the correct predictions far outweight the incorrect ones with values in the 100s on the diagonal and values in the single digits or tens off the diagonal.}
\end{figure}  



The column plots the predicted item label by the model while the row plots the actual correct label. Down the diagonal in green are the number correctly labeled items. Off the diagonal we find the number of items incorrectly labeled as the various labels. For all of your Keyword Spotting models Item A is ``silence'', Item B is ``unknown'' and then Items C, D, etc. are the various words you chose to train on in the order in which you define them in the Colab. For example in the pretrained model we had:

\PYTHON{WANTED\_WORDS = "yes,no"}

Which means that Item C was ``yes'' and Item D was ``no'' and the matrix was a 4x4 matrix. If we had instead set:

\PYTHON{WANTED\_WORDS = "yes,no,left"}

Then it would be a 5x5 matrix and Item C would still be ``yes'' and Item D would still be ``no'' and then there would be an Item E that was ``left''.

We hope that between TensorBoard and the DEBUG output you'll be able to understand how well your training is going and gain insights to improve your results!

\section{Assignment: Training your own Keyword Spotting Model}

Now that you know how to train a new Keyword Spotting model it's time for your assignment -- pick your own keywords from the Speech Commands dataset and train your own model! You’ll get to test it with your own audio input when you’re done! We hope you have a lot of fun with this one! As a quick warning, with the default values this should take ~2 hours to train on a GPU runtime once you launch the training cell -- so feel free to keep working through the rest of this section while it trains!

\url{https://colab.research.google.com/github/tinyMLx/colabs/blob/master/3-5-18-TrainingKeywordSpotting.ipynb}

\section{Assignment Solution}

For this assignment you needed to input a few hyper-parameters for the model training. A possible solution is below but any combination of 2-4 words from the list provided in the comment should have done reasonably well (and more words might have worked too although after a point the small model may have started to struggle). For training steps, the longer you trained, the better the accuracy should have gotten, up to a point, and then you may have found that the model started to overfit. As for the learning rate, if it was too high the training most likely failed due to divergence, and if it was too low the training likely failed because it didn’t make much progress. However, values somewhat close to the default settings should have worked reasonably well.

\begin{lstlisting}
WANTED_WORDS = "yes,no"
TRAINING_STEPS = "12000,3000"
LEARNING_RATE = "0.001,0.0001"
\end{lstlisting}

This assignment was mostly designed to give you experience with the full training pipeline and flow to prepare you for Course 3 where we will be training on your own custom data and deploying the model to your own hardware. Now is the time to make sure you understand the full flow well so you are in a better position to debug later when things get more complicated.

We hope you had some fun picking some words and getting to see the model learn based on your choices! We hope you also had fun trying to trick it with custom input! As always you’ll find the link to the assignment again below in case you want to explore it a bit more!

\url{https://colab.research.google.com/github/tinyMLx/colabs/blob/master/3-5-18-TrainingKeywordSpotting.ipynb}

\section{Metrics}

VIJAY JANAPA REDDI: Welcome back from Training the Keyword Spotting Model.
At this point, I'm hoping that you had a bit of fun
training the neural network model, changing a couple of parameters,
and seeing how the TinyCom model actually performs on the dataset
that we have.
At this point, you've gone through many different stages
already in this machine learning workflow.
You've gone from understanding how to collect the data we're using,
speech command data set.
Then we talked a lot about how to pre-process the data.
Then we thought about how to make the model small by inherently picking
a neural network architecture that is confined enough that it would fit
on our microcontroller-based device.
Then you've gone through the exercise of actually training the model.
And now, we're at the stage where we want to understand how
to assess the keyword spotting model.
Obviously, it's a neural network.
So you will immediately think of accuracy.
Well, of course, I'm going to look at what my accuracy of the model
is and see how well it performed.
But I want you to know that in the TinyML context,
efficiency is also extremely important.
Latency, power consumption, storage/memory
are extremely critical aspects.
This is exactly why we picked the tiny com model.
But in addition to both accuracy and efficiency,
I also want to touch upon the aspect of looking
at metrics that extend beyond just model performance or resource constraints.
When you build an application, there is more than just raw accuracy.
We'll touch on this soon.
But I'm going to try and get you to understand a different way of looking
at machine learning system behavior.
That's the critical message I want to get across in this particular video.
Let's talk about accuracy first.
First and foremost, when you talk about accuracy for something
like keyword spotting, what's the task here?
When I say, Alexa, I want the machine to wake up,
is it waking up only when I say it, or does it accidentally also wake up?
Why would it accidentally wake up?
Remember, it's got a microphone.
It's constantly listening.
And remember, sounds can be quite similar.
So sometimes, someone might be saying something
that might happen to feel like it's a yes or a hiss,
and that immediately causes your machine to think that perhaps you're saying
yes, which might be a keyword for it.
So that's this notion of false positive.
You want the machine to not have too many false positives.
In other words, if you did not say the keyword,
then if the device ends up triggering, you've got a false positive.
If you did not say the keyword, the machine is not supposed to wake up.
It's fairly simple.
But you also have another situation where, for instance, you
might have a false negative.
That is, you did say, OK, Google, or you did say, Alexa,
and the machine does not respond back to you.
That's pretty bad.
Often, that will lead to a lot of user frustration and dissatisfaction.
You don't want that because quality of service or quality of experience
here is king.
So you've got both false positives where the machine might wake up,
and you've got false negatives where the machine might not wake up.
So how do you sort of balance?
This is a fundamental or critical aspect that we
need to understand in terms of really understanding how
keyword spotting needs to be deployed.
That leads into this fundamental notion of what's called
an operating curve or our ROC curve.
Specifically, it's called a receiver operating
curve, which is what I'm showing you on the right-hand side.
It's a graphical representation or illustration
where it's allowing you to understand for a binary classifier that's
a yes or no how the system's actually going to behave.
It gives you a diagnostic sort of perspective
into how the system is behaving.
So for instance, on the x-axis, I'm showing you the false accepts per hour.
How many times are you falsely waking up the machine,
or how many times is a machine thinking that, oh, I think they said Alexa?
You don't want to wake up because that means
you have a rather crappy piece of device that's always
waking up thinking you said something.
And maybe when it wakes up and if it doesn't get the follow on words
that you express--
for example, like dim the lights--
then it might say, I'm sorry, I did not catch that.
Wouldn't that be annoying if your machine kept waking up,
and it kept saying, sorry, I did not catch that, when you are not even
communicating with the machine.
So that's the false accept problem.
You want to stay very much on the left-hand side.
You want to decrease the number of false accepts you have on the x-axis.
But at the same time, you have this problem of false rejects.
So false rejects are basically those false negatives that you have.
Now, the false negatives are where you actually are saying Alexa,
and the machine is not waking up, and that
proves to be problematic when you want the machine to respond to you.
Now, there is a very tight coupling between the false accept
and the false rejects.
That's why they're plotted together in this ROC curve.
What happens is that as you decrease the false accepts,
you're making the system more resilient to the false accept possibility.
But in doing that, you also increase the chances
that your machine might not actually wake up when you need it to wake up.
And that's why as you move more towards the left-hand side,
the curve starts picking up.
And you really want to be able to pick a critical junction where
you want to sit at where you have a right balance
between false accepts and false rejects, and that's called the operating point.
Every keyword spotting architecture out there effectively
relies on this operating point.
Someone has to tune this operating point,
and you have to control that operating point very carefully because depending
on different environments, that operating point might have to shift.
How do you do that?
We'll get into that in a little bit.
But that's that critical messages is that you're
walking this line between false accepts and false rejects.
And you have to walk that line carefully.
So given that you end up deploying these devices in many different locations,
you cannot have a single fixed operating point.
It makes it virtually impossible to be able to detect precisely when you
want the machine to wake up.
Why?
Because there's so many different environments and billions
of these devices floating around there.
How can you possibly preconfigure them?
This is why, for instance, when you look at your Google Assistant
on your device, it actually has the capability
to dynamically adjust that operating point, which is called a threshold.
That operating point threshold can be moved around.
So you can increase the sensitivity.
You can decrease the sensitivity.
It depends on how your device is used because, for instance, your phone might
be in your pocket or something, but you still want it to wake up.
Well, that might mean that it has to be a bit more sensitive.
Or for instance, you might put your device
somewhere near a TV in the living room where it's noisy.
So perhaps you don't want it to be overly sensitive.
Maybe you want it to be actually less sensitive.
So you can adjust this.
And of course, there's the aspect about just how you speak the words that
are coming out of your mouth and whether that accent is something
that the network has been trained for, and maybe if you just
relaxed it a little bit, you adjusted the threshold a little bit,
and allowed it to take a few more signals coming through,
then perhaps it would actually work better.
So these are just examples.
There are many other examples, but I'm trying to get across the critical point
that it's very hard to fix a particular point
when you're shipping the keyword spotting model.
So this is why typically devices will ask you to fine tune it
once you start it up.
It'll ask you, would you like to customize the device?
Or can you say, OK Google five times?
It allows it to customize a little bit.
And part of it is kind of adjusting to this.
And you can also manually adjust.
That's the critical message.
So they are empowering the users to be able to control
how the device should behave.
And you will pick the operating point that is best convenient for you,
and you can dynamically adjust it as needed.
Now, let's talk a little bit about the efficiency side of things.
So remember, on the accuracy, there's no fixed point.
You have to move it, and the operating point is the critical take away.
Now let's talk about efficiency here.
Obviously you know that latency is extremely critical because the machine
should respond to you quickly.
Otherwise, it defeats the whole purpose.
You might think that, sure, the model works really well,
but I don't care if it takes three seconds
to wake up before it can do something.
It has to be almost instantaneous.
So that's extremely important for end user satisfaction.
So the responsiveness is key.
And also, it has to run efficiently because you don't
want it to use up too much battery.
So the latency aspect, again, this is why you want to keep the model small.
So these are all factors you're going to consider.
There's, of course, the other aspect of memory efficiency all together.
We're resource constrained, remember?
This is why we went with the smaller TinyCom model that only
occupies about 17 kilobytes of memory.
And so therefore, we talked about doing quantization and getting it down
from 70 kilobytes in float to just 17 kilobytes in int8 format.
So at this point, you're very familiar with this.
So I'm not going to beat the point too much into you.
That said, I do want you to understand something important.
Now, when we're talking about models, you are focusing on the model size
alone.
Take a step back and think about what's really
happening in the TinyML application.
There is more than just the model, and it's extremely important
you keep that in mind because TinyML is resource constrained.
Let me give you an example here by showing you the memory usage.
These are just totally made up numbers.
We get the general gist.
First, there is some bit of application code.
It's not just the neural network that's running.
There is an application code that is actually
getting the data from the sensor, doing the pre-processing,
massaging things, saving something someplace.
All of that is code that's written.
That takes a little bit of memory.
Then you have the actual inference engine.
In our case, right now, we're using TF Light, but as we go into course 3
and deploy it onto the microcontroller, we'll
introduce you to TF Light Micro, which has got its own APIs
and is an extension of TF Light.
That engine alone takes a little bit of memory because it has to actually run.
Then, of course, you have your model that you have
worked so hard to design and optimize.
But then on top of that, you also have some amount of memory
that is actually needed for hosting all the dynamic variables and so forth.
And last but not least is the audio buffer.
Remember, when you are capturing those signals that are coming in,
you have to store that data someplace.
You have to store that data and then process that data.
Now, this is just an illustrative example.
So I'm not saying that the audio buffer is the maximum size or the biggest
memory consumer, I'm trying to make the fundamental point
that while your model is a small portion of this big equation,
there are many of the components that we have to keep in mind.
And it's extremely important, and I am always
surprised that machine learning engineers do not
realize this fundamental aspect.
As you go into TinyML, it becomes extremely important.
And this is one of the key things that separates TinyML engineers who
are cognizant about these kinds of overall resource constraints
as compared to most of the people who just focus on just the machine learning
model.
Keep that in mind.
Also, there are things beyond just system efficiency.
So we talked about model accuracy and how
to look at the model accuracy from a false accept
and a false reject sort of perspective.
Then we talked about memory latency aspects.
But there are other things.
What are the other things that actually affect?
Think about the overall quality of experience
that you have when you have a TinyML keyword spotting application.
What affects that?
Think about it.
We touched on this right at the beginning
when we talked about the speech command dataset.
The dataset that you train with can't have biases.
It might only work for a certain type of accent
if the people who contributed into that particular dataset
were from a very particular region.
Maybe they're from Europe, or maybe they're only from America.
Maybe they're from India.
Maybe they're from Singapore.
If you have those sorts of clusters and your dataset only comes from them,
well, all of your data, no matter how much pre-processing you do
and so forth, is going to result in a network
that only recognizes those kinds of sounds.
So we have to be very cognizant about what goes into that keyword dataset.
So how diverse are the users is something that we actually
have to test against, not just assume that the data is going
to be distributed just because perhaps I said,
oh, I got the data from everywhere.
We'll go much deeper into this aspect of dataset engineering,
and we'll look at it very carefully because it's fundamental.
Then, of course, to be tested in different backgrounds
in different environments really, right, because no one environment
has got the exact same sort of effect.
The sound that you're hearing right now is bouncing off of the different walls,
and that is having a certain type of effect.
So what you're hearing, you might be able to compensate
for it because your brain, well, it's a remarkable neural network
with lots and lots and lots of neurons.
Well, that's not the case in TinyML.
So we got to make sure that we're actually creating the datasets such
that it's actually going to work for the situation
where we're going to deploy the keyword spotting model.
And do we have noise?
As I mentioned again, earlier when we were talking about speech command,
I said we manually inject noise in there because you actually
have to simulate that noisy behavior.
And that's going to be crucial in order to build
a resilient neural network that's capable of doing keyword spotting.
All right, so we talked a lot about this sort of pipeline.
Hopefully you get a taste for the fact that there
are many different ways of looking at what
or how your keyword spotting model is doing.
It's beyond accuracy.
It's beyond efficiency.
It's about the overall keyword spotting experience that end users have.
And this is the sort of challenge that big organizations like Google or Amazon
or anyone else who is building these applications that
need to get deployed in billions of different devices are dealing with.
Now, are there ways we can actually improve our keyword spotting model?
The answer is absolutely yes.
What we can do is actually look at more than just a keyword.
In fact, we can look at actually the overall experience that you expect
to have, which is I might be having a streaming audio,
for instance, just something that is constantly being talked,
and then you got to be able to pick up that keyword amongst that conversation.
I'll touch on that in the next video, but it's
a very interesting real-world application of keyword spotting.
So far what we talked about is in a very nicely sort of created environment,
and we'll try and make it a little bit more realistic in the next video.
I'll see you there.

\section{Streaming Audio}

VIJAY REDDI: In this video, we're going to focus
on understanding how keyword spotting is really put into production scenarios.
So far when we've been talking about keyword spotting,
we've been talking about it in a very clean, cookie-cutters
of an environment.
Now I'm going to generalize this to help you understand what it really
looks like in the field.
We're not going to actually go ahead and do implementation.
This is much more from an informative standpoint.
In course three, we'll actually pick up on this notion
and go deeper with streaming audio when we deploy things onto our device,
but for now I want you to kind of think about it from the context of the model.
Now, where are we in the grand pipeline?
Well, you've come a very long way.
You're at the final stages about trying to understand how to deploy the model.
Now, in the context of deployment, I'm going
to focus primarily on the model aspect, how
we should think about the audio pipeline that we have built so far.
Specifically when we are looking at keyword spotting,
there's a notion of what we call a single shot versus streaming,
the fundamentally two different ways we can analyze how a keyword spotting
system actually works.
In the examples that we've been working with,
we have been very much focused on assessing
the accuracy of the model when it's only given the key word and nothing
else right, so only when you say, "OK, Google," and only when you say,
"Hello, World," for instance.
That's great.
That's a very fair assessment, but it's the minimum denominator basically
because everything needs to build on top of that.
Why do I say that?
Think of how you might actually be using the device.
Here's an example.
Hello, world.
It's a very nice day to be taking this course, and, oh, OK, Google,
turn up the volume.
Ah now I'm going back to my course.
What just happened here?
Specifically, the important aspect is that the keyword is actually
embedded in the context of a streaming audio sequence
that you're going to be getting.
So it's like trying to find a needle in a haystack, effectively.
You've got a whole bunch of different signals coming through,
and you've got to try and filter and find that keyword
in that continuous track of audio.
Instead, what we did previously was try and pretend
that everything was nice and clean, like this, where we are cleanly
given this nice audio file where there's only a single keyword, effectively,
listed in there.
But in reality, that's not going to be the case.
In reality, what you've got is a streaming audio.
You've got a whole bunch of conversation,
like what I'm saying right now, just kind of happening,
and then what you're looking at here is that continuous stream
that's actually going by you on the machine
that you've actually got to be processing.
We'll talk about this in much more detail in the next course,
and we'll learn how to build a classifier model that can actually
handle these one-second audio snippets on a continuous stream.
Now here is an example from a real patent out from Google,
which talks about how to handle streaming audio for keyword spotting.
The fundamentals that we'll touch on in the next course
will build on this actual production pipeline scenario
where you've got audio sequence that's completely coming through,
and then you're, on the fly, doing the feature generation using
many of the same principles that we actually discussed already
and things that you already practiced in the Colab, in fact.
And from there, we'll see how we can actually generate the keyword detector
so that it's actually able to deal with this continuous audio stream that's
coming through, and then that'll be a pretty awesome experience
because at this point, you should be happy that you're able to detect a word
or you're able to train the network to be able to detect your own word.
The next step is actually being able to detect it in a very complex sequence.
We'll get there, but it's one step at a time or one tiny bit of time.
I'll see you in the next video, and we'll pick up there.

\section{Cascade Architectures}

VIJAY REDDI: In this video, I'm going to wrap keyword spotting up
by talking about yet another interesting paradigm for a real world deployment.
We've talked previously about streaming audio
and how that's actually important.
Now, I'm going to talk about how to efficiently deploy keyword spotting
in a production situation.
Remember that a lot of our problems really stem from the fact
that TinyML is extremely resource constrained,
though it can do this with extreme efficiency from a power perspective.
There's a trade-off we make, and therefore, we
end up with these very, very small little models that we can train,
that are efficient in terms of the number of parameters they have.
But we end up with making many trade-offs.
Think for a second about the kind of trade-offs that we end up making.
What happens to the number of words that you can actually
get the keyword spotting model to recognize?
It's a very small network, so you really can't get it to remember too much.
The other thing is that even if we have only a handful of keywords,
getting that accuracy to be really high is difficult,
because the network is really simple.
Ultimately, these two things add up to limited user experience
or quality of experience for the end user
because the model is extremely simple.
Of course, you can make enhancements to increase the model slightly and so
forth.
But you get the general idea that no matter what you do,
you have that strong confinement in which you have to live in.
So is there something that we can do to alleviate this sort of a problem?
Well, think about the overall system as a potential solution.
By the system, I mean everything from the TinyML device, to your smartphone,
or to your laptop, all the way up to the cloud computing infrastructure that
is on the other end of the extreme.
So you have this massive computing infrastructure that's available to you.
What can we do?
How can we use that to improve the vocabulary or improve the accuracy
and effectively improve the user experience?
Well, let's take a look at the big picture.
When we start off, we're right at the microcontroller.
From the microcontroller, which is running the neural network,
we can do a little bit of processing there.
But if we need more capabilities, we need
to somehow be able to move the data up to the cloud
where there are more computational resources that
can run much bigger neural networks.
In that way, we can effectively do some bits and pieces on the local device
and then do more on the remote device.
So if you remember that GIF that I keep showing you, which is the Alexa.
And then it says, hi, Vijay.
Then, I say, dim the lights-- that animation that I've shown you.
The key thing that's happening there is that in a typical world,
that Alexa is what wakes up the machine.
And the dim the lights is actually something
that requires a bit more processing, so that's typically
offloaded to the cloud.
That's one way we have effectively increased the size or the processing
capabilities of the network in terms of its overall vocabulary.
So you've got to think about keyword spotting in the context of the bigger
picture, not just the keyword.
There are applications with just the keyword.
But there are also bigger applications that you
can enable through keyword spotting.
Now, instead of always going to the cloud, which
can cost a lot of energy, which I mentioned right at the beginning,
are there other ways to mitigate this?
Are these the only two extreme endpoints?
No, right?
Because you go from your TinyML microcontrollers.
You might have a phone, or you might have a watch, which
might have a slightly better processor.
And then, you might be able to actually offload to your PC,
or maybe then, send it to the cloud.
So there's a lot of this computational hierarchy that you have.
And the question is, how do you leverage that?
And that's where the cascading comes in.
The cascading is this idea that you have a little model
do something and a bigger model do something else in order
to unlock the overall capabilities that you actually want for an application.
Now, instead of going to the cloud, one of the things that you can obviously do
is change that step in the middle.
We can use a secondary larger model on a larger local device, like your phone.
So run a very small model on your microcontroller,
and if it detects something that you think is interesting
and maybe it's not very confident, you can also run that in nearby element.
That communication energy between the microcontroller and your phone
is going to be much less than the communication energy that
is needed between your microcontroller and trying
to get it to talk to a cloud server, just from a pure infrastructure
standpoint.
So you can intelligently try and do a system architecture
to try and get that TinyML application to be
a bit more powerful than what you want.
Let's take a real world example, for instance,
of how this is actually done on a smartphone.
A smartphone has a lot of capabilities, but it also has TinyML in it.
Think about the fact that, for instance, when you say, hey Siri,
it's something that your phone is constantly
able to recognize throughout the day.
As long as it's got battery, it's able to respond to, hey Siri.
But how do you think they're able to do that by having you continuously listen
to an audio stream?
Turning the mic on and doing the signal processing actually
costs a fair bit of energy that drains from your battery.
So can you architect the system to be efficient, is the key question.
So in the example that we kind of just looked at, hey Siri comes in.
You have some mechanism, the detector, that wakes the phone up,
that allows you to get more capability that's up in the cloud.
For instance, when you say, hey Siri, play my favorite song,
well, OK, that's pretty complicated.
So that requires, hey Siri to activate your phone.
The phone then has to recognize that there's something else you have said,
which is, play my favorite songs, which means
you've got to figure out what your favorite songs are in the cloud
and then start playing things.
So there's a fairly complicated pipeline.
But how can we enable this to be continuous, always on ML?
Well, that's where the cascading aspects really come in.
What you want to do, is that you have a TinyML subsystem, or an always
on processor module that is physically on the chip, a computing subsystem that
is separate from your main application processor, that is constantly
listening to this audio that is coming in,
and then using a TinyML model that is continuously running.
And then, if it detects something interesting,
then it wakes up a bigger processing engine.
And that bigger processing engine, which is much more accurate,
can then figure out, for instance, with much greater confidence, if it should
really be invoking the cloud, which is very energy
inefficient if you get it wrong.
And of course, you also don't want to have false accepts.
So this is the fundamental idea behind cascade architectures.
And in fact, this is exactly how Siri works.
And in fact, this is almost also how OK Google really
works on your smartphone devices.
TinyML is part of virtually every single phone that's out there.
So taking a step back, the TinyML aspects, the always on processor
is the kind of stuff that's running exactly
like on the microcontroller that we'll be using, which is a Nordic processor.
While your board is about the length of my finger,
the actual chip, the microprocessor itself, is really, really small.
And that actual microprocessor is found in many phones.
Now, by comparison, the bigger network, which has much more capability
and is much more accurate, is actually running on a full on Apple
processor, which can consume anywhere from 1 to 3 watts.
That's a lot of energy compared to the milliwatts of energy
that a TinyML device typically consumes.
And that's the critical thing.
A cascading architecture is a coupling of the different models
working together on different physical subsystems.
That's the key takeaway here.
And that's what cascade architectures are all about.
And the reason we end up doing cascade architectures
is because they give us two ROC curves, which at this point,
you should know what they mean.
And so on the x-axis, we have the false accepts.
On the y-axis, we have the false reject rates.
And this is something, by having two models,
you can have effectively two operating points
simultaneously that you can move around.
And that gives you much more flexibility in the space.
When we get into course 3, we'll go much deeper
into the cascade architectures as well as the streaming audio.
With that said, the next time you say, hey Siri,
I want you to think about all the brilliant mechanics,
from the machine learning, to the embedded systems,
to the streaming audio, to the cascading--
all of this stuff that's actually happening behind the scenes.
It's beautifully executed.
And that's what really gets me excited about TinyML.
And again, remember, this is just the tip of the iceberg.
There's much more exciting stuff that's coming down.
Stay tuned for that.

\section{Keyword Spotting in the Big Picture}

\subsection{Keyword Spotting in the Big Picture}

\begin{figure}
    \includegraphics[width=\textwidth]{TinyML/2 Applications/3-5-24-1.png}
    
    \caption{A high level flowchart showing all of the steps in possible keyword spotting applications and how the keyword spotting task is a small part of a much large speech recognition application architecture. We see how KWS can set off a speech recording leading to a speech to text and then natural language processing model which will determine the request from a server which responds with speech via a text to speech model.}
\end{figure}  


In this section, we have focused on the task of keyword spotting for tiny machine learning systems. We have seen how to develop algorithms for performing keyword spotting on a small set of keywords, and the challenges associated with developing these algorithms. We have also seen how audio can be streamed from a microphone to allow for low-power and yet online monitoring of keywords, and how this can be linked in a cascade architecture as a gateway to more complex and power-hungry operations such as natural language processing.

Keyword spotting has become increasingly popular in recent years, such as the introduction of software-based personal assistants like Siri (Apple), Cortana (Microsoft), Alexa (Amazon), and Google Assistant (Google). Similarly, smartphones utilize keyword spotting even while the CPU is not running through the use of an onboard digital signal processor, which is able to ‘wake up’ the CPU when a keyword is detected. These devices also utilize a cascade architecture, allowing the user to wake the device and subsequently ask questions which can be parsed and analyzed using natural language processing methods via an interface to the Cloud.

The prevalence of keyword spotting in industry will likely continue to increase as voice-recognition and natural language processing algorithms become increasingly accurate and powerful. We are already beginning to see voice-enabled TV’s, alarm clocks (e.g., \href{https://www.hatch.co/?gclid=Cj0KCQiAqdP9BRDVARIsAGSZ8AnmIZ51ebceO9mQPDPBQaTHYdHLfXs569h-iQyGx9mQUc8o4-Jb1YgaAkPZEALw_wcB}{Hatch}), lighting systems (e.g., \href{https://www.philips-hue.com/en-us/explore-hue/works-with/voice-control}{Philips Hue}), thermostats (e.g., \href{https://www.honeywellstore.com/store/products/wi-fi-smart-thermostat-with-voice-rth9590wf1003u.htm}{Honeywell}), and even smart vacuum cleaners (e.g., LG SmartThinQ). Based on the work you have done in this section, you are now equipped with many of the tools to develop your own systems for keyword spotting applications.

% 1.6: Data Engineering for TinyML Applications


\section{Introduction to Data Engineering}


\subsection{What is Data Engineering?}

Data engineering is a critical component of supervised learning, and consists of defining requirements, recording data, processing it, and improving the dataset. The quality and quantity of collected data determines the tractability of a machine learning objective. Training examples need to include enough salient features, along with representative noise induced by the surrounding environment (for example, day or night in images, or quiet and loud background noise for audio), for an ML algorithm to accurately distinguish between classes when deployed into the real world.

Data engineers need a rigorous problem definition in order to know what data should be collected, and must identify the potential sources of data. Data might come from on-device sensors, product users, or paid or unpaid contributors, and each may introduce potential licensing or privacy restrictions. This data must be labeled, and this usually requires manual effort by individual workers. It may also require domain expertise, for example, when labeling medical images. Mislabeled or garbled data may also need to be filtered out through manual inspection. Data engineers must also manage changing needs for a dataset, for example, in order to support additional languages.

\subsection{What's in this Module?}

All machine learning tasks begin with datasets. Datasets are often handed to us, and many of us never think about where they came from. But that’s a big mistake. Not being careful about dataset engineering can have serious consequences down the road. We are going to discover some of those pitfalls so that we don’t fall subject to those same potholes. We will use existing datasets as examples to ground the principles you will learn, specifically around identifying dataset requirements, data collection, data refinement, and sustainability. 

Specifically we will touch upon these items: 

\begin{itemize}
  \item The need for datasets as benchmark standards
  \item How to collect datasets for TinyML use cases
  \item How to be wary of bias and incorrect labeling in the datasets
  \item How to broaden the dataset collection process to accommodate diverse populations
  \item How we can repurpose existing (large) datasets for TinyML use cases
  \item How to be responsible when collecting datasets
\end{itemize}

\section{What's Data Engineering and Why It's Important}

VIJAY JANAPA REDDI: In this video and the coming couple of series of videos,
I'm going to be focusing on data engineering.
I mentioned right at the beginning that data engineering is something
that we have to pay attention to as we try and train machine learning models
because the machine learning models rely on data.
Without data, there is no machine learning model.
Well, because a supervised AI system is trained on a corpus of training data,
the fundamental question becomes, where is that original data set coming from?
Think about it right.
Typically, you almost always start off with that original data set,
some original data set.
You don't even think about where it came from.
Then you say, OK, I'm going to split this original data set into a training
data set and a test data set.
Remember, the training data set is what we actually
use to iteratively tune the hyperparameters of the network
in effect and the test data set is held out in order to actually give
an unbiased perspective on how well the trained model is really doing.
Training data set is typically further split up
into the actual training set and a subclass that's held up for validation.
The validation data set is effectively there
in order to assess how well your model is picking up
its skill in terms of whether you're making changes
in the hyperparameters for instance.
So the model training process is really the training
set plus the validation set.
And then at some point, you will say, I have a very effective model
that works well on my validation data set,
so now I'm ready to put it out there.
And well, when you're ready to assess it,
well, that's when you put in your test data set
and that's the final accuracy that often you report.
Now this is the very typical data processing part
that we all tend to focus on in the machine learning side.
But the key question that I often ask is where did this original data
set come from?
It's extremely important that you learn about this process
because without the data, you don't really have anything to train on.
And there are lots of nuances around building this original data set,
and that's what I want to get through in the next couple of videos
and the content that will follow.
So, data engineering is all about data sets.
I said this earlier in a previous video.
Now let's talk about this.
Good data is absolutely necessary for us.
Why?
Because if you have good data, then it'll give you a good accuracy.
But in order to get that good accuracy, you
have to ask yourself the fundamental question.
What problem are you trying to solve?
What does that mean?
First and foremost, your data has to contain the useful features.
If you're obviously dealing with something like images,
for instance, well, if the images are blurry, is that any good?
What happens if your microphone is not picking up much of the sound
that it's supposed to be picking up and you have poor quality data?
That data is not really useful.
The other part is can a human, for instance,
actually distinguish between the examples of each class?
Now the reason this is extremely important
is because someone has to label the data.
So the human is the actual expert who needs to know whether there's a cat
or whether there's a dog in the picture.
If they don't know the difference between a cat and dog,
then you've got a problem.
Similarly, think about it.
For instance, if you are trying to label a cancer data set
and you have to know the structure of a particular cell, can you do that?
I can't do that because I don't have the domain knowledge.
So, you got to think about what problem you're trying
to solve and find the right people.
Then, you've got to ask yourself the question
of how exactly are you going to assess the performance
of that particular model.
We talk a lot about accuracy in the context of keyword spotting
and I told you, yes, you can look at it from a model perspective,
but then there are also aspects of system level efficiency
and there also this quantity of experience
that the end user has overall.
So you've got to really articulate all those requirements.
After you kind of figure that out, then there's
a notion of how much data do you need and what quantity
does that data really need to have?
So for example, if you are trying to collect
data that's in the context of vision and speech,
you need to think about what are the fundamental attributes of those sensors
that need to be reflected carefully in terms of the quality of those sensors,
in terms of the quality of the data that's coming out of those sensors.
And how much of that do you need?
First and foremost, you need to have a wide distribution of different training
samples.
So for example, if there's a keyword, "yes" and there's a keyword, "no,"
you need to make sure you have several thousands of those examples
in order to get a wide distribution, so that the network is
resilient to any one particular person's way of speaking the keyword.
Then you've got to get accurate labels.
Well, think about it.
If your data is mislabeled, no magical machine learning model
is going to be able to compensate for that,
because we are training on the labels, because we're
using them as the ground truth.
Then we've got to make sure we have a sufficient number
of the different items that we want to tag, uniformly
balanced throughout the data set.
In other words, if you want to classify "yes" and "no,"
you need to make sure that you have "yes" and "no" equally balanced
in the data set.
So let's think about this more systematically,
because in order to go through those questions,
there are four key stages you'll walk through.
First one, starting with requirements.
What are the requirements in terms of building your data set?
We talk a little bit about the problem definition,
and of course, you have to think about what input format do you
want to capture the data in?
Whether it's like an MP3 or a PNG, or--
it's a simple process, but nonetheless, you need to think about it.
Now one of the most important things that I want to emphasize
here is that you've got to focus on permissions and rights.
Often we don't think about that.
Today, we take it for granted, for instance, that I can just go on Google
and I can do an image search and I will get loads of images.
But you got to ask yourself the question,
is that data actually accessible?
Meaning that it actually opens, you can directly use it.
Or do I have to license it for a fee?
Now, there's also other subtle issues about,
is a person-- for instance, in that image, is that person identifiable
and is that OK?
Or is it OK just to have a person in that picture?
These are subtle nuances and we typically
don't think about these things.
But if you are a TinyML engineer, and you come to me, and you said,
hey, I trained this model.
One of my first questions is going to be,
where did you get your data set from?
And that is an important question.
If you answer to me is, I don't know, more than likely,
I will not be hiring you.
That's because that model that you've trained, that IP,
that intellectual property that we created--
well, it might be tainted if it was trained on data for which we did not
explicitly have permissions.
This is extremely often overlooked and it's an extremely serious problem.
Now, let me give you a very specific example.
Let's say you're taking a picture.
You've got your camera-- you went and bought it.
You took a nice little picture and you're very proud of it.
And you're going to put it into your data set.
Depending on how you're using your model--
which might be trained on that data--
you might get into a pickle.
Because if that data for instance, has got critical logos, brand images,
or some sort of property in there, and your model is actually
providing some sort of end user value in order
to be able to pick those things up, well, you need to be careful.
Because those models, properties, brands, for instance,
they're all copyrighted properties.
And so, which means you don't have permission
for effectively accessing that.
You have to make sure you actually have the right permissions set up.
Only then can you actually use the data because the image
actually has information in there.
While the image is yours, that information that's inside it though,
might not belong to you.
So that might seem crazy to you, but it's totally true.
So you have to be very careful.
Similarly in speech for instance, when you're collecting the data,
you've got to make sure that the people who are giving you the data
have also agreed, in order to use that data,
in order to train a model for instance.
If it's not written down in the license file,
then you're going to get into a serious problem
when you actually use their speech to train a model that's later
on doing inferences and deployment.
So many different kinds of licenses exist out there.
I'm just giving you some examples of Apache, BSD, CC licenses.
You've got to understand which of these licenses
actually matter for you and you've got to make sure that the data has
the corresponding license.
In fact, as homework, I want you to go and look
at what's the license for the speech commands data set.
Then let's say you get all these requirements down right,
the second step is that now you've got to gather the data.
How do you go about getting the data?
You figured out you've got to have the right permissions and copyright
things put in place, now how do I tell the people that I
want this sort of data?
You've got to get the right people.
I already gave you an example.
For instance, if you are trying to tag cancer cells,
you better not hire me because you're not
going to get any good data out of me.
You've got to hire the right domain expert.
Then you've got to think about how am I actually going
to go and collect the data.
How do I actually ensure that I'm actually getting
the right labeling put in there?
What are the kind of labels that I actually need?
Is it just a cat and a dog, or is it a cat, dog, and a mouse?
And also.
You've got to think about the data sources from which you actually
want to collect all of this data.
Let me talk about that a little bit because it is relevant to TinyML
specifically.
We're using sensors.
Now you've got to think about, do you have enough of these sensors that
are deployed everywhere, that you are actually going
to get a rich collection of that data.
You also got to think about where these sensors are going to be.
Let's say you want to collect volcanic activity because you're
training some sort of TinyML application that does some good around nature.
Well, if you're collecting sensor data, you've
got to think about, well, how often or how much power and energy does
it consume?
Because that might have implications in terms of the cost of building
that data set.
Because if that sensor needs you to constantly go and replace batteries
in the computing system because it's using too much power,
well, that's not a very effective way of sourcing your data.
Crowdsourcing is a different way of collecting data
from different contributors.
And so contributors you have being people, like you and me.
But then, do we have enough contributors volunteering to generate the data set?
We'll talk about this in the context of speech commands again, soon.
And let's say for a particular product, do you have the consent?
Do you have the right privacy issues figured out?
As well as, got this issue about, like, maybe I
can pay people in order to actually contribute data for me?
And we'll talk about this as in the context of Mechanical Turk later on.
Now that's the whole gathering stage.
Let me go into the refinement stage.
You've got to process the data.
You've already seen some of the processing,
for instance, that we looked at earlier and of course,
one with doing data augmentation.
For instance, you also want to do some sort of basic clean up of the data--
checking for things like, is the mic actually
working when you're looking at the recordings?
Or are these recordings basically empty because the mic was not working?
Occasionally, it totally happens that you might get data
because you're collecting it at the scale,
you might get data that is not very of high quality.
So you need to validate that data.
You cannot just use the raw data.
That's the key message that I'm trying to convey out.
So let's talk a little bit about what data is actually useful.
You've collected the data now, now you're
trying to validate/verify the data.
Well, you got to keep in mind that verifying the data
is very tedious because it's manual in some sense,
that it'll cost you time, a lot of money, as well.
Because you might have to hire individuals to go and listen
to every single one of those keywords, for instance,
and make sure that every file actually has a right keyword.
Maybe there's a sound in it.
Well, that's great.
The mic was working.
But did they actually say the right key word
when they were supposed to say the right key word?
What if someone was being malicious and having a little bit of fun,
and sort of saying "yes," they are saying "no."
So you're actually polluting your data set.
You've got to be careful.
You can certainly use some level of automation.
For example, to detect if there was a sound or no sound in a recording,
well, that's pretty easy.
You can automate it.
And of course, as I was hinting earlier, you've
got to have the right domain expertise.
For instance, you might be thinking that, look, I've
hired person x as a domain expert and they labeled the data like this.
Well, does everyone agree with that particular person's
way of doing the labeling, or is there some nuance way
of doing the labeling slightly differently?
So that might lead into disputes or disagreements.
So you've got to be careful about this.
I know it sounds like a lot, but it's true.
It is a lot of work and this is why most people don't talk about the data set.
But actually, data engineering for data sets
is extremely important and very useful to know.
The final state is really about sustainability.
You've got to store the data someplace.
You've got to worry sometimes in terms of where you are actually
storing the data.
Are you putting it in some sort of cloud service?
And if it is in some sort of cloud service,
does everyone have fair and equal access to that data?
You might not think about this, but certain regions or certain countries
in this world have very strong restrictions
about what data is actually publicly accessible and what data is not.
Then of course, is the data secure?
Depending on what kind of data you're talking about,
sometimes you have to worry about the security of that data.
For example, if it's my medical records and a machine learning model
is being trained on it, well, I hope no one
is going to accidentally put it up on Google
and then be making it completely accessible by mistake.
And of course, then there's a chance that they're going to be errors.
And of course, there's going to be the issue of versioning.
You collect a bunch of data, but that's not it.
Data engineering is a continuous process.
You keep updating the data, which means you've
got to remember what data you originally trained your version one model on.
What data did the version two model trained on?
Because maybe you improved the model, but then there's
something wrong with that model, which means you need to be able to go back.
It's very similar to code versioning, which
we're used to where you always use GitHub or something in order
to keep track of the versions.
You've got to do data versioning as well.
So why would your data evolve?
Well, it's because the demographics might change, for instance.
Let's say for example, speech commands, today you're
learning about speech commands.
And soon you'll realize, well, the volunteers who
contributed into speech commands, well, it's
from a very small subset of the entire planet.
The way you and I say words might be a little different.
What that means is as you are learning about speech commands,
now you need to have your own sounds actually recorded.
So because you're getting informed, well, now we
need a bigger distribution of that data.
And of course, the expanding user base, more number of people.
So we need to just have more amount of data,
so we're much more resilient when we're training our models.
So I've talked about the requirements, gathering, refinement, and sustainment.
So this is one of the critical workflows you
need to keep in mind in terms of building the data set.
Often, again as I've said, people don't talk about it
but it's extremely important that you go through these stages.
And the bold ones that I highlighted earlier
are all extremely important ones and so pay extra attention to those.
Now generally, building data sets requires an insane amount of work
as you can see.
Typically, you're used to using ImageNet.
ImageNet has over 14 million images in there.
It's one of the big reasons that we're all here today,
it's because when ImageNet was created, it
set off this competition where everybody was
competing against being able to classify images correctly
from this ImageNet database.
Which is an awesome thing to do, but you need a lot of images.
ImageNet had 14 million, COCO--
which is for detecting objects--
is just as big.
It's 2.5 million segmented images.
So I'm trying to give you a sample for the scale of the data set.
If you look at driving data sets, like Waymo's open data set or KITTI 360,
well again here, you can see that it's a huge volume of data.
Waymo for instance, has nearly two thousand 20-second driving segments
of information that comes from cameras, LIDARS.
That's every 20 seconds is like 200,000 frames of data that is in there.
So it's a massive amount of data we're talking about.
1,950 might look small, but if you look at it from the total amount of frames,
it's a ginormous data set.
Similarly, KITTI is another very popular data set
that's used in autonomous driving.
It's got data that's about 73 kilometers of driving data
that comes from many different sensors that include laser scans
as well as static images coming in from cameras.
A large volume of data.
Again, as a summary here, look at the scale of these data.
Common Voice-- something that's for spoken audio--
has over 5,000 hours of speech data, which is a huge.
Common Objects, ImageNet, Waymo--
all of these data sets are big.
So what is involved in building these data sets?
How would you go about building your own data set?
How big should the data set be, is a question in and of itself
because wisdom generally is that more data is always better.
So let's start understanding a little bit
about how do you build your own data sets for a TinyML?
We'll pick up in the readings and the recordings
next in trying to get you to understand how you go about this process.


\section{Your New TinyML Applications}

Now that you’ve learned more about dataset design, how would you change your approach to dataset collection for the keyword spotting application you came up with in the previous section? How might you change the application to make data collection easier?

Remember we considered many questions including:

\begin{itemize}
  \item Who are the anticipated end users?
  \item In what environments will the users employ the application?
  \item What are the goals for using the application? How will this impact the requirements of the ML model's performance?
\end{itemize}


\section{Dataset Standards: Speech Commands}

VIJAY JANAPA REDDI: In this video, I'm going
to talk about data engineering using speech commands as a specific example.
In the previous video what we learned about
was that data engineering is a crucial process
and that there are a set of steps that we have to execute.
And I introduced each of those steps in abstract terms, as general concepts.
What I want to do now is actually ground that.
I want to talk about speech commands in each of those different buckets.
So recall that speech command is the dataset
that Pete Warden put together in 2018.
Now this data set pretty much had to go through these fundamental steps.
Any dataset goes through these steps.
You've got a set of requirements that motivate
you to kind of build this data set.
Then you have to figure out how you're actually going to gather the data.
Then you're going to have to clean the data.
And then you're going to have to continue improving the data.
So let's go through each and every single one of these.
Let's talk about what actually motivated him to actually build the data set.
First and foremost, keywood spotting was just
starting to take off back in 2017-18.
But the problem was that it was largely limited to big organizations.
There was nothing out in the public that you and I could, use
be it for research, for education, or for commercial purposes.
There was no public data set.
So that was the seed that kind of jumpstarted the whole process.
There were bigger speech data sets, like LibriSpeech and so forth.
But they're not focused on keyword spotting,
which is the critical requirement here.
Recall that keyword spotting is specifically
looking at one word, not a whole sentence being spoken.
Also, while Google, Amazon, and Apple were putting out
products that had this quote spotting capability,
they were all using proprietary data sets
because they have the ability to build those private data
sets by hiring people to actually get the data that they need.
But there was nothing out there for the public good.
And that's what Pete was going after is trying to build a public data set.
And so the goals were to make sure that the data was being collected
in such a way that it was actually permissible for distribution,
be it for research, education, or commercial purposes.
So these are some of the fundamental requirements.
And ultimately, if you get these fundamental requirements,
a data set can become a new standard.
So speech commands today is effectively the de facto standard
for keyword spotting applications.
I'll talk more about this later on.
So what were some of the additional requirements
that Pete identified in terms of what needs to go into the data set?
This fundamentally ties down to what are the common keyword application
use cases, right?
Yes and no as an example.
That's very common because you might get a machine asking you a question.
And you want to say yes or no.
For robotics for instance, if you want to control
some small autonomous vehicle, well, typical commands
are usually left, right, go, stop.
Or for numbers, maybe you're on a telephone and you often say digits.
The machine responds.
Well, that's key with spotting, for instance.
So you go from 0, 1, 2, 3, 4, 5, 6, 7, 8, 9.
So those are the requirements.
So next step is how do we actually go about gathering the data?
What are the sources?
Who are the people who contributed?
How do you get the data labeled?
Well, it was built on the backs of people.
People contributed their own time and effort.
So over 2,600 volunteers consented to have
their voices recorded and distributed.
They consented.
This was part of the licensing thing.
They understood why or what Pete was doing
and agreed to have their data be shared.
So therefore, we have this access to this data today.
And because people were contributing, we naturally
had a variety of different accents being captured in that data set.
And remember, you need to have a large number of samples.
So there are over 1,000 samples for each keyword in there.
Obviously, some words are more popular than others.
But that's OK because it really depends on what
your application is going to be.
This is why I was saying earlier that you
need to make sure you know what your application
task is because that means you need to have the right data set at hand.
And therefore, you need to know the characteristics
of the data that's inside the data set.
In order to make it accessible and to be able to collect very easily,
simply set up a web browser.
Basically, that was a URL that you could open up your browser on.
And then you would be able to record your voice and contribute.
Nothing was needed.
No app installs were needed.
Nothing custom.
It was just a website.
Next, after getting all this data we said we have to refine the data.
So how did he go about that?
Some data is simply not useful.
We talked about ideas like looking at audio signals
and saying are they too quiet, so that you can't really
extract useful spectrograms from them.
Maybe instead of saying yes, someone was messing around and saying
the opposite word for everything.
Instead of yes, they were saying, no.
Instead of no, they were saying yes.
People like to have a little bit of fun.
Now you can do some of these things through automation.
For example, you can check the volume of the particular clip
and then be able to remove it automatically.
You have methods of doing that.
You also need to be able to effectively identify
the critical part of the entire recording
because when someone says record, it just
starts recording until they stop the recording.
So you might have a 10 second clip that needs
to be cut out just down to the critical piece that
actually has the key words in it.
Then someone, of course, had to manually validate that data set.
So you've got over 100,000 of these utterances that
needed to be manually validated.
And that was again, done through crowdsourcing,
which effectively means people volunteer to actually do the validation.
The final part comes down to sustaining or improving the data set,
making it accessible.
So the original data set only had 10 keywords in it.
Now the data set is over 35 different keywords in it.
And so it's starting to expand as this is
the effective standard that everybody's building keyword spotting models on.
So as the needs become more and more, the data set
naturally continues to grow.
Now the next thing is about having scripts that automatically do
the train validation and test splits.
Remember, I said earlier that when you are doing these splits,
they're not arbitrary splits.
The distributions of the data between, let's say,
the training set and the test set needs to be thought out quite carefully.
It's not some random sampling of the data.
So there are scripts that actually do this.
Keep that in mind when you are building your data sets in the future.
And of course, there's the ultimate question
of how do you expand these new languages.
Speech commands is only focused on English.
What about your language?
What's your native language?
My native language is actually Telugu.
I also speak Hindi.
I also speak a little bit of Tamil.
None of these have keywood spotting data sets
that I really wish we had access to.
There are other things that are coming up
in order to try and fill this gap, this niche, this need
for new languages and keyword spotting.
So we'll talk about that in the coming couple of videos.
And you'll be very excited to see how you can potentially
build the next major data set for your language.
And with that, I hope you picked up a couple of key points.
The key thing that we want to do is that we
want to build a data set that's actually useful by everybody.
Specifically, we want to be able to compare
different models that are trained.
But the models need to be working off of the same standard data set.
Speech commands is that standard dataset.
So that we can actually compare two different models
and see which one is actually performing better.
It becomes a benchmark.
Benchmarks are crucial because TinyML has very strong constraints, remember?
So you want to be able to use the same data set
to benchmark different systems, so we can actually
assess the forward progress we're making, not just from the model
standpoint, but also from the system efficiency standpoint
because now you're running this Tiny model that has been trained
on the exact data that you know.
And then you're running it on a system.
So you can also benchmark the system, if you have a standard dataset.
Finally, there's a process of being able to improve on the data set
that we started off with.
So speech commands is this original standard.
And now we can continue to build on top of it.
The key thing that I'm trying to get do is
that if you want to move the needle forward,
if we all want to enable TinyML in the future,
we need to be able to do apples to apples comparisons from the data set
perspective because data sets are the rocket fuel for machine learning
models.
So we want to make sure we have a standard data set.
If we have a standard data set, and people are building models,
then we can also use that coupling to benchmark different systems.
When we get into course 3, we'll talk much more in detail
about benchmarking and doing apples to apples
comparisons of different systems.
But let's hold on to that thought.
With that said, I hope you understood what data engineering is
and how speech commands, as an example, did all the things that
are needed to be done properly.
And that's why it has become the industry standard today
for applications such as keyword spotting.

\section{Speech Commands Paper}

In this reading assignment, we want you to read through the ``Speech Commands'' ( \url{https://arxiv.org/pdf/1804.03209.pdf}) paper that Pete Warden (one of the course instructors) himself wrote. The paper provides additional information on some of the details that we discussed in the video lectures. It provides a bit more context and detail around the following important topics:

\begin{itemize}
  \item What are Speech Commands?
  \item What was Pete’s motivation behind building Speech Commands?
  \item How is Keyword Spotting different from traditional speech recognition models? 
  \item What are the important metrics in speech recognition for KWS?
\end{itemize}

\section{Crowdsourcing Data for the Long Tail}

VIJAY JANAPA REDDI: Hopefully, at this point
you are excited about keyword spotting, not
from general keyword spotting in English,
but perhaps keyword spotting in your native language.
That's an exciting thought that you and I should be thinking about.
How do we do that?
What's the problem in being able to do that today?
It's the data set.
Do we have a speech command data set for your language?
They sure don't have it from my language, which is Telugu.
Now, I'm going to talk about what's the real problem here,
which is called the long tail.
Let me explain what I mean by the long tail here.
The long tail gets at this point that if you
want to collect data that's really diverse,
which means you have to consider many, many different kinds of people
and their accents and the nuances of how they say things,
then the data will diminish really fast.
But if you're talking about just English, for instance,
there's a large wealth of data out there that we can tap into.
That's effectively the long tail problem.
The long tail really becomes an interesting problem
if you want to do something like keyword spotting
for many different kinds of languages.
So I'm showing you a picture on the left that you will recognize immediately
in terms of the different languages that are there.
And their population clearly indicates the size of the circle.
And you'll see that for some languages, there's
a large number of people who you can tap into to get the data.
But for the many other languages, we don't have that many people.
There are over 7,000 different languages spoken on this planet.
Question is this.
How can we enable speech commands for all of them?
Wouldn't it be cool if you could say, Alexa or OK
Google in some sort of native language and then the machine responds?
That would be a very different experience.
We might be able to come with all kinds of really interesting use cases.
I just talked about Alexa and Google again,
but those are just voice assistance.
I'll talk about one very specific example that really keeps me up
at night as to why we should enable speech commands for the planet.
Think about this in the context of a developing country
where, for instance, you don't have massive computing
infrastructure and rich internet connections and things like that.
But let's say there's some sort of problem.
Let's say there's a virus out there that is spreading across these villages.
And the one base technology that almost all different places have
is access to radio, a very primitive sort of technology
that is pervasive across the planet.
Now, imagine if you were able to train a keyword spotting in a certain language
and get it to detect certain keywords.
And let's say we want to figure out how much of a problem
we think it is that there is a certain virus in a certain village
and build a map-- overall map for a whole region,
for a whole collection of different villages.
Imagine if we could train the keyword spotting model
and put it in front of the radio and tune
in to different regions or different villages
where the signals are coming from.
Then the keyword spotting model is going to be
able to pick up the keywords, let's say, the virus or COVID or something
like that.
And when it picks it up, you can start building and figuring things
out about maybe there is actually a problem in this village--
oh, there's a problem in this village, oh, there's a problem in this village--
without explicitly having to have doctors and nurses
and everybody go out across the entire place,
which is very costly and extremely hard to do,
especially in developing countries.
So that's a very creative example of trying
to use keyword spotting for something outside of voice assistance.
This is why I think it's extremely important to build speech commands
sorts of data sets, small little data sets
for the variety of different languages that we have across the planet.
Now, how do we go about doing this?
Obviously, at this point you know that we can do one of two things.
One is you can either pay in order to get someone to actually do
the labeling, contribute that voice.
Or we can do crowdsourcing the way Pete did for speech commands
and rely on volunteers who actually contribute the data.
What's the difference between the two of them?
Well, if you want to be able to pay for all the different languages,
well, that's costly, and it's got limited scale.
Someone needs to be philanthropically interested in contributing
for the whole planet.
Maybe that's the Gates Foundation.
But that's still a lot to ask from any one organization.
So more than likely, the left-hand side approach is not going to work.
Paying in order to try and build data sets
has very limited applicability unless it's a business proposition around it.
The example I gave you is not so much as a business proposition.
It's much more about doing social good.
So we have to rely on the community to actually help us build the data sets,
because if we can rely on the community the way Pete did
and give them a good cause or good motivation,
it's more than likely that we might be able to build it.
That's where I want to introduce Common Voice to you.
Common Voice is effectively a crowdsourcing platform
that allows people to actually contribute speech.
The way you do that is by simply going to a website,
and you click on Speak on the top left.
And a certain sentence or a certain keyword is given to you,
and you record it.
This is very similar to what Pete did.
It's exactly what I was talking about previously when I said he put up
a website and people contributed.
And how do you validate it?
There's a listen mode where there's a certain statement that
is supposed to have been recorded.
Then you click Play, and whoever previously contributed it,
their voice will be playing.
And you will be listening to it and looking at the text
and saying "check" or "no, this is not right."
So that's the validation stage.
This is effectively what Common Voice does.
It's an infrastructure that is out there that
helps us build data sets for speech, like what we are looking at right now
for English.
But you can actually extend this to many different people
in many different languages.
Common Voice is a platform that's out there for the public good.
Over 50,000 volunteers have already contributed speech in there.
Now, you might be wondering, why am I suddenly talking
about another platform?
Because we already got speech commands, so
what's so special about Common Voice?
Well, the interesting thing is that it's there for 54 different languages.
These are all different languages that people contributed.
Maybe you and I should start contributing to this Common Noise
infrastructure, because that's how a new language gets recorded.
One of the interesting things here is that we can enable speech recognition
for all kinds of different languages because this
is a community-driven effort.
And that's what really gets me excited about Common Voice.
What I'm trying to do here is introduce you to two things.
One is a very specific instantiation of a platform
that you can actually use to build data sets.
Two, I'm trying to get you to understand how we can build data sets.
You need to know the right tools and mechanisms that are out there.
And Common Voice is one repository from which you can get data.
So what Common Voice has is valid data sets.
So people have contributed into it.
So someone has to go and verify that it's correct.
So they have validated some of it.
The way that validation is done is very similar to the way
Pete did his speech commands data set validation.
Some people listen to it, and then they say, yes, this
is exactly what it's supposed to be.
There are going to be recordings in there that are invalid.
So some people said, yes, this is not actually correct.
And finally, there are others where some people have said, yes, this is correct,
and others have said, no, this is not correct.
So that's the other bucket.
I'm just trying to give you a little flavor
for how things are actually structured or what kind of situations
will come up.
It's not always going to be yes.
It's not always going to be no.
You might have a mixture.
And that's the key point as to why we're talking about how we actually
end up with this other bucket.
Some of the very important and interesting aspects
are that it has a permissive license, which means then you
and I can use it for commercial purposes or research-related purposes
and so forth.
There are many contributors, as I said earlier.
It comes with metadata, which is extremely important.
Think about it.
Metadata for a data set is like a nutrition sticker
that is on the back of every food item that you
buy, which tells you how much sugar there is,
how much fat there is, and so forth.
Metadata is very similar to that, but it's specific to the data set.
It actually talks about the internals of the data set.
We'll get into that more in a bit.
Then, of course, there is a specific subset
of Common Voice that is extremely relevant to us, which
is the single word target segment.
It is almost exactly speech commands, which is
something that just started recently.
Pete pioneered the idea about two years ago before this.
The interesting aspect of this is that it supports multiple languages, not
just English, which is the key thing that I
was trying to get to in this video about talking
about different languages for people across the planet.
There are over 18 different languages today,
and it's continuing to grow remarkably.
The languages are effectively capturing a set of keywords.
In this case, it happens to be yes, no, hey, Firefox.
Firefox is a web browser.
And the reason they are capturing this is because, well, the Common Voice
project is part of the Mozilla Foundation, which
actually supports the project, and therefore Firefox is captured.
And it also captures digits, 0 through 9.
So one of the richest things about this is the fact
that, again, because Common Voice has a permissive license,
we can actually work with it, which is really powerful.
Again, here I'm showing you a screenshot of how you would record it, very
similar to what Pete pioneered.
It's also very easy to use because all you need is a browser.
I strongly encourage you to actually go out and contribute
a couple of keywords, because that would help us.
And one of the interesting things here is
that you might actually be able to contribute in starting a new language.
So if you're excited about this, get yourself
and get your friends and everybody to go to Common Voice
and potentially contribute, because we can make the world a bit better
by having additional data that specifically ties to new languages.
One of the important aspects is generally
about the structure of the data that we need to have.
So Common Voice has yes, no.
There's a notion of prompt.
What's the prompt that we want to have?
Then, of course, is a feature of recording
that data that's actually important.
How do we actually record the data?
What are the attributes that are actually captured inside that data?
How do we validate that data?
We've got to know which of the recordings
are actually up, which are the recordings are actually down,
and then use that to create the data splits for the test, the training,
and the validation as well.
I'm talking about all of this in the context
of explaining how speech commands came about as a data set that you are using.
I'm trying to emphasize that there is more to just building machine learning
models.
You need to understand that machinery models are comprised of data.
And now, how do we go about collecting the data
is something that is extremely crucial, and we
need to be systematic in that process.
We cannot just do an ad hoc process.
So hopefully, in this video what you've understand is that, hey,
we started with speech commands, which was just for English,
but now there are other mechanisms to build data sets for our own languages
that we desire.
And there's a method to the madness.
And Common Voice is one systematic process
which uses the fundamental building block steps that we talked about.
And so we can all contribute into it, and then potentially we
can use that data set to train custom keyword spotting models.

\section{Giving Back to the Open Source Community}

We have explored community driven data collection through the examples of both the Speech Commands Dataset and the Common Voice Project. We discussed the benefits and limitations of these approaches. Given that we believe that these projects are important for the TinyML community, we thought it would be a good idea (and fun) to give you a chance to contribute to such a project and fortunately, the Common Voice project is still actively collecting data and looking for help with validating recordings from other users!

This is 100\% optional. You are in no means required to contribute and your contribution will have no impact or bearing on your completion of the course. If you would like to contribute to the Common Voice project the link can be found below. Please consult the data privacy documents and warnings listed on the website before participating.

\url{https://commonvoice.mozilla.org/}

Update: The Common Voice project has decided to combine their single world data collection and validation with their full sentence data collection and validation. As such, you will see a mix of single words and full sentences if you choose to contribute.

\section{Reusing and Adapting Existing Datasets}

VIJAY REDDI: So far I've talked a lot about building data sets from scratch,
where we have understood the process of outlining the requirements,
figuring out how to go about gathering the data,
then figuring out how to process that data effectively to validate that data,
and then finding a means of sustaining that data to keep improving on it.
I wanted to walk you through that whole pipeline,
because it's extremely crucial if you're going to start from scratch.
I wanted to identify all the pitfalls you might have.
Now, what I want to complement that knowledge with,
is saying that you don't always have to start from scratch.
You might be able to potentially reuse data sets.
And the key question there is, why?
Well, TinyML is beautiful because it often focuses on one particular task.
And because it does one particular task, you
might be able to leverage broad swaths of data that is out there
and be able to subset that data down for that one particular task
you want to do.
And I'll talk about that very specifically later on.
So why do we want to do something like that?
Well, building data sets is costly, and it's tedious, and it's error prone.
And therefore, we don't want to do it.
We, instead, want to reuse any data set that's available that has already
gone through that iterative process.
So then, you've got to ask yourself the question of what's actually available,
and what's missing?
If there is a data set that I might be able to repurpose,
then where can I get that data?
And if something's missing, then do I have
to go build that data set from scratch, like Pete did for speech commands?
Well, you've got to identify the answers to these questions.
So let me talk about what's available.
Here's an example from TensorFlow where they have a data set catalog which
has many different areas that are covered, all the way from audio,
to image classification, to object detection,
to question and answering, text summarization--
many different kinds of different data sets
that are readily available that you can actually use.
Before you jump on and start using them, though, make
sure you check the licenses and see what they say about how can you use them?
I'm sure you probably use them for research and education.
But can you use them for commercial purposes?
That's a question for you to answer.
So let me talk about one specific example or way in which
we can actually do subsetting of a large data set
to create a TinyML specific data set.
I had introduced MS COCO, which is specifically
focused on object detection, for instance, where given a picture,
it's got to identify critical bounding boxes.
So the labels are all basically bounding boxes that identify, ah, here's
a person; here's a cat; here's a dog; here's a bus; here's a truck,
and so forth.
Now, this data set is pretty comprehensive.
It covers many different areas, if you remember the details
from the previous video.
Now in TinyML, for example, if I want to do a person detection module,
I want to have a simple camera sensor that the TinyML
model is going to be using in order to detect a person.
For example, maybe I wish I had a smart light in my room,
where whenever I leave the room, I want the light to automatically turn off.
Why are the lights always on when I'm not in the room?
Well, today, I have to manually flip the switch and turn it off.
I wish there was a little TinyML person detection
model, because if it was running very efficiently, then
whenever I leave the room, the lights automatically turn off.
Whenever I come into the room, the lights automatically turn on.
That's a good use case.
There are also security application use cases
of being able to detect if someone is there
when they're not supposed to be there.
So that's a very good application.
What do you want to do if you want to build that application?
Well, do you go out and build the data set from scratch?
You don't have to, because we're lucky to have this MS COCO data
set that's available to us.
We can effectively take that large data set
and break that data set down for only the parts where people are involved.
This data set might have a car or a truck
or bus or whatever, along with people.
But we don't care about that.
We only care about whether there are people or there are no people.
And that, we can effectively do.
So we can take a data set that's already curated and subset it.
So we don't have to start everything from scratch.
Now, similar to the way we're talking about data sets
and we don't have to restart from scratch;
we just need to write the scripts to filter through,
we can also apply the same notion for how we actually train models.
We don't always have to start from scratch.
For example, if, let's say, today you were able to drive a car--
you got your license, and you were able to drive a car, then tomorrow, you
upgrade your car.
Do you go back and get a new license, or do you just
transfer the learning that happened when you first got your car
and when you first got your license to the new car?
Why can you do the transfer?
Well, because, sure, the cars might be a little different,
but they're still effectively a car.
But if it's something like a car versus trying to fly a plane,
well then, you've got to go back and get your certification.
So we'll be getting more into transfer learning later on,
but I'm trying to give you an idea that not everything in machine learning
needs to be started from scratch.
So we might be able to reuse some pretrained models, models
that have already trained on some of the existing data,
and adapt them a little bit by using transfer learning
and then be able to shrink the time it takes for you to actually get
the new application that you want to build really fast off the ground.
Now, you don't always have to work with the data that's
only coming from real-world pictures, for instance, or real-world sounds.
You can also get your data from simulations--
physics simulations, game engines that are becoming extremely
good at simulating the real world.
So you can actually produce a lot of data from simulations
and be actually able to use it directly inside your machine learning model.
Also, if you have heard of GANs, machine learning models
are using GANs to be able to produce new pictures.
So that's also another source of potential data.
Of course, I'm going to include a reading that's
going to tell you about some of the nuances and pitfalls
about using simulation data or synthesized data versus using
real data, because there are some nuances there
that are extremely important.
So you might be thinking, ah, I have solved
the problem of my having a balanced data set or a large data
by using a little bit of real data, a little bit of simulation data,
a little bit of ML models that have generated the data.
And I've created this mega data set.
Well, that data set model, you've got quantity.
You lack this other part-- quality.
So we'll learn a little bit about the nuances
there, because there are some pitfalls there you can get yourself into.
Now, we'll be going into some of the transfer
learning aspects both from a machine learning model perspective, as well
as the data set perspective once we get into the image sensor part,
which is coming up next after we finish up data engineering.
So stay tuned for that.
And after that, we'll be able to extend out to anomaly detection, which
also uses something quite similar.

\section{Responsible Data Collection}

VIJAY JANAPA REDDI: In this video, we're going
to be wrapping up the data engineering section of this course
by focusing on responsible data collection.
Specifically, I'm talking about biases.
Biases are a serious problem in machine learning
and I'm going to use speech recognition as an example
here to explain how we can try and be cognizant about biases creeping in.
What are some sources of bias that we have?
For one, it could be simply from the fact that we have gender diversity.
It is well known that much of the speech data that's actually available today
is really biased towards male speakers, because they're
the ones who are contributing the most volume of it
and female speakers contribute less.
But because of that, the datasets that are public and available to us
are naturally going to end up leading us to have
models that are more tuned for male speakers than for female speakers.
We don't want that, because when we build a TinyML application,
we want that keyword spotting model to work generally for everybody.
Same thing on age groups.
The way the words sound when you're older versus when you younger
tend to be slightly different.
So we've got to make sure we catch the right demographics all over the place.
Yet again, when we go to accents, if you'll recall,
I said that many of the keyword spotting systems that are out there
do not work well for my mom, for instance,
because she happens to have a thick Indian accent, even
though she speaks perfect English.
So we need to make sure that we're cognizant about accents, for instance,
from not creeping in, because that's a source of bias.
Her accent happens to be more Indian because she's coming from India.
Now for a native speaker, that might not be the problem.
So we have to make sure we capture a wide dataset.
Also, many languages tend to have different kinds of dialects
that are very close to one another.
But you've got to be cognizant about the fact
that languages have different dialects, because some dialects are
more pronounced than other dialects.
Which means your data might be, again, biased towards some of those dialects
and less for the other dialects.
Which means the models are not going to work well generally.
Finally, there's the aspect of languages in itself,
which we've talked about quite a bit in terms
of the long tail in the previous videos.
So overall, my point is that there are many different sources of bias.
So let's say because you took this class now you're cognizant about this bias
and you're going to be very careful when you're building your dataset.
Does that solve the problem?
Well, let's see what happens in the real world.
You and I might want to build a dataset because we want to build
some sort of keyword spotting model.
So how do we go about it?
It's highly unlikely that we are going to go out and manually
collect the dataset ourselves.
More often than not, we're going to have to rely on some help,
because the amount of data that you and I can directly collect
is going to be limited.
So we're going to likely hire some sort of third party company who is actually
going to do the labeling.
They're going to go out and hire some set of people
who will be doing the labeling.
That allows us to naturally scale up because given the volume of data
we want, we can tell this intermediary company how much dataset that we want
and they'll hire the right number of people.
But we have to be very careful that the biases that we're
cognizant about, well, they don't necessarily
cover the intermediary services that we're hiring.
Let's be specific, for example.
Amazon Mechanical Turk is a widely used crowdsourcing website
where businesses can effectively hire crowdworkers
to perform certain types of tasks.
Usually people hire them to do a lot of labeling for machine learning datasets.
Well, one of the interesting things here is
that even though every worker comes in being totally unbiased,
they're just being themselves, one of the problems
is that when you're asking someone to do labeling and you hire them,
you do not know where these people come from.
For example, here I'm showing you the long tail of the crowdworkers.
The x-axis is showing the different countries from where they come from
and the y-axis is showing the distribution, effectively.
And you can see that a large portion are actually coming from a few countries
and there's a long tail.
Which means that depending on the kind of labeling you're asking them to do,
their own bias might actually creep into the dataset,
even though you are cognizant about the process of building
a proper dataset without biases.
So we have to keep this in mind, because you and I cannot directly build
a dataset.
We are almost always going to be hiring a third party to build the dataset.
But if we are aware of these kinds of systemic issues that exist out there,
we might be able to build a better dataset in the future.
With that said, the question you should ask yourself
is, how can you work to avoid biases inside your dataset?
Biases are just one of the aspects that we've touched on the data engineering
section of the TinyML course.
There are many other aspects.
We have talked about how we actually go about getting an accurate
model, not just from a moral standpoint, but from overall quality of experience
sort of standpoint.
We will also talk about what that means in terms
of the efficiency of the hardware.
So, long story short, building a dataset is quite complicated
because all of these things.
It's a ripple effect.
If you get the dataset wrong, you can get a model that works just well,
but unfortunately it's not working in the right way.
So just because your Colab says you've got
a certain accuracy does not mean that it's actually doing its job well
from a TinyML application standpoint.
And that's that critical message that I'm going to wrap this up with.
So in the next section, we'll be picking up with the new TinyML application,
where we'll be reusing many of the fundamental concepts
that we have covered here.
I'll see you there.

\section{Forum}

Do you think open-source and crowdsourced data collection are good ideas? What concerns do you have? How would you encourage participation and contributions to the dataset? How would you ensure the data collected is reliable and validated? What strategies might reduce bias in data collection?

\section{Summary: Data Engineering}


\subsection{What is Data Engineering?}

Data engineering is a critical component of supervised learning, and consists of defining requirements, recording data, processing it, and improving the dataset. The quality and quantity of collected data determines the tractability of a machine learning objective. Training examples need to include enough salient features, along with representative noise induced by the surrounding environment (for example, day or night in images, or quiet and loud background noise for audio), for an ML algorithm to accurately distinguish between classes when deployed into the real world.
    
    Data engineers need a rigorous problem definition in order to know what data should be collected, and must identify the potential sources of data. Data might come from on-device sensors, product users, or paid or unpaid contributors, and each may introduce potential licensing or privacy restrictions. This data must be labeled, and this usually requires manual effort by individual workers. It may also require domain expertise, for example, when labeling medical images. Mislabeled or garbled data may also need to be filtered out through manual inspection. Data engineers must also manage changing needs for a dataset, for example, in order to support additional languages.
    
\subsection{Speech Commands}

    Speech Commands is a keyword spotting (KWS) dataset developed by Pete Warden at Google, and we will consider it as a practical example of the steps required in TinyML data engineering.
    
    \textbf{Requirements:} It established a new standard for public, comparable keyword spotting research - previous datasets for KWS were often restrictively licensed or proprietary to individual companies. Speech Commands is available for anyone to use, and allows ML researchers to compare their ML algorithm’s performance on the same data. 
    
    \textbf{Collection:} Speech Commands contains thousands of recorded examples for 35 keywords, collected by over 2,600 volunteers with a variety of accents. All volunteers agree to have their voices redistributed. Data collection was done in the browser, since requiring installation of an app might discourage contributors.
    
    \textbf{Refinement:} As with any dataset, some data collected will be unusable (for example, if an incorrect word is spoken or the microphone gain was too quiet). Some automated techniques were applied (removing low volume recordings and extracting the loudest 1sec from 1.5sec examples), but the remaining 105,829 recordings were manually verified through crowdsourcing.
    
    \textbf{Sustainment:} Speech Commands has been expanded to 35 keywords in Version 2, from the original 25 words in Version 1. Care was also taken to ensure the same recordings remain in the same train, validation, and test splits across the two versions.
    
    The careful construction of the Speech Commands dataset has allowed keyword spotting researchers to compare the performance of different TinyML neural architectures on the same reference data. This benefits many aspects of KWS research, such as reproducibility.
    
\subsection{Crowdsourcing Data for the Long Tail}

    Speech Commands only contains English data, but many languages are spoken across the world. Paying contributors to record and verify keyword data for every language rapidly becomes cost-limited. Companies which pay to collect speech data may prefer to keep this data in-house and proprietary. An alternative approach is crowdsourcing speech data. 
    
    Community contributions have led to the success of other open-source projects such as Linux and TensorFlow, and this model can also work for dataset generation. \href{https://commonvoice.mozilla.org/en}{CommonVoice} is a Mozilla-led effort which seeks to attract community contributions for speech data. So far, over 50,000 volunteers have contributed speech data in 54 languages to CommonVoice. A key draw for potential contributors is the promise of bringing modern advancements in speech processing to underserved languages. 
    
    CommonVoice data is permissively licensed for anyone to use. The majority of data in Common Voice consists of full sentences read aloud by volunteers, and verified to be discernible by other volunteers through a voting process. CommonVoice also contains a single-word target segment dataset (in the same style as Speech Commands) which contains recordings in 18 different languages. 
    
    CommonVoice is an ambitious project to bring automatic speech recognition, voice-based interfaces, and other speech processing technology to the whole planet. Importantly, users can add support for new languages without needing software engineering expertise, as CommonVoice also provides tools for translating their data collection interface to new languages. For data engineers, CommonVoice provides a useful example of how to expand data collection to a worldwide scale.
    
\subsection{Repurposing Existing Datasets for TinyML}

    We’ve seen how challenging it can be to collect brand new datasets. ML research has exploded in popularity and there are many datasets available for a wide variety of tasks already. TinyML research can be accelerated by taking advantage of these existing datasets and repurposing them for embedded tasks. A useful overview of available datasets is provided by TensorFlow’s Datasets Catalog, which covers a wide gamut of machine learning problems, and are ready-to-use for training with the TensorFlow API.
    
    In the next section we will develop a vision-based TinyML tool for person-detection. The dataset for this task, Visual Wake Words, was developed from an existing popular dataset, Common Objects In Context (COCO). By specializing an existing dataset, you can avoid the need to collect data from scratch.
    
    We will also discuss transfer learning in the next section. Transfer learning allows researchers to avoid training models from scratch, since the same features are shared across many tasks in speech or vision. It is faster to transfer features than to train from scratch, and transfer learning also greatly reduces the amount of data which needs to be collected for a new task. Pretrained models can also be used to help curate data for new tasks - ImageNet was built in part using results from Google image search. You may even be able to generate data for your needs using a simulator.
    
    \subsection{Responsible Data Collection}
    
    Data engineers must consider how to reduce bias when collecting datasets. We have seen several potential sources of bias in speech recognition - for example, your data may not contain enough examples from various age groups, accents, or gender. Efforts such as CommonVoice attempt to tackle the problem of underserved languages. Data engineers must determine how to debias data collection efforts for a given ML task, and additionally ensure these debiasing requirements are maintained if an intermediary is used in data collection (for example, paying gig workers to collect data on Amazon’s MechanicalTurk)
    
% 1.7: Visual Wake Words  

\section{Introduction to Visual Wake Words (VWW) Application}


  You've now seen the full TinyML flow in the context of a Keyword Spotting application, noting particular challenges with preprocessing audio data, training TinyML models, and designing effective end-to-end systems and metrics. You then explored the many challenges with (and solutions for) creating an effective TinyML dataset.
  
  \begin{figure}
      \includegraphics[width=\textwidth]{TinyML/2 Applications/3-7-1-1.png}
      \caption{Machine learning workflow and pipeline. Begins with collection and preprocessing of data then the design and training of a model. These steps are done in TensorFlow. Then a model is optimized, converted and deployed using TensorFlow Lite. Finally inferences are made using TensorFlow Lite Micro.}
  \end{figure}  
  
  
  
  \subsection{What's the Focus in this Module?}
  
  In this section we are going to build on that learning by exploring the TinyML flow and data engineering again in the context of a Visual Wake Words application, focusing on some unique challenges presented by this computer vision application. Visual Wake Words represents a common TinyML visual use case of identifying whether an object (or a person) is present in the image or not.
  
  \subsection{What's New and Different?}
  
  Following what we did for Keyword Spotting (KWS), we will look at several interesting perspectives on the end-to-end VWW TinyML application pipeline. We will understand the characteristics of the camera sensor and how that affects our preprocessing pipeline. We will look into the volume of data generated by the sensor and how that affects our downstream ML workflow. We’ll explore a fundamentally different way of doing neural network calculations to make things tiny so they are more efficient on MCUs. Specifically, we will introduce Mobilenets, a new class of models that leverage Depthwise Separable Convolutions to use less memory and reduce the total number of computations. We’ll then discover a whole new way to train models, Transfer Learning, which builds off of an existing pre-trained model to drastically reduce training time. Finally we’ll explore the different end-to-end systems challenges and metrics present in this novel application. We hope you enjoy this section!
  
    \begin{figure}
      \GRAPHICS{1.0}{1.0}{TinyML/2 Applications/3-7-1-2n.png}
      \GRAPHICS{1.0}{1.0}{TinyML/2 Applications/3-7-1-3n.png}
      \caption{The Arduino Nano microcontroller and Arducam we will use in Course 3 for this application.The Arduino Nano microcontroller and Arducam we will use in Course 3 for this application.}
  \end{figure}  
  
  
  
  \section{What are Visual Wake Words (VWW)?}
  
  VIJAY JANAPA REDDI: Hi, In this video and the coming series of videos,
  we're going to start dipping our toes into our second application, which
  is Visual Wake Words.
  Visual Wake Words represents a common microcontroller vision use case
  where the goal is to identify whether a person or an object of interest
  is actually present in the image or not.
  As we did with keyword spotting, I'm going to start off
  by talking about the sensor.
  In this particular case, it's going to be the image sensor, which
  is going to be generating RGB channels of input data
  that we will be feeding into our neural network.
  At this point, I'm going to make sure that we
  go a little deep into neural network design.
  So far, we have been training on neural networks from scratch.
  In this particular module, I'm going to show you
  how we can shortcut that process in order to accelerate the training
  time by reusing some of the information that the network has
  retained from a prior experience.
  More on that later.
  Finally, we'll get into metrics and learning
  how to assess this particular application.
  Specifically, we will look at it from a system standpoint,
  as well as the overall quality of experience standpoint.
  As I've said before and as I'm going to say it again,
  we're going to focus on the end to end application level design,
  not just on the neural network part.
  Now let's get concrete and say what's an application of visual wake words.
  Here you're looking at a smart device where the machine is actually
  waking up only when it recognizes me.
  Now you can't see me in the picture, but that's the implicit assumption.
  Another example of this might actually be
  when you have a smart device that's perhaps
  sitting on the outside of your house.
  It's a smart doorbell.
  And the goal is this.
  It recognizes when someone shows up at a door,
  or perhaps it does not recognize someone who shows up at a door.
  In which case, it might send you an alert.
  Now clearly, both of these examples that we're talking about here
  are about processing the information that's coming in.
  Now these are applications you might already
  be familiar with in some capacity.
  What's interesting, of course, here, is that we're
  talking about untethering the device, which
  has remarkable potential for real world deployments.
  Think about trying to deploy a TinyML device in an office which
  is looking for occupancies.
  If there's no one who is occupying the room,
  perhaps you want the system to turn off the lights.
  Now that seems like a very simple thing to do.
  Well, let me tell you why TinyML version is different from a non-TinyML version,
  again, from a real world perspective.
  In the real world, if you want to put a device in any particular building,
  you often have to get permits, you have to get electrical certifications,
  you have to get permits in order to be able to actually drill
  the hole in there, because that building is not yours.
  So you're trying to install a new piece of equipment.
  So there's all these mechanics that kind of
  come into play that actually cost a lot, which has got nothing
  to do with the actual technology itself of simply, is there anyone in the room.
  That's where TinyML becomes extremely powerful from a practical standpoint,
  because I can take a small little device like this and I can just stick it.
  And then I can set it up to automatically communicate
  to a server on Wi-Fi and then turn off the lights.
  It's so much simpler, and this is why people are actually
  excited about TinyML even though some of these applications
  are things that you could typically do by tethering them to the wall.
  Let's go ahead.
  What's another interesting application?
  Smart glasses are a remarkably powerful application for TinyML.
  You've got to have all this remarkable sensing capability where
  you're processing the visual cues that are coming
  in where you're triggering for interesting actions
  to happen when that data comes in.
  For instance, here perhaps, the camera is actually
  picking up this tag of 30% off sale, for instance, and it triggers,
  and then the lens automatically focuses in on the image that's
  actually of interest in order to guide you there
  for effective retail shopping.
  Or perhaps you are using these lenses in your navigating through.
  Instead of having to look at your phone, your smart glasses
  are detecting critical waypoints and automatically projecting information
  into your lens.
  Remarkably powerful mechanisms that are going to transform the way we
  actually experience things around us.
  And TinyML is right at the forefront of those.
  And these are just a handful of example applications.
  So, we're going to learn about some of the challenges
  associated with trying to achieve that sort of future.
  But we're also going to discover some of the opportunities
  that we have around there.
  Specifically around challenges.
  As we did previously, we'll touch on performance-related aspects,
  like latency and bandwidth.
  We'll talk about model capabilities, in terms of accuracy and personalization.
  We'll talk about the criticality of the data itself,
  how do we keep it secure and private, as well as
  how do we deploy models so that they don't consume too much energy
  and make sure that they fit into resource-constrained devices.
  All of the same things that we've touched on for keyword spotting,
  however, I'm going to focus on the aspects that
  are actually quite unique and relevant to this particular application.
  Finally, we're interested in trying to figure out
  how we can get this model to actually run on device,
  and what are the capabilities that we can unlock if we do that.
  We're not going to deploy the model on the device yet.
  We're going to try and shrink the model down through all our techniques
  so that it's ready to fit on device which we will do in course 3.
  So that said, what are we going to learn?
  We're going to learn a little bit about the data collection process, which
  I'm not showing here, because we're not emphasizing it as much, because we've
  already spent a lot of time talking to you about what it means to actually get
  the requirements for your data, then gather it, then validate it, and then
  actually sustain it.
  So I don't want to harp too much on that.
  Hopefully, you know a lot about it by now.
  But we'll focus on some of the interesting aspects
  around training the model and learning how to optimize it for deployment.
  We're not going to focus on the aspects in deployment
  from inference standpoint, yet.
  We'll get there.
  For now, what we are going to focus on is primarily
  about getting the model down with the right characteristics.
  And we'll do this in Google CoLab.
  And finally, when we actually finish all of this,
  we'll get around to deploying things onto the real hardware.
  With that said, let's get our feet wet and get
  started with learning how to build a TinyML application around Visual Wake
  Words.
  
  \section{Forum}
  
  What unique challenges do you think will be present in the Visual Wake Words application? What makes this different from Keyword Spotting? Do you think it will be harder? Easier? Why?
  
  \section{Visual Wake Words Challenges}
  
  VIJAY JANAPA REDDI: What are some of the interesting challenges that
  come about with visual wake words as compared to keyword spotting?
  That's the key thing that I'm going to try and focus on
  in this particular video.
  Specifically, I want to highlight the interesting attributes
  that are there around the sensor and what
  the sensor is going to do in terms of feeding the network the data
  and how there's a tight coupling between the two
  that you need to understand in order to deploy a practical TinyML
  application that deals with visual data.
  Let's go back and do a quick recall.
  This is a figure that you might recognize
  from keyword spotting where there's a core audio block that
  is running a voice trigger detector for something like Alexa,
  or OK Google, or something to that effect.
  In this particular case, it happened to be Hey Siri.
  And what it's doing is it's actually using a cascade architecture--
  I hope you remember what that means--
  in order to figure out whether to send the data to the cloud.
  There's a TinyML model that's running on the microcontroller which,
  if it's got enough confidence, might actually
  wake up a much larger model in order to get a much more
  precise or accurate inference.
  And then, based on that decision, together they
  will decide whether to go to the cloud.
  Let's take the scenario where you actually
  want to use the cloud for visual information
  and see how things work out in the context of an image sensor.
  That's going to be the key focus here.
  Let's focus on offloading the work to the cloud.
  And let's try and understand what the implications would be.
  If we look at a simple little image, let's
  say we want to build a detector that can detect cats.
  That's what you smart doorbell is going after.
  I want to know if there's a cat in front of my house.
  OK?
  Well, an image in neural networks is around 224
  by 224, sometimes 300 by 300 pixels.
  That's typically the resolution that most of the endpoint
  related TinyML applications are dealing with.
  So you tend to have a camera that has a high resolution image,
  but you down scale it down into a very small image.
  Because neural networks are very good at picking up
  features, even at those small dimensions.
  If it's a 224 by 224 image, that means you have 224 by 224 pixels times 3.
  Because there are three channels, since it's a color image.
  There's a red channel, there's a green channel and a blue channel
  for each one of those primary colors, which then creates our compounded
  picture here.
  And if we assume that each of those pixels
  is occupying 4 bytes, like a floating point value,
  then it comes out to be around 602,000 bytes.
  Now let's talk about what this means in the context
  of uploading this data to the cloud.
  OK.
  So here are some very typical characteristics
  of what a typical network might look like.
  You have a ping of about 25 milliseconds.
  This is the latency just to be able to send something to your local gateway
  and be able to get a response back.
  That's a fixed sort of latency.
  Then you've got the download speed and the upload speed.
  Download speed is often much faster for internet than your upload speed.
  Either way, download speed's about 35 megabits per second.
  Upload speed is about 4.62 megabits per second.
  Now let's go back to our image.
  We have about 602,000 bytes of data that we would like to actually upload.
  If we do the math on the upload speed of about 4.6 megabits per second,
  it's around 570 kilobytes per second.
  Given that we have 602,000 bytes and we've got 570,000 bytes,
  that's going to take about one second in order
  to transfer this little image across.
  Now let's kind of look at that from an end to a network perspective.
  You've got to upload the data, which is going to take you about one second.
  Then you've actually got to do inference,
  which tends to be really small once you're inside the data center,
  inside the cloud.
  Then you've got to pull that response back.
  For example, the question here is that I take an image,
  I want to know whether there's a cat in the image or not.
  Well, it's going to take one second to upload, 10
  milliseconds to do the inference, another 100 milliseconds
  to come back to the machine to say yes or no.
  That's about 1.2 seconds.
  Now this can be pretty bad, because in certain cases,
  if the object of interest moves away within the one second, well,
  by the time you've gotten a decision and you say,
  OK, take a picture, because there is a cat in the--
  well, the image might have changed.
  OK, fine.
  The cat might not move so fast.
  But in certain other cases, depending on your application use case,
  you could totally imagine that this latency is not tolerable
  and therefore you would end up with poor quality of service due to the lag.
  Let's kind of look at this in the perspective of keyword spotting.
  For instance, when we were previously looking
  at the spectrogram, which is shown on the left here,
  it was a much smaller image.
  It's a 49 by 40 two-dimensional image.
  While we're showing you a color image here, when you did your Colab,
  you're only dealing with one channel.
  And let's assume that you were having 4 bytes,
  because you're using floating point values in order to train the network.
  Well, that's still pretty small.
  Because that's only about 7,840 bytes, as
  compared to the cat image, where it's about an order of magnitude larger.
  So the concern here is that if you look at it from a sensor standpoint,
  where you're looking at an audio signal versus a video signal, two
  different types of sensors.
  The volume of data that an image sensor is going to produce
  is going to be substantially more.
  And that's where we're going to get into trouble.
  Lesson number one is that the sensor and the data that it produces
  are crucial to understand.
  Because the more amount of data you have coming
  from a sensor, the more implications it has on terms of having always
  On machine learning models
  Specifically, more data means higher latency
  if you are going to go to the cloud.
  It also means higher power consumption, because you're
  going to have to transmit the data.
  So that's not a good thing.
  Because either way, you end up with lower user satisfaction.
  At that point, you killed a end to end application use case.
  So what I'm trying to get you to understand here
  is about how the sensor plays a very critical role in terms of the data
  it's producing for us to be able to think about the networks.
  All right.
  So let's say we don't go to the cloud.
  And instead, we do on device processing, TinyML style.
  Then what are some of the issues we're going to run into?
  We obviously have to worry about the processing latency
  because, as you know, we're running things
  on a microcontroller, which means that we have to make
  sure our models are kind of small.
  Now remember that, as compared to the data
  centers, the amount of computational horsepower
  that you have on your microcontroller is orders of magnitude smaller.
  And furthermore, you have a huge latency to be able to get to that server
  if you want to be able to run the model there.
  So you might have the beautiful, nice, big, precise model there.
  But by the time you get there, you've already lost a lot.
  Because that latency to get to that actual computational horsepower
  is too high.
  So you want to do things on device if it's possible.
  The second thing, obviously, is around memory.
  At this point, you know that we don't have too much memory.
  I know you know that.
  But what I do want to point out is in the context of visual data,
  where you are building visual models now to recognize things
  or to identify things, typically, the visual models are much larger.
  Which is why I'm highlighting some of the state of the art models
  here and showing you that they tend to be in the order of megabytes
  in memory size versus your microcontroller
  is typically around kilobytes of memory.
  So we have to somehow bridge this gap.
  So we need to learn how to actually train small, little models, which
  is not necessarily an easy thing to do.
  But we'll see how to do it.
  OK.
  And finally, the last thing is yet again about the accuracy of the model
  and how well you were able to perform your task.
  You ideally want to be able to run the model and achieve good accuracy.
  But you have to still worry about subtle things,
  like false positives and negatives.
  This is an image showing you about how the Dazzle Club in London
  is actually using interesting shapes and lines in their faces
  in order to trick cameras from recognizing
  that there's a human in front of them.
  This happens because the camera and the neural network
  running there don't know any better.
  They're just looking at raw pixels.
  So as long as, like, the network is not able to tell
  that there's a certain sort of distance between the eyes
  or what the distance is between your nose and your eyes and so forth,
  you can totally do that masking by just adding
  new lines to your faces and the cameras are not able to perform well.
  So you have to think about the false positives and negatives, which we
  talked a lot about in keyword spotting.
  The same things come about here, but just
  from a totally different representation perspective.
  As I've always said before, any kind of computing system
  has a whole bunch of things that it's running.
  It's not just in the neural network model.
  It's got a full stack.
  Of course, this stack is quite comprehensive,
  because it's got an operating system kernel, it's got middleware,
  and so forth.
  None of which are available or there in the context of a TinyML device.
  That said, though, you still have to think
  about things more than just the model.
  Because you have to think about what does it
  mean for the inference engine that is supporting
  this particular visual model, how much memory do
  I need in order to house the visual images that
  are coming in from the camera.
  Because all of that makes your TinyML application.
  So you've got to think about the model as well as the whole rest
  of the stack that needs memory.
  And the same thing is also true when it comes to latency.
  Because you've got to run the whole stack, not just your neural network
  model in order to build a proper application.
  So I've touched on a couple of interesting challenges
  that we have that are unique to visual wake words when
  we look at it from a sensor standpoint as an image sensor compared
  to an acoustic or audio sensor which we talked about previously.
  So take a moment to kind of reflect on some of these crucial differences.
  Because I will be testing you on them.
  With that said, we're going to next move on
  into understanding how do we actually start
  building out a TinyML application.
  
  \section{Visual Wake Words Dataset}
  
  VIJAY JANAPA REDDI: In this section, we're
  going to kick off the TinyML pipeline for visual Wake Words
  by starting off with data collection.
  As I've always, said data is king.
  So we're going to start talking about data collection
  and pre-processing before we go on to the rest of the parts.
  As with any data, always remember to identify the key requirements
  of the license file that you need to have
  that corresponds to how you intend for the model
  to actually be used in the field.
  I've already talked about this, so I'm not going to re-harp it,
  but it's an extremely critical thing.
  Why?
  Because if you don't read the licenses properly for your data source,
  you can get into a lot of trouble.
  Let me be specific.
  Here's a person who went out and scraped Tinder for about 40,000 images
  to create a facial data set for AI experimentation.
  Sounds like an awesome thing-- you're reusing
  some existing data that's out there, and then you're building a custom data set.
  Well, this is what Tinder had to say officially.
  "This person has violated our terms of service,
  and we're going to take appropriate action and investigate further."
  Well, that's not a good thing because what
  this means is that if anyone happened to build anything that
  is interesting with the tiny ML model into an application that happened
  to be using this data set, well, they can no longer be making
  any business because the data set was polluted, which means you
  don't have the right to use the model.
  Again, I keep harping on this point because it's extremely crucial that you
  understand that the data is king, and that the network, whatever it picks up,
  is dependent upon that data.
  So make sure your data is clean.
  So far, we've talked a lot about data collection.
  We had an entire module on that, and I said,
  if you want to be able to build a good data set,
  you've got to go through a set of stages.
  I said, you got to gather the data after you've identified the requirements,
  and then you've got to figure out how to validate that data.
  And then, make sure you can sustain that data collection pipeline.
  OK and I said that that's really costly, if you want to do it at scale,
  or they have to rely on a whole bunch of crowdsourcing-based methods
  like we did with Common Voice, for instance,
  or like we did with speech commands.
  Well, there's another one that I left out intentionally
  that I wanted to bring up in this particular part of the class, which
  is that you can actually reuse existing data in order
  to create a subset of the data.
  So let's assume, for instance, I have ImageNet,
  which has got a whole bunch of different kinds of classes in it,
  about a thousand different classes, but in TinyML,
  given that you're typically caring about one particular domain, so for example,
  in our doorbell example, it was a cat. .
  So what you really need is just cat pictures In that particular case,
  I don't need a whole bunch of label data that's
  irrelevant to my application use-case.
  I might be able to reuse that existing data with only images of cats
  and balance it out with other images that are not of cats.
  So there's this potential to reuse the data.
  Now this is not new.
  This is exactly how, in fact, Pete actually
  helped build the visual Wake Words data set.
  And the visual Wake Words data set has a single goal,
  which is effectively looking at trying to identify a person in a picture.
  That's what the whole visual Wake Word's data set is all about.
  But the critical thing is really about using publicly available data sets
  in order to create new data sources.
  All right, let's take a concrete example.
  Here's an image that comes from the MS COCO data set.
  What this data set has is a whole bunch of different kinds
  of annotations for many different people, objects, and so forth.
  And specifically, by looking at this image, what the data set
  tells you is that there is a person in here,
  which has got an annotation around him, with a bounding box.
  Well, if you know this exact location in the bigger image,
  then you can crop this image out to just create yourself a new, very specific
  image with one person.
  But hey, it happens that there's another person in here.
  So you can crop that person out and then label that person as them.
  So this way, you are effectively, dynamically creating a new data set out
  of the original data set.
  Now this is a remarkably powerful thing.
  It's a very simple concept, but as long as the licenses are OK,
  you are good to go.
  Because you are effectively now paying the cost up front for a very large data
  set, but then you are able to dynamically create specific instances.
  For example, here we can create a person detection data set.
  We might be able to create a cat detection data set and so forth.
  There are some challenges you have to watch out for, though.
  Think about it.
  The original data set was collected with something in mind.
  There was a very specific goal in mind when that data set was collected.
  Is that data set really going to meet the needs of your particular TinyML
  application?
  It may, but it also may not.
  For example, if the data set that you need is something about planets,
  being able to detect planets and so forth, well, if you're
  imaged in that data set is all about cars, trucks,
  and so forth-- about all the objects that are on this particular planet,
  it's not necessarily going to help you with looking at planets.
  Because that data simply doesn't exist in this data set.
  So you have to worry about that, whether the data actually is there.
  The second thing is whether there's a bias
  because you have to look at the distribution of the samples
  and make sure that the data that you are interested in
  is actually going to be balanced.
  We talked a lot about this in the data engineering
  section of this particular course.
  And make sure you've got all the relevant images,
  and also make sure you've got the right quality and quantity of the images.
  So the point I'm trying to emphasize here is that be careful.
  Don't just jump in and say, hey, I've learned about this new method
  about just using this existing data set.
  Read the licenses.
  Make sure it's got the right data.
  Do your analysis on the data before you start creating the subset of the data
  in order to train your TinyML application.
  We were effectively going to be using the visual Wake Words data
  set which shortcuts the whole process of building a person detection data set.
  And that's why I'm introducing that particular data set in this course.
  With that said, let's go on to the second stage of building up
  the TinyML application.
  
  \section{Reading}
  
    \begin{figure}
      \includegraphics[width=\textwidth]{TinyML/2 Applications/3-7-1.png}
      \caption{Data Privacy with Images.}
  \end{figure}  

  
  The growth of social media has heralded an era of online data sharing, wherein the majority of daily interactions are conducted between technological intermediaries. While this provides huge benefits in terms of productivity, communication, and accessibility, it also presents important privacy challenges. 
  
  \subsection{A Brief History of Data Privacy}
  
  Since the human rights movement following the events of World War II, strict ethical requirements were developed which must be adhered to by social scientists and medical practitioners when performing experiments on humans. The prime example of these is the Belmont Report, created in 1978. This is true regardless of their impact, be it physical, psychological, or emotional. Although data science and machine learning practitioners do not directly experiment with human subjects, their impact can be commensurate with those of the social science community. Every time an e-commerce site changes the layout of their site using an A/B test, they are conducting a large-scale social experiment. 
  
  One of the key concerns of such requirements is the right to individual privacy. The breadth of data that can be obtained about human subjects can be highly valuable to businesses as well as the research community, but requires informed consent of the individual to be used. Today, companies like Google and Facebook have huge amounts of data based on the way users interact with their services. This information is often then marketed to companies and subsequently used for targeted advertising campaigns. The right of the company to use this information is embedded within their terms of use for their respective services. Whilst these uses might seem unethical, especially to privacy advocates, they are entirely legal and abide by the ethical conventions of Belmont Report, as well as more modern ethical requirements developed during the age of big data, such as the Menlo Report. 
  
  However, the principle of informed consent has become increasingly difficult to ensure in the modern world. A good example is genetic information. By uploading genetic information to online platforms such as 23AndMe, an individual is voluntarily waiving their right to privacy. However, genetic information is largely similar between family members, and thus by uploading this information, the privacy of relatives is also violated. A famous example of this was the arrest of the “Golden State Killer”, who went uncaught for decades until DNA information from an online genetic database was linked to a relative of the murderer using DNA from the crime scene. Although this demonstrates a positive use of such data, its potential power is unsettling.
  
  \subsection{Data Privacy in Images}
  
  Image data is one of the most vulnerable mediums of online data. This is troubling as it is also one of the most commonly shared. Datasets may contain images of people curated from online resources wherein the user did not obtain informed consent from the individuals who own or are present in the images. Perhaps the most alarming feature is that of specific types of image which store geographical coordinates of the location a photo was taken, such as GeoTIFF or EXIF.
  
  Most individuals would shrug this off as being harmless, since it has very minor ramifications to their personal life. However, with sufficient knowledge this information can be used to obtain a great deal of personal information. Images uploaded by individuals which are then voluntarily shared can present privacy violations to other individuals present in the image. This information can be used to find personal associations, and combined with image data to determine commonly visited locations or the homes of an individual or their relatives and friends.
  
  This issue is well highlighted by a somewhat amusing and yet concerning study by two Harvard undergraduates looking at images on the Dark Web. By studying images of drugs scraped from websites on the Dark Web (websites only accessible using onion routing on a specialized web browser called Tor), the undergraduates were able to produce a map and isolate the approximate locations of drug dens. Most smartphones today attach this geographical information to pictures automatically, unbeknownst to the average user, and is easily extractable from a scraped image. 
  
  \subsection{Relevance to TinyML}
  
  As data scientists and machine learning practitioners, it is important for us to find ways to uphold the ethical foundations set by our forebearers. One of the best ways for this to be done is to understand the ethical ramifications of our actions, and to try and safeguard against situations where such transgressions may manifest. For example, when curating a publicly available dataset, ensuring that images are within the public domain or by entering into a third party agreement with the original data provider. If the dataset is to be used for computer vision, the geographical data provides no utility, and thus should be removed in order to minimize potential harm to the individuals that produced the original data. Similarly, if images contain personal or revealing information about an individual, they should not be used in a dataset. As a rule of thumb, the minimum amount of data should be utilized to obtain the desired result. If the provenance of a dataset is uncertain or the licensing ambiguous, it is important to err on the side of caution and either confirm the provenance or disregard the data.
  
  This is true not just for during model training but also inference. In our TinyML visual wake words application, we use a camera which is able to take images in real-time. These images are transferred to a framebuffer which then performs inference on the newly obtained image. In the majority of TinyML applications, it is infeasible to obtain informed consent for each person that may potentially fall inside the image, presenting an ethical dilemma. For example, if an entity would like to save images from a smart doorbell for use at a later time, either for future training data or for diagnostics, it may violate the above-mentioned ethical principles. These images may also be saved with metadata, which may further compromise privacy. Thus, we must be mindful when developing these systems to ensure that the obtained images are only for real-time detection and are not archived, or are covered by additional provisions.
  
\section{Neural Network Architectures for Visual Wake Words}

VIJAY REDDI: In this video, we're going to learn what it means
to design a visual wake words model.
Now before we do that, let's take a step back and look at where we are.
We talked a little bit about collecting data and preprocessing
the data for an image-based task.
Now I'm going to talk about designing a model before we
jump into training the model.
Why do I want to do this?
I want to do this because when you look at image classification with a task,
there's lots of different neural networks out there.
This graph is showing you the variety of different models
that are out there that give you accuracy versus performance latency
trade off curve.
So for example, if I take ResNet-152, I can have very good or very high
accuracy, but then I have to pay a fair bit
in terms of the computational requirements,
because it demands over 20 giga operations
to make a single forward pass.
The reason I'm talking about this is just a few years ago,
there was this big challenge that we faced,
which was about how do we take these models
and shrink them down so that they were actually
running efficiently not on some big cloud computing server,
but on smart handheld devices.
Because if we try and run these big models on these small devices,
they're going to run out of battery really fast.
So I want to take you on a little bit of history
so that we can learn something from the history
in order to be able to design models better for TinyML-specific tasks.
So if you want to shrink models down, there are two fundamental things.
You want to be able to shrink down the number of parameters
that are actually inside the model.
The size of the circle here corresponds to the number of parameters.
So for example, VGG-19 has over 155 million parameters.
You want a network that is small and parameters.
The other thing you want to do is you also
want to make sure it runs efficiently.
You don't want it to need a lot of competition.
So those are the two fundamental legs.
There's a number of model parameters, and then there's
a number of hardware operations you can perform per second.
These two together will actually give you
a small network that is actually run efficiently on a mobile device
or TinyML device.
So let's try and understand how we can break a complex model down
so we can try and shrink things down.
I got to take you back to the fundamentals.
When you look at an image-based task, you spend a lot of time
doing convolutional operations.
Almost all the different networks that I showed you,
they all use convolutions at the heart of it.
This is a slide from course 1, where we looked
at applying a filter over the image on the left,
and you end up with the image on the right.
What is really happening behind the scenes
when you do this convolution is extremely important to understand
because it is the foundational building block that has consequences
down the road.
So if you want to be able to shrink models down
or you want to scale things down, you really
have to understand this fundamental building block.
If you look at a Standard Convolution operator, what you're really doing
is that you've got a filter, the yellow here, that
is scanning through the original image, input image on the left,
and it's producing the corresponding blue pixel
on the right for the set of pixels that it's looking at on the left.
I have a 3 channel RGB image, a 3 by 3 by 3 filter
that is scanning over all of those different channels at the same time
and outputting a single point.
To be more specific, input feature map is 8 in width, 8 in height,
and it's got three channels.
The kernel that's running over it is a 3 by 3 by 3 kernel.
That's that yellow kernel that we're looking at.
And the output is that single orange dot that we're
seeing as the kernel is going through the image on the left.
The question is this-- this standard convolution is quite complicated.
But effectively, is there some way to break the convolution down
into simpler operations rather than doing one massive compound convolution?
The reading that I have breaks down the math that tells you
the exact number of multiplications and additions
you're going to be performing for this standard convolution.
And it's quite strongly dependent upon the number of challenge you have,
the number of filters you have, and so forth.
Please pay attention to the math in there.
The thing that I'm leading into is can we break this standard convolution down
into smaller pieces so that's much more efficient to run on the hardware?
So that's where depthwise convolution comes in.
The depthwise convolution is this fundamental idea
that instead of doing a compound convolution,
what if I kind of took the individual channels out
and I did a convolutional on each channel individually?
And after that, I do a pointwise convolution
to fold all of those output feature maps into a single point, which
is exactly where the original standard convolution would have gotten us?
So I'm trying to combine depthwise and pointwise separable convolutions
together.
So instead of doing one compound convolution,
I'm taking two simpler steps--
a depthwise convolution across the channels and then a pointwise
convolution to produce the output single value.
So if we do this, one of the key things that we get out of it
is that we end up reducing the number of multiplications and additions
that we do.
Now I'm not going to get into the math here.
The reading does a beautiful job of explaining this,
so please pay attention to it because it's extremely important that you
have an appreciation for the difference between a standard convolution
and a depthwise convolution or depthwise separable convolution,
because you will see that it can dramatically cut down
the amount of computations needed.
And that is the fundamental foundational block
of making improvements to reducing model size
and improving the latency both at the same time.
Here is a ratio that you would have if you
were to compare a depthwise separable convolution to a standard convolution,
just to put things into perspective.
If, for example, I had 100 filters that I was applying
and I had for each of those filters, the kernel size
was effectively 512 dimensions by 512 dimensions,
just to give you an example, it would be 100 times more
efficient to do a depthwise separable convolution
than a standard convolution.
Do the math for the examples and the reading,
and then you'll see why the difference is so large.
And this is exactly what came out of Google a couple of years
ago when a couple of these engineers there sat down
with the original convolutions, and they were scratching their head, and saying,
how can we reduce this competition?
And voila.
They said, what if we could split this standard convolution up
into dealing with the individual channels,
and then do a coupled pointwise convolution operator at the end?
That's what led to mobile nets, which is a very specific type
of convolutional neural network that has taken off remarkably,
and is now being deployed widely across many different devices,
because it's dramatically reduced the computational requirements, as shown
here.
You can see that the size of the bubble is small,
because it has fewer parameters inside the network.
It also is very much to the left, which means
it is because it requires little computation,
and it achieves a fairly good amount of accuracy.
In recent years, we've made a number of improvements over it.
But the point is that it actually achieves remarkably good efficiency.
That was for a state of the art smartphones back in the day.
So if you compare this baseline, which is still fairly large, because it's
16 megabytes compared to our dinky little kilobyte amount of memory
that we have on our microcontroller, we're about 64 times off.
So we need to do something in order to shrink this model, which is already
quite small, much, much smaller.
That's where there's some subtle optimizations that we can
do on the baseline mobile nets model.
For example, we can play around with what's called a depthwise multiplier.
The depthwise multiplier is a way to reduce the total amount of computation.
This equation, which is presented in greater detail in the reading,
explains the convolution as it's performed in two steps.
The left-hand side corresponds to the depthwise aspect of the convolution,
the right-hand side is doing the pointwise convolution,
as I mentioned earlier.
You're combining these two individual steps,
and therefore, you have an addition in the middle.
Now within each of these individual steps, you can play a couple of games.
As you're doing the convolution, you can control the number of channels
that you're looking at within each layer.
By controlling that, you are effectively going
to be able to control the number of different computations
that you will have to do across the entire network.
Reading has more detail, but the critical message
is you can control this parameter alpha between 0 and 1.
And this has great benefit in terms of cutting down both the size of the model
as well as the computation requirements from the model.
This table here shows you a variety of different pieces of data.
Let me walk you through it.
So the first row is showing you the baseline model for mobile nets.
What this is showing is that the baseline achieves about 70.7%
in accuracy, does really well.
And what's interesting, though, is that if I change the alpha down by scaling
it down by about 25%, I can drop a little bit on accuracy,
but then I can get a real bang for the buck
in terms of the number of parameters that are actually inside the network.
By reducing the number of parameters, I'm effectively
able to shrink the network down.
And as you look at different values of alpha, what you see
is that you can make a network down quite small
as long as you're willing to trade off a certain level of accuracy.
In addition, you can also see that the amount of computation
associated by scaling down alpha starts decreasing.
And this is, again, because you are limiting the number of challenge
you're effectively looking at.
So this is one way of controlling the knob.
Another way of controlling the knob is actually deciding how big or small
is the image input going to be.
Previously, I showed you a cat image, which
is 224 by 224 Pixels we can shrink that down to be even smaller, 128 by 128.
Of course, you're going to take a hit on accuracy, but look at the number of max
or multiply accumulates that you have to perform.
It's a remarkable drop in there.
And that's huge compared to having the same alpha.
That's the critical message that I'm going after here.
Now please keep in mind, though, these are all different kinds
of optimizations.
There are a whole variety of different optimizations you
can perform across a number of channels you applied,
the number of filters you're applying, what the alpha value is,
and the whole bunch of different hyper parameters you can tune.
And that's effectively what's called neural architecture search.
Neural architecture search is a way of trying
to find the best combination of parameters
through the sea of different parameters to achieve the smallest model that
gives you the best accuracy.
Now I don't have an answer or a magic bullet that kind of tells you
what that perfect type of parameter search scheme is going to be.
People have come up with many different kinds of schemes.
But the critical message here really is this.
If you want to design a model for neural networks
that fit on the resource-constrained devices, number one,
I want you to think about more than just the number
of layers you're putting together.
I want you to try and understand that you can actually go into it
and understand what the mechanics that are determining the parameter sizes.
So how can I shrink the number of parameters there?
How do I also make it run more efficiently
by understanding the operators that are actually being used?
Now, this is a bit more advanced, there's a lot more to it.
And as we go into course 3, I'll try and give you a little bit more
detail into these kinds of things.
Here, I'm trying to give you a flavor for carefully designing
the models more than just looking at the high level structure
of the overall neural network.
And there are lots of nuanced parameters we can play with.
And that is key in terms of making the model smaller.
With that said, in the next video, let's pick up
with getting ready to train the neural network model.

\section{The Math Behind MobileNets Efficient Computation}

\subsection{Introduction}

Now that you’ve seen that MobileNets drastically reduce the size of computer vision neural networks let’s dive a little deeper into how their key innovation works: Depthwise Separable Convolutions.

\subsection{Standard Convolutions}

In order to better understand how depthwise separable convolutions are different from standard convolutions let’s first re-examine standard convolutions. In particular, we are going to quantify the number of multiplication operations and parameters in a standard convolution.

In the last course, when we considered filters in detail we only looked at filters applied to grayscale images like the vertical line filter shown below:

\begin{figure}
    \includegraphics[width=\textwidth]{TinyML/2 Applications/CNN}
    \caption{A convolutional filter to recognize vertical lines and the result of applying it to an image of people walking up an outdoor staircase to highlight the vertical lines in the image.}
\end{figure}  


These images can also be called single-channel images as they only have one set of pixel values. In contrast, most color images are three-channel RGB images. Where there is a value representing how much red, green, and blue is in the color. As such the filters (also often referred to as kernels) used are not simply a matrix but instead a tensor as it needs to multiply all three channels at the same time. This tensor operation is shown below. In this example, we have a 3x3x3 kernel being convolved with a 9x9x3 image to produce a 7x7x1 output. The highlighted squares are the inputs and outputs of the last convolution operation.

\begin{figure}
    \includegraphics[width=\textwidth]{TinyML/2 Applications/CNN2}
    \caption{A visual depiction of the convolutional process showing that a 3x3x3 filter can be applied to a 9x9 by 3 channel RGB image resulting in one output in a 7x7 output feature. This is essentially the same process as before when we applied convolutions to a 2D grayscale image, we just now also sum across channels.}
\end{figure}  


Note that in general, each convolution requires  multiplications due to the size of the kernel and the number of channels in the image. Then to produce a single output we need to do these operations    times to produce the full output. Also in general we don’t just use a single kerel, we use N kernels. Multiple kernels are referred to as a filter. A filter is a concatenation of multiple different kernels where each kernel is assigned to a particular channel of the input. As such the total number of multiplications will be: 

$$D_K*D_K*M*D_F*D_F*N=D_K^2*M*D_F^2*N$$


\subsection{Depthwise Separable Convolutions}

Depthwise Separable Convolutions instead proceed in a two-step process. First, each channel is treated independently as if they were separate single-channel images and filters are applied to them creating multiple outputs. This is referred to as a Depthwise Convolution. Next, a Pointwise Convolution is applied to those outputs using a 1x1xC filter to compute the final output. This process can be seen below where we again take our 9x9x3 image and apply three separate 3x3 filters to it depthwise to produce a 7x7x3 output. We then take that output and apply a 1x1x3 pointwise filter to it to produce our final 7x7x1 output.

\begin{figure}
    \includegraphics[width=\textwidth]{TinyML/2 Applications/3-7-9.png}
    \caption{A visual depiction of the depthwise seperable convolution. Showing that a 9x9 by 3 channel RGB input image when applied to a 3x3x3 filter now first produces a 7x7 by 3 channel output feature through the depthwise convolution and then that is filtered with a 3x1 pointwise convolution to produce the 7x7 output feature. This two step process differs from the standard one described above as it uses the second pointwise filter to do the summation of the various 2D convolutions.}
\end{figure}  



Note that in general now each depthwise convolution requires M filters of $D_K \cdot D_K$ multiplications and to produce a single output we need to do these operations $D_F \cdot D_F$ times. Therefore we need $D_K \cdot D_K \cdot M \cdot D_F \cdot D_F$  multiplications for that stage.

For the pointwise convolution, we now use a 1x1xM filter  $D_F \cdot D_F$ times. In general, just like regular convolutions, we don’t use a single filter we will use multiple filters. During Depthwise Separable Convolutions those multiple filters occur in the pointwise step. Therefore if we had N pointwise filters we will then need $M \cdot D_F \cdot D_F \cdot N$ total multiplications for this stage.

Summing the number of multiplications we will need in both stages we find that in total we need: 

$$D_K*D_K*M*D_F*D_F+M*D_F*D_F*N = M*D_F^2*(D_K^2+N)$$

Comparing the two kinds of Convolutions

We can compare the two kinds of convolutions through a ratio of the number of multiplications required for each. Placing standard convolutions on the denominator we get:

$$\frac{\text{Depthwise Separable}}{\text{Standard}} = \frac{M*D_F^2*(D_K^2+N)}{D_K^2*M*D_F^2*N}$$

$$\frac{\text{Depthwise Separable}}{\text{Standard}} = \frac{D_K^2+N}{D_K^2*N}$$

$$\frac{\text{Depthwise Separable}}{\text{Standard}} = \frac{1}{N} + \frac{1}{D_K^2}$$

This means that the more filters we use and the larger the kernels are, the more multiplications we can save. If we use our example from above where $D_K=3$ and we use conservatively only N=10 filters we will find that the ratio becomes 0.2111 meaning that by using Depthwise Separable Convolutions we save almost 5x the number of multiplication operations! This is far more efficient and can greatly improve latency.

Also, note that in the case of standard convolution we have $D_K^2 \cdot M \cdot N$  learnable parameters in our various filters/kernels. In contrast in the Depthwise Separable case, we have $D_K^2 \cdot M + M \cot N$. Again if we take the ratio of the two we find that:

$$\frac{\text{Depthwise Separable}}{\text{Standard}} = \frac{1}{N} + \frac{1}{D_K^2}$$

This means that we also have a much smaller memory requirement as we have far fewer parameters to store!

There is a tradeoff however, in improving our latency and memory needs we have reduced the number of parameters that we can use to learn with. Thus our models are more limited in their expressiveness. This is usually sufficient for TinyML applications but is something to consider when using Depthwise Separable Convolutions in general!

Finally, if you’d like to read more detail about MobileNets you can check out \href{https://arxiv.org/pdf/1704.04861.pdf}{the paper describing them here}.

\section{Transfer Learning for VWW}

VIJAY JANAPA REDDI: In this video, we're going
to focus on what it means to train a visual wake words model.
Specifically, we're going to look at what it means to train a visual wake
words model in an efficient way, in a way that we can cut down the training
time.
Before we go there, let's take a look at where we are right now.
We talked about collecting data.
We talked about pre-processing and designing a model.
No matter how well you design the model, be it big or small,
you've got to train the model.
Training the model can always be tedious because often,
you're dealing with lots of data.
For example, ImageNet has 1,000 classes, and for each of those classes,
there are 1,000 images, which means there's lots of data there.
And you've got to crunch through all of that data.
And that requires a massive amount of computational horsepower,
both from memory, the compute requirements, the server requirements.
And that costs both time and energy.
And typically, when you have a large dataset,
you have to rinse through all of that data over and over and over.
And epoch here is one complete pass through all of the data,
and then you have to repeat it over multiple epochs.
That's typically how you train the network.
At the end of the day, you've got a network
that has learned a key set of weights.
That's what you're really trying to get to.
Now, we want to ask ourselves this question about,
do we have to train the network from scratch every single time?
Before we get into the neural network mechanics,
I want to focus on some very basic intuition.
Our brain has a neural network, right?
The question is this, let's say the first time
we learn to drive a car, what did you do?
You went through the tedious mechanics of learning what a car looks like,
how to operate it, what are the rules that you have to follow, and so forth.
And after a lot of time and practice, you went and got your driver's test.
Now, let's say a couple of years down the road, you upgrade your car.
Do you go back and get another driver's test?
Or do you somehow just seemingly move and adapt to that new car?
It's different, but why is it that you don't need a driver's test?
Because it's got four wheels.
It's got a steering wheel.
It's the same exact mechanics.
Sure, it's got a little bit of its own quirks,
but that's something you can easily adapt to.
When you look at a neural network, that's
the same kind of thing that's going on.
The first part of the network's generally
learned some sort of general features, like in the case of a car,
that it's got wheels and so forth.
And the second part is really the one that kind of picks up
the very task-specific features.
So the question is this--
as we need to train the models, do we need to retrain the whole thing,
or do we just sort of take a pre-trained network that
has gotten exposed to a wide variety of class
and then specialize it for something that is of interest, or in our case,
maybe just detect a person or a cat or dog and so forth?
Well, that's where transfer learning really kicks in.
Transfer learning is built on that observation
that the first part of the network generally
learns these basic features that are typically independent of the task.
So we can reuse those weights that are in there from a pre-trained network.
And then, for the task-specific features,
we effectively allow the network to train.
So the first part of the network layers, we freeze all of them.
We freeze the general feature extraction aspects.
And then we only train the last couple of layers.
That allows us to very quickly take a pre-trained network
and dramatically cut down the training time because you
don't have to train the whole first part of it,
and you only train the last layers, which
allows you to very quickly pivot for a different task.
We're going to do something like this in the Google Colab that's coming up,
which exposes you not only for training for a particular visual wake words
task, but it also teaches you how to do this efficiently.
So there are two things that I'm trying to knock out here for you.
That said, I do want you to keep in mind that transfer learning is a technique
that might not always work well.
For example, if you know how to drive a car, that does not
imply that you can transfer that knowledge to flying a plane.
You still have to go to school and learn everything over
about that particular task.
So keep that in mind.
I have a reading that explains this sort of fundamental difference,
and there are a couple of pitfalls and myths around transfer learning
that are important to appreciate in order
to know when and how it is actually a useful technique to deploy.
With that said, have a little bit of fun in Colab, and I'll see you soon.

\section{Assignment: Transfer Learning in Colab}

Now that you have explored vision applications for TinyML, the MobileNet model, and transfer learning, we are going to put it all together and train a model to help fight COVID-19. We are going to detect if people are wearing masks or not! 

In this Colab you will see through a hands on example how to leverage a pre-trained model as a feature extractor, and how to add new layers and train them to use for classification on a new task. You will also be challenged to find the right number of epochs when using transfer learning. Hint: it's going to be a pretty small number!

\url{https://colab.research.google.com/github/tinyMLx/colabs/blob/master/3-7-11-Assignment.ipynb}

\section{Assignment Solution}

You should have found that you only needed VERY FEW epochs to get the model to train very accurately. Depending on how lucky you were with initialization it could have been in the single digits! We found that by setting EPOCHS = 10 our models were always sufficiently trained, and in some cases we were even able to get \PYTHON{EPOCHS = 2} to work! 

We hope you enjoyed exploring the power of transfer learning! As always if you would like to go back and try the assignment again, it can be found here:

\url{https://colab.research.google.com/github/tinyMLx/colabs/blob/master/3-7-11-Assignment.ipynb}


\section{Common Myths and Pitfalls about Transfer learning}

\subsection{The Pitfalls of Transfer Learning}

\begin{figure}
    \includegraphics[width=\textwidth]{TinyML/2 Applications/3-7-13-1.png}
    \caption{A high level diagram describing transfer learning showing re-training of the last few layers and freezing/reusing pre-trained initial layers. This is shown visually by graying out the first set of layers, which take in the input and have already learned to produce general features from the input, and then highlighting the last few layers that take those features and produce the output layers and noting that those are re-trained.}
\end{figure}  




Hopefully, you have now been convinced of the power of transfer learning. Transfer learning allows us to use a model generated for one task to be fine-tuned for a different but related task. In doing this, we can repurpose models for use in new tasks, thereby saving energy that would be required in computation. Furthermore, models with superior performance can be developed with a relatively small amount of data, which would normally cause overfitting in a network trained from scratch. Consequently, transfer learning provides a trifecta of benefits: reduced energy usage, faster convergence, and higher asymptotic accuracy. However, transfer learning can be difficult to implement and is not always guaranteed to be successful.

\section{Task Similarity}

Transfer learning is only viable when the task of interest is similar in scope to the original task. Image recognition is the archetypal example of this, wherein the convolutional filters close to the input can be ported to many imaging tasks due to their generality (e.g., spotting lines, edges, shapes). The further into the convolutional neural network we get, the more specific the filters become and the less transferrable they become. 

Tasks that are too dissimilar may not receive any benefit from transfer learning. In fact, the use of transfer learning may even be detrimental to the final model performance. For example, performing transfer learning on a neural network used for language translation may perform well if done between two similar languages, such as two Romanic languages (e.g., French and Spanish), but may perform poorly if done between two dissimilar languages (e.g., Mandarin and French). This situation is commonly referred to as negative transfer.

\section{Fragile Co-Adaptation}

Transfer learning involves the use of neuron weights from one trained network to pre-initialize neuron weights in a second network for a related task. Often, the weights of several of the first layers are frozen so that they cannot be changed during learning. The rationale for this is to prevent noise in the new dataset from detrimentally altering the earlier layers. However, freezing weights in itself can be detrimental to model performance as it can lead to fragile co-adaptation of neurons between neighboring frozen and unfrozen layers. This occurs when neurons in the former layers are fixed and the latter layers alone are unable to adapt effectively to the new data. In practice, this is often solved by not freezing any layers and instead using a significantly smaller learning rate used when training the original network. This procedure helps to ensure former layers are not altered significantly while reducing the possibility of fragile co-adaptation.

\section{Fixed Architecture}

A key disadvantage of transfer learning is the restriction on neural architecture. Transfer learning is often performed from well-known models developed by commercial entities such as Inceptionv3, ResNet, and NASNet. These models are trained with specific neural architectures which are reflected in the model weights. Model layers cannot be modified, removed, or added without these changes propagating across the network. As such, if the network is altered, we cannot have confidence that the remaining parameters are still able to accurately model the data. This presents important limitations for model tuning and, more importantly for tiny machine learning applications, also for model size. These well-known architectures are large and often cannot be compressed sufficiently to fit within the hardware constraints of embedded systems. Thus, the use of transfer learning for tiny machine learning applications is fundamentally limited by existing model architectures unless the resources are available to train bespoke models for transfer learning using our own pre-defined architecture.

In summary, here is a general rule of thumb of when transfer learning works well when applying transfer learning to your situation where you have a new dataset that you wish to fine-tune your network upon. 

\begin{tabular}{lll}
    &New dataset is similar to original dataset &
New dataset is (less) similar to original dataset \\
New dataset is small &
Best case scenario for transfer learning&
Not the best scenario for transfer learning\\
New dataset is large &
Transfer learning will work&
Training from scratch might yield better accuracy\\
\end{tabular}

Despite some pitfalls, it is clear that transfer learning is a hugely powerful concept and, although difficult to implement in practice, presents many exciting opportunities for tiny machine learning applications.

\section{Metrics for VWW}

VIJAY JANAPA REDDI: We're almost at the end of the visual wake word's
application.
I'm going to wrap up this video and the series
by focusing on the metrics that really matter.
Looking at the end to end pipeline, you've
come a long way from collecting data to designing
the model to understanding how to train or cut down the training
time to now evaluating a model.
When it comes to evaluation of a model, there are quantitative measures
and there are qualitative measures.
Both have pros and cons, and let's try and understand what it really means.
For quantitative measures, we have things like accuracy.
How well does a model do in doing its classification?
There are also other metrics that actually
matter that are outside of the model, which is, what's
the overall latency the application?
What's the amount of energy?
Does the entire application with the model inside fit in the resource
constraints?
Then, of course, there are qualitative aspects
about how well the overall application does for the end user.
We touched on these in the context of keyword spotting.
Here, I'm going to highlight things that are specific to the visual wake words
model.
Talking about accuracy, I want to quickly touch on the metrics
that we have.
Often, we use Top-1 Accuracy.
Top-1 Accuracy means that the exact softmax classification
score that you get has got to have the highest probability so that it matches
what the original ground truth is.
In this particular case, it is Firefox.
So that's great.
Occasionally, you'll hear about what's called Top-5.
That simply says that if Firefox shows up
in anywhere of the top five probabilities,
then that is also completely acceptable answer.
It really depends on what you're looking at.
Often, we use Top-1 as the baseline so keep that in mind.
Sometimes you hear Top-5, and that's the only reason
I'm bringing up that notion.
Now, one of the key things that is important
is that you got to think about what data you have in order to determine
whether you are actually going to get good accuracy because if you don't have
data, you have data that's irrelevant really, then your accuracy
is going to be bad.
So no matter what your efficiency is, it does not
matter because the end user experience is going to be terrible.
Another thing that you could do is you could obviously
trade-off accuracy for latency, which we've talked about previously
when we looked at designing a model.
For example, I can have a really accurate model that is really slow.
How will that work?
So if you look at it purely from an accuracy standpoint,
it might look like a fantastic choice because it gives you
the highest accuracy in detecting people or something to that effect.
But if it takes 10 seconds to be able to make a classification
and tell you that, yes, there is a person in front of the camera,
is that any good?
Will you be able to deploy that on a smart device, for instance?
No, you won't be.
So that means your user experience is going to be terrible.
So that's not really going to work.
But you might be able to compromise that model quality a little bit
to get a faster model, which is exactly what mobile nets did.
That reduced the complexity of the model through depth
by separable convolutions, which is why we talked about it, in order
to improve the overall performance of the model,
and thereby improve the user experience.
Another thing that we can do, depending on the use cases,
if you want high quality model, then sometimes you
might have to offload it to the cloud.
If you're dealing with applications like medical situations,
the end user might be willing to tolerate long latency
as long as the model is accurate.
So it really depends on the application you're building and deploying.
So that's the message that I'm trying to get across here.
Last but not least, fairness--
diversity in the data is extremely important because if you build a model,
it's got to work across the diverse situations in which it's
going to be exercised.
Specifically, if the particular group that you're looking for
is actually the majority group of your training data, well,
then you're naturally going to get good accuracy,
and everything is going to work fine.
In other words, if I'm building a visual wake words model that
recognizes my skin tone, that's great.
OK, well, I'm in the majority group of that training data
so the model is going to work well.
But what happens now if there's a certain skin tone which the network has
not seen in the training data, someone who's
got a much lighter color face or someone who's got a much darker skin tone?
In that case, if there's something very specific about the skin tone,
for instance, that's getting encoded into the network,
that's going to prove really problematic because even if the model runs really
well, the overall user experience is going
to be bad because it's not generalizing to the characteristics that are needed.
So the point I'm trying to get across here
is that we need to think about diverse set of metrics.
We need to identify the metrics that are important carefully.
Just the way we identified the requirements for a dataset carefully,
we need to identify all the right metrics
and then work through that process.
Though it is a last step, it actually greatly influences
how the entire pipeline is kind of engineered.
And with that said, I hope you learned a couple of new things
in the TinyML visual wake words application.
We're going to be moving forward into yet another interesting application
area where we're going to learn about new fundamentals
as we go through the same process of dealing with sensors and then
the network and then how we actually evaluate that particular scenario.
Stay tuned for that.


  \section{Summary Reading}
  
  In this section you explored the TinyML flow and data engineering in the context of a Visual Wake Words application, focusing on some unique challenges presented by this computer vision application. Visual Wake Words represents a common TinyML visual use case of identifying whether an object (or a person) is present in the image or not.
  
  \begin{figure}
      \includegraphics[width=\textwidth]{TinyML/2 Applications/3-7-16-1.png}
      \caption{Machine learning workflow and pipeline. Begins with collection and preprocessing of data then the design and training of a model. These steps are done in TensorFlow. Then a model is optimized, converted and deployed using TensorFlow Lite. Finally inferences are made using TensorFlow Lite Micro.}
  \end{figure}  
  

  
  
  \subsection{Challenges}
  
  You explored how latency is a strong constraint for VWW applications as images are large (much larger than spectrograms from KWS) and take a long time to send to the cloud. You also explored how memory is also a strong constraint as images, and their traditional models, take up a lot of memory. Finally, you saw how false positives and false negatives can be particularly challenging in the context of computer vision.
  
  \subsection{Datasets}
  
  In the context of datasets your explored three main topics: licensing, privacy, and reuse. Licensing is a major issue for computer vision applications as while images are readily available, they cannot often be used freely without violating copyright. In terms of privacy, you explored how image metadata can often leak significant information about people. Finally, you explored how the VWW dataset reuses the COCO dataset by sub-sampling images using bounding boxes.
  
  \subsection{MobileNets}
  
  You explored how MobileNets use Depthwise Separable Convolutions to reduce the model size and number of operations, at the cost of a decrease in maximal expressiveness. This reduction is dependent on the size of the kernel and the number of kernels used according to the following formula:
  
$$  \frac{\text{Depthwise Separable}}{\text{Standard}} = \frac{1}{N} + \frac{1}{D_K^2}$$
  
  You also explored how reducing the depth multiplier and the input image size can also be used to trade off model memory and latency with accuracy.
  
  \subsection{Transfer Learning}
  
  You next explored how transfer learning can be used to drastically reduce the training time and data requirements. Transfer learning works by reusing the first few layers of a pre-trained model for a task using a similar dataset and re-training only the later layers. This works if the two datasets are close enough together as the first few layers learn to extract generic features from the data while the latter layers learn to extract task specific features. You also learned, however, that transfer learning is not a panacea and is ultimately also limited by the structure of the pre-trained model. You then got to explore using transfer learning to train your own VWW application in Colab!
  
  \subsection{Metrics}
  
  Finally, you explored both qualitative and quantitative metrics for VWW applications. In particular we discussed the latency accuracy tradeoffs often found in VWW applications and the challenges of getting fair results across data represented by both the majority and minority of training samples.
  
  
  % 1.8: Anomaly Detection
  
  \section{Introduction to Anomaly Detection}
  
   \begin{figure}
      \includegraphics[width=\textwidth]{TinyML/2 Applications/3-8-1-1.png}
      \caption{Machine learning workflow and pipeline.}
  \end{figure}  

  
  \subsection{What’s the Focus in this Module?}
  
  In this section we are going to explore a very different kind of TinyML application that is showing great promise for many different application domains, Anomaly Detection. 
  
  \subsection{What's New and Different?}
  
  Anomaly Detection is different from the previous TinyML applications and datasets that we’ve explored because we are not looking to match some particular training data, or to use a particular sensor. Instead, we are focused on finding data that we explicitly have not seen before. We’ll explore what kinds of data we can use for these challenging and important applications. We’ll then explore a new paradigm of machine learning, unsupervised learning, that can be used to solve these problems through the use of both traditional machine learning algorithms, such as K-Means, and neural network based approaches, such as autoencoders. Finally, we’ll train one of these approaches, learn how to evaluate performance with new  metrics and explore how to select a decision threshold when designing functional applications. We hope you enjoy this section!
  
  \section{What Is Anomaly Detection}
  
  VIJAY REDDI: Welcome to your third and final TinyML application.
  In this particular series of videos, we'll
  be learning about anomaly detection.
  Let me first and foremost define what I mean by anomaly detection.
  Anomaly detection is a specific form of data analysis, where
  you are looking to identify interesting and rare items, events,
  and observations, which raise suspicions by deferring significantly
  in their behavior from the majority of the data that you would expect.
  What this means, that occasionally, you will find some sort of a needle
  in a haystack that kind of sticks out, and you're like, well,
  why is this sort of happening here, while the rest of the data
  is typically tightly clustered around some given point?
  Now there are many different interesting applications of TinyML
  and specifically, anomaly detection.
  There are applications in heath?
  There are applications in industrial.
  There are applications in security.
  Before we dig into any one of these, let me just talk about anomaly detection
  as a general application outside of the TinyML context,
  just set the stage for you so that you have a broader
  understanding about the field.
  Let me take security as a specific example of anomaly detection.
  Typically, on smartphones, for instance, there's a lot of malware out there.
  Software that should not be accessing things
  could be accessing things on your phone and leaking it up to some hackers.
  So in cyber security, for instance, anomaly detection
  plays a very big role.
  Specifically for instance, you might be surfing the web.
  And typically, when you're surfing the web on your Android phone or your iOS
  phone, rarely would you ever need to access the contacts.
  So if, let's say for instance, your browser got hacked somehow,
  and then when you're reusing the browser, for some reason,
  every time you run the secret browser, it's constantly accessing the contacts
  and it's sending all your contact information out, well, typically,
  in an anomaly detection system, what we would
  do is that it would actually detect that some kind of unorthodox behavior
  happening that is actually causing your browser to be sending out
  data which it should not be sending
  out.
  Now that's very useful to know, because then suddenly, the system
  might flag and say, hey is this an expected behavior?
  So that's a very simple, yet clean example of anomaly detection.
  Something that's actually quite widely deployed in the larger ecosystem.
  Now we're going to be talking about anomaly detection in this course
  specifically from the perspective of sensors--
  no surprises.
  So in the health case for instance, you might have an EEG sensor.
  My iWatch, for instance, actually has anomaly detection mechanisms
  built into it that kind of help me in terms of my health monitoring.
  For instance, how does a watch know when some old person has fallen over?
  There are detectors and there which are looking for some anomalies that happen,
  which can then raise sort of a flag, and inform the families and so
  forth-- that kind of thing.
  So you have a smartwatch.
  In the case of industry, for instance, we'll be going pretty deep.
  We'll be looking at anomaly detection in the case
  of what happens if a motor fails, where typically, you
  expect the motor to be running normally, and then something goes wrong,
  and then it fails.
  Well, what if I could have predicted that the motor was
  going to fail because I could tell that something's actually
  off in terms of the way the motor is actually sounding when it's running?
  And we've already talked about security.
  There are many different applications of security based on the kind of sensors
  you might have.
  So what's the key thing about anomaly detection?
  It all comes down to having information that's coming through sensors.
  Specifically, what's going to be interesting about anomaly detection
  as we talk about it in this class is that we're
  looking at temporal data streams.
  Now we're looking at time series data, which
  is a sequence of measurements, where the temporal relationship is actually
  a critical component to understanding the data.
  Now this is a new thing that we're actually introducing to you,
  so keep that in mind.
  We're looking at time series data, so we're flipping up what kind of data
  we're looking at.
  Previously, we looked at an audio acoustic input
  and converting that into an image.
  Then we talked about visual wake words, where
  we're looking at just a static image and sort of thinking
  about how do we sort of compress that down
  effectively for our neural networks.
  Now we're talking about a very different kind of data stream,
  which is a time series scale.
  I'm going to quiz you about that.
  So in terms of use data, what you're doing is
  you're kind of scanning through the data as it's
  coming through from your sensors, which might be more than one sensor,
  for sure.
  And you're trying to find interesting patterns in there.
  And based on that pattern of activity, you
  want to determine if there's an anomaly or there's no anomaly.
  An anomaly meaning that the behavior is outside of the expected pattern.
  For instance, in this particular chart, you're looking at that red line,
  and that red line seems to have a nice up and down
  sort of a sawtooth kind of curve.
  An anomaly here would be perhaps when you don't have that up and down sort
  of sawtooth curve--
  perhaps something like the behavior towards the beginning of the curve
  where it kind of just dips a little bit, and it doesn't really
  dip all the way down.
  Perhaps that was an anomaly.
  So you're going to be looking at times use data that's coming in from sensors
  and trying to understand how to process it
  in order to detect if something is off.
  There are three fundamental building blocks
  for any kind of anomaly detection application or scenario.
  There's an input signal that's coming in--
  for instance, like the blue line that's on the right hand side.
  Then you're making a prediction.
  For instance, the predictor here is simply sitting in the value 0.
  It's a straight line.
  It's expecting the input signal to always be a nice strong 0.
  But as you can see, the input signal is varying, while the prediction is flat,
  which means that there's an error here, which
  is what the shaded area under the blue curve is capturing.
  With these three fundamental building blocks,
  you can pretty much build an anomaly detector,
  which is what we will be doing.
  But the reason I'm touching on these three things
  is that they're paramount for the fundamental building blocks
  of any anomaly detection solution.
  So for example here is a poor predictor or poor prediction,
  because your input signal and your prediction are not tracking together.
  But here's an example of a good prediction,
  because your predictor is effectively tracking that expected behavior,
  and it's minimizing the error.
  The error is key here.
  The anomalies are typically detected when there's a large error.
  So for example, perhaps on the left hand side,
  there's continuous error as you're constantly
  detecting that something is pretty off, because the error is pretty large most
  of the time.
  For instance, on the right hand side, the error
  is actually quite small, which means perhaps the machine is actually
  performing as it's supposed to.
  Now why do we want to talk about all of this in the context of TinyML?
  Why anomaly detection in TinyML?
  I gave you some application of the use cases,
  now let me kind of talk about some of the technical things.
  We'll go deep into this in the next couple of readings and videos.
  But real briefly, one of the key things is
  that it's actually really hard to have the ability
  to transmit the data that is at the sensors out,
  because typically, these sensors are sitting
  close to some sort of electrical circuit or someplace
  where it's very hard to get connectivity.
  And we'll look at a ball bearing example in a motor,
  for instance, where it's very deep inside the machine
  that it's actually hard to get the signals out nice and cleanly,
  because you have a lot of electromagnetic interference.
  Then, of course, there's the aspect of latency.
  You've got to track the data really quick,
  because there's a sensor data that's streaming fast,
  so you've got to be able to do that real quickly.
  And you can do that if you're right by the sensor,
  and you're doing on device ML.
  And of course, you want to keep the energy budget low,
  because when you're doing anomaly detection,
  you're talking about normally detection, many different use
  cases where you don't want the battery to be drained out.
  Because often, these anomalies are tethered to sensors,
  and if there are sensors, that means they have very little energy
  associated with them.
  One of the other really important things that I'm
  excited about teaching in this particular aspect of the course
  is that while most of the time we've been talking about machine learning,
  and when we've been going through this sort of a pipeline,
  we have been very heavily focused in one very particular single type of machine
  learning, which is deep learning-- using neural networks.
  What I want to do is to try and take a step back and raise your awareness
  about the level of machine learning methods that are out there.
  For example, neural networks are just one way of doing classification.
  There are also other methods like support vector machines.
  And just the way we can do classifications,
  you might have a nearest neighbor sort of method,
  you might have a clustering-based method, and so forth.
  So there are many different ways of actually unlocking machine learning.
  In the context of anomaly detection, one of the important things to keep in mind
  is that the community or the industry at large
  is trying to figure out which is the best way of going about doing things.
  It has not been yet effectively proven that neural networks are the only way
  of doing things like anomaly detection.
  There are other methods, for instance, like K-Means,
  which are actually quite popular for doing machine learning
  for anomaly detection, because they're very lightweight in terms
  of the compute needed.
  There are also other mechanisms like decision trees that can be used.
  They also require very little computer.
  One of the things that we're going to do when we talk about something
  like K-means where we're actually doing a little bit of clustering
  is understanding when these very basic sort of methods
  are actually remarkably good for doing things like anomaly detection,
  but at the same time, we're also going to understand when they're not good.
  And then we'll talk about how neural network-based methods might
  be good for things where traditional machine learning methods are actually
  not good.
  And the reason I want to do this is to kind of pull together
  your complete knowledge around TinyML and get
  you an unbiased sort of perspective in terms
  of how to look at machine-learning-based solutions
  when you're thinking of applications.
  It's not all about neural networks.
  You've got to take a step back and understand
  that it might be useful to use other types of machine learning methods.
  That said, we will go deep into autoencoders.
  Now as always we will start off by talking
  about what are some of the challenges in actually deploying anomaly detection.
  Then we'll go into our anomaly detection machine learning pipeline
  and sort of understand the full flow as we have always
  done with every single application.
  And then you'll get some hands on experience
  by learning how to train the anomaly detectors
  and then test it out on Google Colab before moving onto course 3,
  where you actually end up deploying it onto the physical device.
  And with that said, I want you to learn some interesting things
  about anomaly detection and keep an open mind
  in terms of what kind of techniques are actually
  useful in the broader ML space.
  
  \section{What are some interesting applications or use cases of Anomaly Detection that come to your mind?}
  
  What are some interesting applications or use cases of Anomaly Detection that come to your mind? Why are you excited about this application/use case? What challenges do you think you would face using TinyML in this context?
  
  \section{Anomaly Detection in Industry}
  
  VIJAY JANAPA REDDI: In this video, we're going
  to start off by talking about what are the challenges with anomaly detection?
  Specifically, I'm going to take a real world example
  scenario where people are really excited about TinyML, which is
  in the context of factory applications.
  In factories or in manufacturing plants, there's
  a lot of interest for doing predictive maintenance.
  Doing predictive maintenance specifically with TinyML
  is a really hot topic.
  Why?
  Predictive maintenance is this idea that you
  want to be able to replace a particular part that may fail,
  or is likely to fail in the near future, before it actually fails.
  And if you can do that proactively, the benefit
  is that you don't have to shut down the factory or the manufacturing plant
  in order to be able to replace that particular part.
  Instead you might be able to schedule ahead of time
  that you want to replace this part perhaps on a weekend,
  or sometime when the manufacturing plant is actually inactive.
  In doing so, you effectively cut down your overall downtime
  for your manufacturing plant, which is awesome, because it dramatically boosts
  your production capabilities, while keeping your overall factory
  or manufacturing plant quite safe.
  Now let's talk about this in the context of what
  it means to actually enable TinyML and how to actually deploy
  TinyML in the context.
  So let's say you have this small little sensor, that is actually
  attached to perhaps a ball bearing.
  And this sensor, like an accelerometer, for instance,
  is collecting a lot of data about the vibrations and things like that,
  that it's picking up from the ball bearing.
  And what it's doing is that you want to be able to look at that data
  and decide if there's some sort of anomalous behavior.
  Now, this gets us into a set of issues, which you might be quite familiar with.
  Latency, limited constraints, and accuracy aspects.
  So keep in mind that you've got a sensor.
  That sensor is actually attached to the device
  and it's going to be looking at a whole bunch of time series data
  that's coming through it.
  How can I processes this efficiently?
  I want to tie things back to whether you should do this in a TinyML context
  or whether you can offload it to the cloud or some sort of large computing
  server that's in the manufacturing plant.
  Let's take a little closer look at this.
  From a latency standpoint, if we look at this particular sensor here,
  this SKF sensor, what the sensor is actually doing
  is it's collecting a bunch of data.
  This is actually a real part that is used in production.
  It has a whole bunch of sensors that tracks the speed, the position,
  and the angle of the ball bearings.
  And this data is continuously available for us to access,
  because this time series data is getting pumped out continuously
  from this machine.
  And what we want to be able to do is we want
  to be able to intelligently process this to see
  if there's some sort of anomaly in the activity of this particular sensor.
  Now, let's say we're collecting data and we
  want to offload it to some sort of nearby computing machine.
  Is that even possible?
  Let's say if every measurement is two bytes of data,
  and there are eight different sensor streams that are coming out
  of this particular device.
  And the data that's coming out is at about 20,000 samples per second,
  or 20 kilohertz.
  That translates to about 320 kilobytes of data per second.
  Is that a lot?
  Let's see.
  Let's say you want to offload this data.
  Let's take, for example, this standard Zigbee
  radio, which you might be able to attach directly next to that sensor which
  is collecting all this data.
  A Zigbee radio is commonly used in wireless sensor networks,
  and it reaches only about 30 kilobytes per second in terms of its ability
  to stream data.
  This basically means that as compared to our 320 kilobytes of data
  that's coming out per second, the Zigbee radio
  is not going to be able to keep up.
  We would be producing data faster than we can actually
  transmit it off the particular device.
  Now, in addition, one of the problems you have
  is that a Zigbee radio consumes about one joule of energy per second.
  Now at that rate, that means a standard coin cell battery would pretty much
  be depleted in about 40 minutes.
  So every hour you're going to be going and replacing batteries.
  So that's not going to work.
  So the point here is that you've got to think about both the amount of data
  that these sensors can actually generate and how do you process all of that data
  locally, rather than trying to move it.
  That's why TinyML is super exciting in this space.
  Therefore a lot of companies are now trying
  to figure out how can we put in the intelligence needed
  real close, right next to the sensors.
  Now remember, we just gave an example where
  we're just talking about a single unit.
  What happens if you have a whole factory of these different machines
  that you want to instrument?
  Think about the volume of data that they're going to be generating.
  So in order to cope with that scale, you need to have TinyML.
  And that's why the industry is super excited about this particular solution,
  if we know how to solve it.
  Now, talking about limited resource constraints,
  you know this all too well.
  Our devices, our TinyML devices, have very tight constraints.
  We only have about 256 kilobytes of RAM, typically,
  on this one particular device.
  Most of the other TinyML devices are just as small or even smaller.
  So, you know this all too well.
  I'm not going to go in too deep into this anymore.
  But let's talk about accuracy for a second.
  What does it really mean to have a model that is actually
  being deployed and doing detection, things like anomaly detection,
  in order to identify dangerous situations like,
  for example, if there's a gas leak.
  Let's say you have this small little clever sensor that's
  tethered to a gas pipe, and its job is supposed
  to be detecting whether there's a gas leak.
  And what happens, for example, if you have a false negative?
  In other words, your machine predicts that there is no gas leak,
  but in fact, there is actually a gas leak and somehow
  your model mispredicted.
  Well, that can have a catastrophic impact.
  Lives could be lost as a result of your model making a misprediction.
  And on the other hand, you might say, OK, well I'm
  going to be a bit more lenient.
  I'm going to allow some false positives to come through.
  Well, be careful if that happens.
  Because if you have too many false positives,
  then people will kind of take that machine to be rather crappy.
  In other words, they're going to say, ah,
  this machine always seems to generate false alarms.
  Just ignore it.
  And that one time when it really matters, well, too bad.
  The other aspect is that if you are generating
  false positives and the concern you have is
  that you might have to stop the overall machines down and then go and actually
  inspect it because you are doing predictive maintenance.
  And if it turns out that most of the time there's really nothing wrong,
  then that is also a penalty in terms of the expenditure.
  So what I'm trying to hint at is that there
  are lots of interesting potential applications for doing
  TinyML in manufacturing plants.
  Specifically around detecting faults and doing predictive maintenance,
  which is going to be our topic of interest in this particular course.
  But I'm trying to give you and highlight about some
  of the critical differences between how you
  need to think about this particular application
  as compared to keyword spotting or as compared to visual wake words.
  Remember, in those kinds of applications,
  if your machine does not wake up when you say, hey
  Alexa, or OK Google, or something like that, it's not the end of the world.
  We'll say it again.
  Sure, you'll be frustrated.
  OK, fine.
  Well, I might not buy that product.
  OK.
  In this particular scenario, if you are relying on these tiny little sensors,
  these intelligent tiny little sensors that we're going to build,
  and you're deploying them, and you're relying on them for safety,
  you have a serious job and you have a big responsibility.
  So keep that in mind as we go through the exercise of building
  TinyML models for things like anomaly detection, which
  are very powerful concepts.
  There are a lot of dangers, but there's a lot of potential good that we can do.
  With that said, in the next video, we can just
  get started with trying to understand what it means to build a TinyML
  model for anomaly detection.
  
  \section{Industry 4.0 and TinyML}
  
   \begin{figure}
      \includegraphics[width=\textwidth]{TinyML/2 Applications/3-8-5-1.png}
      \caption{Industry 4.0 and TinyML.}
  \end{figure}  
  
  

  
  Possibly the greatest impact of tinyML is likely to occur in industrial sectors. The technology is likely to be both disruptive and enabling, providing large benefits to those who are able to leverage it for effective use in business operations. Some industry professionals prognosticate that the introduction of “intelligent” internet of things (IoT)  devices afforded by tinyML is akin to a new industrial revolution, leading some to refer to its future potential as Industry 4.0.
  
  \subsection{What is Industry 4.0?}
  
  The reason for this name is relatively simple and refers to the fourth industrial revolution. The industrial revolutions refer to specific transformative technologies that revolutionized the industry. 
  
  The first industrial revolution started around the 1770’s with the large-scale use of steam and water power, moving from traditional human labor on farms to the first factories and resulting in an exponential increase in productivity. 
  
  This continued until the introduction of electricity at the end of the 19th century, when electricity began to be introduced, prompting the second industrial revolution. This revolution kick-started mass production via assembly lines and automation and further increased the productivity of industrial processes, giving us the notion of a factory as we know it today. 
  
  This continued until the 1950’s when the digital revolution prompted the third industrial revolution. The introduction of computation allowed industrial processes to become controlled and monitored with greater ease, allowing some of the blue-collar work to be automated. 
  
  In recent years, we have started to see a move towards another industrial revolution focused on the interconnectedness of industrial processes. Instead of considering standalone industrial processes, they are considered as part of a network of interconnected processes. Essentially, these processes are in a constant state of communication. This has been enabled by the introduction of intelligent IoT devices tightly coupled with cyber-physical systems, affording capabilities for real-time monitoring and control such as predictive maintenance.
  
    \begin{figure}
      \includegraphics[width=\textwidth]{TinyML/2 Applications/3-8-5-2.png}
      \caption{Four industrial revolutions. First mechanization and steam and water power. Second, mass production using electricity. Third, automated production, computers and IT systems. Fourth, the smart factory using autonomous systems, internet of things, and machine learning.}
  \end{figure}  
  

  
  
  \subsection{Where does TinyML fit in?}
  
  Industry 4.0 encompasses multiple novel technologies, including IoT, blockchains, and tinyML. TinyML, specifically, will come to play a very important role in Industry 4.0. 
  
  One reason for this is a reduced network load. In a network of interconnected devices, it is easy for the network to become saturated with information. Such a situation might occur if all data for monitoring purposes has to be continually transmitted to a central repository for monitoring purposes. Instead, TinyML enables individual devices to manage their own data and make intelligent decisions based on this information, thus limiting the communication load on the network.
  
  Another important factor is energy efficiency. As we have discussed previously, transmitting data across devices is energy-intensive. By performing inference on-device and minimizing external communications, we maximize energy efficiency. Given the importance of sustainability in modern industry and the projected exponential growth in IoT usage in the coming decade, this provides an important stimulus for using tinyML.
  
  Finally, performing inference on-device minimizes latency. Nowadays, it is common to find systems that upload data from an edge device to a cloud server for processing and then return the processed information to the edge device. This message passing across a network is not only energy-intensive and clogs the network, but is also slow. By performing inference on-device, we are not limited by network speeds and only by the hardware constraints of our devices.
  
  Hopefully, it is now clear that with the introduction of tinyML, the wealth of information procured by edge devices can now be leveraged for intelligent control of industrial processes and business operations. Industry 4.0 essentially provides the capability to perform data-driven optimization of business processes in real-time, allowing inefficiencies to be ironed out almost entirely. TinyML will play a fundamental role for the large-scale use of machine learning in the industry.
  
  \section{Anomaly Detection Datasets}
  
  VIJAY JANAPA REDDI: In this section, we're
  going to get started with the anomaly detection machine learning pipeline.
  As always, I like to start right at the beginning talking
  about data collection.
  Because I'm going to say this yet again, you get the data collection part wrong,
  everything else down the pipe is a wasted effort.
  So let's try and understand a little bit about what it means
  to collect data for anomaly detection.
  Let's roll back a bit and reflect on some of the data set principles
  that we learned about previously.
  Here's an example of the speech command data set,
  which we use in keyword spotting.
  Previously, when we wanted to do keyword spotting,
  we knew what the requirements were for the data set implicitly.
  Because you knew what the keywords were going to be,
  you could then go about collecting data for those keywords.
  Sure, you got to be careful about gathering the right keywords
  under the right situations, and making sure that the data is correct,
  and making sure you can sustain that pipeline.
  All of that good stuff that we learned, yes, they're all important.
  But there's one very important fundamental assumption, or baseline,
  that you can take for granted there.
  It's that you know your keywords.
  That makes a world of difference because when it comes to anomaly detection,
  you don't know your keyword.
  What do I mean by that?
  Specifically, you don't know the numerous ways
  in which a machine can fail.
  So when you're talking about anomaly, you're
  talking about finding in 1 of 10 ways, 100 ways,
  1,000 ways a particular machine can fail.
  Compare that and contrast it with keyword spotting or the Visual Wake
  Words data set that we talked about.
  We knew exactly what we were looking for.
  In keyword spotting, we knew the keywords.
  In the Visual Wake Words, we know the particular item
  of interest in those images.
  Here, you don't know which way that the machine is going to fail.
  All you know is that you're trying to predict a failure.
  So that presents a real problem.
  That's one.
  Number two, it is quite hard to capture a wealth, or a volume, of data
  for something that is inherently rare.
  Recall that you're trying to collect anomalous data.
  What that means is that you're trying to collect a lot of data for something
  that happens extremely rarely.
  So it's like a really hard problem that you're trying to solve.
  And if you wanted to collect that data very fast,
  then you have to go about buying machines and intentionally breaking
  them.
  Well, that doesn't sound like an easy thing to do.
  Or it does not sound like a cheap thing to do.
  Usually, you're trying to use TinyML and anomaly detection on costly machines
  because you want to predictively do maintenance on these costly machines
  so that they don't break down at critical times.
  Well, you're not intentionally going to go and break down a jet engine
  or break down 1,000 or 100,000 jet engines just
  because you want to be able to do predictive maintenance on one or two
  jet engines when it's timely.
  Think about that.
  So this presents a really hard problem for us
  because normal data is usually a closed set, meaning the normal operation.
  You know what it looks like.
  And then therefore you can easily capture the data
  because you have a lot of normal data.
  Abnormal data presents a really big problem
  because abnormal data is an open set.
  It's a set where a system can fail in any number
  of different and unpredictable ways.
  It's the unknown unknowns.
  Therefore, we need some way to train a model
  to detect something that we can't actually capture in our data set.
  Think about the difference here.
  In the case of keyword spotting or Visual Wake Words,
  you would just go and get that data.
  But here, it's very hard to get that data.
  And even if you could get that data, it's going to cost a lot.
  So you need to be clever about how you would actually engineer the system
  in order to detect these abnormalities.
  One of the biggest problem is that you have
  unbalanced data in anomaly detection.
  In the context of traditional detector cat versus dog and all the data
  engineering that we taught you, there was lots of that data
  that you could create in order to create a balanced data set.
  Here, because anomalies are really rare, the problem you're going to have
  is that you're going to have very little data available for you.
  So one of the things that has been going on
  is that there's been a lot of interest in the community
  to try and build data sets for anomaly detection, specifically around sound.
  I'm actually introducing sound in here because we've done stuff around sound.
  So here's MIIMII data set that you should know about,
  which is a data set which is specifically
  designed for malfunctioning industrial machine investigation.
  It's a collection of sounds.
  And what they had done is they had gone out
  and instrumented things like pumps, fans, valves,
  and so forth for different kinds of sounds
  that they make, both during normal operation as well as abnormal operation
  in real factory settings.
  And so the spectrograms here are kind of showing you the visualization
  for different kinds of sounds.
  In fact, you're going to have a deep reading that's coming up,
  which is specifically on the MIIMII data set,
  so I'm not going to go too much deeper into this.
  The reading will do so.
  Now, this is in the context of a real factory setting.
  But what about getting creative in order to get lots of that data, that hard
  to get data?
  Well, there's a ToyADMOS, which is a really interesting way
  to think about collecting data.
  The ToyADMOS data set is basically a data set of miniature machine operating
  sounds.
  This data set happens to have 540 hours of sound for anomaly detection,
  again based on sound.
  It's based on toys--
  car, conveyor, trains-- and it's about the sounds that they make,
  and you're trying to detect anomalies from them.
  So the way they kind of set it up is that there's a track here
  that you're seeing where there's small little trains--
  yes, they're toys-- that are actually going
  around doing the job of continuously doing what a toy does.
  And there are different microphones that are positioned in different places,
  and their sounds are recorded.
  And one of the key things that the authors have done
  is that they deliberately damaged a toy.
  So you can collect different kinds of sounds.
  So they're damaged in many different ways, for instance.
  And that allows you to have a diverse portfolio
  of sounds inside the data set.
  Looking one level deeper, the data set is
  spread across three different types of toys, as I said earlier--
  toy car, a toy conveyor which is actually fixed,
  versus a toy train which is actually continuously moving around tracks.
  So you get very different kinds of sounds coming from them.
  And you're collecting different machine types of sounds.
  And you also got data in there that's actually quiet,
  which is the background noise, or the environmental noise
  that you have of that operating situation,
  as well as the data when the machine is actually running.
  Now with that, you can actually think about how
  to collect the data intelligently because, as I said earlier, the four
  different microphones they use.
  Why?
  Because you want to look at the data from different angles.
  So if you use one microphone alone, it does not
  give you a full representation of what the sound would look like.
  So there are a lot of nuances about how you actually collect sounds.
  And when you're building a model, for instance, you can decide,
  am I going to use the data that's coming from all four,
  or I'm actually going to use all three, or only two, and so forth.
  Point is, there's a lot of flexibility in this data set.
  They also have multiple machines.
  For example, they have multiple cars.
  They have multiple conveyors.
  They have multiple trains for the exact same class.
  And that leads to variations that are inherently there
  [? between ?] these different machines, which
  gives you different kinds of sounds for the exact same class.
  And the data is recorded several times.
  So the anomalous sounds you end up picking up
  tend to have diverse characteristics, which is a good thing.
  Ultimately, you have about 180 hours of normal machine operating sounds
  and over 4,048 kilohertz samples of anomalous behavior.
  So it gives you a fairly different behavioral pattern
  about what the entire data set looks like.
  The reason I'm talking about this is because I'm
  trying to highlight two key differences [? between ?] the way
  the MIIMII data set is kind of done from a real world factory
  setting and the ToyADMOS, which is kind of done with the mock environment.
  While you can train models on both of them,
  one of the key things I want you to keep in mind
  is the issue about transferability.
  Let's say you come up with this magical model,
  which we'll get into next, that can actually do
  anomaly detection on a given data set.
  We talked about transfer learning in these kinds of things.
  I want you to know that, in the context of anomaly detection,
  it's actually quite hard to easily be able to transfer the application
  environments.
  For example, if you're collecting sounds for something
  like about how the airplane landing gear works,
  the data that you collect there is not easily repurposable for something
  like an oil rig.
  So just kind of keep that in mind because sometimes anomaly detection
  models tend to overfit to their training data
  within that particular sound context.
  And this is just one of the pitfalls.
  So keep that in mind because that is extremely
  tied to your data engineering aspects.
  So with that said, I want you to go on and do your Colab.
  And you're going to be doing a little bit of hands-on programming.
  Specifically, you'll be touching up on things around some of the thresholds
  that we can change, and what are the differences
  k-means versus anomaly detection, and so forth.
  So we'll be getting pretty deep into that.
  With that said, I'm going to next talk about how do we
  actually build an anomaly detection model
  because that's also going to be new.
  It's different from the other types of networks
  that you've actually played with so far.
  And I'll see you there.
  
  \section{MIMII DATASET}
  
  In this reading assignment, we want you to read through the “MIMII Dataset” paper. The paper provides additional information on some of the details that we discussed in the video lectures. It provides a bit more context and detail around the following important topics:
  
  \begin{itemize}
    \item What is the MIMII Dataset?
    \item   What was the motivation behind building this dataset?
    \item   How was the dataset collected?
    \item   Why did the creators also include sample benchmark models?
  \end{itemize}
  
  \section{Real vs. Synthetic Data}
  
  As we have explored throughout this course, data is at the heart of machine learning. Oftentimes, data is the limiting factor and in turn dictates model performance. A scarcity in the amount of data points or number of features hinders the predictive power of our model. The solution seems obvious: collect more data, but this is not always feasible. So what can an ML engineer do?
  
  Fortunately, with recent advances in deep generative models such as generative adversarial networks (GANs) and variational autoencoders (VAEs), very realistic synthetic data can be generated from a small number of known examples. This is achieved by approximately “learning” the data distribution and drawing random samples from the latent space of the generative model. This can be done on any type of data, allowing us to generate time series information (such as voice or sensor data), as well as photorealistic images. Using these new techniques, ML engineers can now bootstrap machine learning with machine learning -- that is, they can train models to generate more data with which to train other models!
  
  \subsection{What is Synthetic Data?}
  
  Imagine a power plant engineer that wants to use data to prevent a nuclear meltdown. They clearly cannot generate sample nuclear meltdowns to obtain relevant data, but they still need to be able to somehow make predictions. So what can they do? Well, if they have a very small amount of data from previous historical meltdowns or a simulation, then they can possibly use one of these techniques to train a model to generate enough data such that it becomes feasible to train an anomaly detection classifier using supervised learning. In this case, synthetic data is the only viable path forward. 
  
  The need to generate anomalous data is a commonplace situation in anomaly detection applications, since often the so-called “anomalies” are undesirable and can be extremely costly or damaging. Perhaps two of the most important applications of anomaly detection are for industrial processes and medical devices, where it can be used for predictive maintenance (fixing something before catastrophic failure or death). 
  
  Predictive maintenance will likely become an important feature of Industry 4.0, reducing costs by predicting component failures in advance and resolving the issue before it becomes more severe. An open-source dataset was created for this aiding with this specific application, referred to as MIMII (Malfunctioning Industrial Machine Investigation and Inspection). The MIMII dataset focuses on sounds exhibited by various components and can be used as a starting point to generate synthetic anomalous data for a variety of industrial components, reducing the need for individuals or organizations to collect these data.
  
  Similarly, in the medical device space, anomaly detection allows for moderately malfunctioning devices to be detected on and replaced before they pose life-threatening issues to the device wearer. The ECG5000 dataset was created for this purpose, and provides 20 hours of ECG data - a type of data used to study heart health. Anomalous data is also provided in this dataset which corresponds to an individual with severe congestive heart failure. The dataset is open-source and can be freely used by anyone working on medical devices, and can be used as a starting point to train an anomaly detection model or to create more anomalous data.   
  
  This may sound incredible - we can effectively train models by simulating new data based on our current data. However, you have already done this in this course! When you use Data Augmentation, you are effectively expanding your dataset with synthetic data! That said, there are concerns when relying heavily on synthetic data that should be highlighted.
  
  \subsection{Quality Concerns of Synthetic Data}
  
  Real-life datasets are often complex and multifaceted. For example, when recording sounds, the actual sound plays an important role, but so does the microphone, background noise, microphone orientation, as well as myriad other factors. Creating synthetic data that is able to mimic this variability is infeasible, and currently there is no universally accepted method for generating high-quality synthetic data. One important reason for this infeasibility is that data must often go through multiple preprocessing stages before the generation of synthetic data. During this preprocessing stage, assumptions are implicitly made about the data which directly influence the synthetic data, meaning it is not purely based on the properties of the unprocessed data.
  
  Similarly, if only a small amount of real-life data is collected, synthetic data generated using it may exhibit a high degree of bias. This can make it difficult to effectively extrapolate to other environments. Thus, it is still necessary to have access to as large and diverse a dataset as is possible given the contextual constraints to improve the generalizability of our synthetic data. The quality of the generated synthetic data must also be assessed, which can be problematic since the “quality” of a dataset can sometimes be very complex or specific, and not readily described by standard metrics.
  
  \section{Unsupervised Learning for Anomaly Detection (with K-Means)}
  
Now that you’ve learned the basics of anomaly detection and the data-based challenges it imposes, we will explore anomaly detection through the lens of a traditional ML technique, K-means clustering. In this Colab you will visualize anomaly detection and unsupervised learning as well as learn and evaluate the strengths and weaknesses of K-means clustering. We will be using both the SKLearn \href{https://scikit-learn.org/stable/modules/generated/sklearn.cluster.KMeans.html}{KMeans} and \href{https://scikit-learn.org/stable/modules/generated/sklearn.manifold.TSNE.html}{TSNE} libraries.

\url{https://colab.research.google.com/github/tinyMLx/colabs/blob/master/3-8-9-K-means.ipynb}

\section{Unsupervised Learning for Anomaly Detection with Autoencoders}

VIJAY JANAPA REDDI: In this video, I'm very
excited to introduce a whole new neural network architecture to you
called the autoencoders.
Before I do that, let's take a step back and look
at what you learned in the Colab.
One of the things that we introduced to you
was K-Means, which is a machine learning algorithm that is specifically
based on clustering, that allows you to sort of segment data
into different clusters as centroids.
Now hopefully by doing machine learning using K-Means,
you learned that one of the key problems that K-Means has as that data
set increases, or the number of observations
you have in your collection actually increases,
K-Means starts to fall apart.
Especially when you have high dimensional data, it does really poor.
But that still leaves us with this fundamental question about what
do we do with anomaly detection.
Remember, with anomaly detection, the problem
is, that we have a lot of data for the normal behavior
and those really important, yet rare events are the ones
that we care about the most.
So how do we go about solving a problem?
Well, let's take a step back and look at it more intuitively.
We have to come to terms with the fact that we don't know what we don't know.
Do I mean by that?
You and I don't know the 1,000 ways in which the machine might fail.
But there is one thing that we do know really well.
We do know how the machine is supposed to work.
You do know what you know.
In other words, if I take you back to this data
and show you the wealth of data that we have
here as to when the machine is working normally,
well then, you know how the machine is supposed to work normally.
You don't know the various anomalies that you might see,
but when everything is good, you know what that pattern is supposed to be.
Can we exploit that, is the key question?
And the answer is, yes.
We can exploit the fact that we have lots of normal data,
forgetting the anomaly itself.
How do we go about doing that?
What we want to do is, effectively, we want
to build in an ideal class, a neural network that
can act as an identity function.
In other words, it's supposed to be able to take an input
and regenerate the exact same input.
Can we train a neural network to do that?
The answer is, absolutely yes.
Because we have lots of normal data, we can certainly
architect a neural network to be able to take
a set of inputs on the left-hand side and split them out
on the right-hand side.
The key question is, what is this architecture in the middle?
And what does it need to look like?
Well, that's where the interesting aspects come in.
When you have a lot of input signals that are coming in,
there's a lot of noise in these signals.
What you really want to do, is you want to compress that input signal down
to its core fundamental elements.
In other words, you want to eliminate the noise that's in the signal
and really get to the essence of that signal.
And that's where the autoencoder comes to bear.
The autoencoder idea is this.
That we create an encoder network on the left-hand side,
whose job is to take this high dimensional data that's
coming in on the input neurons and then compress it down
into this core compressed latent space representation, which we can then
deconstruct via the decoder, reconstruct the original signal.
Effectively, this would mean that you have
a set of neurons that are on the left-hand side for the encoder, a set
of neurons on the right-hand side for the decoder with this sort
of architecture where you are compressing the information down
on the encoder side by pushing it through fewer and fewer and fewer
neurons, so that the network effectively learns to figure out the hidden code.
And then from that hidden code, you train the decoder side
of the network to learn how to correctly reconstruct the original image.
Now because we have a wealth of normal data,
we can certainly do the back propagation process
and repeatedly train this neural network through all the traditional mechanisms
that you have learned so far.
And in doing so, because you have so much of that normal data,
it is possible to get a reconstructed network that
minimizes this delta between what the input signal looks like
and what that output signal needs to look like.
In effect, you want to minimize the reconstruction loss.
To be more specific, let's say I give you this example here.
And I want this network to be trained repeatedly
until it's able to reproduce this output on the right-hand side.
At every stage, I'm going to take the input on the left,
I'm going to take the output on the right,
and I'm going to compare them and do a diff
and see what's the error, what's the reconstruction loss
or what's the Mean squared error between these two.
And I'm going to push the network repeatedly through the training process
until it minimizes that Mean Squared Error.
At that point, my network has learned for a given set of input
how to reconstruct that input when it sees that signal.
That's great, because this means I have really good low reconstruction error.
Now here's the critical thing.
Now that your network has actually learned
how to reconstruct your wealth of normal data, when
you throw in an abnormal data, like something
that's coming from MNIST Fashion data set instead of the numeric data set,
now suddenly when you ask it to reconstruct,
it's going to reconstruct something like this
that does not look anything similar to what the original left-hand image looks
like.
At this point, if you were to compare the two, what you would end up finding
is that the left-hand image does not match the right-hand image,
and therefore, you have a high error.
And that is the key nugget for anomaly detection.
Because voila!
Suddenly you have trained this network to raise a flag or say, hey, here's
an error that I have.
I have a high error on this, because I'm not able to reconstruct it.
That's a beautiful thing.
Because unlike traditional supervised-based learning methods
where you typically want that error to be as small as possible, here,
it's almost like you want the network to try and generate that image.
But it's not going to be able to generate that original image.
And when you do the diff between the two of them,
it's going to be significantly larger, because it only
knows how to reconstruct the normal data that it's trained on.
So when it sees an abnormal data, the error
is going to be large in its reconstruction.
And that, my friend, is the key for doing anomaly detection
using autoencoders.
And with that said, congratulations.
I'm super excited you learned about autoencoders, because they're
a very powerful piece of technology that everybody is quite excited about,
and it's got many applications.
Because anomaly detection is a very fundamental thing.
It can be used in so many creative ways to find all kinds of anomalies
that can be used to signal for predictive maintenance
or other kinds of applications.
So I'm really excited that you got the chance
to learn about anomaly detection.
One of the things that you'll be going through in your Colabs
is learning how to tune the anomaly detector so
that you have a deeper understanding of how
to calibrate the thresholds before you flag as something being anomalous
or not.
With that said, congratulations.
There are a couple other things you're going to pick up in this class,
and I'll see you back when I'm about to wrap up the course 2.

\section{Autoencoder Model Architecture}

  \begin{figure}
    \includegraphics[width=\textwidth]{TinyML/2 Applications/3-8-12-1.png}
    \caption{Autoencoder Model Architecture.}
\end{figure}  



You’ve previously been introduced to the basic concept of an autoencoder model. At its core, an autoencoder is trained to reconstruct data after compressing it into an embedding. You’ve learned how autoencoders can be applied to anomaly detection by training only on normal data and then using the reconstruction error to determine if a new input is similar to the normal training data. In this reading, we will examine autoencoders in more detail through a high level overview describing the strengths of different types of autoencoders. While autoencoders have many uses, including denoising, compression, and data generation, in this reading we will continue to focus on time-series anomaly detection. If you’d like to examine more applications of autoencoders check out this article.

\subsection{Types of Autoencoders}

It turns out that there are many kinds of autoencoders that are used in the wild. In this rest of this article we will survey some of the most popular types. To see how to implement all of these different autoencoders in keras check out this article.

F\subsection{ully Connected Autoencoder}

A deep fully-connected autoencoder has multiple fully connected layers with some small embedding in the middle. You’ve already seen an example of these in the previous video. They are the most basic form of an autoencoder and are simple to implement and deploy.

\subsection{Convolutional Autoencoder}

A convolutional autoencoder operates under the same principle as the deep fully-connected autoencoder but uses convolutions for the encoder layers and transpose convolutions for the decoder layers. Due to the spatial nature of convolutions, these autoencoders are particularly useful for images and spectrograms (which are essentially images). Unfortunately, at this time, transpose convolutions are less commonly supported by tinyML frameworks, so deploying convolutional autoencoders to microcontrollers is more challenging. That said, hopefully they will become more accessible soon. In the meantime, if you would like to try using a convolutional autoencoder for anomaly detection in Colab, check out this keras example. 

\subsection{LSTM Autoencoders}

Long short term memory networks are a type of recurrent neural network. They are useful for processing longer sequences while capturing the temporal structure of the data like the time-series data we have explored thus far. This can lead to a more detailed analysis of longer sequences which leads to higher accuracy in some cases. However, these RNN layers generally use more working memory as they maintain and modify a state over long sequences. This means they often consume more SRAM, which is a precious resource in tinyML. Additionally, it is less common for tinyML frameworks to support LSTMs or RNNs in general. For more detail on LSTM autoencoders check out this article.

\subsection{Variational Autoencoders}

Unlike the previous examples, the unique aspect of variational autoencoders is not that they use a different type of layer. Instead, variational autoencoders impose an additional constraint during training, which forces the model to map the input onto a distribution, which means the normal latent vector of an autoencoder is replaced by a mean vector and a standard deviation vector from which a latent vector is sampled. For anomaly detection, we can use the reconstruction probability that the model produces instead of the reconstruction error used in normal autoencoders, which can at times lead to better predictions. 

 \begin{figure}
    \includegraphics[width=\textwidth]{TinyML/2 Applications/3-8-12-2.png}
    \caption{The variational autoencoder model architecture showing how inputs are encoded into a mean and standard deviation before being encoded into the low dimensional latent space and then decoded back into the full input space. Training tries to ensure that the decoded reconstruction matches the input.}
\end{figure}  





\subsection{Visualizing the Latent Space of Variational Autoencoders}

There are multiple ways of visualizing the learned low-dimensional representation, or latent space, in an autoencoder. Due to the additional constraints imposed on the latent representation in variational autoencoders, visualizing the latent space is more straightforward than with traditional autoencoders. 

When variational autoencoders are trained on images, a common technique to visualize the latent space is to randomly draw values from the learned distributions and use these values as inputs to the decoder (ignoring the encoder entirely). This produces a novel image sampled from the latent space. For example, after training on MNIST digits, the resultant images will usually resemble specific digits. Note that these generated images are not present in the training set - they are representations of what the network has learned about encoding digits. In fact, by running linearly-spaced sample values through the decoder, you will notice that the output images appear to interpolate between digits. We are thus visualizing samples along a continuous latent space (i.e., a manifold):

 \begin{figure}
    \includegraphics[width=\textwidth]{TinyML/2 Applications/3-8-12-3.png}
    \caption{Visualizing the clustering of digits when mapped to the learned latent space. This clustering allows us to tell apart the digits as they are mapped into discrete areas of the latent space. This particular example show that in this particular dataset there are 9 main clusters and when mapped into two dimensions we can see that most of these clusters are quite well seperated although there are also a decent amount of overlaps.}
\end{figure}  





If your training images belong to specific classes, as MNIST digits do, you may also observe potential clustering within the latent space. You might expect each digit to be spatially separated from other digits, and indeed this is the case. The variational autoencoder in this article (which we used for inspiration for this article) encodes MNIST digits in a two-dimensional latent space. In the below figure, each test image from the MNIST dataset is fed through the encoder (the decoder is ignored). The two-dimensional mean vector of the encoder’s output is plotted on the x-y plane, and each digit (indicated by the legend on the right using different colors) is fairly well-clustered:


\begin{figure}
    \includegraphics[width=\textwidth]{TinyML/2 Applications/3-8-12-4.png}
    \caption{Visualizing the clustering of digits when mapped to the learned latent space. This clustering allows us to tell apart the digits as they are mapped into discrete areas of the latent space.}
\end{figure}  


For variational autoencoder networks with larger latent spaces (more than 2 dimensions), common dimensionality reduction algorithms like PCA, t-SNE, or UMAP may help with visualization, though these are outside the scope of our discussion.

Unfortunately, variational autoencoders are more complicated to train and deploy, therefore in the remainder of this course we will focus on standard autoencoders. However, since you may run into these new and powerful techniques in the future we wanted to make sure to introduce you to them now. Next, you will train a fully connected autoencoder to identify abnormal heart rhythms from ECG data and learn how to evaluate an anomaly detection model.
  
\section{Training and Metrics}

Now that you have been introduced to the basics of autoencoder model architectures, you will explore the process of training one to detect anomalous ECG data. This will get you familiar with the unique process of training only on normal data and allow you to play with the architecture of a fully connected autoencoder. Additionally, you will learn the two important metrics, the receiver operating characteristic (ROC) curve and the area under the ROC curve (AUC).

In this Colab you will be using TensorFlow to define and train the autoencoder and \href{https://scikit-learn.org/stable/modules/model_evaluation.html}{SKLearn} to evaluate its performance.

\url{https://colab.research.google.com/github/tinyMLx/colabs/blob/master/3-8-13-Autoencoders.ipynb}

\section{Picking a Threshold}

For the application of anomaly detection that is most interesting to you, what kinds of thresholds would you select? What factors influenced the sensitivity you select? How might you test that sensitivity? Do you expect it might be too conservative? Too aggressive? Why?

\section{Assignment: Training an Anomaly Detection Model}

Now that you know how to train an autoencoder for anomaly detection using only normal data, it’s time to test your skills. You will be repeating the ECG anomaly detection with a two key differences:

\begin{enumerate}
  \item The training set now includes a low percentage of anomalous data. This reflects a real-world scenario where you don’t have access to experts to label input data and therefore can not ensure that the training data contains only normal data. Since the anomalous portion of the training data is small, the model can still perform quite well as long as we don’t overfit.

  \item You will be tasked with determining the size of the autoencoder’s embedding layer and picking the error threshold to obtain the best performance given the application.
\end{enumerate}

\url{https://colab.research.google.com/github/tinyMLx/colabs/blob/master/3-8-16-Assignment.ipynb}

\section{Assignment Solution}

The size of the encoding layer should be less than 8 and more than 0 as we need to have a bottleneck from the previous layer (which was size 8). We have found: \PYTHON{EMBEDDING\_SIZE = 2} seems to work the best! The encoding layer size is a critical factor in how much information our model can encode. Remember, we want the encoding layer (and the model in general) to be large enough that it can encode the normal signals that represent 90\% of our training data but we want it to be small enough so that it can't learn to reproduce the 10% of our training data that is anomalous. For this dataset, the model and encoding layer can be very small, but it's important to remember that this is a very controlled and simple example.

The threshold depends on the training and on the encoding layer size but somewhere in the range of threshold = 0.037 works best. When picking the threshold the goal is to maximize the accuracy, precision, and recall. The ‘knee’ of the ROC curve is a good place to start as it represents a good balance between precision and recall. For this application, it might be more important to properly classify an abnormal rhythm as such, instead of maximizing the accuracy.

As always if you would like to go back and try the assignment again, it can be found here: 

\url{https://colab.research.google.com/github/tinyMLx/colabs/blob/master/3-8-16-Assignment.ipynb}

\section{Summary Reading}

In this section you explored the TinyML flow and data engineering in the context of anomaly detection applications, focusing on some unique challenges presented by unsupervised learning. Anomaly detection represents a common TinyML use case of classifying some time-series data as ‘normal’ or ‘abnormal’. Anomaly detection can be applied to data coming from a wide variety of sensors in a number of different application domains.



\begin{figure}
    \includegraphics[width=\textwidth]{TinyML/2 Applications/3-8-1.png}
    \caption{Machine learning workflow and pipeline. Begins with collection and preprocessing of data then the design and training of a model. These steps are done in TensorFlow. Then a model is optimized, converted and deployed using TensorFlow Lite. Finally inferences are made using TensorFlow Lite Micro.}
\end{figure} 


\subsection{Anomaly Detection in Industry}

You explored an example of anomaly detection in an industrial setting where you learned about the primary constraints of tinyML anomaly detection. First, anomaly detection systems must be fast enough to keep up with the rate at which sensors produce data. Often the sensors produce too much data to be sent over a wireless network therefore computation must be done locally. To accomplish this, the anomaly detection model must fit within the tight memory constraints of a microcontroller. Finally, an anomaly detection system must mitigate false negative and false positives, which, depending on the application, can have dire consequences. 

\subsection{Data and Datasets}

Anomaly detection is unique in that it is impractical to collect real examples of anomalous data. Collecting data on a failure is often too expensive since it usually involves intentionally breaking something. Additionally, anomalies are inherently rare and can take many different forms therefore it is difficult to train anomaly detection models in a normal supervised manner. That all said, you explored the MIMII and ToyADMOS datasets as well as the concept of generating synthetic data to understand some of the processes used to try to collect an anomaly detection dataset.

\subsection{Unsupervised Learning: K-Means and Autoencoders}

To avoid some of the issues with supervised learning you then explored a classical and neural network based approach to unsupervised learning. For the classical approach, you explored anomaly detection through the lens of K-means clustering and learned about its advantages (ease of implementation and use) and disadvantages (scalability). You then learned about autoencoders and how they can be applied to many domains, among which is anomaly detection. Autoencoders are trained to reconstruct normal data after compressing it into an embedding. By comparing the error between the input and the output of the autoencoder we are able to determine how similar the input is to our normal training data. We perform anomaly detection by selecting a threshold above which we classify the input as an anomaly. You also learned about how the model architecture impacts the autoencoder and how metics like ROC and AUC help us evaluate our performance. Critically, you learned that anomaly detection models often have low transferability and therefore have to be trained for their specific application.  
  
% 1.9: Responsible AI Development

\section{Data Collection}

SUSAN KENNEDY: In an effort to build responsible AI,
we want to think about some issues that might arise during the development
stage.
The question we'll focus on now is, what data
will be collected to train the model?
Well, data offers great benefits from the perspective of those designing ML.
It also carries great risks to the individuals
who are contributing their data.
So perhaps unsurprisingly then, there are emerging laws and regulations
surrounding data collection and use.
For example, in 2018, the European Union passed the General Data Protection
Regulation, a legislative act meant to protect individuals' privacy
and prescribe ways for organizations to handle personal data.
But no matter what country you live in, it's
worth taking a closer look at these regulations
because it's predicted that most countries will eventually
either adopt the GDPR or create legislation similar to it
in the near future.
There are roughly three key features to data protection, identifiability, data
minimization, and notice and consent.
So we'll go over each in a bit more detail now.
Identifiability refers to a requirement for data collectors
to reduce and safeguard identifiable components of data as much as possible.
To reduce identifiable components, we want
to think about de-identification techniques.
So we can either fully anonymize data or we
could employ a pseudonymization, which creates
an alias to hide the subject's identity but where we retain
a key to reverse this if necessary.
To safeguard data, there are three things to keep in mind.
Data Encryption, which encodes information
so it can only be accessed with permission.
Secure Servers and the Storage Location.
For instance, storage in the cloud offers benefits.
But in what nation and under what laws is the cloud operating in?
It's also important to distinguish between different classes of data
because it will determine how careful we need to be.
With non-personal data, like aggregated statistics,
you're in the clear to use it and there's no regulations.
But with personal data, which might include someone's name, location,
IP address, and so forth, we need to be more careful.
And sensitive data, like genetic, biometric, and health
data, or data that reveals race and sex, you'll
need to exercise extreme caution.
Another feature of data protection is minimization.
Here the goal is to limit data collection and the duration of storage
to only what's required to fulfill a specific purpose.
It used to be that companies would want to retain data indefinitely.
But as the influx of data increased dramatically and the cost of storage
became an issue, data minimization has become
more important both practically speaking and in terms of laws and regulation.
Questions to ask yourself are, how long you'll
need the data to achieve the purpose, whether there's
unnecessary data that can be deleted, and how often you'll review and delete
what isn't needed.
Related to data minimization is the right to be forgotten
or the right to erasure outlined in the GDPR.
This states that data subjects have the right
to request that their data be erased as soon as possible.
But this right might be overridden in some cases.
For example, when the data is needed to comply with other legal obligations,
or achieve some purpose in the public interest.
So what does this mean for data collectors?
How will it change your workflow?
You'll want to provide subjects with clear information and practical ways
to make a request for data erasure.
The third feature will turn our attention to is notice and consent.
The aim here is to provide notice to subjects about how
their data is planning to be used.
And get their consent so that subjects have the option
to choose if they want to be involved or not.
So what counts as valid consent?
Well, this concept is very similar to informed consent that
was established in the medical field.
We can distinguish a couple of different components.
The decision has to be informed, meaning the subject has sufficient knowledge
and comprehension to make a decision.
And this rules out lies, deceit, or partial disclosure.
The decision must be voluntary, where the subject freely
chooses to give consent, ruling out coercion or inappropriate pressure
or influence.
And the subject must be competent, having the decisional capacity required
to offer consent, which rules out children or adults deemed mentally
incompetent.
Now we'll turn to a related question of how the data will be collected.
Earlier in this course, you were introduced
to some mechanisms for data collection, like crowdsourcing, product users,
and paid contributors.
As you're probably aware, good representative data is crucial.
So it's important to anticipate possible disadvantages of a chosen method.
As Kate Crawford, principal researcher at Microsoft
and professor at NYU Tandon School of Engineering
has said, "We need to ask which people are excluded.
Which places are less visible?
What happens if you live in the shadow of big data sets?"
For example, consider a case of data collected from product users.
An app called Street Bump was designed to identify and report
potholes that require repair.
But since the app was deployed on smartphones,
whose users tend to be young with a higher income,
the data didn't accurately reflect the number
of potholes in poor neighborhoods or those with elderly residents.
Problems with data collection are closely related to issues of bias.
So we'll talk about that more in an upcoming video.
Finally, collecting the right data can be complicated, costly,
and time consuming.
That's why it's especially inspiring to see a movement in the ML community
towards sharing open source, publicly available data
sets, like the Mozilla Common Voice data set,
for instance, which even offers a breakdown of demographics.
This information is practically useful for people
intending to use the data set.
And it also makes clear where the data set could be improved.
Not only does this make working on ML easier than ever before,
but this increased accessibility is great
because it fosters innovation and healthy competition.
And this is how we'll end up with better products that
improve the world we live in.
  
\section{The Many Faces of Bias in ML}

One of the promises of machine learning is that it can potentially offer a superior alternative to human decision-making. Algorithms offer objective interpretations of data and can be trained and developed to produce consistent results, unlike human decision-making which can oftentimes be subjective and inconsistent. Thus, an advantage of machine learning is that it isn’t susceptible to ‘human error’ or mistakes. Even more impressive is the potential for machine learning to predict outcomes with a level of accuracy that surpasses that of the world’s leading experts.

While these are exciting possibilities, one of the major challenges of machine learning is the problem of bias. The presence of bias can be the downfall of an otherwise perfect model, resulting in skewed outcomes and low accuracy. This can be a frustrating experience for those who dedicate their time and energy into creating a model, only to discover that it doesn’t work very well. Even more concerning, however, is the possibility that if bias finds a foothold in machine learning, the model may effectively function to automate bias and inequality and deploy it at scale. When all goes well, machine learning can be consistently accurate. But the presence of bias can cause machine learning to be consistently inaccurate in such a way that leads to discriminatory or unfair outcomes.

Recent discussions of bias have primarily focused on the use of machine learning in high-stakes contexts where inaccurate outputs can have major consequences. For instance, when algorithms are used to inform decisions regarding bail in the criminal justice system, the granting of loans in financial institutions, or the hiring of job applicants, the presence of bias can result in some individuals or groups being systematically disadvantaged in a way that seriously impacts their life prospects. As one might expect, biased algorithms in high stakes contexts have made national news headlines, as demonstrated by Amazon’s hiring algorithm that revealed gender bias and the COMPAS recidivism tool that revealed racial bias.

While the negative impact of bias is most apparent in high stakes contexts, bias remains a serious challenge for any application of machine learning—including tinyML. This may be somewhat surprising when one considers the kinds of tasks with respect to inference that tinyML is meant to tackle with the use of microcontrollers. For instance, if bias leads a keyword spotting or activity detection application to be inaccurate, how problematic could this be? Even though tinyML will not be responsible for complex inference tasks in high-stakes contexts, the unique features of tinyML raise a different kind of concern. 

TinyML is at the cutting-edge in terms of low latency, low power, and small size. This is significant because it enables tinyML to be pervasive and find applications across a wide range of contexts. The number of IoT devices has already grown considerably in recent years, and with the emergence of tinyML this number is expected to surge upwards. It’s precisely because tinyML is uniquely capable of being everywhere, surrounding us, that the problem of bias remains a major concern. If one were to consider the impact of a single device that is inaccurate due to bias, it may not seem especially problematic. But when one considers the likelihood that tinyML will soon become a feature of a wide variety of our daily interactions with machine learning, it’s easy to see how the presence of bias could amount to a more serious concern. For example, if an individual fails to trigger the wake word of a personal voice assistant, this may seem like a minor inconvenience. But if this individual also fails to trigger the wake word for their light switch, thermostat, coffee maker and so forth, then this individual would seem to be systematically disadvantaged insofar as they cannot enjoy the benefits of machine learning applications.

In order for tinyML to make our lives easier and more accessible than ever before, designers will need to remain vigilant in rooting out bias from their applications. So how does bias find its foothold in machine learning algorithms? One reason this can happen is because humans are still responsible for creating machine learning applications, so whatever biases are present in the data scientists may be carried over into the project. The most common reason machine learning becomes biased, however, is owing to the underlying data. 

For instance, a dataset may be biased if it is incomplete. If the dataset doesn’t fully represent the diversity of end users, then this can lead the model to be less accurate for certain individuals or groups. Researchers have shown that voice recognition tools struggle to identify African American Vernacular English, leading widely popular personal voice assistants to work less well for black individuals. Similarly, research shows that voice recognition struggles to identify non-native English speakers or those with speech impairments. After recognizing the importance of making products accessible to everyone, Google launched “Project Euphonia” in 2019 with the goal of collecting more data from individuals with speech impairments or heavy accents to fill the gaps in current voice datasets. 

However, tackling the problem of bias isn’t always as simple as collecting more data from underrepresented users. Rather, bias has many faces and can present at any stage of the data lifecycle, including creation, collection, sampling and processing. A general piece of advice is to pay close attention to the data that is being used to develop an algorithm. It is worth reflecting on questions like the following: From whom is the data being collected and is it representative of the end-user population? What device is being used to collect the data? What decisions have been made about how to process and label the data?

Carefully analyzing the data is a huge step towards mitigating bias, as identifying the exact form of bias is necessary for making a determination about how to address it. For instance, if the dataset is incomplete, then the solution may be collecting additional data that better represents the end-user population. Alternatively, if the data is faulty due to labeling errors then the solution may be conducting a review of the data to correct these errors. In the next video, you will learn to identify some common forms of bias. In addition, you will learn about some potential solutions that are being developed on the industry level, as well as what you can do as an individual designer to tackle the problem of bias in machine learning.

\section{Biased Datasets}

SUSAN KENNEDY: Another question to address during the development stage
is whether the data set is biased.
First, we'll want to be clear about what we mean when we talk about bias.
In general, it's best to understand bias as a deviation
in a predictable direction, not a random one.
This is important because it helps us distinguish between those errors that
aren't attributable to bias.
So in the context of ML, bias refers to systematic errors.
This is oftentimes because of bias data sets,
meaning faulty, poor, or incomplete data,
which can lead to a distribution mismatch between the data set
and reality.
Moreover, bias can not only lead to inaccurate results, but in some cases
discriminatory or unfair results, which we'll
discuss more in an upcoming video.
Now I'm sure you've heard how important data is for machine learning.
If we have bad data inputs, we'll end up with bad results.
And that's even if we have an otherwise perfect model.
This is sometimes referred to as the "garbage in, garbage out" problem.
But how does our data become biased in the first place?
Well, there are actually a lot of ways that this
can happen, ranging from conditions under which the data was created,
how it was collected and sampled, and how it was processed.
Unfortunately, bias can basically enter the picture
at any point in the data lifecycle, making it a big challenge for ML.
In order to tackle this challenge, I'll introduce you
to some common forms of bias in data so you know what to look out for.
One way bias enters the picture is owing to how we define the target variable.
This gives our model an objective, but it can sometimes
be a rather subjective matter, making it more
of an art form than a perfect science.
For example, if we're designing an algorithm
to make promotion recommendations, how do you define a good employee
and how could you put this in terms that the computer could understand?
Well, you have to come up with some quantifiable measures, maybe
the number of sick days, number of times late to work,
number of sales per quarter, and the length of employment.
But defining the target variable in this way might unintentionally disadvantage
some class of people.
For example, the number of sick days might
disadvantage pregnant women or individuals
with disabilities missing work.
Or the number of times late to work might disadvantage those people
who rely on public transportation.
How might this apply to cases of Tiny ML?
Well imagine you're using biometric sensors for a health wearable device.
How should you define healthy?
Maybe you could focus on heart rate, blood pressure, or the number of steps.
While these might seem like innocuous measures, focusing on blood pressure
might disadvantage the elderly, who tend to have higher blood pressure even when
they have heart healthy habits.
Another way bias enters the picture is through the labeling of data.
The reason is that the labels we apply to the training data
must serve as the ground truth.
So as far as our model can tell, the labels perfectly match reality.
But what happens when there's mislabeled data?
For example, if we build a horse versus human classifier
and some of the images of horses are actually mislabeled as human.
Even if it's obvious to us that the picture is of a horse,
the computer only knows what's in the picture based off
of the label we gave it.
And so the model will start to misclassify images of horses as human.
While this example might seem easy to spot with a careful review of the data,
bias in labeling can sometimes be much more difficult to address.
Suppose we're labeling medical images for a disease classifier
and we present a set of patient images to train
graders who label the data for us.
While the two extremes might show a clear pattern,
if we look at this graph, we'll see that there's
disagreement among the experts about which label best
applies to a particular image, in which case, we
have to make a rather controversial decision about how to label
these images and train our model.
Another common form of bias relates to how
we sample the data that's used to train the model.
Imagine we're training a model to respond to a wake word
and our data set looks like this with mostly white individuals.
Now let's suppose that this data doesn't accurately
represent reality, where our users are significantly more diverse.
What will end up happening is that our model
will perform less well for these individuals who aren't represented
in the training data.
Bias might also be a result of prejudice or historical injustices reflected
in the data.
For example, the company Northpointe relied on an algorithm
to predict recidivism risk.
And despite equal accuracy in terms of correct predictions,
black defendants were almost twice as likely as whites
to be labeled as higher risk, but not actually reoffend.
And the opposite mistake was made among whites,
who were much more likely to be labeled lower risk
but then go on to commit other crimes.
This is attributed to historical injustice
in policing practices that's been disproportionately focused
on minority populations, resulting in higher arrest
rates among these populations, which then gets reflected in the data set.
This type of bias might also occur when stereotypes skew the data set.
For example, the stereotype that women are primarily
responsible for cooking and other household chores
might result in women being overrepresented in a data set that
includes images of people cooking.
This could then cause the algorithm to be
inaccurate by predicting women or cooking a majority of the time, even
when this actually isn't the case.
Bias can also be attributed to measurement distortion, which
most often occurs when there's a problem with the device that's
making a measurement or an observation.
To put it simply, hardware matters.
For example, optic sensors and cameras have
problems detecting darker skin tones due to the way light is reflected
and this has led to notable issues with automatic soap dispensers
and tools for activity detection and facial recognition.
Additionally, this can happen when the device used
to collect data for training differs from the device that's
used when the model's deployed.
If our data set contains images from a DSLR camera,
but the device our model is deployed on is a mobile phone,
as these images show, the value distortion
will cause the model to generate bad results.
So how can we fix biased data?
Well given the scope of the problem it most importantly
requires us to remain vigilant and keep a lookout for potential bias.
But fortunately, there's some industry wide solutions being proposed
as well as questions to keep in mind to help guide individual designers.
One industry solution called Datasheets for Datasets was published in a paper
by researchers from Google and Microsoft as well as academics.
They offer a list of questions for data set creators
to reflect on during the key stages of the data set lifecycle, example,
the motivation, composition, collection, distribution,
and maintenance of the data set.
Similarly, there are data nutrition labels,
a standard label that highlights the key ingredients of the data set,
including provenance, metadata, missing units, and so forth.
This helps someone at a glance get a better sense of what the data contains
and make a more informed decision about whether it's an appropriate data set
to use for their project.
There are also biased testing toolkits coming out.
For example IBM offers its AI Fairness 360 tool
that makes biased testing a part of your development cycle.
And similar tools are being offered by Google and Microsoft.
Now as an individual designer, there are some helpful tips
to keep in mind to mitigate bias.
You want to carefully research users in advance
and be aware of potential outliers.
Ensure your team of data scientists and data labellers are diverse.
Where possible, combine inputs from multiple sources
to ensure data diversity.
Create a gold standard for data labeling and seek out domain experts
to review your data.
Now that you're familiar with the problem of biased data sets,
you're one step closer to developing reliable and accurate tiny ML.

\section{Forum}

Now that you’ve become acquainted with different forms of bias and potential solutions, reflect on the following questions and discuss your response with other students taking this course!

Which form of bias did you find most surprising or interesting? Which form of bias do you think poses the most serious challenge to TinyML? Why? 

As a reminder, these are some common forms of bias discussed in the video:

\begin{itemize}
 \item Defining the target variable
 \item Labeling the data
 \item Sampling the data
 \item Prejudice reflected in the data
 \item Measurement distortion
\end{itemize}

Which solutions for bias (either industry or individual designer) do you think are most promising? Are there any obstacles you can foresee for the solutions you think are most promising? 

As a reminder, industry solutions included datasheets for datasets, data nutrition labels, and bias testing toolkits. Individual designer solutions included recommendations such as carefully researching users in advance, ensuring your team of data scientists and data labelers are diverse, combining inputs from multiple sources to ensure data diversity, creating a gold standard for data labeling, and seeking out domain experts to review your data.

\section{Fairness}

SUSAN KENNEDY: The lesson question we'll consider during the development stage
is how we can ensure that our model is fair.
While this topic is still being debated by ethicists and computer scientists,
a general way to think about unfairness in ML
is a model that exhibits discriminatory biases,
perpetuates inequality, or performs less well
for historically disadvantaged groups.
But it's important to realize how the term discriminate can
be used in different ways.
Technically speaking, all ML discriminates, since this
just means to recognize the distinction or differentiate.
But what we're talking about in the context of fairness
is wrongful discrimination.
Wrongful discrimination typically involves disadvantaging people
in protected classes, based on characteristics
that you can see listed here, like age, disability, medical history, sex, race,
and so forth.
In US law, wrongful discrimination falls into two categories--
disparate treatment, where membership in a protected class is used as an input
to the model, and decisions are differentiated
on that basis in a way that disadvantages members of the protected
class, and disparate impact--
where outcomes of the model disproportionately disadvantage members
of the protected class.
However, given the unique features of AI,
it's proven difficult to show how cases of unfairness
might fit into these categories.
So instead of focusing on how AI might be considered unfair,
we'll now look at different metrics that have
been devised to help promote fairness.
One metric is group unawareness, where sensitive attributes,
like race or gender are not included as features of the data.
Benefit of this metric is that it avoids disparate treatment.
But a disadvantage is that there might be highly correlated features that
serve as proxies for the sensitive attribute
and end up producing discriminatory results nonetheless.
Another metric is called group threshold,
where we counteract historical biases in data
by adjusting confidence thresholds independently for each group.
So if group A has typically been advantaged relative to group B,
we might decrease the confidence threshold for group B,
so our model doesn't unfairly exclude members of this group.
A potential problem for this metric of fairness
is reaching agreement about the existence of historical bias,
as well as how the confidence thresholds should be adjusted accordingly.
A third metric is demographic parity.
Here, we're concerned with the positive rate being equal across groups.
Suppose we have a health wearable device that predicts if someone's healthy,
and an employer is using this on the basis of offering bonuses to employees.
In this case, we might want all individuals to have
an equal rate of positive results.
But now, let's take a closer look.
If one group is more qualified, and so more
likely to be healthy in our example than the other group,
then achieving equal positive rates might
mean that we have to exclude some qualified members
of the advantaged group.
And this raises concerns about reverse discrimination,
as you may be familiar with from debates about affirmative action policies.
The fourth metric of fairness is equal opportunity,
where we want to make sure that among the people who are actually qualified--
and so actually healthy, in our example--
that they have an equal chance of being accurately classified as such.
So we're primarily focused on the true positive rate,
or the relationship between true positives and false negatives.
For group A, there are nine true positives and three false negatives.
So our true positive rate is 9 out of 12, or 75%.
For group B, we have three true positives to one false negative.
And so our true positive rate is three out of four, or again, 75%.
But if we look closer, we might notice that we've
introduced a lot of false positives for group B.
And this might be problematic in certain contexts.
The last metric of fairness we'll discuss
is equal accuracy, or the percentage of correct classifications as actually
healthy or not, should be the same for all individuals.
This refers to the overall accuracy of the model.
A problem that might arise is that even when
the accuracy rate is equal for group A and group B,
the errors that end up occurring might be different for each.
For example, group A may have a higher rate,
of false negatives, while group B has a higher rate of false positives.
If you remember from the previous video, when
we talked about bias and Northpointe's recidivism tool,
this is what ended up happening when black defendants had a higher
rate of false positives, while white defendants enjoyed
a higher rate of false negatives.
When the model's results informed decisions
about the distribution of some benefit or harm,
differences between false negatives and false positives
can make a big difference.
While you might think that the right thing to do
is satisfy all of these metrics the best that we can,
researchers in computer science have shown
that we can't do this, a problem known as the impossibility theorem.
For example, Group Unawareness, where we don't consider the sensitive attribute,
is incompatible with Group Threshold, where
we set the confidence threshold independently,
based on the sensitive attribute.
An Equal Opportunity, or the true positive rate,
is equal across groups, is incompatible with Equal Accuracy,
where the accuracy rate is equal across groups.
Now that we've covered different metrics for fairness,
we'll focus on some other tips and tricks for mitigating unfairness in ML.
One common pitfalls referred to as the Framing Trap.
Typically, designers operate within the algorithmic frame,
where they consider questions like whether the properties of the output
match the input, and whether the algorithm provides
good accuracy on unseen data.
But this is only one important lens through which
you should evaluate your model.
You'll also want to consider the Data Frame,
as we talked about earlier, where you consider whether bias has been removed
from the training data, and whether the demographic information
requires optimization through one of the fairness metrics we just discussed.
Lastly, you'll want to consider the Sociotechnical Frame.
How does the model operate when considered
as part of a system of humans and social institutions?
For example, you might design a fair algorithm for predicting health.
But if the employers who use this to inform their decisions about handing
out bonuses routinely ignore the recommendations of the algorithm,
then your model might indirectly result in unfairness,
despite how you've optimized it.
Another common pitfall is the portability trap.
It's commonly regarded as a benefit whenever
we can reuse an algorithm from one context to another.
For example, maybe we design a tiny ML algorithm
to add water conservation efforts in California
and optimize it for fairness.
And since the model is successful in this context,
you might think that we can similarly use it in Nigeria, Africa.
But context matters.
Repurposing algorithmic solutions may not preserve fair outcomes.
So you'll always want to do your research before deploying
your model in a new context.
The last common pitfall we'll cover is the Formalism Trap.
If you're working on ML, you might think that the ultimate question that
needs to be answered is which mathematical formulation of fairness
you should choose.
But another important question to ask is which
social groups you'll listen to and include in your deliberation,
as focusing on one social group might give you
a conflicting result about the best metric for fairness
than if you had focused on a different group.
While issues of fairness can be tricky, and researchers are still
making progress in this area, it's becoming far easier
for individual designers to tackle this issue now,
thanks to publicly available tools like Google's What-If tool.
This allows you to explore the different metrics
and help you optimize for fairness in your model.

\section{Google’s What-If Tool}

Now that you’ve gained familiarity with some of the common metrics for fairness in machine learning, you will have the chance to explore Google’s What-If Tool (WIT). In this Colab we will walk you through exploring a real-life dataset and identifying sources of unfairness in the data across different social groups. In this case, we will train a DNN on the UCI census problem (predicting whether a person earns more than \$50K from their census information), and visualize different ways to optimize the model for fairness using Google's What-If Tool.

\url{https://colab.research.google.com/github/tinyMLx/colabs/blob/master/3-10-7-WIT}.ipynb

\section{Forum}

When optimizing for Equal Accuracy in Google’s What-If tool for fairness metrics, we get the following results:

\begin{figure}
    \includegraphics[width=\textwidth]{TinyML/2 Applications/3-10-8-1.png}
    \caption{Overall results showing the threshold for females (0.21) being much lower than the threshold for males (0.4) in order to have equal accuracy. This results in a much higher false negative rate for males (19.5\% vs. 3.9\%) and a much high false positive rate for females (21.7\% vs. 2,6\%).}
\end{figure}




Confusion Matrix for Males


\begin{figure}
    \includegraphics[width=\textwidth]{TinyML/2 Applications/3-10-8-n.png}
    \caption{2x2 Confusion Matrix as described in the earlier section (e.g., on diagonal are correct and off-diagonal are incorrect results) for Males showing the high false negative rate (21\%) for males.}
\end{figure}



Confusion Matrix for Females



\begin{figure}
    \includegraphics[width=\textwidth]{TinyML/2 Applications/3-10-8-n1.png}
    \caption{2x2 Confusion Matrix as described in the earlier section (e.g., on diagonal are correct and off-diagonal are incorrect results) for Females showing the high false positive rate (20\%) for females.}
\end{figure}



Take a closer look at the distribution of inaccurate predictions for males vs females - What do you notice? 

Now let’s suppose we are approving loans on the basis of positive predictions. How might optimizing the model for Equal Accuracy be problematic? Do you think it would be better to avoid false positives more than false negatives or vice versa? Why? In which situations do you think one would be more costly than the other?
  
% 1.10: Summary

\section{Summary}

VIJAY JANAPA REDDI: Congratulations on making this far through Course 2
and completing it because you have learned quite a bit
about TinyML and TinyML applications.
In this video, I'm going to walk you through a general overview
of the big core concepts that we touched on.
We touched on many of them.
So I want to make sure you feel good about learning a lot of things.
Specifically, I walked you through this pipeline
here, which goes all the way from collecting data to making inferences.
Sure, we did much of this in the Google Colab environment but that said,
we learned all the fundamental building blocks
that we need to actually be ready to deploy things onto a real device.
By the way, you actually used the real tools
that real engineers out in the world care about, so good for you.
Now one of the things that we did as we went through this entire workflow
is that I hope you recognize the machine learning code is just
one portion of a bigger puzzle that we are dealing with.
Nowadays, industries really want you to know that machine learning code
itself is part of the big infrastructure that they have to deploy.
Which means you have to know all the core building
blocks that are around there.
Therefore I introduced this lifecycle of ML, which starts from data collection,
goes into data ingestion, data analysis and curation, data labeling, data
validation, data preparation, model training, model evaluation
and so forth.
It's extremely key that you understand this big picture
because every machine learning application has to go through this.
And I spent a lot of time talking about data
and I will do that yet again in this video
because if you get the data part wrong, everything else you have done
is a waste of time.
So one thing we did was to kind of step through this huge gamut
of all these different blocks through the use of different sensors.
Remember, TinyML is all about sensors and data that's coming in from sensors.
So we spent a little bit of time understanding
what are the characteristics of the sensor.
Each of the sensors that we discussed, the microphone, the camera
sensor, the accelerometer sensor, everything that we did
has a unique property and that is the basis for why we chose these sensors.
So hopefully you picked something up interesting there.
Then we talked about different kinds of networks
and then we talked about different metrics to evaluate.
I'll go through each of these in a little bit of detail
as a quick recap of the specific things you
learned but the key thing on this slide is that you learned the entire pipe
and that's important as you go forward as an ML engineer.
Let's kind of slice and dice through this vertical
stack here from data engineering to building the product
and let's try and see what are the core technical nuggets that you picked up.
I'm going to use several example slides that come from all the material
that we built previously so that it serves as a reminder in your head.
That said, I'm not going to cover them in a lot of detail
and I'm only going to touch on the highlighted key messages.
But please, if something does not resonate, go back and look
at the original slides.
So the first thing was to actually look at data collection.
So we used Pete's speech commands data set for keyword spotting.
There were a lot of things that we discussed
around this, which specifically involve the motivations for building a data set
and how to go about engineering the data set.
Why did Pete build a data set all by himself?
Well, because there was no public data set that was available.
And in that process we laid out the groundwork
for what are the systematic steps that you should take when you're
trying to build a data set, identify the key requirements,
figure out how you're going to gather the data,
understand how you're going to refine the data
and decide how you're going to sustain that data
so that you can continue improvements over it.
There are several details that we stepped through,
all these different kinds of bullets here,
but specifically, one of the things that I really want to focus on
is around the licensing, make sure you know where your data is coming
from because you can build an awesome application at the end of the day,
but if your data is polluted, you can get into a lot of trouble.
So if you're building a neural network, know your data source.
Also, while there are a lot of challenges in building
the data set from scratch why not ask this question that if we can actually
repurpose an existing data set and that's
something that we introduced to you.
We asked you to think about the fact that in the case of Visual Wake Words,
we learned that we can actually take an existing data set which
has a wealth of different kinds of data in it
and subset that data so that we can use it
for Tiny ML, where TinyML is specifically
focusing on very specific tasks.
For example, we looked at the MS COCO data set which Visual Wake Words uses
and then it subsets that data set, specifically picking out
these individual peoples being able to detect for person detection application
use cases.
It is a remarkably quick way to shorten the whole data processing pipeline
as long as you have the right licenses.
Of course, you can never have enough data.
One of the things that I talk to you in the context of speech
was that there's a long tail, that is if you
want to be able to support a diverse population as you
try to increase the diversity, you get less and less data.
In the case of languages, we don't have that much data
for a good number of languages.
So how can we go about addressing this problem?
Sure we can either build the data set from scratch and we can hire people
or we can do something like Common Voice did, where we're actually
building a platform that allows people to contribute
their own languages into the community so
that we can build a community data set for different languages for speech.
And even after doing all of this there are
things that we have to be careful about, specifically around biases.
We spent a lot of time about understanding responsible AI
development and in the context of building
data sets it comes down to understanding the sources of different biases
in the data set.
And even if you understand these biases well, you have to worry about the fact
that often when you're trying to build a large data set,
you are hiring third party companies, so biases might implicitly
creep into your data set and you have to be very careful
and it is your responsibility to make sure that biases are not in your data
set.
Then how do we pre-process the data?
For example, in the context of the analog signal that's
coming in from the microphone in keyword spotting,
we talked about how we don't necessarily want to use the raw audio signal.
Instead we want to figure out how to extract or do feature engineering
to extract out the critical and salient parts of that signal
and represent that data in a form that is actually easy for our networks
to consume.
So for instance, we took an audio signal,
converted that into a visual spectrogram and then used
a convolutional neural network in order to be able to recognize the keywords.
Then we went into model engineering or designing a model.
One of the key challenges that we get into when we're designing models
is that because we're dealing with Tiny ML devices
they tend to be really small.
Therefore, we need to think about how can we shrink down
the number of parameters and the number of operations
required to compute a particular model.
I said that we will take one particular example, which
happened to be mobilenet and we looked into it
and we learned how by changing the fundamental operations that
are being performed from convolutions to depth
wise separable convolutions, that includes point wise convolutions,
we could dramatically cut down as much as 100x factor reduction in terms
of the number of computations that need to be performed,
which is a huge improvement in terms of model parameter size
and the computations required.
So sometimes you need to look not at just the overall network architecture
but actually go deep in understanding the fundamental building
blocks that actually repeat over and over throughout the whole network.
And so it's important to look at the algorithms that are there.
We also want to think about training time for models.
This requires you to kind of look at not just
do I always train the network from scratch or can
I do something interesting.
Remember I said that when you are learning
how to drive a car for the first time, you spend a lot of time
learning how to drive and you go and get a test
and you do all of that tedious process, but after you've
learned how to drive the car, when you upgrade your car,
you simply adapt to it.
The same concept was introduced to you with transfer learning,
where we said that you don't have to retrain the whole network from scratch.
Instead you can hold certain general parts of the network locked in.
And then only retrain the path specific layers for your given particular task.
So instead of being able to say it's a cat versus a dog,
well because they both have ears, eyes, nose and four legs
it's easy to be able to freeze the general characteristics
and then just fine tune them for being able to tell what's a cat versus a dog.
So you learned a whole new method that's called transfer
learning in order to shrink down the training time quite dramatically.
Also I introduced a different kind of neural network
in the anomaly detection section where specifically I
introduced what we were calling autoencoders,
where the autoencoders are dealing with unsupervised learning based issues
where you have a wealth of data for the common case
but then for those rare events that you need to pick up,
you don't have enough data.
So instead, you train your network on the common data
and then when the common data pattern does not
get reconstructed on the output side then you
can detect that something is off and that effectively is the anomaly flag.
In addition, we learned how to evaluate these models that we're building.
Specifically, we learned that if we want to make the models run fast then
we should do quantization.
Quantization is something that is going to be key for every TinyML engineer.
So hopefully you remember a little bit about post training quantization
and quantization aware training and what are the differences between the two.
Also when you're doing quantization I hope
you remember why the size of the model actually goes down,
it happens because you're going from a floating point value to an integer
value to represent the weight.
And then how does the performance of the network actually improve?
That's because we're doing integer to integer multiply accumulates.
Together that reduces the size as well as improves
the performance which is perfect for TinyML sorts of devices.
So make sure you know about quantization and what are all the nuances around it.
Then we went into converting the models that we
are training in TensorFlow into appropriate formats for deployment.
Specifically, we learned from Lawrence that the TensorFlow ecosystem
has many, many different ways of building and deploying things.
We specifically focused on TensorFlow versus TensorFlow Lite.
The key reason is that TensorFlow is both a training
and an inference engine versus TensorFlow Lite
is just an inference framework that is designed for running things efficiently
on edge end point devices like mobile devices
and TensorFlow Lite Micro is specifically tuned
for microcontrollers.
More on that in the next course.
But one of the reasons we have these differences
is because there are aspects of whether we're training the model
or we're just running the model that matter
because we can remove a lot of code if we're not training the model
Similarly, there are differences in terms
of what support is needed between TensorFlow and TensorFlow Lite
from a software standpoint.
A lot of devices, for instance in the embedded space,
do not have an operating system.
Well that has consequences on which of the frameworks you would use.
Also, the hardware sometimes might support or might not
support floating point for example.
You need to understand some of these fundamental reasons
as to why you need different sorts of frameworks.
So as a TinyML an engineer I expect you to know these differences
because if you're going to try and deploy something,
you need to be using the right tools to be
able to get to that deployment stage.
Next the question is about deployment itself, which goes into metrics.
So the metrics can be looked at in terms of quantitative perspective
as well as qualitative perspective.
Quantitative perspectives are what we're usually used to.
For example, many of us will look at how well a particular model is
doing but remember, I said that the model alone
is not enough to actually determine if it's good for deployment, why?
For example, you might have a very accurate model
but because it's accurate and it's big it could be really slow,
which means a user is not going to be happy.
For example, if I say OK Google and it takes three seconds
before the machine responds, that's not good, even if it's 100% accurate.
Well, therefore you might want to make some compromises,
you drop the quality by using a smaller model but it's faster.
So perhaps you can respond more quickly to me
and that allows you to actually enable your TinyML application.
Also there is no one fixed calibration point
that makes it suitable for deployment.
In the context of keyword spotting we understood
that we have this operating point that we
have to move along in the false accepts versus false rejects case
and that's important because you are deploying
this device for a diverse number of users.
So there is no one fixed static point on this curve
that is actually good for deployment.
You would have to tune it depending on where you're deploying it.
Also we introduce this notion of streaming audio and cascade
architectures to enable powerful capabilities
and that TinyML is part of a bigger picture of how you
deploy a machine learning application.
So hopefully you remember this particular slide
where we talked about how Siri play my music,
for instance would actually work.
Where there's a mix of streaming audio coupled
with cascade architectures and then finally being able to use the cloud.
This whole system works together.
And as a TinyML engineer, you should know this
because it has implications in terms of how you actually design and deploy
your application.
Finally, there is this aspect of product analytics
where we actually talk about making inferences
and how that is actually integrated into dashboards where
you're getting feedback.
We didn't go too far and too deep into this.
That said, keep it in mind because it'll come up in Course 3.
In Course 3 it's all about hands on deployment,
you're going to be getting yourself a small little TinyML kit,
stay tuned for all the instructions that are coming up
and make sure you order your kit and what
we will be focusing on there is building on all the applications
that we have picked up here and learning how to deploy them.
Deploy them does not mean you simply take the model
and you just push a button, there is a fair bit
of nuanced science involved in actually knowing how to integrate something
into a real device.
Don't be worried, if you don't know much about embedded systems
we will be taking you through every single step of the way
to build your entire TinyML application on a real physical device.
That said, I want you to know that you should be proud of yourself
because you learned a lot of different things in this course.
In Course 1 you understood the language of machine learning
and some of the interesting aspects of embedded systems.
In Course 2 we looked at three really diverse applications that
have very different sensor modalities and we integrated all of that
into the single unified pipeline that we were able to repeat through.
And in Course 3, we're going to learn how
to deploy all of that on a real device.
Congratulations again.
I'll see you in Course 3.
  
\section{Kit for Course 3}

We hope you enjoyed Course 2 and we hope you’ll join us for Course 3!
Looking forward to course 3 you’ll need everything from Course 2, a laptop with a microphone and camera and an internet connection, as well as the course kit as we will be deploying our models onto microcontrollers.

You can buy the course kit directly from Arduino at this link:

\url{https://store.arduino.cc/usa/tiny-machine-learning-kit}

October 14, 2021 update: Due to the COVID-19 crisis and the resulting semiconductor shortage, the microcontroller used by the kit (and therefore the kit) is out of stock globally. We’ve been working with Arduino on this and are expecting kits to be in stock in November. We will update this page and email learners when they are available.

The course kit includes a proprietary shield which will make the camera module much more reliable. However, if you would like to purchase individual components instead of the full kit you will need to purchase the following:

\begin{itemize}
  \item Arduino Nano BLE 33 Sense WITH HEADERS
  \item ArduCam OV7675 
  \item Jumper wires and optionally a breadboard for connecting the camera to the Arduino
\end{itemize}

Keep in mind that if you don't use the Arduino kit, we have found that the camera connection is often unreliable over jumper wires, and it will be hard for your fellow classmates and others to provide guidance if you experience any issues.

\begin{figure}
    \includegraphics[width=\textwidth]{TinyML/2 Applications/akx00028_00.default_1.jpg }
    \caption{Ardunino Nano Kit}
\end{figure}

      
  